% =========================================================================
% Figura (clon de «Tema 2. Figuras»): equipotenciales y líneas de campo del
% dipolo puntual (izquierda) y del dipolo físico (derecha), como StreamPlot
% de Mathematica. Datos: figuras/datos/dip_{puntual,fisico}_*.tex.
% =========================================================================
\begin{figure}[ht]
\centering
\begin{minipage}[t]{0.49\textwidth}\centering
\begin{tikzpicture}[x=0.29cm, y=0.29cm]
  \panelMMA{10}{mmafondo}{1.2}{%
    % Dipolo puntual - equipotenciales
% Generado por scripts/generar_datos.py (no editar a mano).
\EQ{(1.224,-0.476) (1.250,-0.522) (1.275,-0.574) (1.296,-0.625) (1.313,-0.675) (1.326,-0.725) (1.337,-0.775) (1.345,-0.825) (1.350,-0.875) (1.353,-0.925) (1.354,-0.975) (1.352,-1.025) (1.347,-1.075) (1.340,-1.125) (1.331,-1.175) (1.319,-1.225) (1.305,-1.275) (1.288,-1.325) (1.269,-1.375) (1.246,-1.425) (1.221,-1.475) (1.194,-1.525) (1.162,-1.575) (1.128,-1.625) (1.089,-1.675) (1.050,-1.721) (1.000,-1.774) (0.950,-1.822) (0.900,-1.865) (0.856,-1.900) (0.800,-1.940) (0.750,-1.972) (0.702,-2.000) (0.656,-2.025) (0.604,-2.050) (0.550,-2.074) (0.500,-2.093) (0.450,-2.111) (0.400,-2.126) (0.350,-2.139) (0.300,-2.151) (0.250,-2.161) (0.200,-2.168) (0.150,-2.174) (0.100,-2.179) (0.050,-2.181) (0.000,-2.182) (-0.050,-2.181) (-0.100,-2.179) (-0.150,-2.174) (-0.200,-2.168) (-0.250,-2.161) (-0.300,-2.151) (-0.350,-2.139) (-0.400,-2.126) (-0.450,-2.111) (-0.500,-2.093) (-0.547,-2.075) (-0.600,-2.052) (-0.650,-2.028) (-0.700,-2.001) (-0.745,-1.975) (-0.800,-1.940) (-0.850,-1.904) (-0.900,-1.865) (-0.947,-1.825) (-0.999,-1.775) (-1.047,-1.725) (-1.089,-1.675) (-1.125,-1.629) (-1.162,-1.575) (-1.194,-1.525) (-1.221,-1.475) (-1.246,-1.425) (-1.269,-1.375) (-1.288,-1.325) (-1.305,-1.275) (-1.319,-1.225) (-1.331,-1.175) (-1.340,-1.125) (-1.347,-1.075) (-1.352,-1.025) (-1.354,-0.975) (-1.353,-0.925) (-1.350,-0.875) (-1.345,-0.825) (-1.337,-0.775) (-1.326,-0.725) (-1.313,-0.675) (-1.296,-0.625) (-1.275,-0.575) (-1.252,-0.525) (-1.225,-0.477)}
\EQ{(1.229,-0.446) (1.263,-0.500) (1.290,-0.550) (1.313,-0.600) (1.333,-0.650) (1.349,-0.700) (1.362,-0.750) (1.373,-0.800) (1.381,-0.850) (1.386,-0.900) (1.389,-0.950) (1.389,-1.000) (1.387,-1.050) (1.383,-1.100) (1.376,-1.150) (1.367,-1.200) (1.355,-1.250) (1.341,-1.300) (1.325,-1.349) (1.306,-1.400) (1.284,-1.450) (1.260,-1.500) (1.233,-1.550) (1.203,-1.600) (1.175,-1.642) (1.132,-1.700) (1.100,-1.740) (1.050,-1.796) (1.000,-1.846) (0.950,-1.891) (0.909,-1.925) (0.850,-1.970) (0.806,-2.000) (0.750,-2.035) (0.700,-2.063) (0.650,-2.089) (0.600,-2.112) (0.550,-2.133) (0.500,-2.152) (0.450,-2.169) (0.400,-2.184) (0.350,-2.197) (0.300,-2.208) (0.250,-2.218) (0.200,-2.225) (0.150,-2.231) (0.100,-2.236) (0.050,-2.238) (0.000,-2.239) (-0.050,-2.238) (-0.100,-2.236) (-0.150,-2.231) (-0.200,-2.225) (-0.250,-2.218) (-0.300,-2.208) (-0.350,-2.197) (-0.400,-2.184) (-0.450,-2.169) (-0.500,-2.152) (-0.550,-2.133) (-0.600,-2.112) (-0.650,-2.089) (-0.700,-2.063) (-0.750,-2.035) (-0.800,-2.004) (-0.843,-1.975) (-0.900,-1.932) (-0.940,-1.900) (-0.996,-1.850) (-1.046,-1.800) (-1.091,-1.750) (-1.125,-1.709) (-1.169,-1.650) (-1.200,-1.604) (-1.233,-1.550) (-1.260,-1.500) (-1.284,-1.450) (-1.306,-1.400) (-1.325,-1.350) (-1.341,-1.300) (-1.355,-1.250) (-1.367,-1.200) (-1.376,-1.150) (-1.383,-1.100) (-1.387,-1.050) (-1.389,-1.000) (-1.389,-0.950) (-1.386,-0.900) (-1.381,-0.850) (-1.373,-0.800) (-1.362,-0.750) (-1.349,-0.700) (-1.333,-0.650) (-1.313,-0.600) (-1.290,-0.550) (-1.263,-0.500) (-1.232,-0.450) (-1.229,-0.446)}
\EQ{(1.244,-0.431) (1.274,-0.475) (1.300,-0.518) (1.325,-0.565) (1.350,-0.619) (1.372,-0.675) (1.388,-0.725) (1.400,-0.775) (1.411,-0.825) (1.419,-0.875) (1.424,-0.925) (1.427,-0.975) (1.427,-1.025) (1.425,-1.075) (1.421,-1.125) (1.414,-1.175) (1.406,-1.225) (1.395,-1.275) (1.381,-1.325) (1.365,-1.375) (1.347,-1.425) (1.326,-1.475) (1.303,-1.525) (1.277,-1.575) (1.250,-1.621) (1.215,-1.675) (1.180,-1.725) (1.141,-1.775) (1.100,-1.822) (1.050,-1.874) (1.000,-1.922) (0.950,-1.965) (0.906,-2.000) (0.850,-2.040) (0.800,-2.073) (0.755,-2.100) (0.709,-2.125) (0.660,-2.150) (0.605,-2.175) (0.550,-2.198) (0.500,-2.216) (0.450,-2.233) (0.400,-2.247) (0.350,-2.260) (0.300,-2.271) (0.250,-2.280) (0.200,-2.287) (0.150,-2.293) (0.100,-2.297) (0.050,-2.299) (0.000,-2.300) (-0.050,-2.299) (-0.100,-2.297) (-0.150,-2.293) (-0.200,-2.287) (-0.250,-2.280) (-0.300,-2.271) (-0.350,-2.260) (-0.400,-2.247) (-0.450,-2.233) (-0.500,-2.216) (-0.550,-2.198) (-0.600,-2.177) (-0.650,-2.155) (-0.700,-2.130) (-0.750,-2.103) (-0.796,-2.075) (-0.850,-2.040) (-0.900,-2.004) (-0.950,-1.965) (-0.997,-1.925) (-1.049,-1.875) (-1.097,-1.825) (-1.141,-1.775) (-1.175,-1.731) (-1.215,-1.675) (-1.248,-1.625) (-1.275,-1.578) (-1.300,-1.531) (-1.325,-1.478) (-1.347,-1.425) (-1.365,-1.375) (-1.381,-1.325) (-1.395,-1.275) (-1.406,-1.225) (-1.414,-1.175) (-1.421,-1.125) (-1.425,-1.075) (-1.427,-1.025) (-1.427,-0.975) (-1.424,-0.925) (-1.419,-0.875) (-1.411,-0.825) (-1.400,-0.775) (-1.388,-0.725) (-1.372,-0.675) (-1.352,-0.625) (-1.330,-0.575) (-1.304,-0.525) (-1.275,-0.477) (-1.244,-0.431)}
\EQ{(1.250,-0.405) (1.285,-0.450) (1.318,-0.500) (1.347,-0.550) (1.373,-0.600) (1.395,-0.650) (1.413,-0.700) (1.429,-0.750) (1.442,-0.800) (1.452,-0.850) (1.460,-0.900) (1.465,-0.950) (1.468,-1.000) (1.468,-1.050) (1.467,-1.100) (1.463,-1.150) (1.457,-1.200) (1.448,-1.250) (1.438,-1.300) (1.425,-1.350) (1.410,-1.400) (1.392,-1.450) (1.372,-1.500) (1.350,-1.550) (1.325,-1.600) (1.300,-1.645) (1.266,-1.700) (1.232,-1.750) (1.200,-1.793) (1.154,-1.850) (1.109,-1.900) (1.059,-1.950) (1.004,-2.000) (0.950,-2.044) (0.909,-2.075) (0.850,-2.116) (0.800,-2.147) (0.751,-2.175) (0.704,-2.200) (0.652,-2.225) (0.600,-2.248) (0.550,-2.268) (0.500,-2.285) (0.450,-2.301) (0.400,-2.315) (0.350,-2.328) (0.300,-2.338) (0.250,-2.347) (0.200,-2.354) (0.150,-2.360) (0.100,-2.364) (0.050,-2.366) (0.000,-2.367) (-0.050,-2.366) (-0.100,-2.364) (-0.150,-2.360) (-0.200,-2.354) (-0.250,-2.347) (-0.300,-2.338) (-0.350,-2.328) (-0.400,-2.315) (-0.450,-2.301) (-0.500,-2.285) (-0.550,-2.268) (-0.600,-2.248) (-0.650,-2.226) (-0.700,-2.202) (-0.750,-2.176) (-0.795,-2.150) (-0.850,-2.116) (-0.900,-2.081) (-0.942,-2.050) (-1.000,-2.003) (-1.050,-1.958) (-1.100,-1.909) (-1.150,-1.854) (-1.195,-1.800) (-1.225,-1.760) (-1.266,-1.700) (-1.297,-1.650) (-1.325,-1.600) (-1.350,-1.550) (-1.372,-1.500) (-1.392,-1.450) (-1.410,-1.400) (-1.425,-1.350) (-1.438,-1.300) (-1.448,-1.250) (-1.457,-1.200) (-1.463,-1.150) (-1.467,-1.100) (-1.468,-1.050) (-1.468,-1.000) (-1.465,-0.950) (-1.460,-0.900) (-1.452,-0.850) (-1.442,-0.800) (-1.429,-0.750) (-1.413,-0.700) (-1.395,-0.650) (-1.373,-0.600) (-1.350,-0.555) (-1.325,-0.512) (-1.285,-0.450) (-1.250,-0.405)}
\EQ{(1.251,-0.374) (1.295,-0.425) (1.325,-0.465) (1.365,-0.525) (1.393,-0.575) (1.418,-0.625) (1.440,-0.675) (1.458,-0.725) (1.473,-0.775) (1.486,-0.825) (1.496,-0.875) (1.504,-0.925) (1.510,-0.975) (1.513,-1.025) (1.514,-1.075) (1.512,-1.125) (1.509,-1.175) (1.503,-1.225) (1.495,-1.275) (1.486,-1.325) (1.473,-1.375) (1.459,-1.425) (1.442,-1.475) (1.423,-1.525) (1.402,-1.575) (1.378,-1.625) (1.352,-1.675) (1.325,-1.721) (1.290,-1.775) (1.255,-1.825) (1.216,-1.875) (1.175,-1.923) (1.127,-1.975) (1.075,-2.025) (1.025,-2.070) (0.975,-2.110) (0.925,-2.147) (0.875,-2.182) (0.825,-2.213) (0.775,-2.242) (0.725,-2.268) (0.675,-2.292) (0.625,-2.314) (0.575,-2.334) (0.525,-2.352) (0.475,-2.369) (0.425,-2.383) (0.375,-2.396) (0.325,-2.407) (0.275,-2.416) (0.225,-2.424) (0.175,-2.430) (0.125,-2.435) (0.075,-2.438) (0.025,-2.440) (-0.050,-2.439) (-0.100,-2.437) (-0.150,-2.433) (-0.200,-2.427) (-0.250,-2.420) (-0.300,-2.412) (-0.350,-2.402) (-0.400,-2.390) (-0.450,-2.376) (-0.500,-2.361) (-0.550,-2.344) (-0.599,-2.325) (-0.650,-2.304) (-0.700,-2.280) (-0.750,-2.255) (-0.800,-2.228) (-0.846,-2.200) (-0.900,-2.165) (-0.950,-2.129) (-1.000,-2.090) (-1.048,-2.050) (-1.100,-2.002) (-1.150,-1.951) (-1.195,-1.900) (-1.236,-1.850) (-1.273,-1.800) (-1.307,-1.750) (-1.337,-1.700) (-1.365,-1.650) (-1.390,-1.600) (-1.413,-1.550) (-1.433,-1.500) (-1.450,-1.453) (-1.466,-1.400) (-1.480,-1.350) (-1.491,-1.300) (-1.500,-1.250) (-1.506,-1.200) (-1.511,-1.150) (-1.513,-1.100) (-1.514,-1.050) (-1.512,-1.000) (-1.507,-0.950) (-1.501,-0.900) (-1.492,-0.850) (-1.480,-0.800) (-1.466,-0.750) (-1.449,-0.700) (-1.429,-0.650) (-1.406,-0.600) (-1.379,-0.550) (-1.350,-0.502) (-1.314,-0.450) (-1.275,-0.401) (-1.251,-0.374)}
\EQ{(1.266,-0.359) (1.305,-0.400) (1.346,-0.450) (1.382,-0.500) (1.414,-0.550) (1.442,-0.600) (1.467,-0.650) (1.488,-0.700) (1.506,-0.750) (1.522,-0.800) (1.534,-0.850) (1.545,-0.900) (1.553,-0.950) (1.559,-1.000) (1.562,-1.050) (1.563,-1.100) (1.563,-1.150) (1.560,-1.200) (1.555,-1.250) (1.547,-1.300) (1.538,-1.350) (1.527,-1.400) (1.513,-1.450) (1.498,-1.500) (1.480,-1.550) (1.460,-1.600) (1.437,-1.650) (1.412,-1.700) (1.385,-1.750) (1.354,-1.800) (1.325,-1.844) (1.285,-1.900) (1.250,-1.944) (1.201,-2.000) (1.154,-2.050) (1.101,-2.100) (1.050,-2.145) (1.000,-2.185) (0.950,-2.222) (0.908,-2.250) (0.850,-2.287) (0.800,-2.315) (0.750,-2.342) (0.700,-2.366) (0.650,-2.388) (0.600,-2.409) (0.550,-2.427) (0.500,-2.443) (0.450,-2.458) (0.400,-2.471) (0.350,-2.483) (0.300,-2.493) (0.250,-2.501) (0.200,-2.508) (0.150,-2.513) (0.100,-2.517) (0.050,-2.519) (0.000,-2.520) (-0.050,-2.519) (-0.100,-2.517) (-0.150,-2.513) (-0.200,-2.508) (-0.250,-2.501) (-0.300,-2.493) (-0.350,-2.483) (-0.400,-2.471) (-0.450,-2.458) (-0.500,-2.443) (-0.550,-2.427) (-0.600,-2.409) (-0.650,-2.388) (-0.700,-2.366) (-0.750,-2.342) (-0.800,-2.315) (-0.850,-2.287) (-0.900,-2.256) (-0.945,-2.225) (-1.000,-2.185) (-1.044,-2.150) (-1.100,-2.101) (-1.150,-2.054) (-1.200,-2.001) (-1.245,-1.950) (-1.285,-1.900) (-1.321,-1.850) (-1.350,-1.807) (-1.385,-1.750) (-1.412,-1.700) (-1.437,-1.650) (-1.460,-1.600) (-1.480,-1.550) (-1.498,-1.500) (-1.513,-1.450) (-1.527,-1.400) (-1.538,-1.350) (-1.547,-1.300) (-1.555,-1.250) (-1.560,-1.200) (-1.563,-1.150) (-1.563,-1.100) (-1.562,-1.050) (-1.559,-1.000) (-1.553,-0.950) (-1.545,-0.900) (-1.534,-0.850) (-1.522,-0.800) (-1.506,-0.750) (-1.488,-0.700) (-1.467,-0.650) (-1.442,-0.600) (-1.414,-0.550) (-1.382,-0.500) (-1.350,-0.456) (-1.305,-0.400) (-1.266,-0.359)}
\EQ{(1.275,-0.337) (1.314,-0.375) (1.350,-0.414) (1.400,-0.475) (1.436,-0.525) (1.467,-0.575) (1.495,-0.625) (1.519,-0.675) (1.540,-0.725) (1.559,-0.775) (1.574,-0.825) (1.587,-0.875) (1.598,-0.925) (1.606,-0.975) (1.612,-1.025) (1.616,-1.075) (1.618,-1.125) (1.618,-1.175) (1.616,-1.225) (1.611,-1.275) (1.605,-1.325) (1.597,-1.375) (1.586,-1.425) (1.574,-1.475) (1.559,-1.525) (1.543,-1.575) (1.524,-1.625) (1.503,-1.675) (1.479,-1.725) (1.454,-1.775) (1.425,-1.825) (1.394,-1.875) (1.360,-1.925) (1.325,-1.973) (1.283,-2.025) (1.250,-2.063) (1.216,-2.100) (1.175,-2.141) (1.125,-2.187) (1.081,-2.225) (1.025,-2.269) (0.983,-2.300) (0.925,-2.339) (0.875,-2.370) (0.825,-2.398) (0.775,-2.425) (0.725,-2.449) (0.675,-2.471) (0.625,-2.492) (0.575,-2.510) (0.525,-2.527) (0.475,-2.542) (0.425,-2.555) (0.375,-2.567) (0.325,-2.578) (0.275,-2.586) (0.225,-2.594) (0.175,-2.599) (0.125,-2.604) (0.075,-2.607) (0.025,-2.608) (-0.050,-2.608) (-0.100,-2.605) (-0.150,-2.602) (-0.200,-2.597) (-0.250,-2.590) (-0.300,-2.582) (-0.350,-2.573) (-0.400,-2.562) (-0.450,-2.549) (-0.500,-2.535) (-0.550,-2.519) (-0.600,-2.501) (-0.650,-2.482) (-0.700,-2.460) (-0.750,-2.437) (-0.800,-2.412) (-0.850,-2.384) (-0.900,-2.355) (-0.946,-2.325) (-1.000,-2.288) (-1.050,-2.250) (-1.100,-2.209) (-1.139,-2.175) (-1.191,-2.125) (-1.239,-2.075) (-1.275,-2.034) (-1.323,-1.975) (-1.360,-1.925) (-1.394,-1.875) (-1.425,-1.825) (-1.450,-1.781) (-1.475,-1.734) (-1.500,-1.681) (-1.524,-1.625) (-1.543,-1.575) (-1.559,-1.525) (-1.574,-1.475) (-1.586,-1.425) (-1.597,-1.375) (-1.605,-1.325) (-1.611,-1.275) (-1.616,-1.225) (-1.618,-1.175) (-1.618,-1.125) (-1.616,-1.075) (-1.612,-1.025) (-1.606,-0.975) (-1.598,-0.925) (-1.587,-0.875) (-1.574,-0.825) (-1.559,-0.775) (-1.540,-0.725) (-1.519,-0.675) (-1.495,-0.625) (-1.467,-0.575) (-1.436,-0.525) (-1.400,-0.475) (-1.360,-0.425) (-1.325,-0.387) (-1.289,-0.350) (-1.275,-0.337)}
\EQ{(1.275,-0.308) (1.323,-0.350) (1.374,-0.400) (1.419,-0.450) (1.450,-0.489) (1.493,-0.550) (1.524,-0.600) (1.550,-0.647) (1.575,-0.698) (1.597,-0.750) (1.616,-0.800) (1.632,-0.850) (1.645,-0.900) (1.656,-0.950) (1.665,-1.000) (1.672,-1.050) (1.676,-1.100) (1.679,-1.150) (1.679,-1.200) (1.678,-1.250) (1.674,-1.300) (1.669,-1.350) (1.662,-1.400) (1.652,-1.450) (1.641,-1.500) (1.628,-1.550) (1.613,-1.600) (1.595,-1.650) (1.576,-1.700) (1.554,-1.750) (1.530,-1.800) (1.504,-1.850) (1.475,-1.900) (1.444,-1.950) (1.410,-2.000) (1.375,-2.047) (1.333,-2.100) (1.289,-2.150) (1.250,-2.191) (1.200,-2.240) (1.162,-2.275) (1.102,-2.325) (1.050,-2.365) (1.001,-2.400) (0.950,-2.434) (0.900,-2.464) (0.850,-2.492) (0.800,-2.518) (0.750,-2.543) (0.700,-2.565) (0.650,-2.585) (0.600,-2.604) (0.550,-2.621) (0.500,-2.636) (0.450,-2.650) (0.400,-2.662) (0.350,-2.672) (0.300,-2.682) (0.250,-2.689) (0.200,-2.696) (0.150,-2.700) (0.100,-2.704) (0.050,-2.706) (0.000,-2.707) (-0.050,-2.706) (-0.100,-2.704) (-0.150,-2.700) (-0.200,-2.696) (-0.250,-2.689) (-0.300,-2.682) (-0.350,-2.672) (-0.400,-2.662) (-0.450,-2.650) (-0.500,-2.636) (-0.550,-2.621) (-0.600,-2.604) (-0.650,-2.585) (-0.700,-2.565) (-0.750,-2.543) (-0.800,-2.518) (-0.850,-2.492) (-0.900,-2.464) (-0.950,-2.434) (-1.000,-2.401) (-1.050,-2.365) (-1.100,-2.327) (-1.150,-2.285) (-1.189,-2.250) (-1.241,-2.200) (-1.289,-2.150) (-1.325,-2.109) (-1.373,-2.050) (-1.410,-2.000) (-1.444,-1.950) (-1.475,-1.900) (-1.504,-1.850) (-1.530,-1.800) (-1.554,-1.750) (-1.575,-1.702) (-1.595,-1.650) (-1.613,-1.600) (-1.628,-1.550) (-1.641,-1.500) (-1.652,-1.450) (-1.662,-1.400) (-1.669,-1.350) (-1.674,-1.300) (-1.678,-1.250) (-1.679,-1.200) (-1.679,-1.150) (-1.676,-1.100) (-1.672,-1.050) (-1.665,-1.000) (-1.656,-0.950) (-1.645,-0.900) (-1.632,-0.850) (-1.616,-0.800) (-1.597,-0.750) (-1.576,-0.700) (-1.552,-0.650) (-1.525,-0.601) (-1.493,-0.550) (-1.458,-0.500) (-1.425,-0.458) (-1.375,-0.401) (-1.325,-0.352) (-1.275,-0.308)}
\EQ{(1.275,-0.280) (1.325,-0.319) (1.375,-0.362) (1.414,-0.400) (1.450,-0.438) (1.500,-0.497) (1.539,-0.550) (1.572,-0.600) (1.600,-0.648) (1.625,-0.696) (1.650,-0.750) (1.670,-0.800) (1.688,-0.850) (1.703,-0.900) (1.716,-0.950) (1.726,-1.000) (1.734,-1.050) (1.741,-1.100) (1.745,-1.150) (1.747,-1.200) (1.748,-1.250) (1.746,-1.300) (1.743,-1.350) (1.738,-1.400) (1.731,-1.450) (1.722,-1.500) (1.711,-1.550) (1.698,-1.600) (1.683,-1.650) (1.667,-1.700) (1.648,-1.750) (1.627,-1.800) (1.604,-1.850) (1.579,-1.900) (1.552,-1.950) (1.525,-1.995) (1.489,-2.050) (1.454,-2.100) (1.416,-2.150) (1.375,-2.200) (1.330,-2.250) (1.282,-2.300) (1.229,-2.350) (1.175,-2.397) (1.125,-2.437) (1.075,-2.475) (1.025,-2.509) (0.975,-2.542) (0.925,-2.571) (0.875,-2.599) (0.825,-2.625) (0.775,-2.649) (0.725,-2.671) (0.675,-2.691) (0.625,-2.710) (0.575,-2.727) (0.525,-2.742) (0.475,-2.756) (0.425,-2.768) (0.375,-2.779) (0.325,-2.789) (0.275,-2.797) (0.225,-2.804) (0.175,-2.809) (0.125,-2.813) (0.075,-2.816) (0.025,-2.817) (-0.050,-2.817) (-0.100,-2.815) (-0.150,-2.811) (-0.200,-2.807) (-0.250,-2.800) (-0.300,-2.793) (-0.350,-2.784) (-0.400,-2.774) (-0.450,-2.762) (-0.500,-2.749) (-0.550,-2.735) (-0.600,-2.719) (-0.650,-2.701) (-0.700,-2.681) (-0.750,-2.660) (-0.800,-2.637) (-0.850,-2.612) (-0.900,-2.586) (-0.950,-2.557) (-1.000,-2.526) (-1.050,-2.492) (-1.100,-2.456) (-1.141,-2.425) (-1.200,-2.376) (-1.250,-2.331) (-1.300,-2.282) (-1.350,-2.228) (-1.396,-2.175) (-1.436,-2.125) (-1.472,-2.075) (-1.506,-2.025) (-1.537,-1.975) (-1.566,-1.925) (-1.592,-1.875) (-1.616,-1.825) (-1.638,-1.775) (-1.658,-1.725) (-1.675,-1.675) (-1.691,-1.625) (-1.705,-1.575) (-1.716,-1.525) (-1.726,-1.475) (-1.734,-1.425) (-1.741,-1.375) (-1.745,-1.325) (-1.747,-1.275) (-1.748,-1.225) (-1.746,-1.175) (-1.743,-1.125) (-1.738,-1.075) (-1.730,-1.025) (-1.721,-0.975) (-1.709,-0.925) (-1.695,-0.875) (-1.679,-0.825) (-1.660,-0.775) (-1.639,-0.725) (-1.614,-0.675) (-1.587,-0.625) (-1.556,-0.575) (-1.525,-0.530) (-1.482,-0.475) (-1.438,-0.425) (-1.400,-0.386) (-1.362,-0.350) (-1.301,-0.300) (-1.275,-0.280)}
\EQ{(1.288,-0.262) (1.342,-0.300) (1.400,-0.347) (1.450,-0.391) (1.500,-0.442) (1.550,-0.499) (1.590,-0.550) (1.625,-0.600) (1.656,-0.650) (1.684,-0.700) (1.708,-0.750) (1.730,-0.800) (1.750,-0.850) (1.767,-0.900) (1.781,-0.950) (1.794,-1.000) (1.804,-1.050) (1.812,-1.100) (1.818,-1.150) (1.823,-1.200) (1.825,-1.250) (1.825,-1.300) (1.824,-1.350) (1.821,-1.400) (1.816,-1.450) (1.810,-1.500) (1.801,-1.550) (1.791,-1.600) (1.779,-1.650) (1.765,-1.700) (1.750,-1.748) (1.732,-1.800) (1.712,-1.850) (1.691,-1.900) (1.667,-1.950) (1.641,-2.000) (1.613,-2.050) (1.583,-2.100) (1.550,-2.149) (1.514,-2.200) (1.476,-2.250) (1.434,-2.300) (1.400,-2.339) (1.366,-2.375) (1.325,-2.416) (1.289,-2.450) (1.250,-2.484) (1.201,-2.525) (1.150,-2.565) (1.101,-2.600) (1.050,-2.634) (1.000,-2.666) (0.950,-2.695) (0.900,-2.722) (0.850,-2.748) (0.800,-2.771) (0.750,-2.793) (0.700,-2.813) (0.650,-2.831) (0.600,-2.848) (0.550,-2.864) (0.500,-2.878) (0.450,-2.890) (0.400,-2.901) (0.350,-2.911) (0.300,-2.919) (0.250,-2.926) (0.200,-2.932) (0.150,-2.937) (0.100,-2.940) (0.050,-2.942) (0.000,-2.943) (-0.050,-2.942) (-0.100,-2.940) (-0.150,-2.937) (-0.200,-2.932) (-0.250,-2.926) (-0.300,-2.919) (-0.350,-2.911) (-0.400,-2.901) (-0.450,-2.890) (-0.500,-2.878) (-0.550,-2.864) (-0.600,-2.848) (-0.650,-2.831) (-0.700,-2.813) (-0.750,-2.793) (-0.800,-2.771) (-0.850,-2.748) (-0.900,-2.722) (-0.950,-2.695) (-1.000,-2.666) (-1.050,-2.634) (-1.100,-2.601) (-1.150,-2.565) (-1.200,-2.526) (-1.250,-2.484) (-1.289,-2.450) (-1.325,-2.416) (-1.366,-2.375) (-1.400,-2.339) (-1.434,-2.300) (-1.475,-2.251) (-1.514,-2.200) (-1.550,-2.150) (-1.583,-2.100) (-1.613,-2.050) (-1.641,-2.000) (-1.667,-1.950) (-1.691,-1.900) (-1.712,-1.850) (-1.732,-1.800) (-1.749,-1.750) (-1.765,-1.700) (-1.779,-1.650) (-1.791,-1.600) (-1.801,-1.550) (-1.810,-1.500) (-1.816,-1.450) (-1.821,-1.400) (-1.824,-1.350) (-1.825,-1.300) (-1.825,-1.250) (-1.823,-1.200) (-1.818,-1.150) (-1.812,-1.100) (-1.804,-1.050) (-1.794,-1.000) (-1.781,-0.950) (-1.767,-0.900) (-1.750,-0.851) (-1.730,-0.800) (-1.708,-0.750) (-1.684,-0.700) (-1.656,-0.650) (-1.625,-0.601) (-1.590,-0.550) (-1.551,-0.500) (-1.508,-0.450) (-1.459,-0.400) (-1.404,-0.350) (-1.350,-0.307) (-1.307,-0.275) (-1.288,-0.262)}
\EQ{(1.300,-0.243) (1.350,-0.275) (1.400,-0.310) (1.450,-0.348) (1.500,-0.391) (1.550,-0.439) (1.600,-0.493) (1.647,-0.550) (1.684,-0.600) (1.717,-0.650) (1.747,-0.700) (1.774,-0.750) (1.798,-0.800) (1.819,-0.850) (1.838,-0.900) (1.855,-0.950) (1.869,-1.000) (1.881,-1.050) (1.892,-1.100) (1.900,-1.150) (1.906,-1.200) (1.911,-1.250) (1.914,-1.300) (1.915,-1.350) (1.914,-1.400) (1.911,-1.450) (1.907,-1.500) (1.901,-1.550) (1.894,-1.600) (1.885,-1.650) (1.873,-1.700) (1.861,-1.750) (1.846,-1.800) (1.830,-1.850) (1.812,-1.900) (1.792,-1.950) (1.770,-2.000) (1.746,-2.050) (1.720,-2.100) (1.692,-2.150) (1.661,-2.200) (1.629,-2.250) (1.593,-2.300) (1.555,-2.350) (1.515,-2.400) (1.475,-2.445) (1.425,-2.498) (1.375,-2.548) (1.325,-2.593) (1.275,-2.636) (1.226,-2.675) (1.175,-2.713) (1.125,-2.747) (1.083,-2.775) (1.025,-2.810) (0.975,-2.838) (0.925,-2.865) (0.875,-2.890) (0.825,-2.913) (0.775,-2.934) (0.725,-2.954) (0.675,-2.972) (0.625,-2.989) (0.575,-3.004) (0.525,-3.018) (0.475,-3.030) (0.425,-3.042) (0.375,-3.052) (0.325,-3.060) (0.275,-3.068) (0.225,-3.074) (0.175,-3.079) (0.125,-3.082) (0.075,-3.085) (0.025,-3.086) (-0.050,-3.086) (-0.100,-3.084) (-0.150,-3.081) (-0.200,-3.076) (-0.250,-3.071) (-0.300,-3.064) (-0.350,-3.056) (-0.400,-3.047) (-0.450,-3.036) (-0.500,-3.024) (-0.550,-3.011) (-0.600,-2.996) (-0.650,-2.980) (-0.700,-2.963) (-0.750,-2.944) (-0.800,-2.923) (-0.850,-2.901) (-0.900,-2.877) (-0.950,-2.852) (-0.999,-2.825) (-1.050,-2.795) (-1.100,-2.764) (-1.150,-2.730) (-1.192,-2.700) (-1.250,-2.656) (-1.300,-2.615) (-1.345,-2.575) (-1.398,-2.525) (-1.448,-2.475) (-1.493,-2.425) (-1.535,-2.375) (-1.575,-2.325) (-1.611,-2.275) (-1.645,-2.225) (-1.675,-2.178) (-1.706,-2.125) (-1.733,-2.075) (-1.758,-2.025) (-1.781,-1.975) (-1.802,-1.925) (-1.821,-1.875) (-1.838,-1.825) (-1.854,-1.775) (-1.867,-1.725) (-1.879,-1.675) (-1.889,-1.625) (-1.898,-1.575) (-1.905,-1.525) (-1.910,-1.475) (-1.913,-1.425) (-1.915,-1.375) (-1.914,-1.325) (-1.913,-1.275) (-1.909,-1.225) (-1.903,-1.175) (-1.896,-1.125) (-1.887,-1.075) (-1.875,-1.025) (-1.862,-0.975) (-1.847,-0.925) (-1.829,-0.875) (-1.809,-0.825) (-1.786,-0.775) (-1.761,-0.725) (-1.733,-0.675) (-1.701,-0.625) (-1.666,-0.575) (-1.627,-0.525) (-1.584,-0.475) (-1.536,-0.425) (-1.482,-0.375) (-1.425,-0.328) (-1.375,-0.292) (-1.325,-0.258) (-1.300,-0.243)}
\EQ{(1.300,-0.216) (1.350,-0.244) (1.400,-0.274) (1.450,-0.308) (1.500,-0.345) (1.550,-0.385) (1.595,-0.425) (1.646,-0.475) (1.692,-0.525) (1.725,-0.565) (1.770,-0.625) (1.800,-0.669) (1.834,-0.725) (1.862,-0.775) (1.887,-0.825) (1.909,-0.875) (1.929,-0.925) (1.947,-0.975) (1.962,-1.025) (1.975,-1.075) (1.987,-1.125) (1.997,-1.175) (2.004,-1.225) (2.011,-1.275) (2.015,-1.325) (2.017,-1.375) (2.018,-1.425) (2.017,-1.475) (2.015,-1.525) (2.011,-1.575) (2.005,-1.625) (1.998,-1.675) (1.989,-1.725) (1.979,-1.775) (1.967,-1.825) (1.953,-1.875) (1.937,-1.925) (1.920,-1.975) (1.901,-2.025) (1.880,-2.075) (1.858,-2.125) (1.833,-2.175) (1.807,-2.225) (1.778,-2.275) (1.750,-2.321) (1.714,-2.375) (1.679,-2.425) (1.641,-2.475) (1.601,-2.525) (1.557,-2.575) (1.510,-2.625) (1.475,-2.661) (1.434,-2.700) (1.378,-2.750) (1.325,-2.794) (1.285,-2.825) (1.225,-2.869) (1.180,-2.900) (1.125,-2.935) (1.075,-2.965) (1.025,-2.993) (0.975,-3.020) (0.925,-3.045) (0.875,-3.068) (0.825,-3.089) (0.775,-3.109) (0.725,-3.128) (0.675,-3.145) (0.625,-3.161) (0.576,-3.175) (0.525,-3.188) (0.476,-3.200) (0.425,-3.211) (0.375,-3.220) (0.325,-3.229) (0.275,-3.236) (0.225,-3.241) (0.175,-3.246) (0.125,-3.249) (0.075,-3.252) (0.025,-3.253) (-0.050,-3.252) (-0.100,-3.251) (-0.150,-3.248) (-0.200,-3.244) (-0.250,-3.239) (-0.300,-3.232) (-0.350,-3.225) (-0.400,-3.216) (-0.450,-3.206) (-0.500,-3.195) (-0.550,-3.182) (-0.600,-3.168) (-0.650,-3.153) (-0.700,-3.137) (-0.750,-3.119) (-0.798,-3.100) (-0.850,-3.079) (-0.900,-3.056) (-0.950,-3.032) (-1.000,-3.007) (-1.050,-2.980) (-1.100,-2.950) (-1.150,-2.919) (-1.200,-2.886) (-1.250,-2.851) (-1.300,-2.814) (-1.348,-2.775) (-1.400,-2.731) (-1.450,-2.685) (-1.486,-2.650) (-1.525,-2.610) (-1.575,-2.555) (-1.621,-2.500) (-1.661,-2.450) (-1.697,-2.400) (-1.731,-2.350) (-1.763,-2.300) (-1.793,-2.250) (-1.820,-2.200) (-1.846,-2.150) (-1.869,-2.100) (-1.891,-2.050) (-1.911,-2.000) (-1.929,-1.950) (-1.945,-1.900) (-1.960,-1.850) (-1.973,-1.800) (-1.984,-1.750) (-1.994,-1.700) (-2.002,-1.650) (-2.008,-1.600) (-2.013,-1.550) (-2.017,-1.500) (-2.018,-1.450) (-2.018,-1.400) (-2.016,-1.350) (-2.013,-1.300) (-2.008,-1.250) (-2.001,-1.200) (-1.992,-1.150) (-1.982,-1.100) (-1.969,-1.050) (-1.955,-1.000) (-1.938,-0.950) (-1.919,-0.900) (-1.898,-0.850) (-1.875,-0.801) (-1.850,-0.753) (-1.825,-0.710) (-1.787,-0.650) (-1.752,-0.600) (-1.713,-0.550) (-1.675,-0.506) (-1.625,-0.454) (-1.575,-0.407) (-1.525,-0.364) (-1.475,-0.326) (-1.425,-0.291) (-1.375,-0.259) (-1.325,-0.230) (-1.300,-0.216)}
\EQ{(1.300,-0.191) (1.350,-0.215) (1.400,-0.241) (1.450,-0.270) (1.499,-0.300) (1.550,-0.335) (1.600,-0.372) (1.650,-0.414) (1.691,-0.450) (1.742,-0.500) (1.788,-0.550) (1.825,-0.594) (1.868,-0.650) (1.900,-0.696) (1.934,-0.750) (1.963,-0.800) (1.989,-0.850) (2.013,-0.900) (2.034,-0.950) (2.053,-1.000) (2.070,-1.050) (2.085,-1.100) (2.098,-1.150) (2.109,-1.200) (2.118,-1.250) (2.126,-1.300) (2.132,-1.350) (2.137,-1.400) (2.139,-1.450) (2.141,-1.500) (2.140,-1.550) (2.138,-1.600) (2.135,-1.650) (2.130,-1.700) (2.123,-1.750) (2.115,-1.800) (2.106,-1.850) (2.095,-1.900) (2.082,-1.950) (2.068,-2.000) (2.052,-2.050) (2.035,-2.100) (2.016,-2.150) (1.995,-2.200) (1.972,-2.250) (1.950,-2.296) (1.925,-2.344) (1.900,-2.389) (1.864,-2.450) (1.831,-2.500) (1.800,-2.546) (1.760,-2.600) (1.725,-2.645) (1.679,-2.700) (1.634,-2.750) (1.600,-2.786) (1.562,-2.825) (1.525,-2.860) (1.481,-2.900) (1.425,-2.947) (1.375,-2.987) (1.325,-3.025) (1.275,-3.060) (1.225,-3.093) (1.175,-3.124) (1.132,-3.150) (1.087,-3.175) (1.040,-3.200) (0.990,-3.225) (0.936,-3.250) (0.879,-3.275) (0.825,-3.297) (0.775,-3.315) (0.725,-3.333) (0.675,-3.349) (0.625,-3.364) (0.575,-3.377) (0.525,-3.390) (0.475,-3.401) (0.425,-3.411) (0.375,-3.420) (0.325,-3.427) (0.275,-3.434) (0.225,-3.439) (0.175,-3.444) (0.125,-3.447) (0.075,-3.449) (0.025,-3.450) (-0.037,-3.450) (-0.100,-3.448) (-0.150,-3.445) (-0.200,-3.442) (-0.250,-3.437) (-0.300,-3.431) (-0.350,-3.424) (-0.400,-3.415) (-0.450,-3.406) (-0.500,-3.395) (-0.550,-3.384) (-0.600,-3.371) (-0.650,-3.356) (-0.700,-3.341) (-0.748,-3.325) (-0.800,-3.306) (-0.850,-3.287) (-0.900,-3.266) (-0.950,-3.244) (-1.000,-3.220) (-1.050,-3.195) (-1.100,-3.168) (-1.150,-3.139) (-1.200,-3.109) (-1.250,-3.077) (-1.300,-3.043) (-1.350,-3.006) (-1.391,-2.975) (-1.450,-2.927) (-1.500,-2.883) (-1.550,-2.836) (-1.587,-2.800) (-1.625,-2.760) (-1.675,-2.705) (-1.721,-2.650) (-1.760,-2.600) (-1.797,-2.550) (-1.831,-2.500) (-1.864,-2.450) (-1.894,-2.400) (-1.922,-2.350) (-1.948,-2.300) (-1.972,-2.250) (-1.995,-2.200) (-2.016,-2.150) (-2.035,-2.100) (-2.052,-2.050) (-2.068,-2.000) (-2.082,-1.950) (-2.095,-1.900) (-2.106,-1.850) (-2.115,-1.800) (-2.123,-1.750) (-2.130,-1.700) (-2.135,-1.650) (-2.138,-1.600) (-2.140,-1.550) (-2.141,-1.500) (-2.139,-1.450) (-2.137,-1.400) (-2.132,-1.350) (-2.126,-1.300) (-2.118,-1.250) (-2.109,-1.200) (-2.098,-1.150) (-2.085,-1.100) (-2.070,-1.050) (-2.053,-1.000) (-2.034,-0.950) (-2.013,-0.900) (-1.989,-0.850) (-1.963,-0.800) (-1.934,-0.750) (-1.903,-0.700) (-1.868,-0.650) (-1.830,-0.600) (-1.788,-0.550) (-1.750,-0.508) (-1.700,-0.459) (-1.650,-0.414) (-1.603,-0.375) (-1.550,-0.335) (-1.500,-0.301) (-1.450,-0.270) (-1.400,-0.241) (-1.350,-0.215) (-1.300,-0.191)}
\EQ{(1.300,-0.165) (1.350,-0.186) (1.400,-0.208) (1.450,-0.233) (1.500,-0.259) (1.550,-0.288) (1.600,-0.319) (1.646,-0.350) (1.700,-0.390) (1.744,-0.425) (1.800,-0.474) (1.850,-0.522) (1.900,-0.575) (1.943,-0.625) (1.975,-0.665) (2.018,-0.725) (2.050,-0.773) (2.082,-0.825) (2.109,-0.875) (2.134,-0.925) (2.157,-0.975) (2.178,-1.025) (2.197,-1.075) (2.213,-1.125) (2.228,-1.175) (2.242,-1.225) (2.253,-1.275) (2.263,-1.325) (2.271,-1.375) (2.277,-1.425) (2.282,-1.475) (2.286,-1.525) (2.288,-1.575) (2.288,-1.625) (2.288,-1.675) (2.285,-1.725) (2.281,-1.775) (2.276,-1.825) (2.269,-1.875) (2.261,-1.925) (2.252,-1.975) (2.241,-2.025) (2.228,-2.075) (2.215,-2.125) (2.199,-2.175) (2.182,-2.225) (2.164,-2.275) (2.144,-2.325) (2.122,-2.375) (2.100,-2.423) (2.075,-2.473) (2.050,-2.520) (2.025,-2.565) (1.989,-2.625) (1.956,-2.675) (1.925,-2.720) (1.885,-2.775) (1.850,-2.820) (1.805,-2.875) (1.761,-2.925) (1.725,-2.964) (1.690,-3.000) (1.650,-3.039) (1.611,-3.075) (1.554,-3.125) (1.500,-3.170) (1.450,-3.209) (1.400,-3.245) (1.358,-3.275) (1.300,-3.313) (1.250,-3.344) (1.200,-3.373) (1.152,-3.400) (1.104,-3.425) (1.053,-3.450) (1.000,-3.475) (0.950,-3.497) (0.900,-3.517) (0.850,-3.536) (0.800,-3.554) (0.750,-3.571) (0.700,-3.587) (0.650,-3.601) (0.600,-3.614) (0.550,-3.626) (0.500,-3.637) (0.450,-3.647) (0.400,-3.656) (0.350,-3.664) (0.300,-3.670) (0.250,-3.676) (0.200,-3.680) (0.150,-3.684) (0.100,-3.687) (0.050,-3.688) (0.000,-3.689) (-0.050,-3.688) (-0.100,-3.687) (-0.150,-3.684) (-0.200,-3.680) (-0.250,-3.676) (-0.300,-3.670) (-0.350,-3.664) (-0.400,-3.656) (-0.450,-3.647) (-0.500,-3.637) (-0.550,-3.626) (-0.600,-3.614) (-0.650,-3.601) (-0.700,-3.587) (-0.750,-3.571) (-0.800,-3.554) (-0.850,-3.536) (-0.900,-3.517) (-0.950,-3.497) (-1.000,-3.475) (-1.050,-3.452) (-1.100,-3.427) (-1.150,-3.401) (-1.197,-3.375) (-1.250,-3.344) (-1.300,-3.313) (-1.350,-3.280) (-1.394,-3.250) (-1.450,-3.209) (-1.494,-3.175) (-1.550,-3.129) (-1.600,-3.085) (-1.638,-3.050) (-1.675,-3.015) (-1.714,-2.975) (-1.750,-2.937) (-1.800,-2.881) (-1.846,-2.825) (-1.885,-2.775) (-1.922,-2.725) (-1.956,-2.675) (-1.989,-2.625) (-2.019,-2.575) (-2.047,-2.525) (-2.074,-2.475) (-2.099,-2.425) (-2.122,-2.375) (-2.144,-2.325) (-2.164,-2.275) (-2.182,-2.225) (-2.199,-2.175) (-2.215,-2.125) (-2.228,-2.075) (-2.241,-2.025) (-2.252,-1.975) (-2.261,-1.925) (-2.269,-1.875) (-2.276,-1.825) (-2.281,-1.775) (-2.285,-1.725) (-2.288,-1.675) (-2.288,-1.625) (-2.288,-1.575) (-2.286,-1.525) (-2.282,-1.475) (-2.277,-1.425) (-2.271,-1.375) (-2.263,-1.325) (-2.253,-1.275) (-2.242,-1.225) (-2.228,-1.175) (-2.213,-1.125) (-2.197,-1.075) (-2.178,-1.025) (-2.157,-0.975) (-2.134,-0.925) (-2.109,-0.875) (-2.082,-0.825) (-2.051,-0.775) (-2.018,-0.725) (-1.982,-0.675) (-1.950,-0.634) (-1.900,-0.575) (-1.853,-0.525) (-1.801,-0.475) (-1.750,-0.430) (-1.700,-0.390) (-1.650,-0.353) (-1.600,-0.319) (-1.550,-0.288) (-1.500,-0.259) (-1.450,-0.233) (-1.400,-0.208) (-1.350,-0.186) (-1.300,-0.165)}
\EQ{(1.300,-0.141) (1.350,-0.158) (1.400,-0.177) (1.450,-0.197) (1.500,-0.220) (1.550,-0.243) (1.600,-0.269) (1.650,-0.297) (1.697,-0.325) (1.750,-0.359) (1.800,-0.394) (1.841,-0.425) (1.900,-0.473) (1.950,-0.517) (2.000,-0.566) (2.050,-0.619) (2.099,-0.675) (2.138,-0.725) (2.174,-0.775) (2.208,-0.825) (2.239,-0.875) (2.267,-0.925) (2.293,-0.975) (2.317,-1.025) (2.339,-1.075) (2.359,-1.125) (2.377,-1.175) (2.394,-1.225) (2.409,-1.275) (2.422,-1.325) (2.433,-1.375) (2.443,-1.425) (2.451,-1.475) (2.458,-1.525) (2.464,-1.575) (2.468,-1.625) (2.470,-1.675) (2.472,-1.725) (2.472,-1.775) (2.470,-1.825) (2.467,-1.875) (2.463,-1.925) (2.458,-1.975) (2.451,-2.025) (2.443,-2.075) (2.434,-2.125) (2.423,-2.175) (2.411,-2.225) (2.398,-2.275) (2.383,-2.325) (2.367,-2.375) (2.350,-2.423) (2.330,-2.475) (2.310,-2.525) (2.288,-2.575) (2.265,-2.625) (2.240,-2.675) (2.213,-2.725) (2.184,-2.775) (2.154,-2.825) (2.125,-2.871) (2.089,-2.925) (2.053,-2.975) (2.015,-3.025) (1.975,-3.075) (1.933,-3.125) (1.888,-3.175) (1.850,-3.215) (1.800,-3.265) (1.763,-3.300) (1.725,-3.335) (1.679,-3.375) (1.625,-3.420) (1.575,-3.459) (1.525,-3.496) (1.475,-3.531) (1.425,-3.565) (1.375,-3.597) (1.328,-3.625) (1.285,-3.650) (1.239,-3.675) (1.191,-3.700) (1.141,-3.725) (1.087,-3.750) (1.031,-3.775) (0.975,-3.798) (0.925,-3.817) (0.875,-3.835) (0.825,-3.852) (0.775,-3.868) (0.725,-3.883) (0.675,-3.897) (0.625,-3.909) (0.575,-3.921) (0.525,-3.932) (0.475,-3.941) (0.425,-3.950) (0.375,-3.958) (0.325,-3.964) (0.275,-3.970) (0.225,-3.975) (0.175,-3.978) (0.125,-3.981) (0.075,-3.983) (0.025,-3.984) (-0.050,-3.984) (-0.100,-3.982) (-0.150,-3.980) (-0.200,-3.977) (-0.250,-3.972) (-0.300,-3.967) (-0.350,-3.961) (-0.400,-3.954) (-0.450,-3.946) (-0.500,-3.937) (-0.550,-3.926) (-0.600,-3.915) (-0.650,-3.903) (-0.700,-3.890) (-0.750,-3.876) (-0.800,-3.860) (-0.850,-3.844) (-0.900,-3.826) (-0.950,-3.808) (-1.000,-3.788) (-1.050,-3.767) (-1.100,-3.744) (-1.150,-3.721) (-1.200,-3.696) (-1.250,-3.669) (-1.300,-3.641) (-1.350,-3.612) (-1.400,-3.581) (-1.448,-3.550) (-1.500,-3.514) (-1.550,-3.478) (-1.600,-3.440) (-1.650,-3.400) (-1.700,-3.357) (-1.750,-3.312) (-1.790,-3.275) (-1.840,-3.225) (-1.875,-3.189) (-1.910,-3.150) (-1.950,-3.105) (-1.995,-3.050) (-2.034,-3.000) (-2.071,-2.950) (-2.100,-2.908) (-2.139,-2.850) (-2.170,-2.800) (-2.199,-2.750) (-2.225,-2.703) (-2.250,-2.654) (-2.275,-2.603) (-2.299,-2.550) (-2.320,-2.500) (-2.340,-2.450) (-2.358,-2.400) (-2.375,-2.351) (-2.391,-2.300) (-2.405,-2.250) (-2.417,-2.200) (-2.429,-2.150) (-2.439,-2.100) (-2.447,-2.050) (-2.455,-2.000) (-2.461,-1.950) (-2.466,-1.900) (-2.469,-1.850) (-2.471,-1.800) (-2.472,-1.750) (-2.471,-1.700) (-2.469,-1.650) (-2.466,-1.600) (-2.461,-1.550) (-2.455,-1.500) (-2.447,-1.450) (-2.438,-1.400) (-2.427,-1.350) (-2.415,-1.300) (-2.401,-1.250) (-2.386,-1.200) (-2.369,-1.150) (-2.350,-1.101) (-2.329,-1.050) (-2.306,-1.000) (-2.281,-0.950) (-2.253,-0.900) (-2.225,-0.852) (-2.192,-0.800) (-2.157,-0.750) (-2.125,-0.708) (-2.078,-0.650) (-2.033,-0.600) (-1.984,-0.550) (-1.931,-0.500) (-1.875,-0.452) (-1.825,-0.413) (-1.775,-0.376) (-1.725,-0.343) (-1.675,-0.312) (-1.625,-0.283) (-1.575,-0.256) (-1.525,-0.231) (-1.475,-0.208) (-1.425,-0.187) (-1.375,-0.167) (-1.327,-0.150) (-1.300,-0.141)}
\EQ{(1.300,-0.117) (1.350,-0.131) (1.400,-0.146) (1.450,-0.163) (1.500,-0.181) (1.549,-0.200) (1.600,-0.221) (1.650,-0.244) (1.700,-0.268) (1.750,-0.293) (1.800,-0.321) (1.849,-0.350) (1.900,-0.382) (1.950,-0.416) (1.997,-0.450) (2.050,-0.492) (2.090,-0.525) (2.145,-0.575) (2.196,-0.625) (2.243,-0.675) (2.287,-0.725) (2.325,-0.772) (2.365,-0.825) (2.400,-0.875) (2.432,-0.925) (2.462,-0.975) (2.490,-1.025) (2.516,-1.075) (2.540,-1.125) (2.562,-1.175) (2.582,-1.225) (2.600,-1.275) (2.617,-1.325) (2.632,-1.375) (2.646,-1.425) (2.658,-1.475) (2.669,-1.525) (2.679,-1.575) (2.687,-1.625) (2.693,-1.675) (2.699,-1.725) (2.703,-1.775) (2.706,-1.825) (2.707,-1.875) (2.708,-1.925) (2.707,-1.975) (2.705,-2.025) (2.701,-2.075) (2.697,-2.125) (2.691,-2.175) (2.684,-2.225) (2.676,-2.275) (2.667,-2.325) (2.656,-2.375) (2.644,-2.425) (2.631,-2.475) (2.617,-2.525) (2.602,-2.575) (2.585,-2.625) (2.567,-2.675) (2.547,-2.725) (2.527,-2.775) (2.504,-2.825) (2.481,-2.875) (2.456,-2.925) (2.430,-2.975) (2.402,-3.025) (2.375,-3.070) (2.341,-3.125) (2.308,-3.175) (2.275,-3.223) (2.237,-3.275) (2.200,-3.323) (2.158,-3.375) (2.125,-3.414) (2.075,-3.470) (2.025,-3.523) (1.975,-3.574) (1.925,-3.621) (1.875,-3.667) (1.836,-3.700) (1.776,-3.750) (1.725,-3.790) (1.677,-3.825) (1.625,-3.862) (1.575,-3.896) (1.530,-3.925) (1.475,-3.959) (1.425,-3.988) (1.375,-4.016) (1.325,-4.043) (1.275,-4.068) (1.225,-4.092) (1.175,-4.115) (1.125,-4.137) (1.075,-4.158) (1.025,-4.177) (0.975,-4.196) (0.925,-4.213) (0.875,-4.229) (0.825,-4.245) (0.775,-4.259) (0.725,-4.272) (0.675,-4.285) (0.625,-4.296) (0.575,-4.307) (0.525,-4.317) (0.475,-4.325) (0.425,-4.333) (0.375,-4.340) (0.325,-4.346) (0.275,-4.351) (0.225,-4.356) (0.175,-4.359) (0.125,-4.362) (0.075,-4.363) (0.025,-4.364) (-0.050,-4.364) (-0.100,-4.363) (-0.150,-4.361) (-0.200,-4.358) (-0.250,-4.354) (-0.300,-4.349) (-0.350,-4.343) (-0.400,-4.337) (-0.450,-4.329) (-0.500,-4.321) (-0.550,-4.312) (-0.600,-4.302) (-0.650,-4.291) (-0.700,-4.279) (-0.750,-4.266) (-0.800,-4.252) (-0.850,-4.237) (-0.900,-4.221) (-0.950,-4.204) (-1.000,-4.187) (-1.050,-4.168) (-1.100,-4.147) (-1.150,-4.126) (-1.200,-4.104) (-1.250,-4.080) (-1.300,-4.056) (-1.350,-4.030) (-1.400,-4.002) (-1.448,-3.975) (-1.500,-3.944) (-1.550,-3.912) (-1.600,-3.879) (-1.642,-3.850) (-1.700,-3.808) (-1.744,-3.775) (-1.800,-3.730) (-1.850,-3.688) (-1.894,-3.650) (-1.948,-3.600) (-1.999,-3.550) (-2.047,-3.500) (-2.093,-3.450) (-2.137,-3.400) (-2.175,-3.354) (-2.218,-3.300) (-2.250,-3.257) (-2.291,-3.200) (-2.325,-3.150) (-2.357,-3.100) (-2.387,-3.050) (-2.416,-3.000) (-2.443,-2.950) (-2.469,-2.900) (-2.493,-2.850) (-2.516,-2.800) (-2.537,-2.750) (-2.557,-2.700) (-2.575,-2.653) (-2.593,-2.600) (-2.610,-2.550) (-2.624,-2.500) (-2.638,-2.450) (-2.650,-2.400) (-2.662,-2.350) (-2.672,-2.300) (-2.680,-2.250) (-2.688,-2.200) (-2.694,-2.150) (-2.699,-2.100) (-2.703,-2.050) (-2.706,-2.000) (-2.707,-1.950) (-2.708,-1.900) (-2.707,-1.850) (-2.704,-1.800) (-2.701,-1.750) (-2.696,-1.700) (-2.690,-1.650) (-2.683,-1.600) (-2.674,-1.550) (-2.664,-1.500) (-2.652,-1.450) (-2.639,-1.400) (-2.625,-1.351) (-2.609,-1.300) (-2.591,-1.250) (-2.572,-1.200) (-2.551,-1.150) (-2.528,-1.100) (-2.503,-1.050) (-2.476,-1.000) (-2.450,-0.954) (-2.416,-0.900) (-2.383,-0.850) (-2.350,-0.805) (-2.308,-0.750) (-2.275,-0.711) (-2.225,-0.655) (-2.175,-0.604) (-2.125,-0.557) (-2.075,-0.513) (-2.029,-0.475) (-1.975,-0.434) (-1.927,-0.400) (-1.875,-0.366) (-1.825,-0.335) (-1.775,-0.307) (-1.725,-0.280) (-1.675,-0.255) (-1.625,-0.232) (-1.575,-0.211) (-1.525,-0.191) (-1.475,-0.172) (-1.425,-0.155) (-1.375,-0.139) (-1.325,-0.124) (-1.300,-0.117)}
\EQ{(1.300,-0.093) (1.350,-0.104) (1.400,-0.116) (1.450,-0.130) (1.500,-0.144) (1.550,-0.159) (1.600,-0.175) (1.650,-0.193) (1.700,-0.211) (1.750,-0.231) (1.800,-0.252) (1.850,-0.275) (1.900,-0.299) (1.950,-0.324) (1.997,-0.350) (2.041,-0.375) (2.100,-0.412) (2.150,-0.444) (2.194,-0.475) (2.250,-0.517) (2.292,-0.550) (2.350,-0.599) (2.400,-0.645) (2.450,-0.693) (2.500,-0.746) (2.548,-0.800) (2.589,-0.850) (2.625,-0.896) (2.664,-0.950) (2.697,-1.000) (2.725,-1.044) (2.759,-1.100) (2.786,-1.150) (2.812,-1.200) (2.836,-1.250) (2.858,-1.300) (2.879,-1.350) (2.898,-1.400) (2.916,-1.450) (2.932,-1.500) (2.947,-1.550) (2.960,-1.600) (2.972,-1.650) (2.983,-1.700) (2.993,-1.750) (3.001,-1.800) (3.008,-1.850) (3.014,-1.900) (3.019,-1.950) (3.023,-2.000) (3.025,-2.050) (3.027,-2.100) (3.027,-2.150) (3.026,-2.200) (3.025,-2.250) (3.022,-2.300) (3.018,-2.350) (3.013,-2.400) (3.006,-2.450) (2.999,-2.500) (2.991,-2.550) (2.981,-2.600) (2.971,-2.650) (2.959,-2.700) (2.947,-2.750) (2.933,-2.800) (2.918,-2.850) (2.902,-2.900) (2.885,-2.950) (2.866,-3.000) (2.847,-3.050) (2.826,-3.100) (2.804,-3.150) (2.781,-3.200) (2.756,-3.250) (2.730,-3.300) (2.703,-3.350) (2.675,-3.399) (2.645,-3.450) (2.613,-3.500) (2.580,-3.550) (2.550,-3.594) (2.509,-3.650) (2.475,-3.696) (2.432,-3.750) (2.400,-3.789) (2.350,-3.846) (2.301,-3.900) (2.253,-3.950) (2.203,-4.000) (2.151,-4.050) (2.100,-4.096) (2.050,-4.139) (2.007,-4.175) (1.950,-4.220) (1.900,-4.258) (1.850,-4.295) (1.807,-4.325) (1.750,-4.363) (1.700,-4.395) (1.651,-4.425) (1.600,-4.455) (1.550,-4.483) (1.500,-4.510) (1.450,-4.536) (1.400,-4.561) (1.350,-4.585) (1.300,-4.607) (1.250,-4.629) (1.200,-4.649) (1.150,-4.669) (1.100,-4.687) (1.050,-4.705) (1.000,-4.722) (0.950,-4.737) (0.900,-4.752) (0.850,-4.766) (0.800,-4.779) (0.750,-4.792) (0.700,-4.803) (0.650,-4.814) (0.600,-4.824) (0.550,-4.833) (0.500,-4.841) (0.450,-4.848) (0.400,-4.855) (0.350,-4.861) (0.300,-4.866) (0.250,-4.870) (0.200,-4.873) (0.150,-4.876) (0.100,-4.878) (0.050,-4.879) (0.000,-4.880) (-0.050,-4.879) (-0.100,-4.878) (-0.150,-4.876) (-0.200,-4.873) (-0.250,-4.870) (-0.300,-4.866) (-0.350,-4.861) (-0.400,-4.855) (-0.450,-4.848) (-0.500,-4.841) (-0.550,-4.833) (-0.600,-4.824) (-0.650,-4.814) (-0.700,-4.803) (-0.750,-4.792) (-0.800,-4.779) (-0.850,-4.766) (-0.900,-4.752) (-0.950,-4.737) (-1.000,-4.722) (-1.050,-4.705) (-1.100,-4.687) (-1.150,-4.669) (-1.198,-4.650) (-1.250,-4.629) (-1.300,-4.607) (-1.350,-4.585) (-1.400,-4.561) (-1.450,-4.536) (-1.500,-4.510) (-1.550,-4.483) (-1.600,-4.455) (-1.650,-4.426) (-1.700,-4.395) (-1.750,-4.363) (-1.800,-4.330) (-1.843,-4.300) (-1.900,-4.258) (-1.944,-4.225) (-2.000,-4.181) (-2.050,-4.139) (-2.095,-4.100) (-2.150,-4.051) (-2.200,-4.003) (-2.250,-3.954) (-2.300,-3.901) (-2.347,-3.850) (-2.390,-3.800) (-2.425,-3.758) (-2.472,-3.700) (-2.509,-3.650) (-2.546,-3.600) (-2.575,-3.558) (-2.613,-3.500) (-2.645,-3.450) (-2.675,-3.400) (-2.700,-3.356) (-2.725,-3.310) (-2.750,-3.262) (-2.775,-3.212) (-2.800,-3.159) (-2.825,-3.102) (-2.847,-3.050) (-2.866,-3.000) (-2.885,-2.950) (-2.902,-2.900) (-2.918,-2.850) (-2.933,-2.800) (-2.947,-2.750) (-2.959,-2.700) (-2.971,-2.650) (-2.981,-2.600) (-2.991,-2.550) (-2.999,-2.500) (-3.006,-2.450) (-3.013,-2.400) (-3.018,-2.350) (-3.022,-2.300) (-3.025,-2.250) (-3.026,-2.200) (-3.027,-2.150) (-3.027,-2.100) (-3.025,-2.050) (-3.023,-2.000) (-3.019,-1.950) (-3.014,-1.900) (-3.008,-1.850) (-3.001,-1.800) (-2.993,-1.750) (-2.983,-1.700) (-2.972,-1.650) (-2.960,-1.600) (-2.947,-1.550) (-2.932,-1.500) (-2.916,-1.450) (-2.898,-1.400) (-2.879,-1.350) (-2.858,-1.300) (-2.836,-1.250) (-2.812,-1.200) (-2.786,-1.150) (-2.759,-1.100) (-2.729,-1.050) (-2.700,-1.004) (-2.664,-0.950) (-2.628,-0.900) (-2.589,-0.850) (-2.550,-0.802) (-2.504,-0.750) (-2.457,-0.700) (-2.406,-0.650) (-2.351,-0.600) (-2.300,-0.557) (-2.261,-0.525) (-2.200,-0.479) (-2.158,-0.450) (-2.100,-0.412) (-2.050,-0.381) (-2.000,-0.352) (-1.951,-0.325) (-1.902,-0.300) (-1.850,-0.275) (-1.800,-0.252) (-1.750,-0.231) (-1.700,-0.211) (-1.650,-0.193) (-1.600,-0.175) (-1.550,-0.159) (-1.500,-0.144) (-1.450,-0.130) (-1.400,-0.116) (-1.350,-0.104) (-1.300,-0.093)}
\EQ{(1.300,-0.070) (1.350,-0.078) (1.400,-0.087) (1.450,-0.097) (1.500,-0.107) (1.550,-0.118) (1.600,-0.130) (1.650,-0.143) (1.700,-0.157) (1.750,-0.171) (1.800,-0.187) (1.850,-0.203) (1.900,-0.220) (1.950,-0.239) (2.000,-0.258) (2.050,-0.279) (2.098,-0.300) (2.150,-0.324) (2.200,-0.348) (2.250,-0.374) (2.298,-0.400) (2.342,-0.425) (2.400,-0.460) (2.450,-0.492) (2.500,-0.525) (2.550,-0.561) (2.600,-0.598) (2.650,-0.638) (2.694,-0.675) (2.750,-0.724) (2.800,-0.772) (2.850,-0.822) (2.899,-0.875) (2.943,-0.925) (2.985,-0.975) (3.024,-1.025) (3.060,-1.075) (3.095,-1.125) (3.125,-1.171) (3.159,-1.225) (3.188,-1.275) (3.216,-1.325) (3.241,-1.375) (3.266,-1.425) (3.289,-1.475) (3.310,-1.525) (3.330,-1.575) (3.349,-1.625) (3.367,-1.675) (3.383,-1.725) (3.398,-1.775) (3.412,-1.825) (3.425,-1.875) (3.437,-1.925) (3.447,-1.975) (3.456,-2.025) (3.465,-2.075) (3.472,-2.125) (3.479,-2.175) (3.484,-2.225) (3.488,-2.275) (3.491,-2.325) (3.494,-2.375) (3.495,-2.425) (3.496,-2.475) (3.495,-2.525) (3.494,-2.575) (3.491,-2.625) (3.488,-2.675) (3.483,-2.725) (3.478,-2.775) (3.472,-2.825) (3.465,-2.875) (3.457,-2.925) (3.448,-2.975) (3.438,-3.025) (3.427,-3.075) (3.415,-3.125) (3.403,-3.175) (3.389,-3.225) (3.374,-3.275) (3.359,-3.325) (3.342,-3.375) (3.325,-3.423) (3.306,-3.475) (3.286,-3.525) (3.265,-3.575) (3.243,-3.625) (3.220,-3.675) (3.196,-3.725) (3.171,-3.775) (3.145,-3.825) (3.118,-3.875) (3.089,-3.925) (3.059,-3.975) (3.028,-4.025) (3.000,-4.068) (2.962,-4.125) (2.927,-4.175) (2.890,-4.225) (2.852,-4.275) (2.812,-4.325) (2.775,-4.370) (2.728,-4.425) (2.683,-4.475) (2.637,-4.525) (2.600,-4.563) (2.563,-4.600) (2.525,-4.637) (2.485,-4.675) (2.430,-4.725) (2.375,-4.773) (2.325,-4.815) (2.281,-4.850) (2.225,-4.894) (2.183,-4.925) (2.125,-4.967) (2.078,-5.000) (2.025,-5.035) (1.975,-5.067) (1.925,-5.098) (1.881,-5.125) (1.825,-5.157) (1.775,-5.185) (1.725,-5.212) (1.675,-5.238) (1.625,-5.262) (1.575,-5.286) (1.525,-5.309) (1.475,-5.331) (1.425,-5.352) (1.375,-5.373) (1.325,-5.392) (1.275,-5.411) (1.225,-5.428) (1.175,-5.445) (1.125,-5.462) (1.075,-5.477) (1.025,-5.492) (0.975,-5.505) (0.925,-5.519) (0.875,-5.531) (0.825,-5.543) (0.775,-5.553) (0.725,-5.564) (0.675,-5.573) (0.625,-5.582) (0.575,-5.590) (0.525,-5.597) (0.475,-5.604) (0.425,-5.610) (0.375,-5.616) (0.325,-5.620) (0.275,-5.624) (0.225,-5.628) (0.175,-5.630) (0.125,-5.632) (0.075,-5.634) (0.025,-5.634) (-0.050,-5.634) (-0.100,-5.633) (-0.150,-5.631) (-0.200,-5.629) (-0.250,-5.626) (-0.300,-5.622) (-0.350,-5.618) (-0.400,-5.613) (-0.450,-5.607) (-0.500,-5.601) (-0.550,-5.594) (-0.600,-5.586) (-0.650,-5.578) (-0.700,-5.569) (-0.750,-5.559) (-0.800,-5.548) (-0.850,-5.537) (-0.899,-5.525) (-0.950,-5.512) (-1.000,-5.499) (-1.050,-5.484) (-1.100,-5.469) (-1.150,-5.453) (-1.200,-5.437) (-1.250,-5.420) (-1.300,-5.401) (-1.350,-5.382) (-1.400,-5.363) (-1.450,-5.342) (-1.500,-5.320) (-1.550,-5.298) (-1.599,-5.275) (-1.650,-5.250) (-1.700,-5.225) (-1.747,-5.200) (-1.793,-5.175) (-1.838,-5.150) (-1.900,-5.114) (-1.950,-5.083) (-2.000,-5.051) (-2.050,-5.019) (-2.100,-4.985) (-2.149,-4.950) (-2.200,-4.913) (-2.249,-4.875) (-2.300,-4.835) (-2.343,-4.800) (-2.400,-4.751) (-2.450,-4.707) (-2.500,-4.661) (-2.538,-4.625) (-2.575,-4.588) (-2.613,-4.550) (-2.650,-4.511) (-2.700,-4.457) (-2.750,-4.400) (-2.792,-4.350) (-2.825,-4.309) (-2.871,-4.250) (-2.908,-4.200) (-2.944,-4.150) (-2.975,-4.106) (-3.012,-4.050) (-3.044,-4.000) (-3.074,-3.950) (-3.100,-3.906) (-3.125,-3.862) (-3.150,-3.816) (-3.175,-3.768) (-3.200,-3.718) (-3.225,-3.665) (-3.250,-3.610) (-3.275,-3.551) (-3.296,-3.500) (-3.315,-3.450) (-3.333,-3.400) (-3.350,-3.351) (-3.366,-3.300) (-3.382,-3.250) (-3.396,-3.200) (-3.409,-3.150) (-3.421,-3.100) (-3.433,-3.050) (-3.443,-3.000) (-3.453,-2.950) (-3.461,-2.900) (-3.469,-2.850) (-3.475,-2.800) (-3.481,-2.750) (-3.486,-2.700) (-3.490,-2.650) (-3.492,-2.600) (-3.494,-2.550) (-3.495,-2.500) (-3.496,-2.450) (-3.495,-2.400) (-3.493,-2.350) (-3.490,-2.300) (-3.486,-2.250) (-3.481,-2.200) (-3.475,-2.150) (-3.469,-2.100) (-3.461,-2.050) (-3.452,-2.000) (-3.442,-1.950) (-3.431,-1.900) (-3.419,-1.850) (-3.405,-1.800) (-3.391,-1.750) (-3.375,-1.700) (-3.358,-1.650) (-3.340,-1.600) (-3.321,-1.550) (-3.300,-1.501) (-3.277,-1.450) (-3.254,-1.400) (-3.229,-1.350) (-3.202,-1.300) (-3.175,-1.252) (-3.144,-1.200) (-3.112,-1.150) (-3.078,-1.100) (-3.042,-1.050) (-3.004,-1.000) (-2.964,-0.950) (-2.925,-0.904) (-2.876,-0.850) (-2.828,-0.800) (-2.777,-0.750) (-2.725,-0.702) (-2.675,-0.659) (-2.634,-0.625) (-2.575,-0.579) (-2.525,-0.543) (-2.475,-0.508) (-2.425,-0.475) (-2.375,-0.444) (-2.325,-0.415) (-2.275,-0.387) (-2.225,-0.361) (-2.175,-0.336) (-2.125,-0.312) (-2.075,-0.290) (-2.025,-0.269) (-1.975,-0.248) (-1.925,-0.230) (-1.875,-0.212) (-1.825,-0.195) (-1.775,-0.179) (-1.725,-0.164) (-1.676,-0.150) (-1.625,-0.137) (-1.575,-0.124) (-1.525,-0.113) (-1.475,-0.102) (-1.425,-0.092) (-1.375,-0.082) (-1.325,-0.074) (-1.300,-0.070)}
\EQ{(1.300,-0.046) (1.350,-0.052) (1.400,-0.058) (1.450,-0.064) (1.500,-0.071) (1.550,-0.079) (1.600,-0.086) (1.650,-0.095) (1.700,-0.104) (1.750,-0.113) (1.800,-0.123) (1.850,-0.134) (1.900,-0.145) (1.950,-0.157) (2.000,-0.170) (2.050,-0.183) (2.100,-0.197) (2.150,-0.212) (2.200,-0.227) (2.250,-0.243) (2.300,-0.260) (2.350,-0.278) (2.400,-0.297) (2.450,-0.317) (2.500,-0.337) (2.550,-0.359) (2.600,-0.381) (2.650,-0.405) (2.700,-0.429) (2.750,-0.455) (2.800,-0.482) (2.850,-0.510) (2.900,-0.539) (2.950,-0.570) (2.998,-0.600) (3.050,-0.635) (3.100,-0.670) (3.141,-0.700) (3.200,-0.745) (3.250,-0.785) (3.298,-0.825) (3.350,-0.871) (3.400,-0.917) (3.450,-0.966) (3.500,-1.017) (3.550,-1.071) (3.598,-1.125) (3.640,-1.175) (3.675,-1.219) (3.718,-1.275) (3.750,-1.319) (3.789,-1.375) (3.822,-1.425) (3.850,-1.470) (3.883,-1.525) (3.912,-1.575) (3.939,-1.625) (3.965,-1.675) (3.990,-1.725) (4.014,-1.775) (4.036,-1.825) (4.057,-1.875) (4.077,-1.925) (4.096,-1.975) (4.114,-2.025) (4.131,-2.075) (4.147,-2.125) (4.162,-2.175) (4.176,-2.225) (4.189,-2.275) (4.201,-2.325) (4.213,-2.375) (4.223,-2.425) (4.232,-2.475) (4.241,-2.525) (4.249,-2.575) (4.256,-2.625) (4.262,-2.675) (4.267,-2.725) (4.271,-2.775) (4.275,-2.825) (4.278,-2.875) (4.280,-2.925) (4.281,-2.975) (4.281,-3.025) (4.281,-3.075) (4.280,-3.125) (4.278,-3.175) (4.275,-3.225) (4.272,-3.275) (4.268,-3.325) (4.263,-3.375) (4.257,-3.425) (4.250,-3.475) (4.243,-3.525) (4.235,-3.575) (4.226,-3.625) (4.217,-3.675) (4.206,-3.725) (4.195,-3.775) (4.184,-3.825) (4.171,-3.875) (4.157,-3.925) (4.143,-3.975) (4.128,-4.025) (4.112,-4.075) (4.096,-4.125) (4.078,-4.175) (4.060,-4.225) (4.041,-4.275) (4.021,-4.325) (4.000,-4.375) (3.979,-4.425) (3.956,-4.475) (3.933,-4.525) (3.908,-4.575) (3.883,-4.625) (3.857,-4.675) (3.830,-4.725) (3.802,-4.775) (3.775,-4.821) (3.743,-4.875) (3.711,-4.925) (3.679,-4.975) (3.650,-5.019) (3.611,-5.075) (3.576,-5.125) (3.539,-5.175) (3.501,-5.225) (3.462,-5.275) (3.425,-5.321) (3.380,-5.375) (3.336,-5.425) (3.300,-5.466) (3.250,-5.520) (3.200,-5.573) (3.150,-5.623) (3.100,-5.672) (3.050,-5.720) (3.000,-5.766) (2.961,-5.800) (2.904,-5.850) (2.850,-5.895) (2.800,-5.935) (2.750,-5.974) (2.700,-6.013) (2.650,-6.049) (2.600,-6.085) (2.550,-6.120) (2.505,-6.150) (2.450,-6.186) (2.400,-6.218) (2.350,-6.249) (2.306,-6.275) (2.250,-6.308) (2.200,-6.336) (2.150,-6.363) (2.100,-6.390) (2.050,-6.415) (2.000,-6.440) (1.950,-6.464) (1.900,-6.488) (1.850,-6.510) (1.800,-6.532) (1.750,-6.553) (1.700,-6.574) (1.650,-6.593) (1.600,-6.612) (1.550,-6.631) (1.500,-6.648) (1.450,-6.665) (1.400,-6.682) (1.350,-6.698) (1.300,-6.713) (1.250,-6.727) (1.200,-6.741) (1.150,-6.754) (1.100,-6.767) (1.050,-6.779) (1.000,-6.791) (0.950,-6.801) (0.900,-6.812) (0.850,-6.821) (0.800,-6.831) (0.750,-6.839) (0.700,-6.847) (0.650,-6.854) (0.600,-6.861) (0.550,-6.868) (0.500,-6.873) (0.450,-6.879) (0.400,-6.883) (0.350,-6.887) (0.300,-6.891) (0.250,-6.894) (0.200,-6.896) (0.150,-6.898) (0.100,-6.900) (0.050,-6.900) (0.000,-6.901) (-0.050,-6.900) (-0.100,-6.900) (-0.150,-6.898) (-0.200,-6.896) (-0.250,-6.894) (-0.300,-6.891) (-0.350,-6.887) (-0.400,-6.883) (-0.450,-6.879) (-0.500,-6.873) (-0.550,-6.868) (-0.600,-6.861) (-0.650,-6.854) (-0.700,-6.847) (-0.750,-6.839) (-0.800,-6.831) (-0.850,-6.821) (-0.900,-6.812) (-0.950,-6.801) (-1.000,-6.791) (-1.050,-6.779) (-1.100,-6.767) (-1.150,-6.754) (-1.200,-6.741) (-1.250,-6.727) (-1.300,-6.713) (-1.350,-6.698) (-1.400,-6.682) (-1.450,-6.665) (-1.500,-6.648) (-1.550,-6.631) (-1.600,-6.612) (-1.650,-6.593) (-1.697,-6.575) (-1.750,-6.553) (-1.800,-6.532) (-1.850,-6.510) (-1.900,-6.488) (-1.950,-6.464) (-2.000,-6.440) (-2.050,-6.415) (-2.100,-6.390) (-2.150,-6.363) (-2.200,-6.336) (-2.250,-6.308) (-2.300,-6.279) (-2.348,-6.250) (-2.400,-6.218) (-2.450,-6.186) (-2.500,-6.154) (-2.543,-6.125) (-2.600,-6.085) (-2.649,-6.050) (-2.700,-6.013) (-2.749,-5.975) (-2.800,-5.935) (-2.844,-5.900) (-2.900,-5.853) (-2.950,-5.810) (-2.990,-5.775) (-3.044,-5.725) (-3.097,-5.675) (-3.148,-5.625) (-3.198,-5.575) (-3.246,-5.525) (-3.292,-5.475) (-3.336,-5.425) (-3.375,-5.380) (-3.421,-5.325) (-3.462,-5.275) (-3.500,-5.227) (-3.539,-5.175) (-3.575,-5.126) (-3.611,-5.075) (-3.646,-5.025) (-3.675,-4.981) (-3.711,-4.925) (-3.743,-4.875) (-3.773,-4.825) (-3.800,-4.778) (-3.825,-4.734) (-3.850,-4.688) (-3.875,-4.641) (-3.900,-4.592) (-3.925,-4.541) (-3.950,-4.488) (-3.975,-4.433) (-4.000,-4.376) (-4.021,-4.325) (-4.041,-4.275) (-4.060,-4.225) (-4.078,-4.175) (-4.096,-4.125) (-4.112,-4.075) (-4.128,-4.025) (-4.143,-3.975) (-4.157,-3.925) (-4.171,-3.875) (-4.184,-3.825) (-4.195,-3.775) (-4.206,-3.725) (-4.217,-3.675) (-4.226,-3.625) (-4.235,-3.575) (-4.243,-3.525) (-4.250,-3.475) (-4.257,-3.425) (-4.263,-3.375) (-4.268,-3.325) (-4.272,-3.275) (-4.275,-3.225) (-4.278,-3.175) (-4.280,-3.125) (-4.281,-3.075) (-4.281,-3.025) (-4.281,-2.975) (-4.280,-2.925) (-4.278,-2.875) (-4.275,-2.825) (-4.271,-2.775) (-4.267,-2.725) (-4.262,-2.675) (-4.256,-2.625) (-4.249,-2.575) (-4.241,-2.525) (-4.232,-2.475) (-4.223,-2.425) (-4.213,-2.375) (-4.201,-2.325) (-4.189,-2.275) (-4.176,-2.225) (-4.162,-2.175) (-4.147,-2.125) (-4.131,-2.075) (-4.114,-2.025) (-4.096,-1.975) (-4.077,-1.925) (-4.057,-1.875) (-4.036,-1.825) (-4.014,-1.775) (-3.990,-1.725) (-3.965,-1.675) (-3.939,-1.625) (-3.912,-1.575) (-3.883,-1.525) (-3.853,-1.475) (-3.825,-1.430) (-3.789,-1.375) (-3.754,-1.325) (-3.718,-1.275) (-3.680,-1.225) (-3.640,-1.175) (-3.600,-1.127) (-3.554,-1.075) (-3.508,-1.025) (-3.459,-0.975) (-3.408,-0.925) (-3.355,-0.875) (-3.300,-0.827) (-3.250,-0.785) (-3.207,-0.750) (-3.150,-0.707) (-3.107,-0.675) (-3.050,-0.635) (-3.000,-0.602) (-2.950,-0.570) (-2.900,-0.539) (-2.850,-0.510) (-2.800,-0.482) (-2.750,-0.455) (-2.700,-0.429) (-2.650,-0.405) (-2.600,-0.381) (-2.550,-0.359) (-2.500,-0.337) (-2.450,-0.317) (-2.400,-0.297) (-2.350,-0.278) (-2.300,-0.260) (-2.250,-0.243) (-2.200,-0.227) (-2.150,-0.212) (-2.100,-0.197) (-2.050,-0.183) (-2.000,-0.170) (-1.950,-0.157) (-1.900,-0.145) (-1.850,-0.134) (-1.800,-0.123) (-1.750,-0.113) (-1.700,-0.104) (-1.650,-0.095) (-1.600,-0.086) (-1.550,-0.079) (-1.500,-0.071) (-1.450,-0.064) (-1.400,-0.058) (-1.350,-0.052) (-1.300,-0.046)}
\EQ{(1.300,-0.023) (1.350,-0.026) (1.400,-0.029) (1.450,-0.032) (1.500,-0.035) (1.550,-0.039) (1.600,-0.043) (1.650,-0.047) (1.700,-0.052) (1.750,-0.056) (1.800,-0.061) (1.850,-0.067) (1.900,-0.072) (1.950,-0.078) (2.000,-0.084) (2.050,-0.091) (2.100,-0.098) (2.150,-0.105) (2.200,-0.112) (2.250,-0.120) (2.300,-0.128) (2.350,-0.137) (2.400,-0.146) (2.450,-0.155) (2.500,-0.165) (2.550,-0.175) (2.600,-0.186) (2.650,-0.197) (2.700,-0.209) (2.750,-0.220) (2.800,-0.233) (2.850,-0.246) (2.900,-0.259) (2.950,-0.273) (3.000,-0.287) (3.050,-0.302) (3.100,-0.318) (3.150,-0.334) (3.199,-0.350) (3.250,-0.367) (3.300,-0.385) (3.350,-0.403) (3.400,-0.422) (3.450,-0.442) (3.500,-0.462) (3.550,-0.483) (3.600,-0.504) (3.646,-0.525) (3.700,-0.550) (3.750,-0.573) (3.800,-0.598) (3.850,-0.623) (3.900,-0.649) (3.949,-0.675) (3.994,-0.700) (4.038,-0.725) (4.100,-0.761) (4.150,-0.792) (4.200,-0.823) (4.250,-0.856) (4.300,-0.889) (4.350,-0.923) (4.400,-0.959) (4.450,-0.996) (4.500,-1.034) (4.550,-1.073) (4.600,-1.113) (4.644,-1.150) (4.700,-1.198) (4.750,-1.243) (4.800,-1.289) (4.838,-1.325) (4.875,-1.362) (4.913,-1.400) (4.950,-1.438) (4.985,-1.475) (5.025,-1.519) (5.074,-1.575) (5.117,-1.625) (5.150,-1.665) (5.198,-1.725) (5.237,-1.775) (5.274,-1.825) (5.310,-1.875) (5.345,-1.925) (5.375,-1.970) (5.411,-2.025) (5.443,-2.075) (5.473,-2.125) (5.500,-2.171) (5.525,-2.214) (5.550,-2.259) (5.575,-2.306) (5.598,-2.350) (5.623,-2.400) (5.648,-2.450) (5.671,-2.500) (5.694,-2.550) (5.716,-2.600) (5.737,-2.650) (5.757,-2.700) (5.777,-2.750) (5.796,-2.800) (5.814,-2.850) (5.831,-2.900) (5.848,-2.950) (5.863,-3.000) (5.879,-3.050) (5.893,-3.100) (5.907,-3.150) (5.920,-3.200) (5.933,-3.250) (5.944,-3.300) (5.956,-3.350) (5.966,-3.400) (5.976,-3.450) (5.985,-3.500) (5.994,-3.550) (6.002,-3.600) (6.010,-3.650) (6.017,-3.700) (6.023,-3.750) (6.029,-3.800) (6.034,-3.850) (6.038,-3.900) (6.042,-3.950) (6.046,-4.000) (6.049,-4.050) (6.051,-4.100) (6.053,-4.150) (6.054,-4.200) (6.054,-4.250) (6.054,-4.300) (6.054,-4.350) (6.053,-4.400) (6.051,-4.450) (6.049,-4.500) (6.047,-4.550) (6.043,-4.600) (6.040,-4.650) (6.035,-4.700) (6.031,-4.750) (6.025,-4.800) (6.019,-4.850) (6.013,-4.900) (6.006,-4.950) (5.998,-5.000) (5.990,-5.050) (5.982,-5.100) (5.973,-5.150) (5.963,-5.200) (5.953,-5.250) (5.942,-5.300) (5.931,-5.350) (5.919,-5.400) (5.906,-5.450) (5.893,-5.500) (5.880,-5.550) (5.866,-5.600) (5.851,-5.650) (5.836,-5.700) (5.820,-5.750) (5.804,-5.800) (5.787,-5.850) (5.769,-5.900) (5.751,-5.950) (5.732,-6.000) (5.713,-6.050) (5.693,-6.100) (5.673,-6.150) (5.652,-6.200) (5.630,-6.250) (5.608,-6.300) (5.585,-6.350) (5.561,-6.400) (5.537,-6.450) (5.512,-6.500) (5.487,-6.550) (5.461,-6.600) (5.434,-6.650) (5.406,-6.700) (5.378,-6.750) (5.350,-6.799) (5.320,-6.850) (5.289,-6.900) (5.258,-6.950) (5.226,-7.000) (5.194,-7.050) (5.160,-7.100) (5.126,-7.150) (5.091,-7.200) (5.056,-7.250) (5.025,-7.292) (4.981,-7.350) (4.950,-7.391) (4.904,-7.450) (4.864,-7.500) (4.825,-7.547) (4.781,-7.600) (4.738,-7.650) (4.700,-7.693) (4.650,-7.749) (4.603,-7.800) (4.555,-7.850) (4.507,-7.900) (4.457,-7.950) (4.407,-8.000) (4.355,-8.050) (4.302,-8.100) (4.250,-8.147) (4.200,-8.192) (4.162,-8.225) (4.104,-8.275) (4.050,-8.320) (4.000,-8.361) (3.952,-8.400) (3.900,-8.441) (3.855,-8.475) (3.800,-8.517) (3.755,-8.550) (3.700,-8.589) (3.650,-8.625) (3.600,-8.659) (3.550,-8.693) (3.501,-8.725) (3.450,-8.758) (3.400,-8.790) (3.350,-8.821) (3.302,-8.850) (3.250,-8.881) (3.200,-8.910) (3.150,-8.939) (3.100,-8.967) (3.050,-8.994) (3.000,-9.021) (2.950,-9.047) (2.900,-9.072) (2.850,-9.097) (2.800,-9.122) (2.750,-9.146) (2.700,-9.169) (2.650,-9.192) (2.600,-9.214) (2.550,-9.236) (2.500,-9.257) (2.450,-9.278) (2.400,-9.298) (2.350,-9.318) (2.300,-9.337) (2.250,-9.356) (2.200,-9.374) (2.150,-9.392) (2.100,-9.410) (2.050,-9.427) (2.000,-9.443) (1.950,-9.459) (1.900,-9.475) (1.850,-9.490) (1.800,-9.504) (1.750,-9.519) (1.700,-9.533) (1.650,-9.546) (1.600,-9.559) (1.550,-9.571) (1.500,-9.583) (1.450,-9.595) (1.400,-9.606) (1.350,-9.617) (1.300,-9.628) (1.250,-9.638) (1.200,-9.647) (1.150,-9.656) (1.100,-9.665) (1.050,-9.674) (1.000,-9.682) (0.950,-9.689) (0.900,-9.696) (0.850,-9.703) (0.800,-9.710) (0.750,-9.716) (0.700,-9.721) (0.650,-9.726) (0.600,-9.731) (0.550,-9.736) (0.500,-9.740) (0.450,-9.743) (0.400,-9.747) (0.350,-9.750) (0.300,-9.752) (0.250,-9.754) (0.200,-9.756) (0.150,-9.757) (0.100,-9.758) (0.050,-9.759) (0.000,-9.759) (-0.050,-9.759) (-0.100,-9.758) (-0.150,-9.757) (-0.200,-9.756) (-0.250,-9.754) (-0.300,-9.752) (-0.350,-9.750) (-0.400,-9.747) (-0.450,-9.743) (-0.500,-9.740) (-0.550,-9.736) (-0.600,-9.731) (-0.650,-9.726) (-0.700,-9.721) (-0.750,-9.716) (-0.800,-9.710) (-0.850,-9.703) (-0.900,-9.696) (-0.950,-9.689) (-1.000,-9.682) (-1.050,-9.674) (-1.100,-9.665) (-1.150,-9.656) (-1.200,-9.647) (-1.250,-9.638) (-1.300,-9.628) (-1.350,-9.617) (-1.400,-9.606) (-1.450,-9.595) (-1.500,-9.583) (-1.550,-9.571) (-1.600,-9.559) (-1.650,-9.546) (-1.700,-9.533) (-1.750,-9.519) (-1.800,-9.504) (-1.850,-9.490) (-1.899,-9.475) (-1.950,-9.459) (-2.000,-9.443) (-2.050,-9.427) (-2.100,-9.410) (-2.150,-9.392) (-2.198,-9.375) (-2.250,-9.356) (-2.300,-9.337) (-2.350,-9.318) (-2.400,-9.298) (-2.450,-9.278) (-2.500,-9.257) (-2.550,-9.236) (-2.600,-9.214) (-2.650,-9.192) (-2.700,-9.169) (-2.750,-9.146) (-2.800,-9.122) (-2.850,-9.097) (-2.900,-9.072) (-2.950,-9.047) (-3.000,-9.021) (-3.050,-8.994) (-3.100,-8.967) (-3.150,-8.939) (-3.200,-8.910) (-3.250,-8.881) (-3.300,-8.851) (-3.344,-8.825) (-3.400,-8.790) (-3.450,-8.758) (-3.500,-8.726) (-3.550,-8.693) (-3.600,-8.659) (-3.650,-8.625) (-3.700,-8.589) (-3.750,-8.553) (-3.800,-8.517) (-3.850,-8.479) (-3.900,-8.441) (-3.950,-8.401) (-4.000,-8.361) (-4.045,-8.325) (-4.100,-8.279) (-4.150,-8.236) (-4.191,-8.200) (-4.247,-8.150) (-4.300,-8.101) (-4.350,-8.055) (-4.400,-8.007) (-4.450,-7.957) (-4.500,-7.907) (-4.550,-7.856) (-4.600,-7.803) (-4.649,-7.750) (-4.694,-7.700) (-4.738,-7.650) (-4.775,-7.607) (-4.823,-7.550) (-4.864,-7.500) (-4.900,-7.455) (-4.943,-7.400) (-4.975,-7.359) (-5.019,-7.300) (-5.050,-7.258) (-5.091,-7.200) (-5.125,-7.152) (-5.160,-7.100) (-5.194,-7.050) (-5.225,-7.002) (-5.258,-6.950) (-5.289,-6.900) (-5.320,-6.850) (-5.349,-6.800) (-5.375,-6.755) (-5.400,-6.711) (-5.425,-6.666) (-5.450,-6.620) (-5.474,-6.575) (-5.500,-6.525) (-5.525,-6.475) (-5.549,-6.425) (-5.573,-6.375) (-5.596,-6.325) (-5.619,-6.275) (-5.641,-6.225) (-5.662,-6.175) (-5.683,-6.125) (-5.703,-6.075) (-5.723,-6.025) (-5.742,-5.975) (-5.760,-5.925) (-5.778,-5.875) (-5.795,-5.825) (-5.812,-5.775) (-5.828,-5.725) (-5.843,-5.675) (-5.858,-5.625) (-5.873,-5.575) (-5.887,-5.525) (-5.900,-5.475) (-5.913,-5.425) (-5.925,-5.375) (-5.936,-5.325) (-5.947,-5.275) (-5.958,-5.225) (-5.968,-5.175) (-5.977,-5.125) (-5.986,-5.075) (-5.994,-5.025) (-6.002,-4.975) (-6.009,-4.925) (-6.016,-4.875) (-6.022,-4.825) (-6.028,-4.775) (-6.033,-4.725) (-6.038,-4.675) (-6.042,-4.625) (-6.045,-4.575) (-6.048,-4.525) (-6.050,-4.475) (-6.052,-4.425) (-6.054,-4.375) (-6.054,-4.325) (-6.055,-4.275) (-6.054,-4.225) (-6.053,-4.175) (-6.052,-4.125) (-6.050,-4.075) (-6.047,-4.025) (-6.044,-3.975) (-6.040,-3.925) (-6.036,-3.875) (-6.031,-3.825) (-6.026,-3.775) (-6.020,-3.725) (-6.013,-3.675) (-6.006,-3.625) (-5.998,-3.575) (-5.990,-3.525) (-5.981,-3.475) (-5.971,-3.425) (-5.961,-3.375) (-5.950,-3.325) (-5.939,-3.275) (-5.926,-3.225) (-5.914,-3.175) (-5.900,-3.125) (-5.886,-3.075) (-5.871,-3.025) (-5.856,-2.975) (-5.839,-2.925) (-5.822,-2.875) (-5.805,-2.825) (-5.786,-2.775) (-5.767,-2.725) (-5.747,-2.675) (-5.727,-2.625) (-5.705,-2.575) (-5.683,-2.525) (-5.660,-2.475) (-5.636,-2.425) (-5.611,-2.375) (-5.585,-2.325) (-5.559,-2.275) (-5.531,-2.225) (-5.503,-2.175) (-5.475,-2.128) (-5.443,-2.075) (-5.411,-2.025) (-5.378,-1.975) (-5.350,-1.933) (-5.310,-1.875) (-5.275,-1.827) (-5.237,-1.775) (-5.200,-1.727) (-5.158,-1.675) (-5.125,-1.634) (-5.075,-1.576) (-5.030,-1.525) (-4.985,-1.475) (-4.950,-1.438) (-4.913,-1.400) (-4.875,-1.362) (-4.838,-1.325) (-4.800,-1.289) (-4.758,-1.250) (-4.702,-1.200) (-4.650,-1.155) (-4.600,-1.113) (-4.553,-1.075) (-4.500,-1.034) (-4.456,-1.000) (-4.400,-0.959) (-4.352,-0.925) (-4.300,-0.889) (-4.250,-0.856) (-4.203,-0.825) (-4.150,-0.792) (-4.100,-0.761) (-4.050,-0.732) (-4.000,-0.703) (-3.950,-0.676) (-3.902,-0.650) (-3.854,-0.625) (-3.805,-0.600) (-3.754,-0.575) (-3.701,-0.550) (-3.650,-0.527) (-3.600,-0.504) (-3.550,-0.483) (-3.500,-0.462) (-3.450,-0.442) (-3.400,-0.422) (-3.350,-0.403) (-3.300,-0.385) (-3.250,-0.367) (-3.200,-0.350) (-3.150,-0.334) (-3.100,-0.318) (-3.050,-0.302) (-3.000,-0.287) (-2.950,-0.273) (-2.900,-0.259) (-2.850,-0.246) (-2.800,-0.233) (-2.750,-0.220) (-2.700,-0.209) (-2.650,-0.197) (-2.600,-0.186) (-2.550,-0.175) (-2.500,-0.165) (-2.450,-0.155) (-2.400,-0.146) (-2.350,-0.137) (-2.300,-0.128) (-2.250,-0.120) (-2.200,-0.112) (-2.150,-0.105) (-2.100,-0.098) (-2.050,-0.091) (-2.000,-0.084) (-1.950,-0.078) (-1.900,-0.072) (-1.850,-0.067) (-1.800,-0.061) (-1.750,-0.056) (-1.700,-0.052) (-1.650,-0.047) (-1.600,-0.043) (-1.550,-0.039) (-1.500,-0.035) (-1.450,-0.032) (-1.400,-0.029) (-1.350,-0.026) (-1.302,-0.023)}
\EQ{(1.300,0.023) (1.350,0.026) (1.400,0.029) (1.450,0.032) (1.500,0.035) (1.550,0.039) (1.600,0.043) (1.650,0.047) (1.700,0.052) (1.750,0.056) (1.800,0.061) (1.850,0.067) (1.900,0.072) (1.950,0.078) (2.000,0.084) (2.050,0.091) (2.100,0.098) (2.150,0.105) (2.200,0.112) (2.250,0.120) (2.300,0.128) (2.350,0.137) (2.400,0.146) (2.450,0.155) (2.500,0.165) (2.550,0.175) (2.600,0.186) (2.650,0.197) (2.700,0.209) (2.750,0.220) (2.800,0.233) (2.850,0.246) (2.900,0.259) (2.950,0.273) (3.000,0.287) (3.050,0.302) (3.100,0.318) (3.150,0.334) (3.199,0.350) (3.250,0.367) (3.300,0.385) (3.350,0.403) (3.400,0.422) (3.450,0.442) (3.500,0.462) (3.550,0.483) (3.600,0.504) (3.646,0.525) (3.700,0.550) (3.750,0.573) (3.800,0.598) (3.850,0.623) (3.900,0.649) (3.949,0.675) (3.994,0.700) (4.038,0.725) (4.100,0.761) (4.150,0.792) (4.200,0.823) (4.250,0.856) (4.300,0.889) (4.350,0.923) (4.400,0.959) (4.450,0.996) (4.500,1.034) (4.550,1.073) (4.600,1.113) (4.644,1.150) (4.700,1.198) (4.750,1.243) (4.800,1.289) (4.838,1.325) (4.875,1.362) (4.913,1.400) (4.950,1.438) (4.985,1.475) (5.025,1.519) (5.074,1.575) (5.117,1.625) (5.150,1.665) (5.198,1.725) (5.237,1.775) (5.274,1.825) (5.310,1.875) (5.345,1.925) (5.375,1.970) (5.411,2.025) (5.443,2.075) (5.473,2.125) (5.500,2.171) (5.525,2.214) (5.550,2.259) (5.575,2.306) (5.598,2.350) (5.623,2.400) (5.648,2.450) (5.671,2.500) (5.694,2.550) (5.716,2.600) (5.737,2.650) (5.757,2.700) (5.777,2.750) (5.796,2.800) (5.814,2.850) (5.831,2.900) (5.848,2.950) (5.863,3.000) (5.879,3.050) (5.893,3.100) (5.907,3.150) (5.920,3.200) (5.933,3.250) (5.944,3.300) (5.956,3.350) (5.966,3.400) (5.976,3.450) (5.985,3.500) (5.994,3.550) (6.002,3.600) (6.010,3.650) (6.017,3.700) (6.023,3.750) (6.029,3.800) (6.034,3.850) (6.038,3.900) (6.042,3.950) (6.046,4.000) (6.049,4.050) (6.051,4.100) (6.053,4.150) (6.054,4.200) (6.054,4.250) (6.054,4.300) (6.054,4.350) (6.053,4.400) (6.051,4.450) (6.049,4.500) (6.047,4.550) (6.043,4.600) (6.040,4.650) (6.035,4.700) (6.031,4.750) (6.025,4.800) (6.019,4.850) (6.013,4.900) (6.006,4.950) (5.998,5.000) (5.990,5.050) (5.982,5.100) (5.973,5.150) (5.963,5.200) (5.953,5.250) (5.942,5.300) (5.931,5.350) (5.919,5.400) (5.906,5.450) (5.893,5.500) (5.880,5.550) (5.866,5.600) (5.851,5.650) (5.836,5.700) (5.820,5.750) (5.804,5.800) (5.787,5.850) (5.769,5.900) (5.751,5.950) (5.732,6.000) (5.713,6.050) (5.693,6.100) (5.673,6.150) (5.652,6.200) (5.630,6.250) (5.608,6.300) (5.585,6.350) (5.561,6.400) (5.537,6.450) (5.512,6.500) (5.487,6.550) (5.461,6.600) (5.434,6.650) (5.406,6.700) (5.378,6.750) (5.350,6.799) (5.320,6.850) (5.289,6.900) (5.258,6.950) (5.226,7.000) (5.194,7.050) (5.160,7.100) (5.126,7.150) (5.091,7.200) (5.056,7.250) (5.025,7.292) (4.981,7.350) (4.950,7.391) (4.904,7.450) (4.864,7.500) (4.825,7.547) (4.781,7.600) (4.738,7.650) (4.700,7.693) (4.650,7.749) (4.603,7.800) (4.555,7.850) (4.507,7.900) (4.457,7.950) (4.407,8.000) (4.355,8.050) (4.302,8.100) (4.250,8.147) (4.200,8.192) (4.162,8.225) (4.104,8.275) (4.050,8.320) (4.000,8.361) (3.952,8.400) (3.900,8.441) (3.855,8.475) (3.800,8.517) (3.755,8.550) (3.700,8.589) (3.650,8.625) (3.600,8.659) (3.550,8.693) (3.501,8.725) (3.450,8.758) (3.400,8.790) (3.350,8.821) (3.302,8.850) (3.250,8.881) (3.200,8.910) (3.150,8.939) (3.100,8.967) (3.050,8.994) (3.000,9.021) (2.950,9.047) (2.900,9.072) (2.850,9.097) (2.800,9.122) (2.750,9.146) (2.700,9.169) (2.650,9.192) (2.600,9.214) (2.550,9.236) (2.500,9.257) (2.450,9.278) (2.400,9.298) (2.350,9.318) (2.300,9.337) (2.250,9.356) (2.200,9.374) (2.150,9.392) (2.100,9.410) (2.050,9.427) (2.000,9.443) (1.950,9.459) (1.900,9.475) (1.850,9.490) (1.800,9.504) (1.750,9.519) (1.700,9.533) (1.650,9.546) (1.600,9.559) (1.550,9.571) (1.500,9.583) (1.450,9.595) (1.400,9.606) (1.350,9.617) (1.300,9.628) (1.250,9.638) (1.200,9.647) (1.150,9.656) (1.100,9.665) (1.050,9.674) (1.000,9.682) (0.950,9.689) (0.900,9.696) (0.850,9.703) (0.800,9.710) (0.750,9.716) (0.700,9.721) (0.650,9.726) (0.600,9.731) (0.550,9.736) (0.500,9.740) (0.450,9.743) (0.400,9.747) (0.350,9.750) (0.300,9.752) (0.250,9.754) (0.200,9.756) (0.150,9.757) (0.100,9.758) (0.050,9.759) (0.000,9.759) (-0.050,9.759) (-0.100,9.758) (-0.150,9.757) (-0.200,9.756) (-0.250,9.754) (-0.300,9.752) (-0.350,9.750) (-0.400,9.747) (-0.450,9.743) (-0.500,9.740) (-0.550,9.736) (-0.600,9.731) (-0.650,9.726) (-0.700,9.721) (-0.750,9.716) (-0.800,9.710) (-0.850,9.703) (-0.900,9.696) (-0.950,9.689) (-1.000,9.682) (-1.050,9.674) (-1.100,9.665) (-1.150,9.656) (-1.200,9.647) (-1.250,9.638) (-1.300,9.628) (-1.350,9.617) (-1.400,9.606) (-1.450,9.595) (-1.500,9.583) (-1.550,9.571) (-1.600,9.559) (-1.650,9.546) (-1.700,9.533) (-1.750,9.519) (-1.800,9.504) (-1.850,9.490) (-1.899,9.475) (-1.950,9.459) (-2.000,9.443) (-2.050,9.427) (-2.100,9.410) (-2.150,9.392) (-2.198,9.375) (-2.250,9.356) (-2.300,9.337) (-2.350,9.318) (-2.400,9.298) (-2.450,9.278) (-2.500,9.257) (-2.550,9.236) (-2.600,9.214) (-2.650,9.192) (-2.700,9.169) (-2.750,9.146) (-2.800,9.122) (-2.850,9.097) (-2.900,9.072) (-2.950,9.047) (-3.000,9.021) (-3.050,8.994) (-3.100,8.967) (-3.150,8.939) (-3.200,8.910) (-3.250,8.881) (-3.300,8.851) (-3.344,8.825) (-3.400,8.790) (-3.450,8.758) (-3.500,8.726) (-3.550,8.693) (-3.600,8.659) (-3.650,8.625) (-3.700,8.589) (-3.750,8.553) (-3.800,8.517) (-3.850,8.479) (-3.900,8.441) (-3.950,8.401) (-4.000,8.361) (-4.045,8.325) (-4.100,8.279) (-4.150,8.236) (-4.191,8.200) (-4.247,8.150) (-4.300,8.101) (-4.350,8.055) (-4.400,8.007) (-4.450,7.957) (-4.500,7.907) (-4.550,7.856) (-4.600,7.803) (-4.649,7.750) (-4.694,7.700) (-4.738,7.650) (-4.775,7.607) (-4.823,7.550) (-4.864,7.500) (-4.900,7.455) (-4.943,7.400) (-4.975,7.359) (-5.019,7.300) (-5.050,7.258) (-5.091,7.200) (-5.125,7.152) (-5.160,7.100) (-5.194,7.050) (-5.225,7.002) (-5.258,6.950) (-5.289,6.900) (-5.320,6.850) (-5.349,6.800) (-5.375,6.755) (-5.400,6.711) (-5.425,6.666) (-5.450,6.620) (-5.474,6.575) (-5.500,6.525) (-5.525,6.475) (-5.549,6.425) (-5.573,6.375) (-5.596,6.325) (-5.619,6.275) (-5.641,6.225) (-5.662,6.175) (-5.683,6.125) (-5.703,6.075) (-5.723,6.025) (-5.742,5.975) (-5.760,5.925) (-5.778,5.875) (-5.795,5.825) (-5.812,5.775) (-5.828,5.725) (-5.843,5.675) (-5.858,5.625) (-5.873,5.575) (-5.887,5.525) (-5.900,5.475) (-5.913,5.425) (-5.925,5.375) (-5.936,5.325) (-5.947,5.275) (-5.958,5.225) (-5.968,5.175) (-5.977,5.125) (-5.986,5.075) (-5.994,5.025) (-6.002,4.975) (-6.009,4.925) (-6.016,4.875) (-6.022,4.825) (-6.028,4.775) (-6.033,4.725) (-6.038,4.675) (-6.042,4.625) (-6.045,4.575) (-6.048,4.525) (-6.050,4.475) (-6.052,4.425) (-6.054,4.375) (-6.054,4.325) (-6.055,4.275) (-6.054,4.225) (-6.053,4.175) (-6.052,4.125) (-6.050,4.075) (-6.047,4.025) (-6.044,3.975) (-6.040,3.925) (-6.036,3.875) (-6.031,3.825) (-6.026,3.775) (-6.020,3.725) (-6.013,3.675) (-6.006,3.625) (-5.998,3.575) (-5.990,3.525) (-5.981,3.475) (-5.971,3.425) (-5.961,3.375) (-5.950,3.325) (-5.939,3.275) (-5.926,3.225) (-5.914,3.175) (-5.900,3.125) (-5.886,3.075) (-5.871,3.025) (-5.856,2.975) (-5.839,2.925) (-5.822,2.875) (-5.805,2.825) (-5.786,2.775) (-5.767,2.725) (-5.747,2.675) (-5.727,2.625) (-5.705,2.575) (-5.683,2.525) (-5.660,2.475) (-5.636,2.425) (-5.611,2.375) (-5.585,2.325) (-5.559,2.275) (-5.531,2.225) (-5.503,2.175) (-5.475,2.128) (-5.443,2.075) (-5.411,2.025) (-5.378,1.975) (-5.350,1.933) (-5.310,1.875) (-5.275,1.827) (-5.237,1.775) (-5.200,1.727) (-5.158,1.675) (-5.125,1.634) (-5.075,1.576) (-5.030,1.525) (-4.985,1.475) (-4.950,1.438) (-4.913,1.400) (-4.875,1.362) (-4.838,1.325) (-4.800,1.289) (-4.758,1.250) (-4.702,1.200) (-4.650,1.155) (-4.600,1.113) (-4.553,1.075) (-4.500,1.034) (-4.456,1.000) (-4.400,0.959) (-4.352,0.925) (-4.300,0.889) (-4.250,0.856) (-4.203,0.825) (-4.150,0.792) (-4.100,0.761) (-4.050,0.732) (-4.000,0.703) (-3.950,0.676) (-3.902,0.650) (-3.854,0.625) (-3.805,0.600) (-3.754,0.575) (-3.701,0.550) (-3.650,0.527) (-3.600,0.504) (-3.550,0.483) (-3.500,0.462) (-3.450,0.442) (-3.400,0.422) (-3.350,0.403) (-3.300,0.385) (-3.250,0.367) (-3.200,0.350) (-3.150,0.334) (-3.100,0.318) (-3.050,0.302) (-3.000,0.287) (-2.950,0.273) (-2.900,0.259) (-2.850,0.246) (-2.800,0.233) (-2.750,0.220) (-2.700,0.209) (-2.650,0.197) (-2.600,0.186) (-2.550,0.175) (-2.500,0.165) (-2.450,0.155) (-2.400,0.146) (-2.350,0.137) (-2.300,0.128) (-2.250,0.120) (-2.200,0.112) (-2.150,0.105) (-2.100,0.098) (-2.050,0.091) (-2.000,0.084) (-1.950,0.078) (-1.900,0.072) (-1.850,0.067) (-1.800,0.061) (-1.750,0.056) (-1.700,0.052) (-1.650,0.047) (-1.600,0.043) (-1.550,0.039) (-1.500,0.035) (-1.450,0.032) (-1.400,0.029) (-1.350,0.026) (-1.302,0.023)}
\EQ{(1.300,0.046) (1.350,0.052) (1.400,0.058) (1.450,0.064) (1.500,0.071) (1.550,0.079) (1.600,0.086) (1.650,0.095) (1.700,0.104) (1.750,0.113) (1.800,0.123) (1.850,0.134) (1.900,0.145) (1.950,0.157) (2.000,0.170) (2.050,0.183) (2.100,0.197) (2.150,0.212) (2.200,0.227) (2.250,0.243) (2.300,0.260) (2.350,0.278) (2.400,0.297) (2.450,0.317) (2.500,0.337) (2.550,0.359) (2.600,0.381) (2.650,0.405) (2.700,0.429) (2.750,0.455) (2.800,0.482) (2.850,0.510) (2.900,0.539) (2.950,0.570) (2.998,0.600) (3.050,0.635) (3.100,0.670) (3.141,0.700) (3.200,0.745) (3.250,0.785) (3.298,0.825) (3.350,0.871) (3.400,0.917) (3.450,0.966) (3.500,1.017) (3.550,1.071) (3.598,1.125) (3.640,1.175) (3.675,1.219) (3.718,1.275) (3.750,1.319) (3.789,1.375) (3.822,1.425) (3.850,1.470) (3.883,1.525) (3.912,1.575) (3.939,1.625) (3.965,1.675) (3.990,1.725) (4.014,1.775) (4.036,1.825) (4.057,1.875) (4.077,1.925) (4.096,1.975) (4.114,2.025) (4.131,2.075) (4.147,2.125) (4.162,2.175) (4.176,2.225) (4.189,2.275) (4.201,2.325) (4.213,2.375) (4.223,2.425) (4.232,2.475) (4.241,2.525) (4.249,2.575) (4.256,2.625) (4.262,2.675) (4.267,2.725) (4.271,2.775) (4.275,2.825) (4.278,2.875) (4.280,2.925) (4.281,2.975) (4.281,3.025) (4.281,3.075) (4.280,3.125) (4.278,3.175) (4.275,3.225) (4.272,3.275) (4.268,3.325) (4.263,3.375) (4.257,3.425) (4.250,3.475) (4.243,3.525) (4.235,3.575) (4.226,3.625) (4.217,3.675) (4.206,3.725) (4.195,3.775) (4.184,3.825) (4.171,3.875) (4.157,3.925) (4.143,3.975) (4.128,4.025) (4.112,4.075) (4.096,4.125) (4.078,4.175) (4.060,4.225) (4.041,4.275) (4.021,4.325) (4.000,4.375) (3.979,4.425) (3.956,4.475) (3.933,4.525) (3.908,4.575) (3.883,4.625) (3.857,4.675) (3.830,4.725) (3.802,4.775) (3.775,4.821) (3.743,4.875) (3.711,4.925) (3.679,4.975) (3.650,5.019) (3.611,5.075) (3.576,5.125) (3.539,5.175) (3.501,5.225) (3.462,5.275) (3.425,5.321) (3.380,5.375) (3.336,5.425) (3.300,5.466) (3.250,5.520) (3.200,5.573) (3.150,5.623) (3.100,5.672) (3.050,5.720) (3.000,5.766) (2.961,5.800) (2.904,5.850) (2.850,5.895) (2.800,5.935) (2.750,5.974) (2.700,6.013) (2.650,6.049) (2.600,6.085) (2.550,6.120) (2.505,6.150) (2.450,6.186) (2.400,6.218) (2.350,6.249) (2.306,6.275) (2.250,6.308) (2.200,6.336) (2.150,6.363) (2.100,6.390) (2.050,6.415) (2.000,6.440) (1.950,6.464) (1.900,6.488) (1.850,6.510) (1.800,6.532) (1.750,6.553) (1.700,6.574) (1.650,6.593) (1.600,6.612) (1.550,6.631) (1.500,6.648) (1.450,6.665) (1.400,6.682) (1.350,6.698) (1.300,6.713) (1.250,6.727) (1.200,6.741) (1.150,6.754) (1.100,6.767) (1.050,6.779) (1.000,6.791) (0.950,6.801) (0.900,6.812) (0.850,6.821) (0.800,6.831) (0.750,6.839) (0.700,6.847) (0.650,6.854) (0.600,6.861) (0.550,6.868) (0.500,6.873) (0.450,6.879) (0.400,6.883) (0.350,6.887) (0.300,6.891) (0.250,6.894) (0.200,6.896) (0.150,6.898) (0.100,6.900) (0.050,6.900) (0.000,6.901) (-0.050,6.900) (-0.100,6.900) (-0.150,6.898) (-0.200,6.896) (-0.250,6.894) (-0.300,6.891) (-0.350,6.887) (-0.400,6.883) (-0.450,6.879) (-0.500,6.873) (-0.550,6.868) (-0.600,6.861) (-0.650,6.854) (-0.700,6.847) (-0.750,6.839) (-0.800,6.831) (-0.850,6.821) (-0.900,6.812) (-0.950,6.801) (-1.000,6.791) (-1.050,6.779) (-1.100,6.767) (-1.150,6.754) (-1.200,6.741) (-1.250,6.727) (-1.300,6.713) (-1.350,6.698) (-1.400,6.682) (-1.450,6.665) (-1.500,6.648) (-1.550,6.631) (-1.600,6.612) (-1.650,6.593) (-1.697,6.575) (-1.750,6.553) (-1.800,6.532) (-1.850,6.510) (-1.900,6.488) (-1.950,6.464) (-2.000,6.440) (-2.050,6.415) (-2.100,6.390) (-2.150,6.363) (-2.200,6.336) (-2.250,6.308) (-2.300,6.279) (-2.348,6.250) (-2.400,6.218) (-2.450,6.186) (-2.500,6.154) (-2.543,6.125) (-2.600,6.085) (-2.649,6.050) (-2.700,6.013) (-2.749,5.975) (-2.800,5.935) (-2.844,5.900) (-2.900,5.853) (-2.950,5.810) (-2.990,5.775) (-3.044,5.725) (-3.097,5.675) (-3.148,5.625) (-3.198,5.575) (-3.246,5.525) (-3.292,5.475) (-3.336,5.425) (-3.375,5.380) (-3.421,5.325) (-3.462,5.275) (-3.500,5.227) (-3.539,5.175) (-3.575,5.126) (-3.611,5.075) (-3.646,5.025) (-3.675,4.981) (-3.711,4.925) (-3.743,4.875) (-3.773,4.825) (-3.800,4.778) (-3.825,4.734) (-3.850,4.688) (-3.875,4.641) (-3.900,4.592) (-3.925,4.541) (-3.950,4.488) (-3.975,4.433) (-4.000,4.376) (-4.021,4.325) (-4.041,4.275) (-4.060,4.225) (-4.078,4.175) (-4.096,4.125) (-4.112,4.075) (-4.128,4.025) (-4.143,3.975) (-4.157,3.925) (-4.171,3.875) (-4.184,3.825) (-4.195,3.775) (-4.206,3.725) (-4.217,3.675) (-4.226,3.625) (-4.235,3.575) (-4.243,3.525) (-4.250,3.475) (-4.257,3.425) (-4.263,3.375) (-4.268,3.325) (-4.272,3.275) (-4.275,3.225) (-4.278,3.175) (-4.280,3.125) (-4.281,3.075) (-4.281,3.025) (-4.281,2.975) (-4.280,2.925) (-4.278,2.875) (-4.275,2.825) (-4.271,2.775) (-4.267,2.725) (-4.262,2.675) (-4.256,2.625) (-4.249,2.575) (-4.241,2.525) (-4.232,2.475) (-4.223,2.425) (-4.213,2.375) (-4.201,2.325) (-4.189,2.275) (-4.176,2.225) (-4.162,2.175) (-4.147,2.125) (-4.131,2.075) (-4.114,2.025) (-4.096,1.975) (-4.077,1.925) (-4.057,1.875) (-4.036,1.825) (-4.014,1.775) (-3.990,1.725) (-3.965,1.675) (-3.939,1.625) (-3.912,1.575) (-3.883,1.525) (-3.853,1.475) (-3.825,1.430) (-3.789,1.375) (-3.754,1.325) (-3.718,1.275) (-3.680,1.225) (-3.640,1.175) (-3.600,1.127) (-3.554,1.075) (-3.508,1.025) (-3.459,0.975) (-3.408,0.925) (-3.355,0.875) (-3.300,0.827) (-3.250,0.785) (-3.207,0.750) (-3.150,0.707) (-3.107,0.675) (-3.050,0.635) (-3.000,0.602) (-2.950,0.570) (-2.900,0.539) (-2.850,0.510) (-2.800,0.482) (-2.750,0.455) (-2.700,0.429) (-2.650,0.405) (-2.600,0.381) (-2.550,0.359) (-2.500,0.337) (-2.450,0.317) (-2.400,0.297) (-2.350,0.278) (-2.300,0.260) (-2.250,0.243) (-2.200,0.227) (-2.150,0.212) (-2.100,0.197) (-2.050,0.183) (-2.000,0.170) (-1.950,0.157) (-1.900,0.145) (-1.850,0.134) (-1.800,0.123) (-1.750,0.113) (-1.700,0.104) (-1.650,0.095) (-1.600,0.086) (-1.550,0.079) (-1.500,0.071) (-1.450,0.064) (-1.400,0.058) (-1.350,0.052) (-1.300,0.046)}
\EQ{(1.300,0.070) (1.350,0.078) (1.400,0.087) (1.450,0.097) (1.500,0.107) (1.550,0.118) (1.600,0.130) (1.650,0.143) (1.700,0.157) (1.750,0.171) (1.800,0.187) (1.850,0.203) (1.900,0.220) (1.950,0.239) (2.000,0.258) (2.050,0.279) (2.098,0.300) (2.150,0.324) (2.200,0.348) (2.250,0.374) (2.298,0.400) (2.342,0.425) (2.400,0.460) (2.450,0.492) (2.500,0.525) (2.550,0.561) (2.600,0.598) (2.650,0.638) (2.694,0.675) (2.750,0.724) (2.800,0.772) (2.850,0.822) (2.899,0.875) (2.943,0.925) (2.985,0.975) (3.024,1.025) (3.060,1.075) (3.095,1.125) (3.125,1.171) (3.159,1.225) (3.188,1.275) (3.216,1.325) (3.241,1.375) (3.266,1.425) (3.289,1.475) (3.310,1.525) (3.330,1.575) (3.349,1.625) (3.367,1.675) (3.383,1.725) (3.398,1.775) (3.412,1.825) (3.425,1.875) (3.437,1.925) (3.447,1.975) (3.456,2.025) (3.465,2.075) (3.472,2.125) (3.479,2.175) (3.484,2.225) (3.488,2.275) (3.491,2.325) (3.494,2.375) (3.495,2.425) (3.496,2.475) (3.495,2.525) (3.494,2.575) (3.491,2.625) (3.488,2.675) (3.483,2.725) (3.478,2.775) (3.472,2.825) (3.465,2.875) (3.457,2.925) (3.448,2.975) (3.438,3.025) (3.427,3.075) (3.415,3.125) (3.403,3.175) (3.389,3.225) (3.374,3.275) (3.359,3.325) (3.342,3.375) (3.325,3.423) (3.306,3.475) (3.286,3.525) (3.265,3.575) (3.243,3.625) (3.220,3.675) (3.196,3.725) (3.171,3.775) (3.145,3.825) (3.118,3.875) (3.089,3.925) (3.059,3.975) (3.028,4.025) (3.000,4.068) (2.962,4.125) (2.927,4.175) (2.890,4.225) (2.852,4.275) (2.812,4.325) (2.775,4.370) (2.728,4.425) (2.683,4.475) (2.637,4.525) (2.600,4.563) (2.563,4.600) (2.525,4.637) (2.485,4.675) (2.430,4.725) (2.375,4.773) (2.325,4.815) (2.281,4.850) (2.225,4.894) (2.183,4.925) (2.125,4.967) (2.078,5.000) (2.025,5.035) (1.975,5.067) (1.925,5.098) (1.881,5.125) (1.825,5.157) (1.775,5.185) (1.725,5.212) (1.675,5.238) (1.625,5.262) (1.575,5.286) (1.525,5.309) (1.475,5.331) (1.425,5.352) (1.375,5.373) (1.325,5.392) (1.275,5.411) (1.225,5.428) (1.175,5.445) (1.125,5.462) (1.075,5.477) (1.025,5.492) (0.975,5.505) (0.925,5.519) (0.875,5.531) (0.825,5.543) (0.775,5.553) (0.725,5.564) (0.675,5.573) (0.625,5.582) (0.575,5.590) (0.525,5.597) (0.475,5.604) (0.425,5.610) (0.375,5.616) (0.325,5.620) (0.275,5.624) (0.225,5.628) (0.175,5.630) (0.125,5.632) (0.075,5.634) (0.025,5.634) (-0.050,5.634) (-0.100,5.633) (-0.150,5.631) (-0.200,5.629) (-0.250,5.626) (-0.300,5.622) (-0.350,5.618) (-0.400,5.613) (-0.450,5.607) (-0.500,5.601) (-0.550,5.594) (-0.600,5.586) (-0.650,5.578) (-0.700,5.569) (-0.750,5.559) (-0.800,5.548) (-0.850,5.537) (-0.899,5.525) (-0.950,5.512) (-1.000,5.499) (-1.050,5.484) (-1.100,5.469) (-1.150,5.453) (-1.200,5.437) (-1.250,5.420) (-1.300,5.401) (-1.350,5.382) (-1.400,5.363) (-1.450,5.342) (-1.500,5.320) (-1.550,5.298) (-1.599,5.275) (-1.650,5.250) (-1.700,5.225) (-1.747,5.200) (-1.793,5.175) (-1.838,5.150) (-1.900,5.114) (-1.950,5.083) (-2.000,5.051) (-2.050,5.019) (-2.100,4.985) (-2.149,4.950) (-2.200,4.913) (-2.249,4.875) (-2.300,4.835) (-2.343,4.800) (-2.400,4.751) (-2.450,4.707) (-2.500,4.661) (-2.538,4.625) (-2.575,4.588) (-2.613,4.550) (-2.650,4.511) (-2.700,4.457) (-2.750,4.400) (-2.792,4.350) (-2.825,4.309) (-2.871,4.250) (-2.908,4.200) (-2.944,4.150) (-2.975,4.106) (-3.012,4.050) (-3.044,4.000) (-3.074,3.950) (-3.100,3.906) (-3.125,3.862) (-3.150,3.816) (-3.175,3.768) (-3.200,3.718) (-3.225,3.665) (-3.250,3.610) (-3.275,3.551) (-3.296,3.500) (-3.315,3.450) (-3.333,3.400) (-3.350,3.351) (-3.366,3.300) (-3.382,3.250) (-3.396,3.200) (-3.409,3.150) (-3.421,3.100) (-3.433,3.050) (-3.443,3.000) (-3.453,2.950) (-3.461,2.900) (-3.469,2.850) (-3.475,2.800) (-3.481,2.750) (-3.486,2.700) (-3.490,2.650) (-3.492,2.600) (-3.494,2.550) (-3.495,2.500) (-3.496,2.450) (-3.495,2.400) (-3.493,2.350) (-3.490,2.300) (-3.486,2.250) (-3.481,2.200) (-3.475,2.150) (-3.469,2.100) (-3.461,2.050) (-3.452,2.000) (-3.442,1.950) (-3.431,1.900) (-3.419,1.850) (-3.405,1.800) (-3.391,1.750) (-3.375,1.700) (-3.358,1.650) (-3.340,1.600) (-3.321,1.550) (-3.300,1.501) (-3.277,1.450) (-3.254,1.400) (-3.229,1.350) (-3.202,1.300) (-3.175,1.252) (-3.144,1.200) (-3.112,1.150) (-3.078,1.100) (-3.042,1.050) (-3.004,1.000) (-2.964,0.950) (-2.925,0.904) (-2.876,0.850) (-2.828,0.800) (-2.777,0.750) (-2.725,0.702) (-2.675,0.659) (-2.634,0.625) (-2.575,0.579) (-2.525,0.543) (-2.475,0.508) (-2.425,0.475) (-2.375,0.444) (-2.325,0.415) (-2.275,0.387) (-2.225,0.361) (-2.175,0.336) (-2.125,0.312) (-2.075,0.290) (-2.025,0.269) (-1.975,0.248) (-1.925,0.230) (-1.875,0.212) (-1.825,0.195) (-1.775,0.179) (-1.725,0.164) (-1.676,0.150) (-1.625,0.137) (-1.575,0.124) (-1.525,0.113) (-1.475,0.102) (-1.425,0.092) (-1.375,0.082) (-1.325,0.074) (-1.300,0.070)}
\EQ{(1.300,0.093) (1.350,0.104) (1.400,0.116) (1.450,0.130) (1.500,0.144) (1.550,0.159) (1.600,0.175) (1.650,0.193) (1.700,0.211) (1.750,0.231) (1.800,0.252) (1.850,0.275) (1.900,0.299) (1.950,0.324) (1.997,0.350) (2.041,0.375) (2.100,0.412) (2.150,0.444) (2.194,0.475) (2.250,0.517) (2.292,0.550) (2.350,0.599) (2.400,0.645) (2.450,0.693) (2.500,0.746) (2.548,0.800) (2.589,0.850) (2.625,0.896) (2.664,0.950) (2.697,1.000) (2.725,1.044) (2.759,1.100) (2.786,1.150) (2.812,1.200) (2.836,1.250) (2.858,1.300) (2.879,1.350) (2.898,1.400) (2.916,1.450) (2.932,1.500) (2.947,1.550) (2.960,1.600) (2.972,1.650) (2.983,1.700) (2.993,1.750) (3.001,1.800) (3.008,1.850) (3.014,1.900) (3.019,1.950) (3.023,2.000) (3.025,2.050) (3.027,2.100) (3.027,2.150) (3.026,2.200) (3.025,2.250) (3.022,2.300) (3.018,2.350) (3.013,2.400) (3.006,2.450) (2.999,2.500) (2.991,2.550) (2.981,2.600) (2.971,2.650) (2.959,2.700) (2.947,2.750) (2.933,2.800) (2.918,2.850) (2.902,2.900) (2.885,2.950) (2.866,3.000) (2.847,3.050) (2.826,3.100) (2.804,3.150) (2.781,3.200) (2.756,3.250) (2.730,3.300) (2.703,3.350) (2.675,3.399) (2.645,3.450) (2.613,3.500) (2.580,3.550) (2.550,3.594) (2.509,3.650) (2.475,3.696) (2.432,3.750) (2.400,3.789) (2.350,3.846) (2.301,3.900) (2.253,3.950) (2.203,4.000) (2.151,4.050) (2.100,4.096) (2.050,4.139) (2.007,4.175) (1.950,4.220) (1.900,4.258) (1.850,4.295) (1.807,4.325) (1.750,4.363) (1.700,4.395) (1.651,4.425) (1.600,4.455) (1.550,4.483) (1.500,4.510) (1.450,4.536) (1.400,4.561) (1.350,4.585) (1.300,4.607) (1.250,4.629) (1.200,4.649) (1.150,4.669) (1.100,4.687) (1.050,4.705) (1.000,4.722) (0.950,4.737) (0.900,4.752) (0.850,4.766) (0.800,4.779) (0.750,4.792) (0.700,4.803) (0.650,4.814) (0.600,4.824) (0.550,4.833) (0.500,4.841) (0.450,4.848) (0.400,4.855) (0.350,4.861) (0.300,4.866) (0.250,4.870) (0.200,4.873) (0.150,4.876) (0.100,4.878) (0.050,4.879) (0.000,4.880) (-0.050,4.879) (-0.100,4.878) (-0.150,4.876) (-0.200,4.873) (-0.250,4.870) (-0.300,4.866) (-0.350,4.861) (-0.400,4.855) (-0.450,4.848) (-0.500,4.841) (-0.550,4.833) (-0.600,4.824) (-0.650,4.814) (-0.700,4.803) (-0.750,4.792) (-0.800,4.779) (-0.850,4.766) (-0.900,4.752) (-0.950,4.737) (-1.000,4.722) (-1.050,4.705) (-1.100,4.687) (-1.150,4.669) (-1.198,4.650) (-1.250,4.629) (-1.300,4.607) (-1.350,4.585) (-1.400,4.561) (-1.450,4.536) (-1.500,4.510) (-1.550,4.483) (-1.600,4.455) (-1.650,4.426) (-1.700,4.395) (-1.750,4.363) (-1.800,4.330) (-1.843,4.300) (-1.900,4.258) (-1.944,4.225) (-2.000,4.181) (-2.050,4.139) (-2.095,4.100) (-2.150,4.051) (-2.200,4.003) (-2.250,3.954) (-2.300,3.901) (-2.347,3.850) (-2.390,3.800) (-2.425,3.758) (-2.472,3.700) (-2.509,3.650) (-2.546,3.600) (-2.575,3.558) (-2.613,3.500) (-2.645,3.450) (-2.675,3.400) (-2.700,3.356) (-2.725,3.310) (-2.750,3.262) (-2.775,3.212) (-2.800,3.159) (-2.825,3.102) (-2.847,3.050) (-2.866,3.000) (-2.885,2.950) (-2.902,2.900) (-2.918,2.850) (-2.933,2.800) (-2.947,2.750) (-2.959,2.700) (-2.971,2.650) (-2.981,2.600) (-2.991,2.550) (-2.999,2.500) (-3.006,2.450) (-3.013,2.400) (-3.018,2.350) (-3.022,2.300) (-3.025,2.250) (-3.026,2.200) (-3.027,2.150) (-3.027,2.100) (-3.025,2.050) (-3.023,2.000) (-3.019,1.950) (-3.014,1.900) (-3.008,1.850) (-3.001,1.800) (-2.993,1.750) (-2.983,1.700) (-2.972,1.650) (-2.960,1.600) (-2.947,1.550) (-2.932,1.500) (-2.916,1.450) (-2.898,1.400) (-2.879,1.350) (-2.858,1.300) (-2.836,1.250) (-2.812,1.200) (-2.786,1.150) (-2.759,1.100) (-2.729,1.050) (-2.700,1.004) (-2.664,0.950) (-2.628,0.900) (-2.589,0.850) (-2.550,0.802) (-2.504,0.750) (-2.457,0.700) (-2.406,0.650) (-2.351,0.600) (-2.300,0.557) (-2.261,0.525) (-2.200,0.479) (-2.158,0.450) (-2.100,0.412) (-2.050,0.381) (-2.000,0.352) (-1.951,0.325) (-1.902,0.300) (-1.850,0.275) (-1.800,0.252) (-1.750,0.231) (-1.700,0.211) (-1.650,0.193) (-1.600,0.175) (-1.550,0.159) (-1.500,0.144) (-1.450,0.130) (-1.400,0.116) (-1.350,0.104) (-1.300,0.093)}
\EQ{(1.300,0.117) (1.350,0.131) (1.400,0.146) (1.450,0.163) (1.500,0.181) (1.549,0.200) (1.600,0.221) (1.650,0.244) (1.700,0.268) (1.750,0.293) (1.800,0.321) (1.849,0.350) (1.900,0.382) (1.950,0.416) (1.997,0.450) (2.050,0.492) (2.090,0.525) (2.145,0.575) (2.196,0.625) (2.243,0.675) (2.287,0.725) (2.325,0.772) (2.365,0.825) (2.400,0.875) (2.432,0.925) (2.462,0.975) (2.490,1.025) (2.516,1.075) (2.540,1.125) (2.562,1.175) (2.582,1.225) (2.600,1.275) (2.617,1.325) (2.632,1.375) (2.646,1.425) (2.658,1.475) (2.669,1.525) (2.679,1.575) (2.687,1.625) (2.693,1.675) (2.699,1.725) (2.703,1.775) (2.706,1.825) (2.707,1.875) (2.708,1.925) (2.707,1.975) (2.705,2.025) (2.701,2.075) (2.697,2.125) (2.691,2.175) (2.684,2.225) (2.676,2.275) (2.667,2.325) (2.656,2.375) (2.644,2.425) (2.631,2.475) (2.617,2.525) (2.602,2.575) (2.585,2.625) (2.567,2.675) (2.547,2.725) (2.527,2.775) (2.504,2.825) (2.481,2.875) (2.456,2.925) (2.430,2.975) (2.402,3.025) (2.375,3.070) (2.341,3.125) (2.308,3.175) (2.275,3.223) (2.237,3.275) (2.200,3.323) (2.158,3.375) (2.125,3.414) (2.075,3.470) (2.025,3.523) (1.975,3.574) (1.925,3.621) (1.875,3.667) (1.836,3.700) (1.776,3.750) (1.725,3.790) (1.677,3.825) (1.625,3.862) (1.575,3.896) (1.530,3.925) (1.475,3.959) (1.425,3.988) (1.375,4.016) (1.325,4.043) (1.275,4.068) (1.225,4.092) (1.175,4.115) (1.125,4.137) (1.075,4.158) (1.025,4.177) (0.975,4.196) (0.925,4.213) (0.875,4.229) (0.825,4.245) (0.775,4.259) (0.725,4.272) (0.675,4.285) (0.625,4.296) (0.575,4.307) (0.525,4.317) (0.475,4.325) (0.425,4.333) (0.375,4.340) (0.325,4.346) (0.275,4.351) (0.225,4.356) (0.175,4.359) (0.125,4.362) (0.075,4.363) (0.025,4.364) (-0.050,4.364) (-0.100,4.363) (-0.150,4.361) (-0.200,4.358) (-0.250,4.354) (-0.300,4.349) (-0.350,4.343) (-0.400,4.337) (-0.450,4.329) (-0.500,4.321) (-0.550,4.312) (-0.600,4.302) (-0.650,4.291) (-0.700,4.279) (-0.750,4.266) (-0.800,4.252) (-0.850,4.237) (-0.900,4.221) (-0.950,4.204) (-1.000,4.187) (-1.050,4.168) (-1.100,4.147) (-1.150,4.126) (-1.200,4.104) (-1.250,4.080) (-1.300,4.056) (-1.350,4.030) (-1.400,4.002) (-1.448,3.975) (-1.500,3.944) (-1.550,3.912) (-1.600,3.879) (-1.642,3.850) (-1.700,3.808) (-1.744,3.775) (-1.800,3.730) (-1.850,3.688) (-1.894,3.650) (-1.948,3.600) (-1.999,3.550) (-2.047,3.500) (-2.093,3.450) (-2.137,3.400) (-2.175,3.354) (-2.218,3.300) (-2.250,3.257) (-2.291,3.200) (-2.325,3.150) (-2.357,3.100) (-2.387,3.050) (-2.416,3.000) (-2.443,2.950) (-2.469,2.900) (-2.493,2.850) (-2.516,2.800) (-2.537,2.750) (-2.557,2.700) (-2.575,2.653) (-2.593,2.600) (-2.610,2.550) (-2.624,2.500) (-2.638,2.450) (-2.650,2.400) (-2.662,2.350) (-2.672,2.300) (-2.680,2.250) (-2.688,2.200) (-2.694,2.150) (-2.699,2.100) (-2.703,2.050) (-2.706,2.000) (-2.707,1.950) (-2.708,1.900) (-2.707,1.850) (-2.704,1.800) (-2.701,1.750) (-2.696,1.700) (-2.690,1.650) (-2.683,1.600) (-2.674,1.550) (-2.664,1.500) (-2.652,1.450) (-2.639,1.400) (-2.625,1.351) (-2.609,1.300) (-2.591,1.250) (-2.572,1.200) (-2.551,1.150) (-2.528,1.100) (-2.503,1.050) (-2.476,1.000) (-2.450,0.954) (-2.416,0.900) (-2.383,0.850) (-2.350,0.805) (-2.308,0.750) (-2.275,0.711) (-2.225,0.655) (-2.175,0.604) (-2.125,0.557) (-2.075,0.513) (-2.029,0.475) (-1.975,0.434) (-1.927,0.400) (-1.875,0.366) (-1.825,0.335) (-1.775,0.307) (-1.725,0.280) (-1.675,0.255) (-1.625,0.232) (-1.575,0.211) (-1.525,0.191) (-1.475,0.172) (-1.425,0.155) (-1.375,0.139) (-1.325,0.124) (-1.300,0.117)}
\EQ{(1.300,0.141) (1.350,0.158) (1.400,0.177) (1.450,0.197) (1.500,0.220) (1.550,0.243) (1.600,0.269) (1.650,0.297) (1.697,0.325) (1.750,0.359) (1.800,0.394) (1.841,0.425) (1.900,0.473) (1.950,0.517) (2.000,0.566) (2.050,0.619) (2.099,0.675) (2.138,0.725) (2.174,0.775) (2.208,0.825) (2.239,0.875) (2.267,0.925) (2.293,0.975) (2.317,1.025) (2.339,1.075) (2.359,1.125) (2.377,1.175) (2.394,1.225) (2.409,1.275) (2.422,1.325) (2.433,1.375) (2.443,1.425) (2.451,1.475) (2.458,1.525) (2.464,1.575) (2.468,1.625) (2.470,1.675) (2.472,1.725) (2.472,1.775) (2.470,1.825) (2.467,1.875) (2.463,1.925) (2.458,1.975) (2.451,2.025) (2.443,2.075) (2.434,2.125) (2.423,2.175) (2.411,2.225) (2.398,2.275) (2.383,2.325) (2.367,2.375) (2.350,2.423) (2.330,2.475) (2.310,2.525) (2.288,2.575) (2.265,2.625) (2.240,2.675) (2.213,2.725) (2.184,2.775) (2.154,2.825) (2.125,2.871) (2.089,2.925) (2.053,2.975) (2.015,3.025) (1.975,3.075) (1.933,3.125) (1.888,3.175) (1.850,3.215) (1.800,3.265) (1.763,3.300) (1.725,3.335) (1.679,3.375) (1.625,3.420) (1.575,3.459) (1.525,3.496) (1.475,3.531) (1.425,3.565) (1.375,3.597) (1.328,3.625) (1.285,3.650) (1.239,3.675) (1.191,3.700) (1.141,3.725) (1.087,3.750) (1.031,3.775) (0.975,3.798) (0.925,3.817) (0.875,3.835) (0.825,3.852) (0.775,3.868) (0.725,3.883) (0.675,3.897) (0.625,3.909) (0.575,3.921) (0.525,3.932) (0.475,3.941) (0.425,3.950) (0.375,3.958) (0.325,3.964) (0.275,3.970) (0.225,3.975) (0.175,3.978) (0.125,3.981) (0.075,3.983) (0.025,3.984) (-0.050,3.984) (-0.100,3.982) (-0.150,3.980) (-0.200,3.977) (-0.250,3.972) (-0.300,3.967) (-0.350,3.961) (-0.400,3.954) (-0.450,3.946) (-0.500,3.937) (-0.550,3.926) (-0.600,3.915) (-0.650,3.903) (-0.700,3.890) (-0.750,3.876) (-0.800,3.860) (-0.850,3.844) (-0.900,3.826) (-0.950,3.808) (-1.000,3.788) (-1.050,3.767) (-1.100,3.744) (-1.150,3.721) (-1.200,3.696) (-1.250,3.669) (-1.300,3.641) (-1.350,3.612) (-1.400,3.581) (-1.448,3.550) (-1.500,3.514) (-1.550,3.478) (-1.600,3.440) (-1.650,3.400) (-1.700,3.357) (-1.750,3.312) (-1.790,3.275) (-1.840,3.225) (-1.875,3.189) (-1.910,3.150) (-1.950,3.105) (-1.995,3.050) (-2.034,3.000) (-2.071,2.950) (-2.100,2.908) (-2.139,2.850) (-2.170,2.800) (-2.199,2.750) (-2.225,2.703) (-2.250,2.654) (-2.275,2.603) (-2.299,2.550) (-2.320,2.500) (-2.340,2.450) (-2.358,2.400) (-2.375,2.351) (-2.391,2.300) (-2.405,2.250) (-2.417,2.200) (-2.429,2.150) (-2.439,2.100) (-2.447,2.050) (-2.455,2.000) (-2.461,1.950) (-2.466,1.900) (-2.469,1.850) (-2.471,1.800) (-2.472,1.750) (-2.471,1.700) (-2.469,1.650) (-2.466,1.600) (-2.461,1.550) (-2.455,1.500) (-2.447,1.450) (-2.438,1.400) (-2.427,1.350) (-2.415,1.300) (-2.401,1.250) (-2.386,1.200) (-2.369,1.150) (-2.350,1.101) (-2.329,1.050) (-2.306,1.000) (-2.281,0.950) (-2.253,0.900) (-2.225,0.852) (-2.192,0.800) (-2.157,0.750) (-2.125,0.708) (-2.078,0.650) (-2.033,0.600) (-1.984,0.550) (-1.931,0.500) (-1.875,0.452) (-1.825,0.413) (-1.775,0.376) (-1.725,0.343) (-1.675,0.312) (-1.625,0.283) (-1.575,0.256) (-1.525,0.231) (-1.475,0.208) (-1.425,0.187) (-1.375,0.167) (-1.327,0.150) (-1.300,0.141)}
\EQ{(1.300,0.165) (1.350,0.186) (1.400,0.208) (1.450,0.233) (1.500,0.259) (1.550,0.288) (1.600,0.319) (1.646,0.350) (1.700,0.390) (1.744,0.425) (1.800,0.474) (1.850,0.522) (1.900,0.575) (1.943,0.625) (1.975,0.665) (2.018,0.725) (2.050,0.773) (2.082,0.825) (2.109,0.875) (2.134,0.925) (2.157,0.975) (2.178,1.025) (2.197,1.075) (2.213,1.125) (2.228,1.175) (2.242,1.225) (2.253,1.275) (2.263,1.325) (2.271,1.375) (2.277,1.425) (2.282,1.475) (2.286,1.525) (2.288,1.575) (2.288,1.625) (2.288,1.675) (2.285,1.725) (2.281,1.775) (2.276,1.825) (2.269,1.875) (2.261,1.925) (2.252,1.975) (2.241,2.025) (2.228,2.075) (2.215,2.125) (2.199,2.175) (2.182,2.225) (2.164,2.275) (2.144,2.325) (2.122,2.375) (2.100,2.423) (2.075,2.473) (2.050,2.520) (2.025,2.565) (1.989,2.625) (1.956,2.675) (1.925,2.720) (1.885,2.775) (1.850,2.820) (1.805,2.875) (1.761,2.925) (1.725,2.964) (1.690,3.000) (1.650,3.039) (1.611,3.075) (1.554,3.125) (1.500,3.170) (1.450,3.209) (1.400,3.245) (1.358,3.275) (1.300,3.313) (1.250,3.344) (1.200,3.373) (1.152,3.400) (1.104,3.425) (1.053,3.450) (1.000,3.475) (0.950,3.497) (0.900,3.517) (0.850,3.536) (0.800,3.554) (0.750,3.571) (0.700,3.587) (0.650,3.601) (0.600,3.614) (0.550,3.626) (0.500,3.637) (0.450,3.647) (0.400,3.656) (0.350,3.664) (0.300,3.670) (0.250,3.676) (0.200,3.680) (0.150,3.684) (0.100,3.687) (0.050,3.688) (0.000,3.689) (-0.050,3.688) (-0.100,3.687) (-0.150,3.684) (-0.200,3.680) (-0.250,3.676) (-0.300,3.670) (-0.350,3.664) (-0.400,3.656) (-0.450,3.647) (-0.500,3.637) (-0.550,3.626) (-0.600,3.614) (-0.650,3.601) (-0.700,3.587) (-0.750,3.571) (-0.800,3.554) (-0.850,3.536) (-0.900,3.517) (-0.950,3.497) (-1.000,3.475) (-1.050,3.452) (-1.100,3.427) (-1.150,3.401) (-1.197,3.375) (-1.250,3.344) (-1.300,3.313) (-1.350,3.280) (-1.394,3.250) (-1.450,3.209) (-1.494,3.175) (-1.550,3.129) (-1.600,3.085) (-1.638,3.050) (-1.675,3.015) (-1.714,2.975) (-1.750,2.937) (-1.800,2.881) (-1.846,2.825) (-1.885,2.775) (-1.922,2.725) (-1.956,2.675) (-1.989,2.625) (-2.019,2.575) (-2.047,2.525) (-2.074,2.475) (-2.099,2.425) (-2.122,2.375) (-2.144,2.325) (-2.164,2.275) (-2.182,2.225) (-2.199,2.175) (-2.215,2.125) (-2.228,2.075) (-2.241,2.025) (-2.252,1.975) (-2.261,1.925) (-2.269,1.875) (-2.276,1.825) (-2.281,1.775) (-2.285,1.725) (-2.288,1.675) (-2.288,1.625) (-2.288,1.575) (-2.286,1.525) (-2.282,1.475) (-2.277,1.425) (-2.271,1.375) (-2.263,1.325) (-2.253,1.275) (-2.242,1.225) (-2.228,1.175) (-2.213,1.125) (-2.197,1.075) (-2.178,1.025) (-2.157,0.975) (-2.134,0.925) (-2.109,0.875) (-2.082,0.825) (-2.051,0.775) (-2.018,0.725) (-1.982,0.675) (-1.950,0.634) (-1.900,0.575) (-1.853,0.525) (-1.801,0.475) (-1.750,0.430) (-1.700,0.390) (-1.650,0.353) (-1.600,0.319) (-1.550,0.288) (-1.500,0.259) (-1.450,0.233) (-1.400,0.208) (-1.350,0.186) (-1.300,0.165)}
\EQ{(1.300,0.191) (1.350,0.215) (1.400,0.241) (1.450,0.270) (1.499,0.300) (1.550,0.335) (1.600,0.372) (1.650,0.414) (1.691,0.450) (1.742,0.500) (1.788,0.550) (1.825,0.594) (1.868,0.650) (1.900,0.696) (1.934,0.750) (1.963,0.800) (1.989,0.850) (2.013,0.900) (2.034,0.950) (2.053,1.000) (2.070,1.050) (2.085,1.100) (2.098,1.150) (2.109,1.200) (2.118,1.250) (2.126,1.300) (2.132,1.350) (2.137,1.400) (2.139,1.450) (2.141,1.500) (2.140,1.550) (2.138,1.600) (2.135,1.650) (2.130,1.700) (2.123,1.750) (2.115,1.800) (2.106,1.850) (2.095,1.900) (2.082,1.950) (2.068,2.000) (2.052,2.050) (2.035,2.100) (2.016,2.150) (1.995,2.200) (1.972,2.250) (1.950,2.296) (1.925,2.344) (1.900,2.389) (1.864,2.450) (1.831,2.500) (1.800,2.546) (1.760,2.600) (1.725,2.645) (1.679,2.700) (1.634,2.750) (1.600,2.786) (1.562,2.825) (1.525,2.860) (1.481,2.900) (1.425,2.947) (1.375,2.987) (1.325,3.025) (1.275,3.060) (1.225,3.093) (1.175,3.124) (1.132,3.150) (1.087,3.175) (1.040,3.200) (0.990,3.225) (0.936,3.250) (0.879,3.275) (0.825,3.297) (0.775,3.315) (0.725,3.333) (0.675,3.349) (0.625,3.364) (0.575,3.377) (0.525,3.390) (0.475,3.401) (0.425,3.411) (0.375,3.420) (0.325,3.427) (0.275,3.434) (0.225,3.439) (0.175,3.444) (0.125,3.447) (0.075,3.449) (0.025,3.450) (-0.037,3.450) (-0.100,3.448) (-0.150,3.445) (-0.200,3.442) (-0.250,3.437) (-0.300,3.431) (-0.350,3.424) (-0.400,3.415) (-0.450,3.406) (-0.500,3.395) (-0.550,3.384) (-0.600,3.371) (-0.650,3.356) (-0.700,3.341) (-0.748,3.325) (-0.800,3.306) (-0.850,3.287) (-0.900,3.266) (-0.950,3.244) (-1.000,3.220) (-1.050,3.195) (-1.100,3.168) (-1.150,3.139) (-1.200,3.109) (-1.250,3.077) (-1.300,3.043) (-1.350,3.006) (-1.391,2.975) (-1.450,2.927) (-1.500,2.883) (-1.550,2.836) (-1.587,2.800) (-1.625,2.760) (-1.675,2.705) (-1.721,2.650) (-1.760,2.600) (-1.797,2.550) (-1.831,2.500) (-1.864,2.450) (-1.894,2.400) (-1.922,2.350) (-1.948,2.300) (-1.972,2.250) (-1.995,2.200) (-2.016,2.150) (-2.035,2.100) (-2.052,2.050) (-2.068,2.000) (-2.082,1.950) (-2.095,1.900) (-2.106,1.850) (-2.115,1.800) (-2.123,1.750) (-2.130,1.700) (-2.135,1.650) (-2.138,1.600) (-2.140,1.550) (-2.141,1.500) (-2.139,1.450) (-2.137,1.400) (-2.132,1.350) (-2.126,1.300) (-2.118,1.250) (-2.109,1.200) (-2.098,1.150) (-2.085,1.100) (-2.070,1.050) (-2.053,1.000) (-2.034,0.950) (-2.013,0.900) (-1.989,0.850) (-1.963,0.800) (-1.934,0.750) (-1.903,0.700) (-1.868,0.650) (-1.830,0.600) (-1.788,0.550) (-1.750,0.508) (-1.700,0.459) (-1.650,0.414) (-1.603,0.375) (-1.550,0.335) (-1.500,0.301) (-1.450,0.270) (-1.400,0.241) (-1.350,0.215) (-1.300,0.191)}
\EQ{(1.300,0.216) (1.350,0.244) (1.400,0.274) (1.450,0.308) (1.500,0.345) (1.550,0.385) (1.595,0.425) (1.646,0.475) (1.692,0.525) (1.725,0.565) (1.770,0.625) (1.800,0.669) (1.834,0.725) (1.862,0.775) (1.887,0.825) (1.909,0.875) (1.929,0.925) (1.947,0.975) (1.962,1.025) (1.975,1.075) (1.987,1.125) (1.997,1.175) (2.004,1.225) (2.011,1.275) (2.015,1.325) (2.017,1.375) (2.018,1.425) (2.017,1.475) (2.015,1.525) (2.011,1.575) (2.005,1.625) (1.998,1.675) (1.989,1.725) (1.979,1.775) (1.967,1.825) (1.953,1.875) (1.937,1.925) (1.920,1.975) (1.901,2.025) (1.880,2.075) (1.858,2.125) (1.833,2.175) (1.807,2.225) (1.778,2.275) (1.750,2.321) (1.714,2.375) (1.679,2.425) (1.641,2.475) (1.601,2.525) (1.557,2.575) (1.510,2.625) (1.475,2.661) (1.434,2.700) (1.378,2.750) (1.325,2.794) (1.285,2.825) (1.225,2.869) (1.180,2.900) (1.125,2.935) (1.075,2.965) (1.025,2.993) (0.975,3.020) (0.925,3.045) (0.875,3.068) (0.825,3.089) (0.775,3.109) (0.725,3.128) (0.675,3.145) (0.625,3.161) (0.576,3.175) (0.525,3.188) (0.476,3.200) (0.425,3.211) (0.375,3.220) (0.325,3.229) (0.275,3.236) (0.225,3.241) (0.175,3.246) (0.125,3.249) (0.075,3.252) (0.025,3.253) (-0.050,3.252) (-0.100,3.251) (-0.150,3.248) (-0.200,3.244) (-0.250,3.239) (-0.300,3.232) (-0.350,3.225) (-0.400,3.216) (-0.450,3.206) (-0.500,3.195) (-0.550,3.182) (-0.600,3.168) (-0.650,3.153) (-0.700,3.137) (-0.750,3.119) (-0.798,3.100) (-0.850,3.079) (-0.900,3.056) (-0.950,3.032) (-1.000,3.007) (-1.050,2.980) (-1.100,2.950) (-1.150,2.919) (-1.200,2.886) (-1.250,2.851) (-1.300,2.814) (-1.348,2.775) (-1.400,2.731) (-1.450,2.685) (-1.486,2.650) (-1.525,2.610) (-1.575,2.555) (-1.621,2.500) (-1.661,2.450) (-1.697,2.400) (-1.731,2.350) (-1.763,2.300) (-1.793,2.250) (-1.820,2.200) (-1.846,2.150) (-1.869,2.100) (-1.891,2.050) (-1.911,2.000) (-1.929,1.950) (-1.945,1.900) (-1.960,1.850) (-1.973,1.800) (-1.984,1.750) (-1.994,1.700) (-2.002,1.650) (-2.008,1.600) (-2.013,1.550) (-2.017,1.500) (-2.018,1.450) (-2.018,1.400) (-2.016,1.350) (-2.013,1.300) (-2.008,1.250) (-2.001,1.200) (-1.992,1.150) (-1.982,1.100) (-1.969,1.050) (-1.955,1.000) (-1.938,0.950) (-1.919,0.900) (-1.898,0.850) (-1.875,0.801) (-1.850,0.753) (-1.825,0.710) (-1.787,0.650) (-1.752,0.600) (-1.713,0.550) (-1.675,0.506) (-1.625,0.454) (-1.575,0.407) (-1.525,0.364) (-1.475,0.326) (-1.425,0.291) (-1.375,0.259) (-1.325,0.230) (-1.300,0.216)}
\EQ{(1.300,0.243) (1.350,0.275) (1.400,0.310) (1.450,0.348) (1.500,0.391) (1.550,0.439) (1.600,0.493) (1.647,0.550) (1.684,0.600) (1.717,0.650) (1.747,0.700) (1.774,0.750) (1.798,0.800) (1.819,0.850) (1.838,0.900) (1.855,0.950) (1.869,1.000) (1.881,1.050) (1.892,1.100) (1.900,1.150) (1.906,1.200) (1.911,1.250) (1.914,1.300) (1.915,1.350) (1.914,1.400) (1.911,1.450) (1.907,1.500) (1.901,1.550) (1.894,1.600) (1.885,1.650) (1.873,1.700) (1.861,1.750) (1.846,1.800) (1.830,1.850) (1.812,1.900) (1.792,1.950) (1.770,2.000) (1.746,2.050) (1.720,2.100) (1.692,2.150) (1.661,2.200) (1.629,2.250) (1.593,2.300) (1.555,2.350) (1.515,2.400) (1.475,2.445) (1.425,2.498) (1.375,2.548) (1.325,2.593) (1.275,2.636) (1.226,2.675) (1.175,2.713) (1.125,2.747) (1.083,2.775) (1.025,2.810) (0.975,2.838) (0.925,2.865) (0.875,2.890) (0.825,2.913) (0.775,2.934) (0.725,2.954) (0.675,2.972) (0.625,2.989) (0.575,3.004) (0.525,3.018) (0.475,3.030) (0.425,3.042) (0.375,3.052) (0.325,3.060) (0.275,3.068) (0.225,3.074) (0.175,3.079) (0.125,3.082) (0.075,3.085) (0.025,3.086) (-0.050,3.086) (-0.100,3.084) (-0.150,3.081) (-0.200,3.076) (-0.250,3.071) (-0.300,3.064) (-0.350,3.056) (-0.400,3.047) (-0.450,3.036) (-0.500,3.024) (-0.550,3.011) (-0.600,2.996) (-0.650,2.980) (-0.700,2.963) (-0.750,2.944) (-0.800,2.923) (-0.850,2.901) (-0.900,2.877) (-0.950,2.852) (-0.999,2.825) (-1.050,2.795) (-1.100,2.764) (-1.150,2.730) (-1.192,2.700) (-1.250,2.656) (-1.300,2.615) (-1.345,2.575) (-1.398,2.525) (-1.448,2.475) (-1.493,2.425) (-1.535,2.375) (-1.575,2.325) (-1.611,2.275) (-1.645,2.225) (-1.675,2.178) (-1.706,2.125) (-1.733,2.075) (-1.758,2.025) (-1.781,1.975) (-1.802,1.925) (-1.821,1.875) (-1.838,1.825) (-1.854,1.775) (-1.867,1.725) (-1.879,1.675) (-1.889,1.625) (-1.898,1.575) (-1.905,1.525) (-1.910,1.475) (-1.913,1.425) (-1.915,1.375) (-1.914,1.325) (-1.913,1.275) (-1.909,1.225) (-1.903,1.175) (-1.896,1.125) (-1.887,1.075) (-1.875,1.025) (-1.862,0.975) (-1.847,0.925) (-1.829,0.875) (-1.809,0.825) (-1.786,0.775) (-1.761,0.725) (-1.733,0.675) (-1.701,0.625) (-1.666,0.575) (-1.627,0.525) (-1.584,0.475) (-1.536,0.425) (-1.482,0.375) (-1.425,0.328) (-1.375,0.292) (-1.325,0.258) (-1.300,0.243)}
\EQ{(1.288,0.262) (1.342,0.300) (1.400,0.347) (1.450,0.391) (1.500,0.442) (1.550,0.499) (1.590,0.550) (1.625,0.600) (1.656,0.650) (1.684,0.700) (1.708,0.750) (1.730,0.800) (1.750,0.850) (1.767,0.900) (1.781,0.950) (1.794,1.000) (1.804,1.050) (1.812,1.100) (1.818,1.150) (1.823,1.200) (1.825,1.250) (1.825,1.300) (1.824,1.350) (1.821,1.400) (1.816,1.450) (1.810,1.500) (1.801,1.550) (1.791,1.600) (1.779,1.650) (1.765,1.700) (1.750,1.748) (1.732,1.800) (1.712,1.850) (1.691,1.900) (1.667,1.950) (1.641,2.000) (1.613,2.050) (1.583,2.100) (1.550,2.149) (1.514,2.200) (1.476,2.250) (1.434,2.300) (1.400,2.339) (1.366,2.375) (1.325,2.416) (1.289,2.450) (1.250,2.484) (1.201,2.525) (1.150,2.565) (1.101,2.600) (1.050,2.634) (1.000,2.666) (0.950,2.695) (0.900,2.722) (0.850,2.748) (0.800,2.771) (0.750,2.793) (0.700,2.813) (0.650,2.831) (0.600,2.848) (0.550,2.864) (0.500,2.878) (0.450,2.890) (0.400,2.901) (0.350,2.911) (0.300,2.919) (0.250,2.926) (0.200,2.932) (0.150,2.937) (0.100,2.940) (0.050,2.942) (0.000,2.943) (-0.050,2.942) (-0.100,2.940) (-0.150,2.937) (-0.200,2.932) (-0.250,2.926) (-0.300,2.919) (-0.350,2.911) (-0.400,2.901) (-0.450,2.890) (-0.500,2.878) (-0.550,2.864) (-0.600,2.848) (-0.650,2.831) (-0.700,2.813) (-0.750,2.793) (-0.800,2.771) (-0.850,2.748) (-0.900,2.722) (-0.950,2.695) (-1.000,2.666) (-1.050,2.634) (-1.100,2.601) (-1.150,2.565) (-1.200,2.526) (-1.250,2.484) (-1.289,2.450) (-1.325,2.416) (-1.366,2.375) (-1.400,2.339) (-1.434,2.300) (-1.475,2.251) (-1.514,2.200) (-1.550,2.150) (-1.583,2.100) (-1.613,2.050) (-1.641,2.000) (-1.667,1.950) (-1.691,1.900) (-1.712,1.850) (-1.732,1.800) (-1.749,1.750) (-1.765,1.700) (-1.779,1.650) (-1.791,1.600) (-1.801,1.550) (-1.810,1.500) (-1.816,1.450) (-1.821,1.400) (-1.824,1.350) (-1.825,1.300) (-1.825,1.250) (-1.823,1.200) (-1.818,1.150) (-1.812,1.100) (-1.804,1.050) (-1.794,1.000) (-1.781,0.950) (-1.767,0.900) (-1.750,0.851) (-1.730,0.800) (-1.708,0.750) (-1.684,0.700) (-1.656,0.650) (-1.625,0.601) (-1.590,0.550) (-1.551,0.500) (-1.508,0.450) (-1.459,0.400) (-1.404,0.350) (-1.350,0.307) (-1.307,0.275) (-1.288,0.262)}
\EQ{(1.275,0.280) (1.325,0.319) (1.375,0.362) (1.414,0.400) (1.450,0.438) (1.500,0.497) (1.539,0.550) (1.572,0.600) (1.600,0.648) (1.625,0.696) (1.650,0.750) (1.670,0.800) (1.688,0.850) (1.703,0.900) (1.716,0.950) (1.726,1.000) (1.734,1.050) (1.741,1.100) (1.745,1.150) (1.747,1.200) (1.748,1.250) (1.746,1.300) (1.743,1.350) (1.738,1.400) (1.731,1.450) (1.722,1.500) (1.711,1.550) (1.698,1.600) (1.683,1.650) (1.667,1.700) (1.648,1.750) (1.627,1.800) (1.604,1.850) (1.579,1.900) (1.552,1.950) (1.525,1.995) (1.489,2.050) (1.454,2.100) (1.416,2.150) (1.375,2.200) (1.330,2.250) (1.282,2.300) (1.229,2.350) (1.175,2.397) (1.125,2.437) (1.075,2.475) (1.025,2.509) (0.975,2.542) (0.925,2.571) (0.875,2.599) (0.825,2.625) (0.775,2.649) (0.725,2.671) (0.675,2.691) (0.625,2.710) (0.575,2.727) (0.525,2.742) (0.475,2.756) (0.425,2.768) (0.375,2.779) (0.325,2.789) (0.275,2.797) (0.225,2.804) (0.175,2.809) (0.125,2.813) (0.075,2.816) (0.025,2.817) (-0.050,2.817) (-0.100,2.815) (-0.150,2.811) (-0.200,2.807) (-0.250,2.800) (-0.300,2.793) (-0.350,2.784) (-0.400,2.774) (-0.450,2.762) (-0.500,2.749) (-0.550,2.735) (-0.600,2.719) (-0.650,2.701) (-0.700,2.681) (-0.750,2.660) (-0.800,2.637) (-0.850,2.612) (-0.900,2.586) (-0.950,2.557) (-1.000,2.526) (-1.050,2.492) (-1.100,2.456) (-1.141,2.425) (-1.200,2.376) (-1.250,2.331) (-1.300,2.282) (-1.350,2.228) (-1.396,2.175) (-1.436,2.125) (-1.472,2.075) (-1.506,2.025) (-1.537,1.975) (-1.566,1.925) (-1.592,1.875) (-1.616,1.825) (-1.638,1.775) (-1.658,1.725) (-1.675,1.675) (-1.691,1.625) (-1.705,1.575) (-1.716,1.525) (-1.726,1.475) (-1.734,1.425) (-1.741,1.375) (-1.745,1.325) (-1.747,1.275) (-1.748,1.225) (-1.746,1.175) (-1.743,1.125) (-1.738,1.075) (-1.730,1.025) (-1.721,0.975) (-1.709,0.925) (-1.695,0.875) (-1.679,0.825) (-1.660,0.775) (-1.639,0.725) (-1.614,0.675) (-1.587,0.625) (-1.556,0.575) (-1.525,0.530) (-1.482,0.475) (-1.438,0.425) (-1.400,0.386) (-1.362,0.350) (-1.301,0.300) (-1.275,0.280)}
\EQ{(1.275,0.308) (1.323,0.350) (1.374,0.400) (1.419,0.450) (1.450,0.489) (1.493,0.550) (1.524,0.600) (1.550,0.647) (1.575,0.698) (1.597,0.750) (1.616,0.800) (1.632,0.850) (1.645,0.900) (1.656,0.950) (1.665,1.000) (1.672,1.050) (1.676,1.100) (1.679,1.150) (1.679,1.200) (1.678,1.250) (1.674,1.300) (1.669,1.350) (1.662,1.400) (1.652,1.450) (1.641,1.500) (1.628,1.550) (1.613,1.600) (1.595,1.650) (1.576,1.700) (1.554,1.750) (1.530,1.800) (1.504,1.850) (1.475,1.900) (1.444,1.950) (1.410,2.000) (1.375,2.047) (1.333,2.100) (1.289,2.150) (1.250,2.191) (1.200,2.240) (1.162,2.275) (1.102,2.325) (1.050,2.365) (1.001,2.400) (0.950,2.434) (0.900,2.464) (0.850,2.492) (0.800,2.518) (0.750,2.543) (0.700,2.565) (0.650,2.585) (0.600,2.604) (0.550,2.621) (0.500,2.636) (0.450,2.650) (0.400,2.662) (0.350,2.672) (0.300,2.682) (0.250,2.689) (0.200,2.696) (0.150,2.700) (0.100,2.704) (0.050,2.706) (0.000,2.707) (-0.050,2.706) (-0.100,2.704) (-0.150,2.700) (-0.200,2.696) (-0.250,2.689) (-0.300,2.682) (-0.350,2.672) (-0.400,2.662) (-0.450,2.650) (-0.500,2.636) (-0.550,2.621) (-0.600,2.604) (-0.650,2.585) (-0.700,2.565) (-0.750,2.543) (-0.800,2.518) (-0.850,2.492) (-0.900,2.464) (-0.950,2.434) (-1.000,2.401) (-1.050,2.365) (-1.100,2.327) (-1.150,2.285) (-1.189,2.250) (-1.241,2.200) (-1.289,2.150) (-1.325,2.109) (-1.373,2.050) (-1.410,2.000) (-1.444,1.950) (-1.475,1.900) (-1.504,1.850) (-1.530,1.800) (-1.554,1.750) (-1.575,1.702) (-1.595,1.650) (-1.613,1.600) (-1.628,1.550) (-1.641,1.500) (-1.652,1.450) (-1.662,1.400) (-1.669,1.350) (-1.674,1.300) (-1.678,1.250) (-1.679,1.200) (-1.679,1.150) (-1.676,1.100) (-1.672,1.050) (-1.665,1.000) (-1.656,0.950) (-1.645,0.900) (-1.632,0.850) (-1.616,0.800) (-1.597,0.750) (-1.576,0.700) (-1.552,0.650) (-1.525,0.601) (-1.493,0.550) (-1.458,0.500) (-1.425,0.458) (-1.375,0.401) (-1.325,0.352) (-1.275,0.308)}
\EQ{(1.275,0.337) (1.314,0.375) (1.350,0.414) (1.400,0.475) (1.436,0.525) (1.467,0.575) (1.495,0.625) (1.519,0.675) (1.540,0.725) (1.559,0.775) (1.574,0.825) (1.587,0.875) (1.598,0.925) (1.606,0.975) (1.612,1.025) (1.616,1.075) (1.618,1.125) (1.618,1.175) (1.616,1.225) (1.611,1.275) (1.605,1.325) (1.597,1.375) (1.586,1.425) (1.574,1.475) (1.559,1.525) (1.543,1.575) (1.524,1.625) (1.503,1.675) (1.479,1.725) (1.454,1.775) (1.425,1.825) (1.394,1.875) (1.360,1.925) (1.325,1.973) (1.283,2.025) (1.250,2.063) (1.216,2.100) (1.175,2.141) (1.125,2.187) (1.081,2.225) (1.025,2.269) (0.983,2.300) (0.925,2.339) (0.875,2.370) (0.825,2.398) (0.775,2.425) (0.725,2.449) (0.675,2.471) (0.625,2.492) (0.575,2.510) (0.525,2.527) (0.475,2.542) (0.425,2.555) (0.375,2.567) (0.325,2.578) (0.275,2.586) (0.225,2.594) (0.175,2.599) (0.125,2.604) (0.075,2.607) (0.025,2.608) (-0.050,2.608) (-0.100,2.605) (-0.150,2.602) (-0.200,2.597) (-0.250,2.590) (-0.300,2.582) (-0.350,2.573) (-0.400,2.562) (-0.450,2.549) (-0.500,2.535) (-0.550,2.519) (-0.600,2.501) (-0.650,2.482) (-0.700,2.460) (-0.750,2.437) (-0.800,2.412) (-0.850,2.384) (-0.900,2.355) (-0.946,2.325) (-1.000,2.288) (-1.050,2.250) (-1.100,2.209) (-1.139,2.175) (-1.191,2.125) (-1.239,2.075) (-1.275,2.034) (-1.323,1.975) (-1.360,1.925) (-1.394,1.875) (-1.425,1.825) (-1.450,1.781) (-1.475,1.734) (-1.500,1.681) (-1.524,1.625) (-1.543,1.575) (-1.559,1.525) (-1.574,1.475) (-1.586,1.425) (-1.597,1.375) (-1.605,1.325) (-1.611,1.275) (-1.616,1.225) (-1.618,1.175) (-1.618,1.125) (-1.616,1.075) (-1.612,1.025) (-1.606,0.975) (-1.598,0.925) (-1.587,0.875) (-1.574,0.825) (-1.559,0.775) (-1.540,0.725) (-1.519,0.675) (-1.495,0.625) (-1.467,0.575) (-1.436,0.525) (-1.400,0.475) (-1.360,0.425) (-1.325,0.387) (-1.289,0.350) (-1.275,0.337)}
\EQ{(1.266,0.359) (1.305,0.400) (1.346,0.450) (1.382,0.500) (1.414,0.550) (1.442,0.600) (1.467,0.650) (1.488,0.700) (1.506,0.750) (1.522,0.800) (1.534,0.850) (1.545,0.900) (1.553,0.950) (1.559,1.000) (1.562,1.050) (1.563,1.100) (1.563,1.150) (1.560,1.200) (1.555,1.250) (1.547,1.300) (1.538,1.350) (1.527,1.400) (1.513,1.450) (1.498,1.500) (1.480,1.550) (1.460,1.600) (1.437,1.650) (1.412,1.700) (1.385,1.750) (1.354,1.800) (1.325,1.844) (1.285,1.900) (1.250,1.944) (1.201,2.000) (1.154,2.050) (1.101,2.100) (1.050,2.145) (1.000,2.185) (0.950,2.222) (0.908,2.250) (0.850,2.287) (0.800,2.315) (0.750,2.342) (0.700,2.366) (0.650,2.388) (0.600,2.409) (0.550,2.427) (0.500,2.443) (0.450,2.458) (0.400,2.471) (0.350,2.483) (0.300,2.493) (0.250,2.501) (0.200,2.508) (0.150,2.513) (0.100,2.517) (0.050,2.519) (0.000,2.520) (-0.050,2.519) (-0.100,2.517) (-0.150,2.513) (-0.200,2.508) (-0.250,2.501) (-0.300,2.493) (-0.350,2.483) (-0.400,2.471) (-0.450,2.458) (-0.500,2.443) (-0.550,2.427) (-0.600,2.409) (-0.650,2.388) (-0.700,2.366) (-0.750,2.342) (-0.800,2.315) (-0.850,2.287) (-0.900,2.256) (-0.945,2.225) (-1.000,2.185) (-1.044,2.150) (-1.100,2.101) (-1.150,2.054) (-1.200,2.001) (-1.245,1.950) (-1.285,1.900) (-1.321,1.850) (-1.350,1.807) (-1.385,1.750) (-1.412,1.700) (-1.437,1.650) (-1.460,1.600) (-1.480,1.550) (-1.498,1.500) (-1.513,1.450) (-1.527,1.400) (-1.538,1.350) (-1.547,1.300) (-1.555,1.250) (-1.560,1.200) (-1.563,1.150) (-1.563,1.100) (-1.562,1.050) (-1.559,1.000) (-1.553,0.950) (-1.545,0.900) (-1.534,0.850) (-1.522,0.800) (-1.506,0.750) (-1.488,0.700) (-1.467,0.650) (-1.442,0.600) (-1.414,0.550) (-1.382,0.500) (-1.350,0.456) (-1.305,0.400) (-1.266,0.359)}
\EQ{(1.251,0.374) (1.295,0.425) (1.325,0.465) (1.365,0.525) (1.393,0.575) (1.418,0.625) (1.440,0.675) (1.458,0.725) (1.473,0.775) (1.486,0.825) (1.496,0.875) (1.504,0.925) (1.510,0.975) (1.513,1.025) (1.514,1.075) (1.512,1.125) (1.509,1.175) (1.503,1.225) (1.495,1.275) (1.486,1.325) (1.473,1.375) (1.459,1.425) (1.442,1.475) (1.423,1.525) (1.402,1.575) (1.378,1.625) (1.352,1.675) (1.325,1.721) (1.290,1.775) (1.255,1.825) (1.216,1.875) (1.175,1.923) (1.127,1.975) (1.075,2.025) (1.025,2.070) (0.975,2.110) (0.925,2.147) (0.875,2.182) (0.825,2.213) (0.775,2.242) (0.725,2.268) (0.675,2.292) (0.625,2.314) (0.575,2.334) (0.525,2.352) (0.475,2.369) (0.425,2.383) (0.375,2.396) (0.325,2.407) (0.275,2.416) (0.225,2.424) (0.175,2.430) (0.125,2.435) (0.075,2.438) (0.025,2.440) (-0.050,2.439) (-0.100,2.437) (-0.150,2.433) (-0.200,2.427) (-0.250,2.420) (-0.300,2.412) (-0.350,2.402) (-0.400,2.390) (-0.450,2.376) (-0.500,2.361) (-0.550,2.344) (-0.599,2.325) (-0.650,2.304) (-0.700,2.280) (-0.750,2.255) (-0.800,2.228) (-0.846,2.200) (-0.900,2.165) (-0.950,2.129) (-1.000,2.090) (-1.048,2.050) (-1.100,2.002) (-1.150,1.951) (-1.195,1.900) (-1.236,1.850) (-1.273,1.800) (-1.307,1.750) (-1.337,1.700) (-1.365,1.650) (-1.390,1.600) (-1.413,1.550) (-1.433,1.500) (-1.450,1.453) (-1.466,1.400) (-1.480,1.350) (-1.491,1.300) (-1.500,1.250) (-1.506,1.200) (-1.511,1.150) (-1.513,1.100) (-1.514,1.050) (-1.512,1.000) (-1.507,0.950) (-1.501,0.900) (-1.492,0.850) (-1.480,0.800) (-1.466,0.750) (-1.449,0.700) (-1.429,0.650) (-1.406,0.600) (-1.379,0.550) (-1.350,0.502) (-1.314,0.450) (-1.275,0.401) (-1.251,0.374)}
\EQ{(1.250,0.405) (1.285,0.450) (1.318,0.500) (1.347,0.550) (1.373,0.600) (1.395,0.650) (1.413,0.700) (1.429,0.750) (1.442,0.800) (1.452,0.850) (1.460,0.900) (1.465,0.950) (1.468,1.000) (1.468,1.050) (1.467,1.100) (1.463,1.150) (1.457,1.200) (1.448,1.250) (1.438,1.300) (1.425,1.350) (1.410,1.400) (1.392,1.450) (1.372,1.500) (1.350,1.550) (1.325,1.600) (1.300,1.645) (1.266,1.700) (1.232,1.750) (1.200,1.793) (1.154,1.850) (1.109,1.900) (1.059,1.950) (1.004,2.000) (0.950,2.044) (0.909,2.075) (0.850,2.116) (0.800,2.147) (0.751,2.175) (0.704,2.200) (0.652,2.225) (0.600,2.248) (0.550,2.268) (0.500,2.285) (0.450,2.301) (0.400,2.315) (0.350,2.328) (0.300,2.338) (0.250,2.347) (0.200,2.354) (0.150,2.360) (0.100,2.364) (0.050,2.366) (0.000,2.367) (-0.050,2.366) (-0.100,2.364) (-0.150,2.360) (-0.200,2.354) (-0.250,2.347) (-0.300,2.338) (-0.350,2.328) (-0.400,2.315) (-0.450,2.301) (-0.500,2.285) (-0.550,2.268) (-0.600,2.248) (-0.650,2.226) (-0.700,2.202) (-0.750,2.176) (-0.795,2.150) (-0.850,2.116) (-0.900,2.081) (-0.942,2.050) (-1.000,2.003) (-1.050,1.958) (-1.100,1.909) (-1.150,1.854) (-1.195,1.800) (-1.225,1.760) (-1.266,1.700) (-1.297,1.650) (-1.325,1.600) (-1.350,1.550) (-1.372,1.500) (-1.392,1.450) (-1.410,1.400) (-1.425,1.350) (-1.438,1.300) (-1.448,1.250) (-1.457,1.200) (-1.463,1.150) (-1.467,1.100) (-1.468,1.050) (-1.468,1.000) (-1.465,0.950) (-1.460,0.900) (-1.452,0.850) (-1.442,0.800) (-1.429,0.750) (-1.413,0.700) (-1.395,0.650) (-1.373,0.600) (-1.350,0.555) (-1.325,0.512) (-1.285,0.450) (-1.250,0.405)}
\EQ{(1.244,0.431) (1.274,0.475) (1.300,0.518) (1.325,0.565) (1.350,0.619) (1.372,0.675) (1.388,0.725) (1.400,0.775) (1.411,0.825) (1.419,0.875) (1.424,0.925) (1.427,0.975) (1.427,1.025) (1.425,1.075) (1.421,1.125) (1.414,1.175) (1.406,1.225) (1.395,1.275) (1.381,1.325) (1.365,1.375) (1.347,1.425) (1.326,1.475) (1.303,1.525) (1.277,1.575) (1.250,1.621) (1.215,1.675) (1.180,1.725) (1.141,1.775) (1.100,1.822) (1.050,1.874) (1.000,1.922) (0.950,1.965) (0.906,2.000) (0.850,2.040) (0.800,2.073) (0.755,2.100) (0.709,2.125) (0.660,2.150) (0.605,2.175) (0.550,2.198) (0.500,2.216) (0.450,2.233) (0.400,2.247) (0.350,2.260) (0.300,2.271) (0.250,2.280) (0.200,2.287) (0.150,2.293) (0.100,2.297) (0.050,2.299) (0.000,2.300) (-0.050,2.299) (-0.100,2.297) (-0.150,2.293) (-0.200,2.287) (-0.250,2.280) (-0.300,2.271) (-0.350,2.260) (-0.400,2.247) (-0.450,2.233) (-0.500,2.216) (-0.550,2.198) (-0.600,2.177) (-0.650,2.155) (-0.700,2.130) (-0.750,2.103) (-0.796,2.075) (-0.850,2.040) (-0.900,2.004) (-0.950,1.965) (-0.997,1.925) (-1.049,1.875) (-1.097,1.825) (-1.141,1.775) (-1.175,1.731) (-1.215,1.675) (-1.248,1.625) (-1.275,1.578) (-1.300,1.531) (-1.325,1.478) (-1.347,1.425) (-1.365,1.375) (-1.381,1.325) (-1.395,1.275) (-1.406,1.225) (-1.414,1.175) (-1.421,1.125) (-1.425,1.075) (-1.427,1.025) (-1.427,0.975) (-1.424,0.925) (-1.419,0.875) (-1.411,0.825) (-1.400,0.775) (-1.388,0.725) (-1.372,0.675) (-1.352,0.625) (-1.330,0.575) (-1.304,0.525) (-1.275,0.477) (-1.244,0.431)}
\EQ{(1.229,0.446) (1.263,0.500) (1.290,0.550) (1.313,0.600) (1.333,0.650) (1.349,0.700) (1.362,0.750) (1.373,0.800) (1.381,0.850) (1.386,0.900) (1.389,0.950) (1.389,1.000) (1.387,1.050) (1.383,1.100) (1.376,1.150) (1.367,1.200) (1.355,1.250) (1.341,1.300) (1.325,1.349) (1.306,1.400) (1.284,1.450) (1.260,1.500) (1.233,1.550) (1.203,1.600) (1.175,1.642) (1.132,1.700) (1.100,1.740) (1.050,1.796) (1.000,1.846) (0.950,1.891) (0.909,1.925) (0.850,1.970) (0.806,2.000) (0.750,2.035) (0.700,2.063) (0.650,2.089) (0.600,2.112) (0.550,2.133) (0.500,2.152) (0.450,2.169) (0.400,2.184) (0.350,2.197) (0.300,2.208) (0.250,2.218) (0.200,2.225) (0.150,2.231) (0.100,2.236) (0.050,2.238) (0.000,2.239) (-0.050,2.238) (-0.100,2.236) (-0.150,2.231) (-0.200,2.225) (-0.250,2.218) (-0.300,2.208) (-0.350,2.197) (-0.400,2.184) (-0.450,2.169) (-0.500,2.152) (-0.550,2.133) (-0.600,2.112) (-0.650,2.089) (-0.700,2.063) (-0.750,2.035) (-0.800,2.004) (-0.843,1.975) (-0.900,1.932) (-0.940,1.900) (-0.996,1.850) (-1.046,1.800) (-1.091,1.750) (-1.125,1.709) (-1.169,1.650) (-1.200,1.604) (-1.233,1.550) (-1.260,1.500) (-1.284,1.450) (-1.306,1.400) (-1.325,1.350) (-1.341,1.300) (-1.355,1.250) (-1.367,1.200) (-1.376,1.150) (-1.383,1.100) (-1.387,1.050) (-1.389,1.000) (-1.389,0.950) (-1.386,0.900) (-1.381,0.850) (-1.373,0.800) (-1.362,0.750) (-1.349,0.700) (-1.333,0.650) (-1.313,0.600) (-1.290,0.550) (-1.263,0.500) (-1.232,0.450) (-1.229,0.446)}
\EQ{(1.224,0.476) (1.250,0.522) (1.275,0.574) (1.296,0.625) (1.313,0.675) (1.326,0.725) (1.337,0.775) (1.345,0.825) (1.350,0.875) (1.353,0.925) (1.354,0.975) (1.352,1.025) (1.347,1.075) (1.340,1.125) (1.331,1.175) (1.319,1.225) (1.305,1.275) (1.288,1.325) (1.269,1.375) (1.246,1.425) (1.221,1.475) (1.194,1.525) (1.162,1.575) (1.128,1.625) (1.089,1.675) (1.050,1.721) (1.000,1.774) (0.950,1.822) (0.900,1.865) (0.856,1.900) (0.800,1.940) (0.750,1.972) (0.702,2.000) (0.656,2.025) (0.604,2.050) (0.550,2.074) (0.500,2.093) (0.450,2.111) (0.400,2.126) (0.350,2.139) (0.300,2.151) (0.250,2.161) (0.200,2.168) (0.150,2.174) (0.100,2.179) (0.050,2.181) (0.000,2.182) (-0.050,2.181) (-0.100,2.179) (-0.150,2.174) (-0.200,2.168) (-0.250,2.161) (-0.300,2.151) (-0.350,2.139) (-0.400,2.126) (-0.450,2.111) (-0.500,2.093) (-0.547,2.075) (-0.600,2.052) (-0.650,2.028) (-0.700,2.001) (-0.745,1.975) (-0.800,1.940) (-0.850,1.904) (-0.900,1.865) (-0.947,1.825) (-0.999,1.775) (-1.047,1.725) (-1.089,1.675) (-1.125,1.629) (-1.162,1.575) (-1.194,1.525) (-1.221,1.475) (-1.246,1.425) (-1.269,1.375) (-1.288,1.325) (-1.305,1.275) (-1.319,1.225) (-1.331,1.175) (-1.340,1.125) (-1.347,1.075) (-1.352,1.025) (-1.354,0.975) (-1.353,0.925) (-1.350,0.875) (-1.345,0.825) (-1.337,0.775) (-1.326,0.725) (-1.313,0.675) (-1.296,0.625) (-1.275,0.575) (-1.252,0.525) (-1.225,0.477)}

    \draw[mmaEQ] (-10,0) -- (10,0);
    % Dipolo puntual - líneas de campo
% Generado por scripts/generar_datos.py (no editar a mano).
\LC{(1.082,0.489) (1.149,0.455) (1.212,0.414) (1.269,0.365) (1.319,0.309) (1.361,0.248) (1.394,0.180) (1.417,0.109) (1.428,0.035) (1.428,-0.040) (1.415,-0.114) (1.392,-0.185) (1.358,-0.252) (1.316,-0.314) (1.265,-0.369) (1.207,-0.417) (1.144,-0.458) (1.077,-0.491)}
\LC{(0.453,1.106) (0.495,1.168) (0.539,1.229) (0.583,1.290) (0.629,1.349) (0.676,1.408) (0.724,1.465) (0.773,1.522) (0.823,1.578) (0.873,1.633) (0.925,1.688) (0.978,1.741) (1.031,1.794) (1.085,1.846) (1.141,1.896) (1.197,1.946) (1.253,1.996) (1.311,2.044) (1.369,2.091) (1.428,2.137) (1.487,2.183) (1.548,2.228) (1.609,2.271) (1.670,2.314) (1.733,2.356) (1.795,2.397) (1.859,2.437) (1.923,2.476) (1.988,2.514) (2.053,2.551) (2.118,2.587) (2.185,2.622) (2.251,2.657) (2.319,2.690) (2.386,2.722) (2.454,2.753) (2.523,2.784) (2.592,2.813) (2.661,2.841) (2.731,2.869) (2.802,2.895) (2.872,2.920) (2.943,2.945) (3.014,2.968) (3.086,2.990) (3.158,3.011) (3.230,3.031) (3.303,3.050) (3.376,3.068) (3.449,3.085) (3.522,3.101) (3.596,3.116) (3.669,3.130) (3.743,3.142) (3.817,3.154) (3.892,3.164) (3.966,3.173) (4.041,3.181) (4.115,3.188) (4.190,3.194) (4.265,3.199) (4.340,3.203) (4.415,3.205) (4.490,3.206) (4.565,3.206) (4.640,3.205) (4.715,3.203) (4.790,3.200) (4.864,3.195) (4.939,3.189) (5.014,3.182) (5.088,3.174) (5.163,3.165) (5.237,3.154) (5.311,3.142) (5.385,3.129) (5.459,3.115) (5.532,3.099) (5.605,3.082) (5.678,3.064) (5.750,3.045) (5.822,3.024) (5.894,3.002) (5.965,2.979) (6.036,2.955) (6.107,2.929) (6.177,2.902) (6.246,2.874) (6.315,2.844) (6.384,2.813) (6.451,2.781) (6.518,2.748) (6.585,2.713) (6.651,2.677) (6.716,2.640) (6.780,2.601) (6.844,2.561) (6.906,2.520) (6.968,2.478) (7.029,2.434) (7.089,2.389) (7.148,2.343) (7.207,2.296) (7.264,2.247) (7.320,2.197) (7.374,2.146) (7.428,2.094) (7.481,2.040) (7.532,1.985) (7.582,1.930) (7.631,1.873) (7.679,1.815) (7.725,1.755) (7.769,1.695) (7.813,1.634) (7.854,1.572) (7.894,1.508) (7.933,1.444) (7.970,1.379) (8.006,1.313) (8.039,1.246) (8.071,1.178) (8.102,1.109) (8.130,1.040) (8.157,0.970) (8.182,0.899) (8.205,0.828) (8.226,0.756) (8.246,0.683) (8.263,0.610) (8.278,0.537) (8.292,0.463) (8.303,0.389) (8.313,0.315) (8.320,0.240) (8.326,0.165) (8.329,0.090) (8.331,0.015) (8.330,-0.060) (8.327,-0.134) (8.323,-0.209) (8.316,-0.284) (8.307,-0.359) (8.297,-0.433) (8.284,-0.507) (8.270,-0.580) (8.253,-0.653) (8.234,-0.726) (8.214,-0.798) (8.192,-0.870) (8.168,-0.941) (8.142,-1.011) (8.114,-1.081) (8.084,-1.150) (8.053,-1.218) (8.020,-1.285) (7.985,-1.352) (7.949,-1.417) (7.911,-1.482) (7.871,-1.546) (7.830,-1.608) (7.787,-1.670) (7.743,-1.731) (7.698,-1.790) (7.651,-1.849) (7.603,-1.906) (7.553,-1.963) (7.502,-2.018) (7.450,-2.072) (7.397,-2.124) (7.342,-2.176) (7.287,-2.226) (7.230,-2.276) (7.173,-2.324) (7.114,-2.370) (7.054,-2.416) (6.994,-2.460) (6.932,-2.503) (6.870,-2.544) (6.806,-2.585) (6.742,-2.624) (6.678,-2.662) (6.612,-2.698) (6.546,-2.733) (6.479,-2.767) (6.412,-2.800) (6.343,-2.831) (6.275,-2.862) (6.206,-2.890) (6.136,-2.918) (6.066,-2.944) (5.995,-2.969) (5.924,-2.993) (5.852,-3.015) (5.780,-3.036) (5.708,-3.056) (5.635,-3.075) (5.562,-3.092) (5.489,-3.108) (5.415,-3.123) (5.342,-3.137) (5.268,-3.149) (5.194,-3.160) (5.119,-3.170) (5.045,-3.179) (4.970,-3.187) (4.895,-3.193) (4.821,-3.198) (4.746,-3.202) (4.671,-3.205) (4.596,-3.206) (4.521,-3.206) (4.446,-3.206) (4.371,-3.204) (4.296,-3.201) (4.221,-3.196) (4.146,-3.191) (4.071,-3.184) (3.997,-3.177) (3.922,-3.168) (3.848,-3.158) (3.774,-3.147) (3.700,-3.135) (3.626,-3.122) (3.552,-3.107) (3.479,-3.092) (3.406,-3.075) (3.333,-3.058) (3.260,-3.039) (3.188,-3.020) (3.116,-2.999) (3.044,-2.977) (2.973,-2.954) (2.901,-2.930) (2.831,-2.906) (2.760,-2.880) (2.690,-2.853) (2.621,-2.825) (2.551,-2.796) (2.483,-2.766) (2.414,-2.735) (2.346,-2.703) (2.279,-2.670) (2.212,-2.637) (2.146,-2.602) (2.080,-2.566) (2.014,-2.529) (1.950,-2.492) (1.885,-2.453) (1.822,-2.413) (1.758,-2.373) (1.696,-2.331) (1.634,-2.289) (1.573,-2.246) (1.512,-2.201) (1.452,-2.156) (1.393,-2.110) (1.335,-2.063) (1.277,-2.016) (1.220,-1.967) (1.164,-1.917) (1.108,-1.867) (1.053,-1.815) (1.000,-1.763) (0.947,-1.710) (0.895,-1.656) (0.843,-1.601) (0.793,-1.545) (0.744,-1.489) (0.696,-1.432) (0.648,-1.373) (0.602,-1.314) (0.557,-1.254) (0.513,-1.194) (0.470,-1.132) (0.442,-1.091)}
\LC{(-0.453,1.106) (-0.495,1.168) (-0.539,1.229) (-0.583,1.290) (-0.629,1.349) (-0.676,1.408) (-0.724,1.465) (-0.773,1.522) (-0.823,1.578) (-0.873,1.633) (-0.925,1.688) (-0.978,1.741) (-1.031,1.794) (-1.085,1.846) (-1.141,1.896) (-1.197,1.946) (-1.253,1.996) (-1.311,2.044) (-1.369,2.091) (-1.428,2.137) (-1.487,2.183) (-1.548,2.228) (-1.609,2.271) (-1.670,2.314) (-1.733,2.356) (-1.795,2.397) (-1.859,2.437) (-1.923,2.476) (-1.988,2.514) (-2.053,2.551) (-2.118,2.587) (-2.185,2.622) (-2.251,2.657) (-2.319,2.690) (-2.386,2.722) (-2.454,2.753) (-2.523,2.784) (-2.592,2.813) (-2.661,2.841) (-2.731,2.869) (-2.802,2.895) (-2.872,2.920) (-2.943,2.945) (-3.014,2.968) (-3.086,2.990) (-3.158,3.011) (-3.230,3.031) (-3.303,3.050) (-3.376,3.068) (-3.449,3.085) (-3.522,3.101) (-3.596,3.116) (-3.669,3.130) (-3.743,3.142) (-3.817,3.154) (-3.892,3.164) (-3.966,3.173) (-4.041,3.181) (-4.115,3.188) (-4.190,3.194) (-4.265,3.199) (-4.340,3.203) (-4.415,3.205) (-4.490,3.206) (-4.565,3.206) (-4.640,3.205) (-4.715,3.203) (-4.790,3.200) (-4.864,3.195) (-4.939,3.189) (-5.014,3.182) (-5.088,3.174) (-5.163,3.165) (-5.237,3.154) (-5.311,3.142) (-5.385,3.129) (-5.459,3.115) (-5.532,3.099) (-5.605,3.082) (-5.678,3.064) (-5.750,3.045) (-5.822,3.024) (-5.894,3.002) (-5.965,2.979) (-6.036,2.955) (-6.107,2.929) (-6.177,2.902) (-6.246,2.874) (-6.315,2.844) (-6.384,2.813) (-6.451,2.781) (-6.518,2.748) (-6.585,2.713) (-6.651,2.677) (-6.716,2.640) (-6.780,2.601) (-6.844,2.561) (-6.906,2.520) (-6.968,2.478) (-7.029,2.434) (-7.089,2.389) (-7.148,2.343) (-7.207,2.296) (-7.264,2.247) (-7.320,2.197) (-7.374,2.146) (-7.428,2.094) (-7.481,2.040) (-7.532,1.985) (-7.582,1.930) (-7.631,1.873) (-7.679,1.815) (-7.725,1.755) (-7.769,1.695) (-7.813,1.634) (-7.854,1.572) (-7.894,1.508) (-7.933,1.444) (-7.970,1.379) (-8.006,1.313) (-8.039,1.246) (-8.071,1.178) (-8.102,1.109) (-8.130,1.040) (-8.157,0.970) (-8.182,0.899) (-8.205,0.828) (-8.226,0.756) (-8.246,0.683) (-8.263,0.610) (-8.278,0.537) (-8.292,0.463) (-8.303,0.389) (-8.313,0.315) (-8.320,0.240) (-8.326,0.165) (-8.329,0.090) (-8.331,0.015) (-8.330,-0.060) (-8.327,-0.134) (-8.323,-0.209) (-8.316,-0.284) (-8.307,-0.359) (-8.297,-0.433) (-8.284,-0.507) (-8.270,-0.580) (-8.253,-0.653) (-8.234,-0.726) (-8.214,-0.798) (-8.192,-0.870) (-8.168,-0.941) (-8.142,-1.011) (-8.114,-1.081) (-8.084,-1.150) (-8.053,-1.218) (-8.020,-1.285) (-7.985,-1.352) (-7.949,-1.417) (-7.911,-1.482) (-7.871,-1.546) (-7.830,-1.608) (-7.787,-1.670) (-7.743,-1.731) (-7.698,-1.790) (-7.651,-1.849) (-7.603,-1.906) (-7.553,-1.963) (-7.502,-2.018) (-7.450,-2.072) (-7.397,-2.124) (-7.342,-2.176) (-7.287,-2.226) (-7.230,-2.276) (-7.173,-2.324) (-7.114,-2.370) (-7.054,-2.416) (-6.994,-2.460) (-6.932,-2.503) (-6.870,-2.544) (-6.806,-2.585) (-6.742,-2.624) (-6.678,-2.662) (-6.612,-2.698) (-6.546,-2.733) (-6.479,-2.767) (-6.412,-2.800) (-6.343,-2.831) (-6.275,-2.862) (-6.206,-2.890) (-6.136,-2.918) (-6.066,-2.944) (-5.995,-2.969) (-5.924,-2.993) (-5.852,-3.015) (-5.780,-3.036) (-5.708,-3.056) (-5.635,-3.075) (-5.562,-3.092) (-5.489,-3.108) (-5.415,-3.123) (-5.342,-3.137) (-5.268,-3.149) (-5.194,-3.160) (-5.119,-3.170) (-5.045,-3.179) (-4.970,-3.187) (-4.895,-3.193) (-4.821,-3.198) (-4.746,-3.202) (-4.671,-3.205) (-4.596,-3.206) (-4.521,-3.206) (-4.446,-3.206) (-4.371,-3.204) (-4.296,-3.201) (-4.221,-3.196) (-4.146,-3.191) (-4.071,-3.184) (-3.997,-3.177) (-3.922,-3.168) (-3.848,-3.158) (-3.774,-3.147) (-3.700,-3.135) (-3.626,-3.122) (-3.552,-3.107) (-3.479,-3.092) (-3.406,-3.075) (-3.333,-3.058) (-3.260,-3.039) (-3.188,-3.020) (-3.116,-2.999) (-3.044,-2.977) (-2.973,-2.954) (-2.901,-2.930) (-2.831,-2.906) (-2.760,-2.880) (-2.690,-2.853) (-2.621,-2.825) (-2.551,-2.796) (-2.483,-2.766) (-2.414,-2.735) (-2.346,-2.703) (-2.279,-2.670) (-2.212,-2.637) (-2.146,-2.602) (-2.080,-2.566) (-2.014,-2.529) (-1.950,-2.492) (-1.885,-2.453) (-1.822,-2.413) (-1.758,-2.373) (-1.696,-2.331) (-1.634,-2.289) (-1.573,-2.246) (-1.512,-2.201) (-1.452,-2.156) (-1.393,-2.110) (-1.335,-2.063) (-1.277,-2.016) (-1.220,-1.967) (-1.164,-1.917) (-1.108,-1.867) (-1.053,-1.815) (-1.000,-1.763) (-0.947,-1.710) (-0.895,-1.656) (-0.843,-1.601) (-0.793,-1.545) (-0.744,-1.489) (-0.696,-1.432) (-0.648,-1.373) (-0.602,-1.314) (-0.557,-1.254) (-0.513,-1.194) (-0.470,-1.132) (-0.442,-1.091)}
\LC{(-1.082,0.489) (-1.149,0.455) (-1.212,0.414) (-1.269,0.365) (-1.319,0.309) (-1.361,0.248) (-1.394,0.180) (-1.417,0.109) (-1.428,0.035) (-1.428,-0.040) (-1.415,-0.114) (-1.392,-0.185) (-1.358,-0.252) (-1.316,-0.314) (-1.265,-0.369) (-1.207,-0.417) (-1.144,-0.458) (-1.077,-0.491)}
\LC{(4.893,-10.000) (4.845,-9.951) (4.797,-9.893) (4.749,-9.835) (4.702,-9.777) (4.655,-9.719) (4.607,-9.661) (4.560,-9.602) (4.514,-9.544) (4.467,-9.485) (4.420,-9.426) (4.374,-9.368) (4.328,-9.309) (4.281,-9.249) (4.235,-9.190) (4.190,-9.131) (4.144,-9.071) (4.099,-9.012) (4.053,-8.952) (4.008,-8.892) (3.963,-8.832) (3.918,-8.772) (3.873,-8.712) (3.829,-8.651) (3.785,-8.591) (3.740,-8.530) (3.696,-8.470) (3.652,-8.409) (3.609,-8.348) (3.565,-8.287) (3.522,-8.226) (3.479,-8.164) (3.435,-8.103) (3.393,-8.041) (3.350,-7.980) (3.307,-7.918) (3.265,-7.856) (3.223,-7.794) (3.181,-7.732) (3.139,-7.670) (3.097,-7.607) (3.056,-7.545) (3.015,-7.482) (2.973,-7.419) (2.932,-7.357) (2.892,-7.294) (2.851,-7.231) (2.811,-7.167) (2.771,-7.104) (2.730,-7.041) (2.691,-6.977) (2.651,-6.913) (2.612,-6.850) (2.572,-6.786) (2.533,-6.722) (2.494,-6.658) (2.456,-6.593) (2.417,-6.529) (2.379,-6.465) (2.341,-6.400) (2.303,-6.335) (2.265,-6.270) (2.228,-6.205) (2.190,-6.140) (2.153,-6.075) (2.116,-6.010) (2.080,-5.945) (2.043,-5.879) (2.007,-5.813) (1.971,-5.748) (1.935,-5.682) (1.899,-5.616) (1.864,-5.550) (1.829,-5.483) (1.794,-5.417) (1.759,-5.351) (1.724,-5.284) (1.690,-5.217) (1.656,-5.151) (1.622,-5.084) (1.588,-5.017) (1.555,-4.950) (1.522,-4.882) (1.489,-4.815) (1.456,-4.747) (1.423,-4.680) (1.391,-4.612) (1.359,-4.544) (1.327,-4.476) (1.296,-4.408) (1.265,-4.340) (1.234,-4.272) (1.203,-4.203) (1.172,-4.135) (1.142,-4.066) (1.112,-3.998) (1.082,-3.929) (1.053,-3.860) (1.024,-3.791) (0.995,-3.721) (0.966,-3.652) (0.938,-3.583) (0.910,-3.513) (0.882,-3.444) (0.854,-3.374) (0.827,-3.304) (0.800,-3.234) (0.773,-3.164) (0.747,-3.094) (0.721,-3.023) (0.695,-2.953) (0.670,-2.882) (0.645,-2.812) (0.620,-2.741) (0.595,-2.670) (0.571,-2.599) (0.547,-2.528) (0.524,-2.457) (0.501,-2.385) (0.478,-2.314) (0.455,-2.242) (0.433,-2.171) (0.412,-2.099) (0.390,-2.027) (0.369,-1.955) (0.349,-1.883) (0.328,-1.811) (0.309,-1.738) (0.289,-1.666) (0.283,-1.642)}
\LC{(-5.042,-10.000) (-4.993,-9.951) (-4.944,-9.894) (-4.895,-9.837) (-4.847,-9.780) (-4.798,-9.722) (-4.750,-9.665) (-4.702,-9.607) (-4.654,-9.550) (-4.606,-9.492) (-4.558,-9.434) (-4.511,-9.376) (-4.463,-9.318) (-4.416,-9.259) (-4.369,-9.201) (-4.322,-9.142) (-4.276,-9.084) (-4.229,-9.025) (-4.183,-8.966) (-4.136,-8.907) (-4.090,-8.848) (-4.044,-8.789) (-3.999,-8.729) (-3.953,-8.670) (-3.908,-8.610) (-3.862,-8.550) (-3.817,-8.490) (-3.772,-8.430) (-3.728,-8.370) (-3.683,-8.310) (-3.639,-8.249) (-3.594,-8.189) (-3.550,-8.128) (-3.506,-8.067) (-3.463,-8.006) (-3.419,-7.945) (-3.376,-7.884) (-3.333,-7.823) (-3.289,-7.761) (-3.247,-7.700) (-3.204,-7.638) (-3.161,-7.576) (-3.119,-7.514) (-3.077,-7.452) (-3.035,-7.390) (-2.993,-7.328) (-2.952,-7.265) (-2.910,-7.203) (-2.869,-7.140) (-2.828,-7.077) (-2.787,-7.015) (-2.747,-6.951) (-2.706,-6.888) (-2.666,-6.825) (-2.626,-6.762) (-2.586,-6.698) (-2.546,-6.634) (-2.507,-6.571) (-2.467,-6.507) (-2.428,-6.443) (-2.389,-6.379) (-2.351,-6.314) (-2.312,-6.250) (-2.274,-6.186) (-2.236,-6.121) (-2.198,-6.056) (-2.160,-5.991) (-2.123,-5.926) (-2.086,-5.861) (-2.049,-5.796) (-2.012,-5.731) (-1.975,-5.665) (-1.939,-5.600) (-1.903,-5.534) (-1.867,-5.468) (-1.831,-5.402) (-1.796,-5.336) (-1.760,-5.270) (-1.725,-5.204) (-1.690,-5.137) (-1.656,-5.071) (-1.621,-5.004) (-1.587,-4.937) (-1.553,-4.870) (-1.520,-4.803) (-1.486,-4.736) (-1.453,-4.669) (-1.420,-4.601) (-1.388,-4.534) (-1.355,-4.466) (-1.323,-4.399) (-1.291,-4.331) (-1.260,-4.263) (-1.228,-4.195) (-1.197,-4.126) (-1.166,-4.058) (-1.135,-3.990) (-1.105,-3.921) (-1.075,-3.852) (-1.045,-3.783) (-1.016,-3.714) (-0.986,-3.645) (-0.957,-3.576) (-0.929,-3.507) (-0.900,-3.438) (-0.872,-3.368) (-0.844,-3.298) (-0.817,-3.229) (-0.790,-3.159) (-0.763,-3.089) (-0.736,-3.019) (-0.710,-2.948) (-0.684,-2.878) (-0.658,-2.808) (-0.633,-2.737) (-0.608,-2.666) (-0.583,-2.595) (-0.559,-2.524) (-0.535,-2.453) (-0.511,-2.382) (-0.488,-2.311) (-0.465,-2.240) (-0.442,-2.168) (-0.420,-2.096) (-0.398,-2.025) (-0.377,-1.953) (-0.356,-1.881) (-0.335,-1.809) (-0.315,-1.736) (-0.295,-1.664)}
\LC{(10.000,-6.289) (9.931,-6.297) (9.856,-6.305) (9.782,-6.311) (9.707,-6.317) (9.632,-6.323) (9.557,-6.328) (9.482,-6.332) (9.407,-6.336) (9.332,-6.339) (9.257,-6.342) (9.182,-6.344) (9.107,-6.345) (9.032,-6.346) (8.957,-6.346) (8.882,-6.346) (8.808,-6.344) (8.733,-6.343) (8.658,-6.341) (8.583,-6.338) (8.508,-6.335) (8.433,-6.331) (8.358,-6.326) (8.283,-6.321) (8.208,-6.315) (8.134,-6.309) (8.059,-6.302) (7.984,-6.295) (7.910,-6.287) (7.835,-6.278) (7.761,-6.269) (7.686,-6.260) (7.612,-6.249) (7.538,-6.239) (7.464,-6.227) (7.390,-6.215) (7.316,-6.203) (7.242,-6.190) (7.168,-6.176) (7.094,-6.162) (7.021,-6.148) (6.947,-6.133) (6.874,-6.117) (6.801,-6.100) (6.728,-6.084) (6.655,-6.066) (6.582,-6.048) (6.509,-6.030) (6.437,-6.011) (6.364,-5.991) (6.292,-5.971) (6.220,-5.951) (6.148,-5.930) (6.076,-5.908) (6.004,-5.886) (5.933,-5.863) (5.862,-5.840) (5.791,-5.816) (5.720,-5.792) (5.649,-5.767) (5.578,-5.742) (5.508,-5.716) (5.438,-5.690) (5.367,-5.663) (5.298,-5.636) (5.228,-5.608) (5.159,-5.579) (5.089,-5.550) (5.020,-5.521) (4.952,-5.491) (4.883,-5.461) (4.815,-5.430) (4.747,-5.399) (4.679,-5.367) (4.611,-5.334) (4.544,-5.301) (4.476,-5.268) (4.410,-5.234) (4.343,-5.200) (4.276,-5.165) (4.210,-5.130) (4.144,-5.094) (4.079,-5.058) (4.013,-5.021) (3.948,-4.984) (3.883,-4.946) (3.819,-4.908) (3.755,-4.869) (3.691,-4.830) (3.627,-4.790) (3.564,-4.750) (3.501,-4.709) (3.438,-4.668) (3.375,-4.627) (3.313,-4.585) (3.251,-4.542) (3.190,-4.499) (3.129,-4.456) (3.068,-4.412) (3.008,-4.368) (2.947,-4.323) (2.888,-4.278) (2.828,-4.232) (2.769,-4.186) (2.710,-4.139) (2.652,-4.092) (2.594,-4.045) (2.536,-3.997) (2.479,-3.948) (2.422,-3.899) (2.366,-3.850) (2.309,-3.800) (2.254,-3.750) (2.198,-3.699) (2.144,-3.648) (2.089,-3.597) (2.035,-3.545) (1.981,-3.492) (1.928,-3.439) (1.875,-3.386) (1.823,-3.332) (1.771,-3.278) (1.720,-3.223) (1.669,-3.168) (1.618,-3.113) (1.568,-3.057) (1.519,-3.001) (1.470,-2.944) (1.421,-2.887) (1.373,-2.829) (1.326,-2.771) (1.279,-2.713) (1.233,-2.654) (1.187,-2.594) (1.141,-2.535) (1.097,-2.474) (1.052,-2.414)}
\LC{(-2.969,2.475) (-3.043,2.485) (-3.118,2.495) (-3.192,2.503) (-3.267,2.509) (-3.342,2.515) (-3.417,2.518) (-3.492,2.520) (-3.567,2.521) (-3.642,2.520) (-3.717,2.518) (-3.792,2.514) (-3.866,2.509) (-3.941,2.502) (-4.016,2.494) (-4.090,2.484) (-4.164,2.472) (-4.238,2.459) (-4.312,2.445) (-4.385,2.428) (-4.458,2.411) (-4.530,2.391) (-4.602,2.370) (-4.674,2.348) (-4.745,2.323) (-4.815,2.298) (-4.885,2.270) (-4.954,2.241) (-5.022,2.210) (-5.090,2.178) (-5.157,2.144) (-5.223,2.108) (-5.288,2.071) (-5.352,2.032) (-5.415,1.991) (-5.477,1.949) (-5.538,1.905) (-5.598,1.860) (-5.656,1.813) (-5.713,1.765) (-5.769,1.714) (-5.823,1.663) (-5.876,1.610) (-5.928,1.555) (-5.977,1.499) (-6.025,1.441) (-6.072,1.382) (-6.116,1.322) (-6.159,1.260) (-6.200,1.197) (-6.239,1.133) (-6.275,1.068) (-6.310,1.001) (-6.343,0.934) (-6.373,0.865) (-6.401,0.796) (-6.427,0.725) (-6.450,0.654) (-6.471,0.582) (-6.490,0.509) (-6.506,0.436) (-6.520,0.362) (-6.531,0.288) (-6.539,0.214) (-6.545,0.139) (-6.549,0.064) (-6.550,-0.011) (-6.548,-0.086) (-6.544,-0.161) (-6.537,-0.236) (-6.528,-0.310) (-6.516,-0.384) (-6.502,-0.458) (-6.485,-0.531) (-6.465,-0.603) (-6.444,-0.675) (-6.419,-0.746) (-6.393,-0.816) (-6.364,-0.885) (-6.333,-0.954) (-6.300,-1.021) (-6.265,-1.087) (-6.227,-1.152) (-6.188,-1.216) (-6.147,-1.279) (-6.103,-1.340) (-6.058,-1.400) (-6.012,-1.458) (-5.963,-1.515) (-5.913,-1.571) (-5.861,-1.625) (-5.808,-1.678) (-5.753,-1.729) (-5.697,-1.779) (-5.639,-1.827) (-5.580,-1.874) (-5.520,-1.918) (-5.459,-1.962) (-5.397,-2.004) (-5.333,-2.044) (-5.269,-2.082) (-5.203,-2.119) (-5.137,-2.154) (-5.070,-2.188) (-5.002,-2.219) (-4.934,-2.250) (-4.864,-2.278) (-4.794,-2.305) (-4.724,-2.331) (-4.653,-2.355) (-4.581,-2.377) (-4.509,-2.397) (-4.436,-2.416) (-4.363,-2.433) (-4.290,-2.449) (-4.216,-2.463) (-4.142,-2.476) (-4.068,-2.487) (-3.994,-2.496) (-3.919,-2.504) (-3.845,-2.511) (-3.770,-2.515) (-3.695,-2.519) (-3.620,-2.521) (-3.545,-2.521) (-3.470,-2.520) (-3.395,-2.517) (-3.320,-2.513) (-3.245,-2.508) (-3.170,-2.501) (-3.096,-2.492) (-3.022,-2.482) (-2.972,-2.475)}
\LC{(0.389,1.769) (0.413,1.840) (0.437,1.911) (0.463,1.982) (0.488,2.052) (0.514,2.123) (0.541,2.193) (0.568,2.263) (0.595,2.333) (0.623,2.402) (0.651,2.472) (0.680,2.541) (0.709,2.610) (0.738,2.679) (0.768,2.748) (0.798,2.817) (0.829,2.885) (0.860,2.954) (0.891,3.022) (0.923,3.090) (0.955,3.158) (0.987,3.225) (1.020,3.293) (1.053,3.360) (1.086,3.427) (1.120,3.494) (1.154,3.561) (1.189,3.628) (1.223,3.694) (1.258,3.760) (1.294,3.827) (1.329,3.893) (1.365,3.958) (1.402,4.024) (1.438,4.089) (1.475,4.155) (1.512,4.220) (1.550,4.285) (1.588,4.350) (1.626,4.414) (1.664,4.479) (1.703,4.543) (1.742,4.607) (1.781,4.671) (1.821,4.735) (1.860,4.798) (1.900,4.862) (1.941,4.925) (1.981,4.988) (2.022,5.051) (2.063,5.113) (2.105,5.176) (2.146,5.238) (2.188,5.301) (2.230,5.363) (2.273,5.424) (2.316,5.486) (2.358,5.548) (2.402,5.609) (2.445,5.670) (2.489,5.731) (2.533,5.792) (2.577,5.853) (2.621,5.913) (2.666,5.973) (2.711,6.033) (2.756,6.093) (2.801,6.153) (2.847,6.213) (2.892,6.272) (2.938,6.331) (2.985,6.390) (3.031,6.449) (3.078,6.508) (3.125,6.566) (3.172,6.625) (3.219,6.683) (3.267,6.741) (3.315,6.799) (3.363,6.856) (3.411,6.914) (3.459,6.971) (3.508,7.028) (3.557,7.085) (3.606,7.142) (3.655,7.198) (3.705,7.255) (3.754,7.311) (3.804,7.367) (3.854,7.423) (3.905,7.478) (3.955,7.534) (4.006,7.589) (4.057,7.644) (4.108,7.699) (4.159,7.754) (4.210,7.808) (4.262,7.863) (4.314,7.917) (4.366,7.971) (4.418,8.025) (4.471,8.078) (4.523,8.132) (4.576,8.185) (4.629,8.238) (4.682,8.291) (4.735,8.344) (4.789,8.396) (4.843,8.449) (4.896,8.501) (4.950,8.553) (5.005,8.605) (5.059,8.656) (5.114,8.708) (5.168,8.759) (5.223,8.810) (5.278,8.861) (5.334,8.912) (5.389,8.962) (5.445,9.013) (5.500,9.063) (5.556,9.113) (5.613,9.163) (5.669,9.212) (5.725,9.262) (5.782,9.311) (5.839,9.360) (5.895,9.409) (5.953,9.457) (6.010,9.506) (6.067,9.554) (6.125,9.602) (6.183,9.650) (6.240,9.698) (6.298,9.745) (6.357,9.793) (6.415,9.840) (6.473,9.887) (6.532,9.933) (6.591,9.980) (6.630,10.000)}
\LC{(-3.797,3.606) (-3.869,3.626) (-3.942,3.645) (-4.014,3.663) (-4.087,3.680) (-4.161,3.696) (-4.234,3.712) (-4.307,3.726) (-4.381,3.740) (-4.455,3.753) (-4.529,3.765) (-4.603,3.775) (-4.678,3.785) (-4.752,3.794) (-4.827,3.802) (-4.901,3.810) (-4.976,3.816) (-5.051,3.821) (-5.126,3.825) (-5.201,3.829) (-5.276,3.831) (-5.351,3.832) (-5.426,3.833) (-5.501,3.832) (-5.576,3.831) (-5.651,3.828) (-5.726,3.825) (-5.800,3.820) (-5.875,3.815) (-5.950,3.808) (-6.025,3.801) (-6.099,3.792) (-6.173,3.782) (-6.248,3.772) (-6.322,3.760) (-6.396,3.748) (-6.470,3.734) (-6.543,3.719) (-6.616,3.703) (-6.689,3.687) (-6.762,3.669) (-6.835,3.650) (-6.907,3.630) (-6.979,3.609) (-7.051,3.587) (-7.122,3.564) (-7.193,3.540) (-7.264,3.514) (-7.334,3.488) (-7.404,3.461) (-7.473,3.432) (-7.542,3.402) (-7.611,3.372) (-7.679,3.340) (-7.746,3.307) (-7.813,3.273) (-7.879,3.239) (-7.945,3.202) (-8.010,3.165) (-8.075,3.127) (-8.139,3.088) (-8.202,3.048) (-8.265,3.006) (-8.326,2.964) (-8.387,2.920) (-8.448,2.876) (-8.507,2.830) (-8.566,2.784) (-8.624,2.736) (-8.681,2.687) (-8.737,2.637) (-8.792,2.587) (-8.846,2.535) (-8.900,2.482) (-8.952,2.428) (-9.003,2.373) (-9.054,2.318) (-9.103,2.261) (-9.151,2.203) (-9.198,2.145) (-9.243,2.085) (-9.288,2.025) (-9.331,1.964) (-9.373,1.902) (-9.414,1.839) (-9.454,1.775) (-9.492,1.711) (-9.529,1.645) (-9.564,1.579) (-9.598,1.513) (-9.631,1.445) (-9.662,1.377) (-9.692,1.308) (-9.720,1.239) (-9.747,1.169) (-9.772,1.098) (-9.796,1.027) (-9.818,0.955) (-9.839,0.883) (-9.858,0.810) (-9.875,0.737) (-9.891,0.664) (-9.905,0.590) (-9.918,0.517) (-9.928,0.442) (-9.938,0.368) (-9.945,0.293) (-9.951,0.218) (-9.955,0.144) (-9.957,0.069) (-9.958,-0.006) (-9.957,-0.081) (-9.954,-0.156) (-9.950,-0.231) (-9.944,-0.306) (-9.936,-0.381) (-9.927,-0.455) (-9.916,-0.529) (-9.903,-0.603) (-9.888,-0.677) (-9.872,-0.750) (-9.855,-0.823) (-9.836,-0.895) (-9.815,-0.967) (-9.792,-1.039) (-9.768,-1.110) (-9.743,-1.180) (-9.716,-1.250) (-9.687,-1.320) (-9.657,-1.388) (-9.626,-1.457) (-9.593,-1.524) (-9.558,-1.591) (-9.522,-1.656) (-9.485,-1.722) (-9.447,-1.786) (-9.407,-1.850) (-9.366,-1.912) (-9.324,-1.974) (-9.280,-2.035) (-9.236,-2.096) (-9.190,-2.155) (-9.143,-2.213) (-9.094,-2.271) (-9.045,-2.327) (-8.995,-2.383) (-8.943,-2.437) (-8.891,-2.491) (-8.837,-2.544) (-8.783,-2.595) (-8.728,-2.646) (-8.671,-2.695) (-8.614,-2.744) (-8.556,-2.792) (-8.497,-2.838) (-8.438,-2.883) (-8.377,-2.928) (-8.316,-2.971) (-8.254,-3.013) (-8.191,-3.055) (-8.128,-3.095) (-8.064,-3.134) (-7.999,-3.172) (-7.934,-3.209) (-7.868,-3.245) (-7.802,-3.279) (-7.735,-3.313) (-7.667,-3.346) (-7.599,-3.377) (-7.530,-3.408) (-7.461,-3.437) (-7.392,-3.465) (-7.322,-3.492) (-7.252,-3.519) (-7.181,-3.544) (-7.110,-3.568) (-7.039,-3.591) (-6.967,-3.613) (-6.895,-3.633) (-6.823,-3.653) (-6.750,-3.672) (-6.677,-3.690) (-6.604,-3.706) (-6.531,-3.722) (-6.457,-3.736) (-6.383,-3.750) (-6.309,-3.762) (-6.235,-3.774) (-6.161,-3.784) (-6.086,-3.794) (-6.012,-3.802) (-5.937,-3.809) (-5.863,-3.816) (-5.788,-3.821) (-5.713,-3.825) (-5.638,-3.829) (-5.563,-3.831) (-5.488,-3.832) (-5.413,-3.833) (-5.338,-3.832) (-5.263,-3.831) (-5.188,-3.828) (-5.113,-3.825) (-5.038,-3.820) (-4.964,-3.815) (-4.889,-3.808) (-4.814,-3.801) (-4.740,-3.793) (-4.665,-3.784) (-4.591,-3.774) (-4.517,-3.763) (-4.443,-3.751) (-4.369,-3.738) (-4.295,-3.724) (-4.221,-3.709) (-4.148,-3.694) (-4.075,-3.677) (-4.002,-3.660) (-3.929,-3.641) (-3.857,-3.622) (-3.809,-3.609)}
\LC{(-0.048,-10.000) (-0.048,-9.947) (-0.047,-9.872) (-0.046,-9.797) (-0.046,-9.722) (-0.045,-9.647) (-0.045,-9.572) (-0.044,-9.497) (-0.044,-9.422) (-0.043,-9.347) (-0.043,-9.272) (-0.042,-9.197) (-0.042,-9.122) (-0.041,-9.047) (-0.041,-8.972) (-0.040,-8.897) (-0.040,-8.822) (-0.039,-8.747) (-0.039,-8.672) (-0.038,-8.597) (-0.038,-8.522) (-0.037,-8.447) (-0.037,-8.372) (-0.036,-8.297) (-0.036,-8.222) (-0.035,-8.147) (-0.035,-8.072) (-0.034,-7.997) (-0.034,-7.922) (-0.033,-7.847) (-0.033,-7.772) (-0.032,-7.697) (-0.032,-7.622) (-0.031,-7.547) (-0.031,-7.472) (-0.030,-7.397) (-0.030,-7.322) (-0.030,-7.247) (-0.029,-7.172) (-0.029,-7.097) (-0.028,-7.022) (-0.028,-6.947) (-0.027,-6.872) (-0.027,-6.797) (-0.026,-6.722) (-0.026,-6.647) (-0.026,-6.572) (-0.025,-6.497) (-0.025,-6.422) (-0.024,-6.347) (-0.024,-6.272) (-0.023,-6.197) (-0.023,-6.122) (-0.023,-6.047) (-0.022,-5.972) (-0.022,-5.897) (-0.021,-5.822) (-0.021,-5.747) (-0.020,-5.672) (-0.020,-5.597) (-0.020,-5.522) (-0.019,-5.447) (-0.019,-5.372) (-0.018,-5.297) (-0.018,-5.222) (-0.018,-5.147) (-0.017,-5.072) (-0.017,-4.997) (-0.017,-4.922) (-0.016,-4.847) (-0.016,-4.772) (-0.015,-4.697) (-0.015,-4.622) (-0.015,-4.547) (-0.014,-4.472) (-0.014,-4.397) (-0.014,-4.322) (-0.013,-4.247) (-0.013,-4.172) (-0.013,-4.097) (-0.012,-4.022) (-0.012,-3.947) (-0.012,-3.872) (-0.011,-3.797) (-0.011,-3.722) (-0.011,-3.647) (-0.010,-3.572) (-0.010,-3.497) (-0.010,-3.422) (-0.009,-3.347) (-0.009,-3.272) (-0.009,-3.197) (-0.008,-3.122) (-0.008,-3.047) (-0.008,-2.972) (-0.007,-2.897) (-0.007,-2.822) (-0.007,-2.747) (-0.007,-2.672) (-0.006,-2.597) (-0.006,-2.522) (-0.006,-2.447) (-0.006,-2.372) (-0.005,-2.297)}
\LC{(-0.488,1.887) (-0.516,1.956) (-0.545,2.025) (-0.574,2.094) (-0.604,2.163) (-0.634,2.232) (-0.665,2.300) (-0.696,2.369) (-0.727,2.437) (-0.759,2.504) (-0.792,2.572) (-0.825,2.639) (-0.858,2.707) (-0.892,2.774) (-0.926,2.840) (-0.961,2.907) (-0.996,2.973) (-1.031,3.039) (-1.067,3.105) (-1.103,3.171) (-1.140,3.236) (-1.177,3.301) (-1.214,3.366) (-1.252,3.431) (-1.290,3.496) (-1.329,3.560) (-1.368,3.624) (-1.407,3.688) (-1.446,3.752) (-1.486,3.816) (-1.526,3.879) (-1.567,3.942) (-1.608,4.005) (-1.649,4.067) (-1.691,4.130) (-1.733,4.192) (-1.775,4.254) (-1.818,4.316) (-1.860,4.377) (-1.904,4.439) (-1.947,4.500) (-1.991,4.560) (-2.035,4.621) (-2.080,4.682) (-2.124,4.742) (-2.169,4.802) (-2.215,4.861) (-2.260,4.921) (-2.306,4.980) (-2.352,5.039) (-2.399,5.098) (-2.446,5.157) (-2.493,5.215) (-2.540,5.273) (-2.588,5.331) (-2.635,5.389) (-2.684,5.447) (-2.732,5.504) (-2.781,5.561) (-2.830,5.618) (-2.879,5.674) (-2.928,5.731) (-2.978,5.787) (-3.028,5.843) (-3.078,5.899) (-3.129,5.954) (-3.179,6.009) (-3.230,6.064) (-3.282,6.119) (-3.333,6.174) (-3.385,6.228) (-3.437,6.282) (-3.489,6.336) (-3.542,6.389) (-3.594,6.443) (-3.647,6.496) (-3.700,6.549) (-3.754,6.602) (-3.807,6.654) (-3.861,6.706) (-3.915,6.758) (-3.969,6.810) (-4.024,6.862) (-4.079,6.913) (-4.134,6.964) (-4.189,7.015) (-4.244,7.065) (-4.300,7.116) (-4.355,7.166) (-4.411,7.216) (-4.468,7.265) (-4.524,7.315) (-4.581,7.364) (-4.638,7.413) (-4.695,7.461) (-4.752,7.510) (-4.809,7.558) (-4.867,7.606) (-4.925,7.654) (-4.983,7.701) (-5.041,7.749) (-5.100,7.796) (-5.158,7.842) (-5.217,7.889) (-5.276,7.935) (-5.335,7.981) (-5.395,8.027) (-5.454,8.073) (-5.514,8.118) (-5.574,8.163) (-5.634,8.208) (-5.694,8.252) (-5.755,8.297) (-5.816,8.341) (-5.876,8.385) (-5.937,8.428) (-5.999,8.472) (-6.060,8.515) (-6.122,8.558) (-6.183,8.600) (-6.245,8.643) (-6.307,8.685) (-6.369,8.727) (-6.432,8.768) (-6.494,8.810) (-6.557,8.851) (-6.620,8.892) (-6.683,8.932) (-6.746,8.973) (-6.810,9.013) (-6.873,9.053) (-6.937,9.092) (-7.001,9.132) (-7.065,9.171) (-7.129,9.210) (-7.193,9.248) (-7.257,9.287) (-7.322,9.325) (-7.387,9.363) (-7.452,9.400) (-7.517,9.437) (-7.582,9.475) (-7.647,9.511) (-7.713,9.548) (-7.778,9.584) (-7.844,9.620) (-7.910,9.656) (-7.976,9.692) (-8.042,9.727) (-8.109,9.762) (-8.175,9.797) (-8.242,9.831) (-8.308,9.866) (-8.375,9.900) (-8.442,9.933) (-8.509,9.967) (-8.577,10.000) (-8.599,10.000)}
\LC{(-1.957,3.712) (-2.008,3.767) (-2.059,3.821) (-2.111,3.876) (-2.163,3.930) (-2.215,3.983) (-2.268,4.036) (-2.322,4.089) (-2.375,4.142) (-2.429,4.194) (-2.484,4.245) (-2.538,4.297) (-2.593,4.348) (-2.649,4.398) (-2.704,4.448) (-2.761,4.498) (-2.817,4.547) (-2.874,4.596) (-2.931,4.645) (-2.988,4.693) (-3.046,4.741) (-3.104,4.789) (-3.162,4.836) (-3.221,4.883) (-3.280,4.929) (-3.339,4.975) (-3.399,5.020) (-3.459,5.065) (-3.519,5.110) (-3.579,5.155) (-3.640,5.199) (-3.701,5.242) (-3.763,5.285) (-3.824,5.328) (-3.886,5.370) (-3.948,5.412) (-4.011,5.454) (-4.073,5.495) (-4.136,5.536) (-4.199,5.576) (-4.263,5.616) (-4.327,5.656) (-4.391,5.695) (-4.455,5.734) (-4.519,5.772) (-4.584,5.810) (-4.649,5.848) (-4.714,5.885) (-4.780,5.921) (-4.845,5.958) (-4.911,5.994) (-4.977,6.029) (-5.043,6.064) (-5.110,6.099) (-5.177,6.133) (-5.244,6.167) (-5.311,6.200) (-5.378,6.233) (-5.446,6.265) (-5.514,6.298) (-5.582,6.329) (-5.650,6.360) (-5.718,6.391) (-5.787,6.422) (-5.855,6.452) (-5.924,6.481) (-5.994,6.510) (-6.063,6.539) (-6.132,6.567) (-6.202,6.595) (-6.272,6.622) (-6.342,6.649) (-6.412,6.676) (-6.482,6.702) (-6.553,6.727) (-6.623,6.753) (-6.694,6.777) (-6.765,6.802) (-6.836,6.825) (-6.908,6.849) (-6.979,6.872) (-7.051,6.894) (-7.122,6.916) (-7.194,6.938) (-7.266,6.959) (-7.338,6.980) (-7.410,7.000) (-7.483,7.020) (-7.555,7.039) (-7.628,7.058) (-7.700,7.076) (-7.773,7.094) (-7.846,7.112) (-7.919,7.129) (-7.992,7.145) (-8.066,7.161) (-8.139,7.177) (-8.212,7.192) (-8.286,7.207) (-8.360,7.221) (-8.433,7.235) (-8.507,7.248) (-8.581,7.261) (-8.655,7.273) (-8.729,7.285) (-8.803,7.297) (-8.877,7.308) (-8.952,7.318) (-9.026,7.328) (-9.100,7.337) (-9.175,7.346) (-9.249,7.355) (-9.324,7.363) (-9.399,7.371) (-9.473,7.378) (-9.548,7.384) (-9.623,7.390) (-9.697,7.396) (-9.772,7.401) (-9.847,7.406) (-9.922,7.410) (-9.997,7.413) (-10.000,7.414)}
\LC{(-10.000,-8.150) (-9.943,-8.140) (-9.868,-8.130) (-9.794,-8.119) (-9.720,-8.107) (-9.646,-8.095) (-9.572,-8.083) (-9.498,-8.070) (-9.424,-8.057) (-9.350,-8.044) (-9.277,-8.030) (-9.203,-8.016) (-9.130,-8.001) (-9.056,-7.986) (-8.983,-7.970) (-8.910,-7.954) (-8.836,-7.938) (-8.763,-7.921) (-8.690,-7.903) (-8.617,-7.886) (-8.545,-7.868) (-8.472,-7.849) (-8.399,-7.830) (-8.327,-7.811) (-8.255,-7.791) (-8.182,-7.771) (-8.110,-7.750) (-8.038,-7.729) (-7.966,-7.708) (-7.895,-7.686) (-7.823,-7.664) (-7.751,-7.641) (-7.680,-7.618) (-7.609,-7.595) (-7.538,-7.571) (-7.467,-7.547) (-7.396,-7.522) (-7.325,-7.497) (-7.255,-7.472) (-7.184,-7.446) (-7.114,-7.420) (-7.044,-7.393) (-6.974,-7.366) (-6.904,-7.339) (-6.834,-7.311) (-6.765,-7.283) (-6.695,-7.255) (-6.626,-7.226) (-6.557,-7.196) (-6.488,-7.167) (-6.420,-7.137) (-6.351,-7.106) (-6.283,-7.075) (-6.214,-7.044) (-6.146,-7.012) (-6.079,-6.980) (-6.011,-6.948) (-5.944,-6.915) (-5.876,-6.882) (-5.809,-6.849) (-5.742,-6.815) (-5.676,-6.780) (-5.609,-6.746) (-5.543,-6.711) (-5.477,-6.675) (-5.411,-6.640) (-5.345,-6.603) (-5.279,-6.567) (-5.214,-6.530) (-5.149,-6.493) (-5.084,-6.455) (-5.020,-6.417) (-4.955,-6.379) (-4.891,-6.340) (-4.827,-6.301) (-4.763,-6.261) (-4.700,-6.222) (-4.636,-6.181) (-4.573,-6.141) (-4.510,-6.100) (-4.448,-6.059) (-4.385,-6.017) (-4.323,-5.975) (-4.261,-5.933) (-4.200,-5.890) (-4.138,-5.847) (-4.077,-5.803) (-4.016,-5.760) (-3.956,-5.715) (-3.895,-5.671) (-3.835,-5.626) (-3.775,-5.581) (-3.716,-5.535) (-3.656,-5.489) (-3.597,-5.443) (-3.539,-5.397) (-3.480,-5.350) (-3.422,-5.302) (-3.364,-5.255) (-3.306,-5.207) (-3.249,-5.158) (-3.192,-5.110) (-3.135,-5.061) (-3.079,-5.011) (-3.023,-4.961) (-2.967,-4.911) (-2.911,-4.861) (-2.856,-4.810) (-2.801,-4.759) (-2.747,-4.708) (-2.692,-4.656) (-2.638,-4.604) (-2.585,-4.551) (-2.531,-4.499) (-2.479,-4.445) (-2.426,-4.392) (-2.374,-4.338) (-2.322,-4.284) (-2.270,-4.230) (-2.219,-4.175) (-2.168,-4.120) (-2.118,-4.064) (-2.068,-4.008) (-2.018,-3.952) (-1.968,-3.896) (-1.919,-3.839) (-1.871,-3.782) (-1.823,-3.724) (-1.775,-3.667) (-1.727,-3.609) (-1.680,-3.550) (-1.634,-3.491) (-1.587,-3.432) (-1.542,-3.373) (-1.496,-3.313) (-1.451,-3.253) (-1.407,-3.193) (-1.363,-3.132) (-1.319,-3.071) (-1.276,-3.010) (-1.233,-2.948) (-1.191,-2.886) (-1.149,-2.824)}
\LC{(3.767,2.687) (3.842,2.687) (3.917,2.685) (3.992,2.683) (4.066,2.678) (4.141,2.673) (4.216,2.665) (4.290,2.657) (4.365,2.647) (4.439,2.635) (4.513,2.622) (4.586,2.608) (4.660,2.592) (4.732,2.574) (4.805,2.555) (4.877,2.535) (4.949,2.513) (5.020,2.489) (5.091,2.464) (5.161,2.438) (5.231,2.410) (5.299,2.380) (5.368,2.349) (5.435,2.316) (5.502,2.282) (5.568,2.246) (5.633,2.209) (5.697,2.170) (5.760,2.130) (5.823,2.088) (5.884,2.045) (5.944,2.000) (6.003,1.954) (6.061,1.906) (6.117,1.856) (6.172,1.806) (6.226,1.753) (6.279,1.700) (6.330,1.645) (6.379,1.588) (6.427,1.531) (6.473,1.472) (6.518,1.411) (6.560,1.350) (6.601,1.287) (6.641,1.223) (6.678,1.158) (6.713,1.092) (6.747,1.025) (6.778,0.956) (6.807,0.887) (6.834,0.817) (6.859,0.747) (6.882,0.675) (6.902,0.603) (6.920,0.530) (6.936,0.457) (6.950,0.383) (6.961,0.309) (6.970,0.235) (6.976,0.160) (6.980,0.085) (6.981,0.010) (6.980,-0.065) (6.977,-0.140) (6.971,-0.215) (6.963,-0.289) (6.953,-0.364) (6.940,-0.437) (6.925,-0.511) (6.907,-0.584) (6.887,-0.656) (6.865,-0.728) (6.841,-0.799) (6.814,-0.869) (6.786,-0.938) (6.755,-1.007) (6.722,-1.074) (6.687,-1.141) (6.651,-1.206) (6.612,-1.270) (6.571,-1.333) (6.529,-1.395) (6.485,-1.456) (6.439,-1.515) (6.392,-1.573) (6.343,-1.630) (6.292,-1.685) (6.240,-1.739) (6.187,-1.792) (6.132,-1.843) (6.076,-1.893) (6.018,-1.941) (5.960,-1.988) (5.900,-2.033) (5.839,-2.077) (5.777,-2.119) (5.714,-2.160) (5.650,-2.199) (5.585,-2.237) (5.519,-2.273) (5.453,-2.308) (5.386,-2.341) (5.318,-2.372) (5.249,-2.402) (5.179,-2.431) (5.109,-2.458) (5.039,-2.483) (4.968,-2.507) (4.896,-2.529) (4.824,-2.550) (4.752,-2.569) (4.679,-2.587) (4.606,-2.604) (4.532,-2.618) (4.458,-2.632) (4.384,-2.644) (4.310,-2.654) (4.236,-2.663) (4.161,-2.671) (4.086,-2.677) (4.011,-2.682) (3.936,-2.685) (3.861,-2.687) (3.786,-2.687) (3.761,-2.687)}
\LC{(10.000,-7.569) (9.949,-7.564) (9.875,-7.558) (9.800,-7.553) (9.725,-7.546) (9.650,-7.539) (9.576,-7.532) (9.501,-7.524) (9.427,-7.516) (9.352,-7.508) (9.278,-7.498) (9.203,-7.489) (9.129,-7.479) (9.055,-7.468) (8.980,-7.457) (8.906,-7.445) (8.832,-7.434) (8.758,-7.421) (8.684,-7.408) (8.611,-7.395) (8.537,-7.381) (8.463,-7.367) (8.390,-7.352) (8.316,-7.337) (8.243,-7.321) (8.170,-7.305) (8.097,-7.288) (8.024,-7.271) (7.951,-7.254) (7.878,-7.236) (7.805,-7.218) (7.733,-7.199) (7.660,-7.179) (7.588,-7.160) (7.515,-7.140) (7.443,-7.119) (7.371,-7.098) (7.299,-7.076) (7.228,-7.054) (7.156,-7.032) (7.085,-7.009) (7.013,-6.986) (6.942,-6.962) (6.871,-6.938) (6.800,-6.914) (6.730,-6.889) (6.659,-6.863) (6.589,-6.837) (6.518,-6.811) (6.448,-6.784) (6.378,-6.757) (6.309,-6.730) (6.239,-6.702) (6.170,-6.673) (6.100,-6.645) (6.031,-6.615) (5.963,-6.586) (5.894,-6.555) (5.825,-6.525) (5.757,-6.494) (5.689,-6.463) (5.621,-6.431) (5.553,-6.399) (5.486,-6.366) (5.418,-6.333) (5.351,-6.300) (5.284,-6.266) (5.218,-6.232) (5.151,-6.197) (5.085,-6.162) (5.019,-6.126) (4.953,-6.091) (4.887,-6.054) (4.822,-6.018) (4.757,-5.980) (4.692,-5.943) (4.627,-5.905) (4.562,-5.867) (4.498,-5.828) (4.434,-5.789) (4.370,-5.749) (4.307,-5.710) (4.244,-5.669) (4.181,-5.629) (4.139,-5.601)}
\LC{(3.620,3.580) (3.691,3.603) (3.763,3.624) (3.835,3.645) (3.908,3.665) (3.980,3.684) (4.053,3.703) (4.126,3.720) (4.199,3.736) (4.272,3.752) (4.346,3.767) (4.420,3.781) (4.494,3.793) (4.568,3.805) (4.642,3.816) (4.716,3.827) (4.790,3.836) (4.865,3.844) (4.940,3.851) (5.014,3.858) (5.089,3.863) (5.164,3.868) (5.239,3.871) (5.314,3.874) (5.389,3.876) (5.464,3.876) (5.539,3.876) (5.614,3.875) (5.689,3.873) (5.764,3.870) (5.839,3.865) (5.914,3.860) (5.988,3.854) (6.063,3.847) (6.137,3.839) (6.212,3.830) (6.286,3.820) (6.360,3.808) (6.434,3.796) (6.508,3.783) (6.582,3.769) (6.655,3.754) (6.728,3.737) (6.801,3.720) (6.874,3.702) (6.947,3.682) (7.019,3.662) (7.091,3.640) (7.162,3.618) (7.233,3.594) (7.304,3.570) (7.375,3.544) (7.445,3.517) (7.514,3.489) (7.584,3.460) (7.652,3.431) (7.721,3.400) (7.788,3.367) (7.856,3.334) (7.922,3.300) (7.989,3.265) (8.054,3.229) (8.119,3.191) (8.184,3.153) (8.247,3.113) (8.311,3.073) (8.373,3.031) (8.435,2.988) (8.496,2.945) (8.556,2.900) (8.615,2.854) (8.674,2.807) (8.731,2.759) (8.788,2.710) (8.844,2.660) (8.899,2.610) (8.953,2.558) (9.007,2.505) (9.059,2.451) (9.110,2.396) (9.160,2.340) (9.209,2.284) (9.257,2.226) (9.304,2.167) (9.350,2.108) (9.394,2.048) (9.438,1.986) (9.480,1.924) (9.520,1.861) (9.560,1.798) (9.598,1.733) (9.635,1.668) (9.671,1.602) (9.705,1.535) (9.738,1.468) (9.769,1.400) (9.799,1.331) (9.828,1.261) (9.855,1.191) (9.880,1.121) (9.904,1.050) (9.927,0.978) (9.947,0.906) (9.967,0.834) (9.984,0.761) (10.000,0.688)}
\LC{(2.538,5.010) (2.587,5.067) (2.637,5.123) (2.687,5.179) (2.737,5.234) (2.788,5.290) (2.839,5.345) (2.890,5.400) (2.941,5.454) (2.993,5.508) (3.045,5.563) (3.097,5.616) (3.150,5.670) (3.203,5.723) (3.256,5.776) (3.309,5.829) (3.363,5.882) (3.416,5.934) (3.470,5.986) (3.525,6.037) (3.579,6.089) (3.634,6.140) (3.689,6.191) (3.745,6.242) (3.800,6.292) (3.856,6.342) (3.912,6.392) (3.968,6.442) (4.025,6.491) (4.081,6.540) (4.138,6.589) (4.196,6.637) (4.253,6.685) (4.311,6.733) (4.369,6.781) (4.427,6.829) (4.485,6.876) (4.544,6.923) (4.602,6.969) (4.661,7.016) (4.720,7.062) (4.780,7.107) (4.840,7.153) (4.899,7.198) (4.959,7.243) (5.020,7.288) (5.080,7.332) (5.141,7.376) (5.202,7.420) (5.263,7.464) (5.324,7.507) (5.385,7.550) (5.447,7.593) (5.509,7.635) (5.571,7.677) (5.633,7.719) (5.695,7.761) (5.758,7.802) (5.821,7.843) (5.884,7.884) (5.947,7.925) (6.010,7.965) (6.074,8.005) (6.137,8.044) (6.201,8.084) (6.265,8.123) (6.329,8.162) (6.394,8.200) (6.458,8.238) (6.523,8.276) (6.588,8.314) (6.653,8.351) (6.718,8.388) (6.783,8.425) (6.849,8.462) (6.915,8.498) (6.980,8.534) (7.046,8.569) (7.113,8.605) (7.179,8.640) (7.245,8.674) (7.312,8.709) (7.379,8.743) (7.446,8.777) (7.513,8.810) (7.580,8.844) (7.647,8.877) (7.715,8.909) (7.782,8.942) (7.850,8.974) (7.918,9.005) (7.986,9.037) (8.054,9.068) (8.123,9.099) (8.191,9.130) (8.260,9.160) (8.329,9.190) (8.398,9.220) (8.467,9.249) (8.536,9.278) (8.605,9.307) (8.674,9.335) (8.744,9.363) (8.813,9.391) (8.883,9.419) (8.953,9.446) (9.023,9.473) (9.093,9.500) (9.163,9.526) (9.234,9.552) (9.304,9.578) (9.375,9.603) (9.446,9.628) (9.516,9.653) (9.587,9.677) (9.658,9.702) (9.729,9.725) (9.801,9.749) (9.872,9.772) (9.943,9.795) (10.000,9.818)}
\LC{(1.312,1.428) (1.381,1.456) (1.451,1.482) (1.523,1.506) (1.594,1.528) (1.667,1.548) (1.739,1.566) (1.813,1.582) (1.886,1.595) (1.961,1.607) (2.035,1.616) (2.110,1.623) (2.184,1.628) (2.259,1.631) (2.334,1.631) (2.409,1.629) (2.484,1.625) (2.559,1.618) (2.633,1.609) (2.708,1.597) (2.781,1.583) (2.854,1.567) (2.927,1.548) (2.999,1.527) (3.070,1.503) (3.140,1.477) (3.209,1.448) (3.278,1.417) (3.345,1.383) (3.410,1.346) (3.474,1.307) (3.537,1.266) (3.597,1.222) (3.656,1.176) (3.713,1.127) (3.768,1.075) (3.820,1.022) (3.870,0.966) (3.918,0.908) (3.963,0.848) (4.004,0.785) (4.043,0.721) (4.079,0.655) (4.111,0.588) (4.140,0.518) (4.165,0.448) (4.187,0.376) (4.205,0.303) (4.219,0.230) (4.229,0.155) (4.235,0.080) (4.238,0.006) (4.236,-0.069) (4.230,-0.144) (4.221,-0.219) (4.207,-0.292) (4.190,-0.365) (4.169,-0.437) (4.144,-0.508) (4.116,-0.577) (4.084,-0.645) (4.049,-0.712) (4.010,-0.776) (3.969,-0.839) (3.925,-0.899) (3.878,-0.958) (3.828,-1.014) (3.776,-1.068) (3.721,-1.119) (3.665,-1.169) (3.606,-1.215) (3.546,-1.260) (3.483,-1.301) (3.420,-1.341) (3.354,-1.378) (3.288,-1.412) (3.220,-1.444) (3.151,-1.473) (3.080,-1.499) (3.009,-1.524) (2.938,-1.545) (2.865,-1.564) (2.792,-1.581) (2.718,-1.595) (2.644,-1.607) (2.570,-1.617) (2.495,-1.624) (2.420,-1.628) (2.345,-1.631) (2.270,-1.631) (2.195,-1.629) (2.121,-1.624) (2.046,-1.617) (1.971,-1.608) (1.897,-1.597) (1.824,-1.584) (1.750,-1.568) (1.677,-1.551) (1.605,-1.531) (1.533,-1.510) (1.462,-1.486) (1.391,-1.460) (1.322,-1.432) (1.299,-1.423)}
\LC{(2.736,-10.000) (2.707,-9.935) (2.678,-9.866) (2.648,-9.797) (2.619,-9.728) (2.591,-9.659) (2.562,-9.589) (2.533,-9.520) (2.504,-9.451) (2.476,-9.382) (2.447,-9.312) (2.419,-9.243) (2.391,-9.173) (2.363,-9.104) (2.335,-9.034) (2.307,-8.964) (2.279,-8.895) (2.252,-8.825) (2.224,-8.755) (2.197,-8.685) (2.169,-8.616) (2.142,-8.546) (2.115,-8.476) (2.088,-8.406) (2.061,-8.336) (2.034,-8.266) (2.008,-8.196) (1.981,-8.125) (1.955,-8.055) (1.929,-7.985) (1.902,-7.915) (1.876,-7.844) (1.851,-7.774) (1.825,-7.704) (1.799,-7.633) (1.773,-7.563) (1.748,-7.492) (1.723,-7.421) (1.697,-7.351) (1.672,-7.280) (1.647,-7.209) (1.623,-7.139) (1.598,-7.068) (1.573,-6.997) (1.549,-6.926) (1.524,-6.855) (1.500,-6.784) (1.476,-6.713) (1.452,-6.642) (1.428,-6.571) (1.405,-6.500) (1.381,-6.429) (1.358,-6.357) (1.334,-6.286) (1.311,-6.215) (1.288,-6.143) (1.265,-6.072) (1.243,-6.000) (1.220,-5.929) (1.197,-5.857) (1.175,-5.786) (1.153,-5.714) (1.131,-5.642) (1.109,-5.571) (1.087,-5.499) (1.066,-5.427) (1.044,-5.355) (1.023,-5.283) (1.001,-5.211) (0.980,-5.139) (0.959,-5.067) (0.939,-4.995) (0.918,-4.923) (0.898,-4.851) (0.877,-4.779) (0.857,-4.707) (0.837,-4.634) (0.817,-4.562) (0.798,-4.490) (0.778,-4.417) (0.759,-4.345) (0.740,-4.272)}
\LC{(-3.539,5.679) (-3.596,5.728) (-3.653,5.776) (-3.711,5.824) (-3.768,5.872) (-3.826,5.920) (-3.885,5.967) (-3.943,6.014) (-4.002,6.061) (-4.061,6.107) (-4.120,6.153) (-4.179,6.199) (-4.239,6.244) (-4.299,6.289) (-4.359,6.334) (-4.420,6.378) (-4.480,6.422) (-4.541,6.466) (-4.602,6.509) (-4.664,6.553) (-4.725,6.595) (-4.787,6.638) (-4.849,6.680) (-4.911,6.722) (-4.974,6.764) (-5.037,6.805) (-5.099,6.846) (-5.162,6.886) (-5.226,6.926) (-5.289,6.966) (-5.353,7.006) (-5.417,7.045) (-5.481,7.084) (-5.545,7.123) (-5.610,7.161) (-5.674,7.199) (-5.739,7.237) (-5.804,7.274) (-5.870,7.311) (-5.935,7.347) (-6.001,7.384) (-6.067,7.420) (-6.133,7.455) (-6.199,7.490) (-6.265,7.525) (-6.332,7.560) (-6.398,7.594) (-6.465,7.628) (-6.532,7.662) (-6.600,7.695) (-6.667,7.728) (-6.735,7.761) (-6.802,7.793) (-6.870,7.825) (-6.938,7.856) (-7.006,7.887) (-7.075,7.918) (-7.143,7.949) (-7.212,7.979) (-7.281,8.009) (-7.350,8.038) (-7.419,8.067) (-7.488,8.096) (-7.558,8.125) (-7.627,8.153) (-7.697,8.180) (-7.767,8.208) (-7.837,8.235) (-7.907,8.261) (-7.977,8.288) (-8.047,8.314) (-8.118,8.339) (-8.188,8.365) (-8.259,8.389) (-8.330,8.414) (-8.401,8.438) (-8.472,8.462) (-8.543,8.485) (-8.615,8.509) (-8.686,8.531) (-8.758,8.554) (-8.829,8.576) (-8.901,8.597) (-8.973,8.619) (-9.045,8.640) (-9.117,8.660) (-9.190,8.680) (-9.262,8.700) (-9.334,8.720) (-9.407,8.739) (-9.479,8.757) (-9.552,8.776) (-9.625,8.794) (-9.698,8.811) (-9.771,8.829) (-9.844,8.845) (-9.917,8.862) (-9.990,8.878) (-10.000,8.883)}
\LC{(8.623,2.249) (8.673,2.193) (8.722,2.136) (8.769,2.078) (8.815,2.019) (8.860,1.959) (8.904,1.898) (8.947,1.836) (8.988,1.773) (9.028,1.710) (9.066,1.646) (9.103,1.580) (9.139,1.514) (9.173,1.448) (9.206,1.380) (9.237,1.312) (9.267,1.243) (9.295,1.174) (9.321,1.103) (9.346,1.033) (9.370,0.961) (9.391,0.890) (9.411,0.817) (9.430,0.745) (9.446,0.671) (9.461,0.598) (9.474,0.524) (9.486,0.450) (9.496,0.376) (9.504,0.301) (9.510,0.226) (9.514,0.151) (9.517,0.076) (9.518,0.001) (9.517,-0.073) (9.514,-0.148) (9.510,-0.223) (9.504,-0.298) (9.496,-0.373) (9.486,-0.447) (9.475,-0.521) (9.462,-0.595) (9.447,-0.668) (9.430,-0.742) (9.412,-0.814) (9.392,-0.887) (9.370,-0.958) (9.347,-1.030) (9.322,-1.101) (9.296,-1.171) (9.268,-1.240) (9.238,-1.309) (9.207,-1.377) (9.174,-1.445) (9.140,-1.512) (9.105,-1.578) (9.068,-1.643) (9.029,-1.707) (8.990,-1.771) (8.948,-1.834) (8.906,-1.896) (8.862,-1.956) (8.817,-2.016) (8.771,-2.075) (8.724,-2.134) (8.675,-2.191) (8.625,-2.247) (8.574,-2.302) (8.522,-2.356) (8.469,-2.409) (8.415,-2.461) (8.360,-2.512) (8.304,-2.562) (8.248,-2.611) (8.190,-2.659) (8.131,-2.705) (8.071,-2.751) (8.011,-2.795) (7.950,-2.839) (7.888,-2.881) (7.825,-2.922) (7.762,-2.962) (7.698,-3.001) (7.633,-3.039) (7.567,-3.076) (7.501,-3.111) (7.435,-3.146) (7.368,-3.179) (7.300,-3.211) (7.232,-3.242) (7.163,-3.272) (7.094,-3.301) (7.024,-3.328) (6.954,-3.355) (6.883,-3.380) (6.812,-3.405) (6.741,-3.428) (6.669,-3.450) (6.597,-3.471) (6.525,-3.491) (6.452,-3.509) (6.379,-3.527) (6.306,-3.544) (6.233,-3.559) (6.159,-3.574) (6.085,-3.587) (6.011,-3.599) (5.937,-3.610) (5.863,-3.620) (5.788,-3.629) (5.714,-3.637) (5.639,-3.644) (5.564,-3.650) (5.490,-3.655) (5.415,-3.658) (5.340,-3.661) (5.265,-3.663) (5.190,-3.663) (5.115,-3.663) (5.040,-3.662) (4.965,-3.659) (4.890,-3.656) (4.815,-3.651)}
\LC{(9.552,-10.000) (9.482,-9.976) (9.412,-9.949) (9.342,-9.921) (9.273,-9.894) (9.203,-9.866) (9.134,-9.837) (9.064,-9.809) (8.995,-9.780) (8.926,-9.751) (8.857,-9.722) (8.788,-9.692) (8.719,-9.662) (8.651,-9.632) (8.582,-9.601) (8.514,-9.571) (8.445,-9.539) (8.377,-9.508) (8.309,-9.476) (8.241,-9.444) (8.174,-9.412) (8.106,-9.380) (8.039,-9.347) (7.971,-9.314) (7.904,-9.281) (7.837,-9.247) (7.770,-9.213) (7.703,-9.179) (7.637,-9.144) (7.570,-9.110) (7.504,-9.075) (7.438,-9.039) (7.372,-9.004) (7.306,-8.968) (7.240,-8.932) (7.175,-8.895) (7.109,-8.859) (7.044,-8.822) (6.979,-8.785) (6.914,-8.747) (6.849,-8.709) (6.784,-8.671) (6.720,-8.633) (6.656,-8.594) (6.592,-8.556) (6.528,-8.516) (6.464,-8.477) (6.400,-8.437) (6.337,-8.398) (6.273,-8.357) (6.210,-8.317) (6.147,-8.276) (6.084,-8.235) (6.022,-8.194) (5.959,-8.152) (5.897,-8.110) (5.835,-8.068) (5.773,-8.026) (5.711,-7.983) (5.650,-7.941) (5.588,-7.897) (5.527,-7.854) (5.466,-7.810) (5.405,-7.766) (5.345,-7.722) (5.284,-7.678) (5.224,-7.633) (5.164,-7.588) (5.104,-7.543) (5.045,-7.497) (4.985,-7.452) (4.926,-7.406) (4.867,-7.359) (4.808,-7.313) (4.750,-7.266) (4.691,-7.219) (4.633,-7.172) (4.575,-7.124) (4.517,-7.076) (4.460,-7.028) (4.402,-6.980) (4.345,-6.931) (4.288,-6.882) (4.232,-6.833) (4.175,-6.784) (4.119,-6.734) (4.063,-6.684) (4.007,-6.634) (3.952,-6.584) (3.896,-6.533) (3.841,-6.482) (3.786,-6.431) (3.732,-6.380) (3.677,-6.328) (3.623,-6.276) (3.569,-6.224) (3.515,-6.172) (3.462,-6.119) (3.409,-6.066) (3.356,-6.013) (3.303,-5.960) (3.251,-5.906) (3.199,-5.852) (3.147,-5.798) (3.095,-5.744) (3.044,-5.689) (2.993,-5.634) (2.942,-5.579) (2.891,-5.524) (2.841,-5.468) (2.791,-5.412) (2.741,-5.356) (2.691,-5.300) (2.642,-5.243) (2.593,-5.187) (2.545,-5.130) (2.496,-5.072) (2.448,-5.015) (2.400,-4.957) (2.353,-4.899) (2.306,-4.841) (2.259,-4.782) (2.212,-4.723) (2.166,-4.664) (2.120,-4.605) (2.074,-4.546) (2.029,-4.486) (1.984,-4.426) (1.939,-4.366) (1.894,-4.305) (1.850,-4.245) (1.807,-4.184) (1.763,-4.123) (1.720,-4.061) (1.677,-4.000) (1.635,-3.938) (1.593,-3.876) (1.551,-3.813) (1.510,-3.751) (1.482,-3.709)}
\LC{(-10.000,-5.600) (-9.942,-5.616) (-9.869,-5.631) (-9.795,-5.646) (-9.721,-5.659) (-9.648,-5.673) (-9.574,-5.685) (-9.500,-5.697) (-9.425,-5.708) (-9.351,-5.719) (-9.277,-5.729) (-9.202,-5.738) (-9.128,-5.747) (-9.053,-5.754) (-8.979,-5.762) (-8.904,-5.768) (-8.829,-5.774) (-8.754,-5.779) (-8.679,-5.784) (-8.605,-5.788) (-8.530,-5.791) (-8.455,-5.793) (-8.380,-5.795) (-8.305,-5.796) (-8.230,-5.797) (-8.155,-5.797) (-8.080,-5.796) (-8.005,-5.795) (-7.930,-5.793) (-7.855,-5.790) (-7.780,-5.787) (-7.705,-5.783) (-7.630,-5.779) (-7.555,-5.773) (-7.480,-5.768) (-7.406,-5.761) (-7.331,-5.754) (-7.256,-5.747) (-7.182,-5.738) (-7.107,-5.729) (-7.033,-5.720) (-6.959,-5.710) (-6.885,-5.699) (-6.810,-5.688) (-6.736,-5.676) (-6.662,-5.663) (-6.589,-5.650) (-6.515,-5.636) (-6.441,-5.622) (-6.368,-5.607) (-6.294,-5.591) (-6.221,-5.575) (-6.148,-5.558) (-6.075,-5.541) (-6.002,-5.523) (-5.930,-5.504) (-5.857,-5.485) (-5.785,-5.465) (-5.713,-5.445) (-5.641,-5.424) (-5.569,-5.403) (-5.497,-5.380) (-5.426,-5.358) (-5.354,-5.335) (-5.283,-5.311) (-5.212,-5.287) (-5.141,-5.262) (-5.071,-5.236) (-5.001,-5.210) (-4.930,-5.184) (-4.860,-5.156) (-4.791,-5.129) (-4.721,-5.100) (-4.652,-5.072) (-4.583,-5.042) (-4.514,-5.012) (-4.446,-4.982) (-4.377,-4.951) (-4.309,-4.919) (-4.242,-4.887) (-4.174,-4.854) (-4.107,-4.821) (-4.040,-4.787) (-3.973,-4.753) (-3.907,-4.718) (-3.841,-4.683) (-3.775,-4.647) (-3.709,-4.611) (-3.644,-4.574) (-3.579,-4.536) (-3.514,-4.498) (-3.450,-4.460) (-3.386,-4.421) (-3.322,-4.381) (-3.259,-4.341) (-3.196,-4.301) (-3.133,-4.259) (-3.071,-4.218) (-3.009,-4.176) (-2.947,-4.133) (-2.886,-4.090) (-2.825,-4.046) (-2.764,-4.002) (-2.704,-3.957) (-2.644,-3.912) (-2.584,-3.867) (-2.525,-3.820) (-2.467,-3.774)}
\LC{(10.000,-4.876) (9.939,-4.901) (9.868,-4.925) (9.797,-4.949) (9.725,-4.972) (9.653,-4.993) (9.581,-5.015) (9.509,-5.035) (9.437,-5.055) (9.364,-5.073) (9.292,-5.092) (9.219,-5.109) (9.145,-5.125) (9.072,-5.141) (8.999,-5.156) (8.925,-5.171) (8.851,-5.184) (8.777,-5.197) (8.703,-5.209) (8.629,-5.220) (8.555,-5.231) (8.481,-5.241) (8.406,-5.250) (8.332,-5.258) (8.257,-5.266) (8.182,-5.273) (8.108,-5.279) (8.033,-5.285) (7.958,-5.289) (7.883,-5.293) (7.808,-5.297) (7.733,-5.299) (7.658,-5.301) (7.583,-5.302) (7.508,-5.303) (7.433,-5.302) (7.358,-5.301) (7.283,-5.300) (7.208,-5.297) (7.133,-5.294) (7.058,-5.290) (6.984,-5.286) (6.909,-5.281) (6.834,-5.275) (6.759,-5.268) (6.685,-5.261) (6.610,-5.253) (6.535,-5.245) (6.461,-5.236) (6.387,-5.226) (6.312,-5.215) (6.238,-5.204) (6.164,-5.192) (6.090,-5.179) (6.017,-5.166) (5.943,-5.152) (5.869,-5.137) (5.796,-5.122) (5.723,-5.106) (5.649,-5.090) (5.576,-5.073) (5.504,-5.055) (5.431,-5.036) (5.358,-5.017) (5.286,-4.998) (5.214,-4.977) (5.142,-4.956) (5.070,-4.934) (4.998,-4.912) (4.927,-4.889) (4.856,-4.866) (4.785,-4.842) (4.714,-4.817) (4.643,-4.791) (4.573,-4.765) (4.503,-4.739) (4.433,-4.711) (4.363,-4.684) (4.294,-4.655) (4.248,-4.636)}
\LC{(-0.445,2.886) (-0.462,2.959) (-0.479,3.032) (-0.497,3.105) (-0.515,3.178) (-0.533,3.251) (-0.551,3.323) (-0.570,3.396) (-0.588,3.469) (-0.607,3.541) (-0.626,3.614) (-0.645,3.686) (-0.665,3.759) (-0.684,3.831) (-0.704,3.904) (-0.724,3.976) (-0.744,4.048) (-0.765,4.120) (-0.785,4.192) (-0.806,4.265) (-0.827,4.337) (-0.848,4.409) (-0.869,4.481) (-0.891,4.552) (-0.912,4.624) (-0.934,4.696) (-0.956,4.768) (-0.978,4.839) (-1.000,4.911) (-1.023,4.983) (-1.045,5.054) (-1.068,5.126) (-1.091,5.197) (-1.114,5.268) (-1.137,5.340) (-1.160,5.411) (-1.184,5.482) (-1.208,5.553) (-1.231,5.624) (-1.255,5.695) (-1.280,5.766) (-1.304,5.837) (-1.328,5.908) (-1.353,5.979) (-1.378,6.050) (-1.403,6.121) (-1.428,6.191) (-1.453,6.262) (-1.478,6.333) (-1.504,6.403) (-1.529,6.474) (-1.555,6.544) (-1.581,6.615) (-1.607,6.685) (-1.633,6.755) (-1.659,6.825) (-1.686,6.896) (-1.712,6.966) (-1.739,7.036) (-1.766,7.106) (-1.793,7.176) (-1.820,7.246) (-1.847,7.316) (-1.875,7.385) (-1.902,7.455) (-1.930,7.525) (-1.958,7.595) (-1.986,7.664) (-2.014,7.734) (-2.042,7.803) (-2.070,7.873) (-2.099,7.942) (-2.127,8.011) (-2.156,8.081) (-2.185,8.150) (-2.214,8.219) (-2.243,8.288) (-2.272,8.357) (-2.301,8.426) (-2.331,8.495) (-2.360,8.564) (-2.390,8.633) (-2.420,8.702) (-2.450,8.771) (-2.480,8.840) (-2.510,8.908) (-2.540,8.977) (-2.570,9.045) (-2.601,9.114) (-2.632,9.182) (-2.662,9.251) (-2.693,9.319) (-2.724,9.387) (-2.755,9.456) (-2.786,9.524) (-2.818,9.592) (-2.849,9.660) (-2.881,9.728) (-2.912,9.796) (-2.944,9.864) (-2.976,9.932) (-3.008,10.000) (-3.019,10.000)}
\LC{(-10.000,-9.569) (-9.950,-9.548) (-9.878,-9.527) (-9.807,-9.505) (-9.735,-9.484) (-9.663,-9.461) (-9.592,-9.439) (-9.520,-9.416) (-9.449,-9.393) (-9.378,-9.369) (-9.307,-9.346) (-9.236,-9.321) (-9.165,-9.297) (-9.094,-9.272) (-9.023,-9.247) (-8.953,-9.222) (-8.882,-9.196) (-8.812,-9.170) (-8.742,-9.143) (-8.672,-9.117) (-8.602,-9.090) (-8.532,-9.062) (-8.462,-9.035) (-8.393,-9.007) (-8.323,-8.978) (-8.254,-8.950) (-8.185,-8.921) (-8.116,-8.892) (-8.047,-8.862) (-7.978,-8.832) (-7.909,-8.802) (-7.841,-8.771) (-7.772,-8.741) (-7.704,-8.709) (-7.636,-8.678) (-7.568,-8.646) (-7.500,-8.614) (-7.432,-8.582) (-7.365,-8.549) (-7.298,-8.516) (-7.230,-8.483) (-7.163,-8.449) (-7.096,-8.415) (-7.030,-8.381) (-6.963,-8.347) (-6.897,-8.312) (-6.830,-8.277) (-6.764,-8.241) (-6.698,-8.206) (-6.633,-8.170) (-6.567,-8.133) (-6.501,-8.097) (-6.436,-8.060) (-6.371,-8.023) (-6.306,-7.985) (-6.241,-7.947) (-6.177,-7.909) (-6.112,-7.871) (-6.048,-7.832) (-5.984,-7.793) (-5.920,-7.754) (-5.857,-7.714) (-5.793,-7.674) (-5.730,-7.634) (-5.667,-7.593) (-5.604,-7.553) (-5.541,-7.512) (-5.478,-7.470) (-5.416,-7.429) (-5.354,-7.387) (-5.292,-7.344) (-5.230,-7.302) (-5.168,-7.259) (-5.107,-7.216) (-5.046,-7.173) (-4.985,-7.129) (-4.924,-7.085) (-4.864,-7.041) (-4.803,-6.996) (-4.743,-6.951) (-4.683,-6.906) (-4.624,-6.861) (-4.564,-6.815) (-4.505,-6.769) (-4.446,-6.723) (-4.387,-6.676) (-4.328,-6.629) (-4.270,-6.582) (-4.212,-6.535) (-4.154,-6.487) (-4.096,-6.439) (-4.039,-6.391) (-3.982,-6.342) (-3.925,-6.294) (-3.868,-6.245) (-3.812,-6.195) (-3.755,-6.146) (-3.699,-6.096) (-3.644,-6.045)}
\LC{(-2.919,-10.000) (-2.888,-9.932) (-2.857,-9.864) (-2.826,-9.795) (-2.795,-9.727) (-2.764,-9.658) (-2.734,-9.590) (-2.703,-9.521) (-2.673,-9.453) (-2.642,-9.384) (-2.612,-9.316) (-2.582,-9.247) (-2.552,-9.178) (-2.523,-9.109) (-2.493,-9.040) (-2.463,-8.971) (-2.434,-8.902) (-2.405,-8.833) (-2.375,-8.764) (-2.346,-8.695) (-2.317,-8.626) (-2.288,-8.557) (-2.260,-8.488) (-2.231,-8.418) (-2.203,-8.349) (-2.174,-8.279) (-2.146,-8.210) (-2.118,-8.140) (-2.090,-8.071) (-2.062,-8.001) (-2.034,-7.932) (-2.006,-7.862) (-1.979,-7.792) (-1.951,-7.722) (-1.924,-7.652) (-1.897,-7.583) (-1.870,-7.513) (-1.843,-7.443) (-1.816,-7.373) (-1.790,-7.302) (-1.763,-7.232) (-1.737,-7.162) (-1.710,-7.092) (-1.684,-7.021) (-1.658,-6.951) (-1.632,-6.881) (-1.607,-6.810) (-1.581,-6.740) (-1.556,-6.669) (-1.530,-6.599) (-1.505,-6.528) (-1.480,-6.457) (-1.455,-6.387) (-1.431,-6.316) (-1.406,-6.245) (-1.381,-6.174) (-1.357,-6.103) (-1.333,-6.032) (-1.309,-5.961) (-1.285,-5.890) (-1.261,-5.819) (-1.238,-5.748) (-1.214,-5.676) (-1.191,-5.605) (-1.168,-5.534) (-1.145,-5.462) (-1.122,-5.391) (-1.099,-5.320) (-1.076,-5.248) (-1.054,-5.176) (-1.032,-5.105) (-1.010,-5.033) (-0.988,-4.961) (-0.966,-4.890) (-0.944,-4.818) (-0.923,-4.746) (-0.902,-4.674) (-0.880,-4.602) (-0.859,-4.530) (-0.839,-4.458) (-0.818,-4.386)}
\LC{(2.151,2.078) (2.223,2.099) (2.296,2.118) (2.368,2.135) (2.442,2.152) (2.515,2.166) (2.589,2.179) (2.663,2.191) (2.738,2.200) (2.812,2.209) (2.887,2.215) (2.962,2.220) (3.037,2.223) (3.112,2.225) (3.187,2.225) (3.262,2.223) (3.336,2.220) (3.411,2.215) (3.486,2.208) (3.560,2.199) (3.635,2.189) (3.709,2.176) (3.782,2.163) (3.856,2.147) (3.929,2.129) (4.001,2.110) (4.073,2.089) (4.144,2.066) (4.215,2.041) (4.285,2.014) (4.355,1.986) (4.423,1.955) (4.491,1.923) (4.558,1.889) (4.624,1.853) (4.688,1.815) (4.752,1.776) (4.815,1.734) (4.876,1.691) (4.936,1.646) (4.994,1.599) (5.051,1.550) (5.107,1.500) (5.160,1.447) (5.213,1.393) (5.263,1.338) (5.311,1.281) (5.358,1.222) (5.402,1.161) (5.445,1.099) (5.485,1.036) (5.523,0.971) (5.558,0.905) (5.592,0.838) (5.622,0.770) (5.651,0.700) (5.676,0.630) (5.699,0.558) (5.719,0.486) (5.736,0.413) (5.751,0.340) (5.763,0.266) (5.772,0.191) (5.778,0.116) (5.781,0.041) (5.781,-0.034) (5.778,-0.109) (5.772,-0.183) (5.764,-0.258) (5.752,-0.332) (5.738,-0.406) (5.721,-0.479) (5.701,-0.551) (5.679,-0.622) (5.653,-0.693) (5.625,-0.763) (5.595,-0.831) (5.562,-0.899) (5.527,-0.965) (5.489,-1.030) (5.449,-1.093) (5.407,-1.155) (5.363,-1.216) (5.316,-1.275) (5.268,-1.332) (5.218,-1.388) (5.166,-1.442) (5.112,-1.494) (5.057,-1.545) (5.000,-1.594) (4.942,-1.641) (4.882,-1.686) (4.821,-1.730) (4.759,-1.771) (4.695,-1.811) (4.630,-1.849) (4.565,-1.885) (4.498,-1.920) (4.430,-1.952) (4.362,-1.983) (4.293,-2.011) (4.222,-2.038) (4.152,-2.063) (4.080,-2.086) (4.009,-2.108) (3.936,-2.127) (3.863,-2.145) (3.790,-2.161) (3.716,-2.175) (3.642,-2.187) (3.568,-2.198) (3.494,-2.207) (3.419,-2.214) (3.344,-2.219) (3.269,-2.223) (3.194,-2.225) (3.119,-2.225) (3.044,-2.224) (2.969,-2.220) (2.895,-2.216) (2.820,-2.209) (2.745,-2.201) (2.671,-2.192) (2.597,-2.180) (2.523,-2.167) (2.449,-2.153) (2.376,-2.137) (2.303,-2.120) (2.230,-2.101) (2.158,-2.080)}
\LC{(-0.002,2.978) (-0.002,3.053) (-0.002,3.128) (-0.002,3.203) (-0.003,3.278) (-0.003,3.353) (-0.003,3.428) (-0.003,3.503) (-0.003,3.578) (-0.003,3.653) (-0.003,3.728) (-0.003,3.803) (-0.003,3.878) (-0.003,3.953) (-0.003,4.028) (-0.004,4.103) (-0.004,4.178) (-0.004,4.253) (-0.004,4.328) (-0.004,4.403) (-0.004,4.478) (-0.004,4.553) (-0.004,4.628) (-0.004,4.703) (-0.004,4.778) (-0.005,4.853) (-0.005,4.928) (-0.005,5.003) (-0.005,5.078) (-0.005,5.153) (-0.005,5.228) (-0.005,5.303) (-0.005,5.378) (-0.005,5.453) (-0.006,5.528) (-0.006,5.603) (-0.006,5.678) (-0.006,5.753) (-0.006,5.828) (-0.006,5.903) (-0.006,5.978) (-0.006,6.053) (-0.006,6.128) (-0.007,6.203) (-0.007,6.278) (-0.007,6.353) (-0.007,6.428) (-0.007,6.503) (-0.007,6.578) (-0.007,6.653) (-0.007,6.728) (-0.008,6.803) (-0.008,6.878) (-0.008,6.953) (-0.008,7.028) (-0.008,7.103) (-0.008,7.178) (-0.008,7.253) (-0.008,7.328) (-0.009,7.403) (-0.009,7.478) (-0.009,7.553) (-0.009,7.628) (-0.009,7.703) (-0.009,7.778) (-0.009,7.853) (-0.010,7.928) (-0.010,8.003) (-0.010,8.078) (-0.010,8.153) (-0.010,8.228) (-0.010,8.303) (-0.010,8.378) (-0.011,8.453) (-0.011,8.528) (-0.011,8.603) (-0.011,8.678) (-0.011,8.753) (-0.011,8.828) (-0.011,8.903) (-0.012,8.978) (-0.012,9.053) (-0.012,9.128) (-0.012,9.203) (-0.012,9.278) (-0.012,9.353) (-0.012,9.428) (-0.013,9.503) (-0.013,9.578) (-0.013,9.653) (-0.013,9.728) (-0.013,9.803) (-0.013,9.878) (-0.013,9.953) (-0.014,10.000)}
\LC{(4.675,4.277) (4.748,4.294) (4.821,4.311) (4.894,4.327) (4.967,4.342) (5.041,4.357) (5.115,4.370) (5.189,4.383) (5.263,4.395) (5.337,4.406) (5.411,4.417) (5.485,4.427) (5.560,4.435) (5.634,4.444) (5.709,4.451) (5.784,4.457) (5.859,4.463) (5.933,4.468) (6.008,4.472) (6.083,4.475) (6.158,4.477) (6.233,4.479) (6.308,4.480) (6.383,4.480) (6.458,4.479) (6.533,4.477) (6.608,4.474) (6.683,4.471) (6.758,4.466) (6.833,4.461) (6.907,4.455) (6.982,4.448) (7.057,4.440) (7.131,4.432) (7.206,4.422) (7.280,4.412) (7.354,4.401) (7.428,4.388) (7.502,4.375) (7.576,4.361) (7.649,4.347) (7.722,4.331) (7.796,4.314) (7.869,4.297) (7.941,4.278) (8.014,4.259) (8.086,4.239) (8.158,4.218) (8.230,4.196) (8.301,4.173) (8.372,4.149) (8.443,4.124) (8.513,4.099) (8.584,4.072) (8.653,4.045) (8.723,4.016) (8.792,3.987) (8.860,3.957) (8.929,3.926) (8.996,3.893) (9.064,3.860) (9.131,3.827) (9.197,3.792) (9.263,3.756) (9.328,3.719) (9.393,3.682) (9.458,3.643) (9.521,3.604) (9.585,3.563) (9.647,3.522) (9.709,3.480) (9.771,3.437) (9.831,3.393) (9.891,3.348) (9.951,3.302) (10.000,3.255)}
\LC{(-1.771,2.911) (-1.827,2.961) (-1.884,3.010) (-1.941,3.058) (-1.999,3.106) (-2.057,3.154) (-2.116,3.200) (-2.175,3.246) (-2.234,3.292) (-2.295,3.337) (-2.355,3.381) (-2.416,3.425) (-2.477,3.468) (-2.539,3.511) (-2.601,3.552) (-2.664,3.594) (-2.727,3.634) (-2.790,3.674) (-2.854,3.714) (-2.918,3.753) (-2.983,3.791) (-3.048,3.828) (-3.113,3.865) (-3.179,3.902) (-3.245,3.937) (-3.311,3.972) (-3.378,4.007) (-3.445,4.040) (-3.512,4.074) (-3.580,4.106) (-3.648,4.138) (-3.716,4.169) (-3.784,4.199) (-3.853,4.229) (-3.922,4.258) (-3.992,4.287) (-4.061,4.315) (-4.131,4.342) (-4.201,4.369) (-4.272,4.394) (-4.342,4.420) (-4.413,4.444) (-4.484,4.468) (-4.556,4.491) (-4.627,4.514) (-4.699,4.535) (-4.771,4.556) (-4.843,4.577) (-4.915,4.596) (-4.988,4.615) (-5.061,4.634) (-5.134,4.651) (-5.207,4.668) (-5.280,4.684) (-5.353,4.700) (-5.427,4.715) (-5.501,4.729) (-5.574,4.742) (-5.648,4.754) (-5.722,4.766) (-5.797,4.777) (-5.871,4.788) (-5.945,4.798) (-6.020,4.806) (-6.094,4.815) (-6.169,4.822) (-6.244,4.829) (-6.318,4.835) (-6.393,4.840) (-6.468,4.844) (-6.543,4.848) (-6.618,4.851) (-6.693,4.853) (-6.768,4.855) (-6.843,4.855) (-6.918,4.855) (-6.993,4.854) (-7.068,4.852) (-7.143,4.850) (-7.218,4.847) (-7.293,4.843) (-7.367,4.838) (-7.442,4.832) (-7.517,4.826) (-7.592,4.819) (-7.666,4.811) (-7.741,4.802) (-7.815,4.792) (-7.889,4.782) (-7.963,4.770) (-8.037,4.758) (-8.111,4.746) (-8.185,4.732) (-8.259,4.717) (-8.332,4.702) (-8.405,4.686) (-8.478,4.669) (-8.551,4.651) (-8.624,4.633) (-8.696,4.613) (-8.769,4.593) (-8.840,4.572) (-8.912,4.550) (-8.984,4.527) (-9.055,4.503) (-9.126,4.479) (-9.196,4.454) (-9.267,4.427) (-9.337,4.400) (-9.406,4.372) (-9.475,4.344) (-9.544,4.314) (-9.613,4.284) (-9.681,4.253) (-9.749,4.220) (-9.816,4.187) (-9.883,4.154) (-9.950,4.119) (-10.000,4.084)}
\LC{(2.995,4.844) (3.051,4.893) (3.108,4.942) (3.166,4.990) (3.223,5.038) (3.281,5.086) (3.339,5.133) (3.398,5.180) (3.456,5.227) (3.516,5.273) (3.575,5.319) (3.635,5.364) (3.694,5.410) (3.755,5.454) (3.815,5.499) (3.876,5.543) (3.937,5.586) (3.998,5.630) (4.060,5.673) (4.121,5.715) (4.183,5.757) (4.246,5.799) (4.308,5.840) (4.371,5.881) (4.434,5.922) (4.497,5.962) (4.561,6.002) (4.624,6.042) (4.688,6.081) (4.753,6.120) (4.817,6.158) (4.882,6.196) (4.947,6.234) (5.012,6.271) (5.077,6.308) (5.143,6.344) (5.208,6.380) (5.274,6.416) (5.341,6.451) (5.407,6.486) (5.474,6.521) (5.540,6.555) (5.607,6.588) (5.675,6.622) (5.742,6.655) (5.809,6.687) (5.877,6.719) (5.945,6.751) (6.013,6.782) (6.082,6.813) (6.150,6.844) (6.219,6.874) (6.288,6.904) (6.357,6.933) (6.426,6.962) (6.495,6.991) (6.565,7.019) (6.634,7.046) (6.704,7.074) (6.774,7.101) (6.845,7.127) (6.915,7.153) (6.985,7.179) (7.056,7.204) (7.127,7.229) (7.198,7.253) (7.269,7.277) (7.340,7.301) (7.411,7.324) (7.483,7.347) (7.554,7.369) (7.626,7.391) (7.698,7.413) (7.770,7.434) (7.842,7.454) (7.914,7.474) (7.987,7.494) (8.059,7.514) (8.132,7.533) (8.204,7.551) (8.277,7.569) (8.350,7.587) (8.423,7.604) (8.496,7.621) (8.569,7.637) (8.642,7.653) (8.716,7.668) (8.789,7.683) (8.863,7.698) (8.937,7.712) (9.010,7.726) (9.084,7.739) (9.158,7.752) (9.232,7.764) (9.306,7.776) (9.380,7.788) (9.454,7.799) (9.529,7.810) (9.603,7.820) (9.677,7.829) (9.752,7.839) (9.826,7.847) (9.901,7.856) (9.975,7.864) (10.000,7.866)}
\LC{(-1.118,1.217) (-1.188,1.245) (-1.258,1.270) (-1.330,1.293) (-1.402,1.314) (-1.474,1.332) (-1.548,1.348) (-1.622,1.362) (-1.696,1.372) (-1.770,1.380) (-1.845,1.386) (-1.920,1.389) (-1.995,1.389) (-2.070,1.387) (-2.145,1.382) (-2.219,1.374) (-2.294,1.363) (-2.367,1.349) (-2.440,1.333) (-2.513,1.313) (-2.585,1.291) (-2.655,1.266) (-2.725,1.238) (-2.793,1.207) (-2.860,1.172) (-2.925,1.135) (-2.988,1.095) (-3.050,1.052) (-3.109,1.006) (-3.166,0.958) (-3.220,0.906) (-3.272,0.852) (-3.321,0.795) (-3.366,0.735) (-3.409,0.673) (-3.447,0.609) (-3.482,0.543) (-3.514,0.474) (-3.541,0.405) (-3.563,0.333) (-3.582,0.260) (-3.596,0.187) (-3.605,0.112) (-3.609,0.037) (-3.609,-0.038) (-3.605,-0.112) (-3.596,-0.187) (-3.582,-0.261) (-3.563,-0.333) (-3.540,-0.405) (-3.513,-0.475) (-3.482,-0.543) (-3.447,-0.609) (-3.409,-0.673) (-3.366,-0.735) (-3.321,-0.795) (-3.272,-0.852) (-3.220,-0.906) (-3.166,-0.958) (-3.109,-1.006) (-3.049,-1.052) (-2.988,-1.095) (-2.925,-1.135) (-2.859,-1.172) (-2.793,-1.207) (-2.724,-1.238) (-2.655,-1.266) (-2.584,-1.291) (-2.513,-1.313) (-2.440,-1.333) (-2.367,-1.349) (-2.293,-1.363) (-2.219,-1.374) (-2.145,-1.382) (-2.070,-1.387) (-1.995,-1.389) (-1.920,-1.389) (-1.845,-1.386) (-1.770,-1.380) (-1.696,-1.372) (-1.621,-1.361) (-1.548,-1.348) (-1.474,-1.332) (-1.401,-1.314) (-1.329,-1.293) (-1.258,-1.270) (-1.188,-1.245) (-1.118,-1.217)}
\LC{(-10.000,-2.969) (-9.946,-3.019) (-9.890,-3.069) (-9.833,-3.118) (-9.775,-3.165) (-9.716,-3.212) (-9.657,-3.258) (-9.597,-3.303) (-9.537,-3.347) (-9.475,-3.391) (-9.413,-3.433) (-9.351,-3.474) (-9.287,-3.515) (-9.224,-3.554) (-9.159,-3.593) (-9.094,-3.630) (-9.029,-3.667) (-8.963,-3.702) (-8.897,-3.737) (-8.830,-3.771) (-8.762,-3.804) (-8.694,-3.836) (-8.626,-3.867) (-8.557,-3.897) (-8.488,-3.926) (-8.419,-3.954) (-8.349,-3.981) (-8.278,-4.007) (-8.208,-4.033) (-8.137,-4.057) (-8.066,-4.080) (-7.994,-4.103) (-7.922,-4.124) (-7.850,-4.145) (-7.778,-4.165) (-7.705,-4.184) (-7.632,-4.201) (-7.559,-4.218) (-7.486,-4.234) (-7.413,-4.249) (-7.339,-4.264) (-7.265,-4.277) (-7.191,-4.289) (-7.117,-4.301) (-7.043,-4.311) (-6.968,-4.321) (-6.894,-4.330) (-6.819,-4.338) (-6.745,-4.345) (-6.670,-4.351) (-6.595,-4.356) (-6.520,-4.361) (-6.445,-4.364) (-6.370,-4.367) (-6.295,-4.368) (-6.220,-4.369) (-6.145,-4.369) (-6.070,-4.369) (-5.995,-4.367) (-5.920,-4.364) (-5.846,-4.361) (-5.771,-4.357) (-5.696,-4.352) (-5.621,-4.346) (-5.546,-4.339) (-5.472,-4.332) (-5.397,-4.323) (-5.323,-4.314) (-5.248,-4.304) (-5.174,-4.293) (-5.100,-4.282) (-5.051,-4.273)}
\LC{(-2.119,1.838) (-2.193,1.851) (-2.267,1.863) (-2.341,1.872) (-2.416,1.880) (-2.491,1.886) (-2.566,1.890) (-2.641,1.892) (-2.716,1.892) (-2.791,1.890) (-2.866,1.885) (-2.940,1.879) (-3.015,1.871) (-3.089,1.861) (-3.163,1.849) (-3.237,1.834) (-3.310,1.818) (-3.383,1.799) (-3.455,1.778) (-3.526,1.756) (-3.597,1.730) (-3.667,1.703) (-3.736,1.674) (-3.804,1.642) (-3.871,1.608) (-3.936,1.573) (-4.001,1.534) (-4.064,1.494) (-4.126,1.452) (-4.186,1.407) (-4.245,1.360) (-4.302,1.311) (-4.357,1.261) (-4.410,1.208) (-4.461,1.153) (-4.510,1.096) (-4.557,1.037) (-4.601,0.977) (-4.643,0.914) (-4.682,0.850) (-4.718,0.785) (-4.752,0.718) (-4.782,0.649) (-4.810,0.580) (-4.835,0.509) (-4.856,0.437) (-4.874,0.364) (-4.889,0.291) (-4.901,0.216) (-4.909,0.142) (-4.914,0.067) (-4.915,-0.008) (-4.913,-0.083) (-4.908,-0.158) (-4.899,-0.232) (-4.887,-0.306) (-4.871,-0.379) (-4.852,-0.452) (-4.830,-0.524) (-4.805,-0.594) (-4.776,-0.664) (-4.745,-0.732) (-4.711,-0.799) (-4.674,-0.864) (-4.634,-0.927) (-4.592,-0.989) (-4.547,-1.050) (-4.500,-1.108) (-4.451,-1.164) (-4.399,-1.219) (-4.346,-1.271) (-4.290,-1.322) (-4.233,-1.370) (-4.174,-1.417) (-4.113,-1.461) (-4.051,-1.503) (-3.987,-1.543) (-3.923,-1.580) (-3.857,-1.616) (-3.789,-1.649) (-3.721,-1.680) (-3.652,-1.709) (-3.582,-1.736) (-3.511,-1.761) (-3.440,-1.783) (-3.367,-1.803) (-3.295,-1.821) (-3.221,-1.838) (-3.148,-1.851) (-3.074,-1.863) (-2.999,-1.873) (-2.925,-1.881) (-2.850,-1.887) (-2.775,-1.890) (-2.700,-1.892) (-2.625,-1.891) (-2.550,-1.889) (-2.475,-1.885) (-2.400,-1.879) (-2.326,-1.871) (-2.251,-1.861) (-2.177,-1.849) (-2.128,-1.840)}
\LC{(-4.280,5.092) (-4.347,5.126) (-4.414,5.159) (-4.482,5.192) (-4.549,5.225) (-4.617,5.257) (-4.685,5.288) (-4.754,5.319) (-4.822,5.349) (-4.891,5.379) (-4.960,5.409) (-5.029,5.437) (-5.099,5.466) (-5.168,5.494) (-5.238,5.521) (-5.308,5.548) (-5.379,5.574) (-5.449,5.600) (-5.520,5.625) (-5.590,5.650) (-5.661,5.674) (-5.733,5.697) (-5.804,5.720) (-5.876,5.743) (-5.947,5.765) (-6.019,5.786) (-6.091,5.807) (-6.163,5.827) (-6.236,5.847) (-6.308,5.866) (-6.381,5.885) (-6.454,5.903) (-6.527,5.921) (-6.600,5.938) (-6.673,5.954) (-6.746,5.970) (-6.819,5.986) (-6.893,6.001) (-6.967,6.015) (-7.040,6.028) (-7.114,6.042) (-7.188,6.054) (-7.262,6.066) (-7.336,6.078) (-7.410,6.088) (-7.485,6.099) (-7.559,6.108) (-7.634,6.117) (-7.708,6.126) (-7.783,6.134) (-7.857,6.141) (-7.932,6.148) (-8.007,6.154) (-8.082,6.160) (-8.156,6.165) (-8.231,6.169) (-8.306,6.173) (-8.381,6.176) (-8.456,6.179) (-8.531,6.181) (-8.606,6.182) (-8.681,6.183) (-8.756,6.183) (-8.831,6.183) (-8.906,6.182) (-8.981,6.180) (-9.056,6.178) (-9.131,6.175) (-9.206,6.172) (-9.281,6.168) (-9.356,6.163) (-9.430,6.158) (-9.505,6.152) (-9.580,6.145) (-9.654,6.138) (-9.729,6.130) (-9.804,6.122) (-9.878,6.113) (-9.952,6.103) (-10.000,6.096)}
\LC{(1.461,2.722) (1.513,2.776) (1.565,2.830) (1.618,2.883) (1.671,2.936) (1.725,2.988) (1.779,3.040) (1.834,3.091) (1.890,3.142) (1.945,3.192) (2.002,3.242) (2.058,3.291) (2.115,3.339) (2.173,3.387) (2.231,3.435) (2.290,3.482) (2.349,3.528) (2.408,3.574) (2.468,3.619) (2.528,3.664) (2.588,3.708) (2.649,3.752) (2.711,3.795) (2.773,3.837) (2.835,3.879) (2.897,3.921) (2.960,3.962) (3.024,4.002) (3.087,4.042) (3.151,4.081) (3.215,4.119) (3.280,4.157) (3.345,4.195) (3.410,4.232) (3.476,4.268) (3.542,4.304) (3.608,4.339) (3.675,4.373) (3.742,4.407) (3.809,4.441) (3.876,4.473) (3.944,4.506) (4.012,4.537) (4.080,4.568) (4.149,4.599) (4.218,4.628) (4.287,4.658) (4.356,4.686) (4.426,4.714) (4.496,4.741) (4.566,4.768) (4.636,4.794) (4.706,4.820) (4.777,4.845) (4.848,4.869) (4.919,4.893) (4.991,4.916) (5.062,4.938) (5.134,4.960) (5.206,4.981) (5.278,5.002) (5.350,5.022) (5.423,5.041) (5.496,5.060) (5.568,5.078) (5.641,5.095) (5.714,5.112) (5.788,5.128) (5.861,5.143) (5.935,5.158) (6.008,5.172) (6.082,5.185) (6.156,5.198) (6.230,5.210) (6.304,5.222) (6.378,5.233) (6.453,5.243) (6.527,5.252) (6.602,5.261) (6.676,5.269) (6.751,5.276) (6.825,5.283) (6.900,5.289) (6.975,5.294) (7.050,5.299) (7.125,5.303) (7.200,5.306) (7.275,5.309) (7.350,5.311) (7.425,5.312) (7.500,5.312) (7.575,5.312) (7.650,5.311) (7.725,5.310) (7.800,5.307) (7.874,5.304) (7.949,5.300) (8.024,5.296) (8.099,5.291) (8.174,5.285) (8.249,5.278) (8.323,5.271) (8.398,5.262) (8.472,5.254) (8.547,5.244) (8.621,5.234) (8.695,5.223) (8.769,5.211) (8.843,5.198) (8.917,5.185) (8.990,5.171) (9.064,5.156) (9.137,5.140) (9.211,5.124) (9.284,5.107) (9.356,5.089) (9.429,5.071) (9.502,5.051) (9.574,5.031) (9.646,5.010) (9.718,4.989) (9.789,4.966) (9.861,4.943) (9.932,4.919) (10.000,4.894)}
\LC{(1.369,4.863) (1.399,4.932) (1.429,5.000) (1.460,5.069) (1.491,5.137) (1.522,5.205) (1.553,5.273) (1.585,5.341) (1.617,5.409) (1.649,5.477) (1.681,5.545) (1.713,5.613) (1.746,5.680) (1.779,5.748) (1.812,5.815) (1.845,5.882) (1.878,5.949) (1.912,6.017) (1.945,6.084) (1.979,6.150) (2.013,6.217) (2.048,6.284) (2.082,6.350) (2.117,6.417) (2.152,6.483) (2.187,6.550) (2.222,6.616) (2.258,6.682) (2.293,6.748) (2.329,6.814) (2.365,6.880) (2.401,6.945) (2.437,7.011) (2.474,7.077) (2.510,7.142) (2.547,7.207) (2.584,7.273) (2.621,7.338) (2.659,7.403) (2.696,7.468) (2.734,7.533) (2.772,7.597) (2.810,7.662) (2.848,7.727) (2.886,7.791) (2.925,7.855) (2.964,7.920) (3.002,7.984) (3.041,8.048) (3.081,8.112) (3.120,8.176) (3.159,8.239) (3.199,8.303) (3.239,8.367) (3.279,8.430) (3.319,8.494) (3.359,8.557) (3.399,8.620) (3.440,8.683) (3.481,8.746) (3.522,8.809) (3.563,8.872) (3.604,8.934) (3.645,8.997) (3.687,9.059) (3.728,9.122) (3.770,9.184) (3.812,9.246) (3.854,9.308) (3.896,9.370) (3.939,9.432) (3.981,9.494) (4.024,9.556) (4.067,9.617) (4.110,9.679) (4.153,9.740) (4.196,9.802) (4.239,9.863) (4.283,9.924) (4.326,9.985) (4.341,10.000)}
\LC{(-1.407,4.702) (-1.439,4.770) (-1.472,4.837) (-1.504,4.905) (-1.537,4.973) (-1.570,5.040) (-1.603,5.107) (-1.636,5.174) (-1.670,5.241) (-1.704,5.308) (-1.738,5.375) (-1.772,5.442) (-1.806,5.508) (-1.841,5.575) (-1.876,5.641) (-1.911,5.708) (-1.947,5.774) (-1.982,5.840) (-2.018,5.906) (-2.054,5.972) (-2.090,6.037) (-2.126,6.103) (-2.163,6.168) (-2.200,6.234) (-2.237,6.299) (-2.274,6.364) (-2.311,6.429) (-2.349,6.494) (-2.386,6.559) (-2.424,6.624) (-2.463,6.688) (-2.501,6.753) (-2.539,6.817) (-2.578,6.881) (-2.617,6.945) (-2.656,7.009) (-2.695,7.073) (-2.735,7.137) (-2.774,7.201) (-2.814,7.264) (-2.854,7.328) (-2.894,7.391) (-2.935,7.454) (-2.975,7.517) (-3.016,7.580) (-3.057,7.643) (-3.098,7.706) (-3.139,7.769) (-3.180,7.831) (-3.222,7.894) (-3.264,7.956) (-3.306,8.018) (-3.348,8.080) (-3.390,8.142) (-3.432,8.204) (-3.475,8.266) (-3.518,8.327) (-3.561,8.389) (-3.604,8.450) (-3.647,8.512) (-3.690,8.573) (-3.734,8.634) (-3.778,8.695) (-3.822,8.756) (-3.866,8.816) (-3.910,8.877) (-3.954,8.937) (-3.999,8.998) (-4.043,9.058) (-4.088,9.118) (-4.133,9.178) (-4.178,9.238) (-4.224,9.298) (-4.269,9.357) (-4.315,9.417) (-4.360,9.476) (-4.406,9.536) (-4.452,9.595) (-4.499,9.654) (-4.545,9.713) (-4.591,9.772) (-4.638,9.831) (-4.685,9.889) (-4.732,9.948) (-4.779,10.000)}
\LC{(0.449,3.419) (0.464,3.492) (0.479,3.566) (0.494,3.639) (0.509,3.713) (0.524,3.786) (0.540,3.860) (0.556,3.933) (0.571,4.006) (0.587,4.080) (0.603,4.153) (0.619,4.226) (0.636,4.299) (0.652,4.373) (0.669,4.446) (0.686,4.519) (0.703,4.592) (0.720,4.665) (0.737,4.738) (0.754,4.811) (0.772,4.884) (0.789,4.957) (0.807,5.030) (0.825,5.102) (0.843,5.175) (0.861,5.248) (0.879,5.321) (0.897,5.394) (0.916,5.466) (0.934,5.539) (0.953,5.612) (0.972,5.684) (0.991,5.757) (1.010,5.829) (1.029,5.902) (1.048,5.974) (1.068,6.047) (1.087,6.119) (1.107,6.191) (1.127,6.264) (1.147,6.336) (1.167,6.408) (1.187,6.481) (1.207,6.553) (1.227,6.625) (1.248,6.697) (1.268,6.769) (1.289,6.841) (1.310,6.913) (1.331,6.986) (1.352,7.057) (1.373,7.129) (1.394,7.201) (1.415,7.273) (1.437,7.345) (1.458,7.417) (1.480,7.489) (1.502,7.561) (1.524,7.632) (1.546,7.704) (1.568,7.776) (1.590,7.847) (1.612,7.919) (1.635,7.991) (1.657,8.062) (1.680,8.134) (1.702,8.205) (1.725,8.277) (1.748,8.348) (1.771,8.419) (1.794,8.491) (1.817,8.562) (1.841,8.633) (1.864,8.705) (1.887,8.776) (1.911,8.847) (1.935,8.918) (1.958,8.989) (1.982,9.060) (2.006,9.132) (2.030,9.203) (2.054,9.274) (2.079,9.345) (2.103,9.416) (2.127,9.486) (2.152,9.557) (2.176,9.628) (2.201,9.699) (2.226,9.770) (2.251,9.841) (2.276,9.911) (2.301,9.982) (2.309,10.000)}
\LC{(10.000,-9.013) (9.929,-8.996) (9.856,-8.978) (9.784,-8.960) (9.711,-8.942) (9.638,-8.924) (9.565,-8.905) (9.493,-8.886) (9.421,-8.866) (9.348,-8.846) (9.276,-8.826) (9.204,-8.805) (9.132,-8.784) (9.060,-8.763) (8.988,-8.741) (8.917,-8.719) (8.845,-8.696) (8.774,-8.674) (8.702,-8.650) (8.631,-8.627) (8.560,-8.603) (8.489,-8.579) (8.418,-8.554) (8.347,-8.529) (8.277,-8.504) (8.206,-8.478) (8.136,-8.453) (8.066,-8.426) (7.996,-8.400) (7.926,-8.373) (7.856,-8.345) (7.786,-8.318) (7.717,-8.289) (7.647,-8.261) (7.578,-8.232) (7.509,-8.203) (7.440,-8.174) (7.371,-8.144) (7.302,-8.114) (7.234,-8.084) (7.165,-8.053) (7.097,-8.022) (7.029,-7.990) (6.961,-7.959) (6.893,-7.927) (6.825,-7.894) (6.758,-7.861) (6.691,-7.828) (6.624,-7.795) (6.557,-7.761) (6.490,-7.727) (6.423,-7.692) (6.357,-7.658) (6.290,-7.623) (6.224,-7.587) (6.158,-7.551) (6.093,-7.515) (6.027,-7.479) (5.962,-7.442) (5.897,-7.405) (5.832,-7.368) (5.767,-7.330) (5.702,-7.292) (5.638,-7.254) (5.573,-7.215) (5.509,-7.176) (5.445,-7.137) (5.382,-7.097) (5.318,-7.057) (5.255,-7.017) (5.192,-6.976) (5.129,-6.935) (5.067,-6.894) (5.004,-6.852) (4.942,-6.810) (4.880,-6.768) (4.818,-6.726) (4.757,-6.683) (4.716,-6.654)}
\LC{(4.277,5.274) (4.342,5.310) (4.408,5.346) (4.475,5.381) (4.541,5.415) (4.608,5.449) (4.675,5.483) (4.742,5.516) (4.810,5.549) (4.878,5.581) (4.946,5.613) (5.014,5.644) (5.082,5.675) (5.151,5.705) (5.220,5.735) (5.289,5.764) (5.358,5.793) (5.427,5.822) (5.497,5.850) (5.567,5.877) (5.637,5.904) (5.707,5.930) (5.777,5.956) (5.848,5.982) (5.919,6.006) (5.989,6.031) (6.061,6.055) (6.132,6.078) (6.203,6.101) (6.275,6.123) (6.347,6.145) (6.418,6.167) (6.490,6.188) (6.563,6.208) (6.635,6.228) (6.707,6.247) (6.780,6.266) (6.853,6.284) (6.926,6.302) (6.999,6.319) (7.072,6.336) (7.145,6.352) (7.218,6.368) (7.292,6.383) (7.365,6.398) (7.439,6.412) (7.513,6.425) (7.587,6.439) (7.661,6.451) (7.735,6.463) (7.809,6.474) (7.883,6.485) (7.957,6.496) (8.032,6.506) (8.106,6.515) (8.180,6.524) (8.255,6.532) (8.330,6.539) (8.404,6.547) (8.479,6.553) (8.554,6.559) (8.629,6.565) (8.703,6.569) (8.778,6.574) (8.853,6.578) (8.928,6.581) (9.003,6.583) (9.078,6.586) (9.153,6.587) (9.228,6.588) (9.303,6.588) (9.378,6.588) (9.453,6.587) (9.528,6.586) (9.603,6.584) (9.678,6.582) (9.753,6.579) (9.828,6.575) (9.903,6.571) (9.977,6.566) (10.000,6.564)}
\LC{(0.444,6.194) (0.452,6.268) (0.460,6.343) (0.468,6.418) (0.477,6.492) (0.485,6.567) (0.493,6.641) (0.502,6.716) (0.510,6.790) (0.518,6.865) (0.527,6.939) (0.535,7.014) (0.544,7.088) (0.553,7.163) (0.561,7.237) (0.570,7.312) (0.579,7.386) (0.588,7.461) (0.596,7.535) (0.605,7.610) (0.614,7.684) (0.623,7.759) (0.632,7.833) (0.641,7.908) (0.650,7.982) (0.660,8.056) (0.669,8.131) (0.678,8.205) (0.687,8.280) (0.697,8.354) (0.706,8.429) (0.715,8.503) (0.725,8.577) (0.734,8.652) (0.744,8.726) (0.753,8.801) (0.763,8.875) (0.773,8.949) (0.782,9.024) (0.792,9.098) (0.802,9.172) (0.812,9.247) (0.822,9.321) (0.831,9.395) (0.841,9.470) (0.851,9.544) (0.861,9.618) (0.871,9.693) (0.881,9.767) (0.892,9.841) (0.902,9.916) (0.912,9.990) (0.915,10.000)}
\LC{(-0.458,5.643) (-0.467,5.717) (-0.476,5.792) (-0.485,5.866) (-0.495,5.941) (-0.504,6.015) (-0.513,6.090) (-0.523,6.164) (-0.532,6.238) (-0.542,6.313) (-0.552,6.387) (-0.561,6.461) (-0.571,6.536) (-0.581,6.610) (-0.591,6.685) (-0.601,6.759) (-0.611,6.833) (-0.621,6.908) (-0.631,6.982) (-0.641,7.056) (-0.651,7.130) (-0.661,7.205) (-0.672,7.279) (-0.682,7.353) (-0.692,7.428) (-0.703,7.502) (-0.713,7.576) (-0.724,7.650) (-0.734,7.725) (-0.745,7.799) (-0.756,7.873) (-0.767,7.947) (-0.777,8.022) (-0.788,8.096) (-0.799,8.170) (-0.810,8.244) (-0.821,8.318) (-0.832,8.392) (-0.843,8.467) (-0.854,8.541) (-0.866,8.615) (-0.877,8.689) (-0.888,8.763) (-0.900,8.837) (-0.911,8.912) (-0.922,8.986) (-0.934,9.060) (-0.945,9.134) (-0.957,9.208) (-0.969,9.282) (-0.980,9.356) (-0.992,9.430) (-1.004,9.504) (-1.016,9.578) (-1.028,9.652) (-1.039,9.726) (-1.051,9.801) (-1.063,9.875) (-1.076,9.949) (-1.088,10.000)}
\LC{(1.579,1.049) (1.654,1.043) (1.728,1.033) (1.802,1.020) (1.875,1.002) (1.947,0.981) (2.017,0.956) (2.086,0.927) (2.154,0.894) (2.219,0.857) (2.282,0.816) (2.342,0.771) (2.399,0.723) (2.453,0.671) (2.503,0.615) (2.549,0.556) (2.591,0.493) (2.628,0.428) (2.660,0.360) (2.686,0.290) (2.707,0.218) (2.722,0.144) (2.731,0.070) (2.733,-0.005) (2.730,-0.080) (2.720,-0.154) (2.704,-0.228) (2.683,-0.299) (2.656,-0.369) (2.623,-0.437) (2.586,-0.502) (2.544,-0.564) (2.497,-0.622) (2.446,-0.678) (2.392,-0.729) (2.334,-0.777) (2.274,-0.822) (2.210,-0.862) (2.145,-0.898) (2.077,-0.931) (2.008,-0.959) (1.937,-0.984) (1.865,-1.005) (1.792,-1.022) (1.718,-1.035) (1.644,-1.044) (1.594,-1.048)}
\LC{(1.474,-10.000) (1.457,-9.936) (1.441,-9.863) (1.425,-9.789) (1.409,-9.716) (1.393,-9.643) (1.377,-9.570) (1.361,-9.496) (1.345,-9.423) (1.329,-9.350) (1.313,-9.276) (1.298,-9.203) (1.282,-9.130) (1.266,-9.056) (1.251,-8.983) (1.235,-8.910) (1.220,-8.836) (1.205,-8.763) (1.189,-8.689) (1.174,-8.616) (1.159,-8.542) (1.144,-8.469) (1.129,-8.395) (1.114,-8.322) (1.099,-8.248) (1.084,-8.175) (1.070,-8.101) (1.055,-8.028) (1.040,-7.954) (1.026,-7.881) (1.012,-7.807) (0.997,-7.734) (0.983,-7.660) (0.969,-7.586) (0.954,-7.513) (0.940,-7.439) (0.926,-7.365) (0.912,-7.292) (0.898,-7.218) (0.884,-7.144) (0.871,-7.070) (0.857,-6.997) (0.843,-6.923) (0.830,-6.849) (0.816,-6.775) (0.803,-6.702) (0.790,-6.628) (0.776,-6.554) (0.763,-6.480) (0.750,-6.406) (0.737,-6.332) (0.724,-6.259) (0.711,-6.185) (0.699,-6.111) (0.686,-6.037) (0.673,-5.963) (0.661,-5.889) (0.648,-5.815) (0.636,-5.741) (0.623,-5.667) (0.611,-5.593) (0.599,-5.519) (0.587,-5.445) (0.575,-5.371) (0.571,-5.346)}
\LC{(-7.036,-10.000) (-6.975,-9.958) (-6.914,-9.914) (-6.853,-9.870) (-6.793,-9.826) (-6.732,-9.782) (-6.672,-9.737) (-6.612,-9.692) (-6.552,-9.647) (-6.492,-9.602) (-6.433,-9.556) (-6.373,-9.510) (-6.314,-9.464) (-6.255,-9.418) (-6.196,-9.372) (-6.137,-9.326) (-6.078,-9.279) (-6.020,-9.232) (-5.961,-9.185) (-5.903,-9.137) (-5.845,-9.090) (-5.787,-9.042) (-5.730,-8.994) (-5.672,-8.946) (-5.615,-8.898) (-5.558,-8.849) (-5.501,-8.800) (-5.444,-8.751) (-5.387,-8.702) (-5.331,-8.653) (-5.274,-8.603) (-5.218,-8.554) (-5.162,-8.504) (-5.107,-8.454) (-5.051,-8.403) (-4.996,-8.353) (-4.940,-8.302) (-4.885,-8.251) (-4.830,-8.200) (-4.776,-8.149) (-4.721,-8.097) (-4.667,-8.045) (-4.613,-7.993) (-4.559,-7.941) (-4.505,-7.889) (-4.452,-7.836) (-4.398,-7.784) (-4.345,-7.731) (-4.292,-7.678) (-4.239,-7.624) (-4.187,-7.571) (-4.134,-7.517) (-4.082,-7.463) (-4.030,-7.409) (-3.978,-7.355) (-3.927,-7.301) (-3.876,-7.246) (-3.824,-7.191) (-3.773,-7.136) (-3.723,-7.081) (-3.672,-7.025) (-3.622,-6.970) (-3.572,-6.914) (-3.522,-6.858) (-3.472,-6.802) (-3.423,-6.745) (-3.374,-6.689) (-3.325,-6.632) (-3.276,-6.575) (-3.227,-6.518) (-3.179,-6.461) (-3.131,-6.403) (-3.083,-6.345) (-3.035,-6.287) (-2.988,-6.229) (-2.941,-6.171) (-2.894,-6.113) (-2.847,-6.054) (-2.800,-5.995) (-2.754,-5.936) (-2.708,-5.877) (-2.662,-5.817) (-2.617,-5.758) (-2.572,-5.698) (-2.527,-5.638) (-2.482,-5.578) (-2.437,-5.517) (-2.393,-5.457) (-2.349,-5.396) (-2.320,-5.356)}
\LC{(10.000,-2.658) (9.952,-2.713) (9.900,-2.767) (9.847,-2.819) (9.793,-2.871) (9.738,-2.922) (9.682,-2.972) (9.625,-3.021) (9.568,-3.069) (9.509,-3.117) (9.450,-3.163) (9.391,-3.208) (9.330,-3.253) (9.269,-3.296) (9.207,-3.339) (9.145,-3.381) (9.082,-3.421) (9.018,-3.461) (8.954,-3.500) (8.889,-3.537) (8.824,-3.574) (8.758,-3.610) (8.691,-3.645) (8.625,-3.679) (8.557,-3.711) (8.489,-3.743) (8.421,-3.774) (8.352,-3.804) (8.283,-3.833) (8.214,-3.861) (8.144,-3.889) (8.073,-3.915) (8.003,-3.940) (7.932,-3.964) (7.860,-3.987) (7.789,-4.010) (7.717,-4.031) (7.645,-4.051) (7.572,-4.071) (7.500,-4.089) (7.427,-4.107) (7.354,-4.124) (7.280,-4.139) (7.207,-4.154) (7.133,-4.168) (7.059,-4.181) (6.985,-4.193) (6.911,-4.204) (6.837,-4.214) (6.762,-4.223) (6.688,-4.231) (6.613,-4.239) (6.538,-4.245) (6.463,-4.251) (6.389,-4.255) (6.314,-4.259) (6.239,-4.262) (6.164,-4.264) (6.089,-4.265) (6.014,-4.265) (5.939,-4.265) (5.864,-4.263) (5.789,-4.261) (5.714,-4.258) (5.639,-4.253) (5.564,-4.248) (5.489,-4.243) (5.415,-4.236) (5.340,-4.228) (5.266,-4.220) (5.191,-4.211) (5.117,-4.201) (5.043,-4.190) (4.969,-4.178) (4.895,-4.165) (4.821,-4.152) (4.747,-4.138) (4.674,-4.123) (4.600,-4.107) (4.527,-4.090) (4.454,-4.073) (4.382,-4.055) (4.309,-4.036) (4.237,-4.016) (4.165,-3.995) (4.093,-3.974) (4.021,-3.952) (3.950,-3.929) (3.879,-3.905) (3.808,-3.880) (3.737,-3.855) (3.667,-3.829) (3.597,-3.802) (3.527,-3.774) (3.458,-3.746) (3.388,-3.717) (3.320,-3.687) (3.251,-3.657) (3.183,-3.625) (3.115,-3.593) (3.048,-3.560) (2.981,-3.527) (2.914,-3.492) (2.848,-3.457) (2.782,-3.422) (2.716,-3.385) (2.651,-3.348) (2.586,-3.310) (2.522,-3.272)}
\LC{(-10.000,-4.265) (-9.942,-4.298) (-9.875,-4.330) (-9.806,-4.362) (-9.738,-4.392) (-9.669,-4.422) (-9.600,-4.451) (-9.530,-4.479) (-9.460,-4.506) (-9.390,-4.532) (-9.320,-4.558) (-9.249,-4.583) (-9.178,-4.606) (-9.106,-4.629) (-9.035,-4.652) (-8.963,-4.673) (-8.891,-4.694) (-8.818,-4.713) (-8.746,-4.732) (-8.673,-4.751) (-8.600,-4.768) (-8.527,-4.784) (-8.454,-4.800) (-8.380,-4.815) (-8.306,-4.829) (-8.233,-4.842) (-8.159,-4.855) (-8.085,-4.867) (-8.010,-4.878) (-7.936,-4.888) (-7.862,-4.897) (-7.787,-4.906) (-7.712,-4.913) (-7.638,-4.920) (-7.563,-4.926) (-7.488,-4.932) (-7.413,-4.936) (-7.339,-4.940) (-7.264,-4.943) (-7.189,-4.946) (-7.114,-4.947) (-7.039,-4.948) (-6.964,-4.948) (-6.889,-4.947) (-6.814,-4.946) (-6.739,-4.944) (-6.664,-4.941) (-6.589,-4.937) (-6.514,-4.933) (-6.439,-4.927) (-6.364,-4.921) (-6.290,-4.915) (-6.215,-4.907) (-6.140,-4.899) (-6.066,-4.890) (-5.992,-4.881) (-5.917,-4.870) (-5.843,-4.859) (-5.769,-4.848) (-5.695,-4.835) (-5.621,-4.822) (-5.548,-4.808) (-5.474,-4.794) (-5.401,-4.778)}
\LC{(-10.000,-7.079) (-9.926,-7.079) (-9.851,-7.078) (-9.776,-7.076) (-9.701,-7.074) (-9.626,-7.072) (-9.552,-7.069) (-9.477,-7.066) (-9.402,-7.062) (-9.327,-7.057) (-9.252,-7.052) (-9.177,-7.046) (-9.102,-7.040) (-9.028,-7.034) (-8.953,-7.027) (-8.878,-7.019) (-8.804,-7.011) (-8.729,-7.002) (-8.655,-6.993) (-8.581,-6.984) (-8.506,-6.973) (-8.432,-6.963) (-8.358,-6.952) (-8.284,-6.940) (-8.210,-6.928) (-8.136,-6.915) (-8.062,-6.902) (-7.988,-6.888) (-7.915,-6.874) (-7.841,-6.860) (-7.768,-6.845) (-7.694,-6.829) (-7.621,-6.813) (-7.548,-6.796) (-7.475,-6.779) (-7.402,-6.762) (-7.329,-6.744) (-7.256,-6.725) (-7.184,-6.706) (-7.111,-6.687) (-7.039,-6.667) (-6.967,-6.646) (-6.895,-6.625) (-6.823,-6.604) (-6.751,-6.582) (-6.680,-6.560) (-6.608,-6.537) (-6.537,-6.514) (-6.466,-6.490) (-6.395,-6.466) (-6.324,-6.441) (-6.253,-6.416) (-6.183,-6.391) (-6.112,-6.365) (-6.042,-6.338) (-5.972,-6.311) (-5.902,-6.284) (-5.833,-6.256) (-5.763,-6.227) (-5.694,-6.199) (-5.625,-6.169) (-5.556,-6.140) (-5.488,-6.110) (-5.419,-6.079) (-5.351,-6.048) (-5.283,-6.016) (-5.215,-5.985) (-5.147,-5.952) (-5.080,-5.919) (-5.013,-5.886)}
\LC{(6.534,-10.000) (6.476,-9.962) (6.418,-9.915) (6.360,-9.868) (6.302,-9.820) (6.244,-9.772) (6.186,-9.724) (6.129,-9.676) (6.072,-9.628) (6.014,-9.579) (5.957,-9.530) (5.901,-9.481) (5.844,-9.432) (5.787,-9.383) (5.731,-9.333) (5.675,-9.284) (5.619,-9.234) (5.563,-9.184) (5.507,-9.134) (5.452,-9.083) (5.396,-9.033) (5.341,-8.982) (5.286,-8.931) (5.231,-8.880) (5.177,-8.828) (5.122,-8.777) (5.068,-8.725) (5.013,-8.673) (4.959,-8.621) (4.906,-8.569) (4.852,-8.517) (4.799,-8.464) (4.745,-8.411) (4.692,-8.358) (4.639,-8.305) (4.586,-8.252) (4.534,-8.198) (4.482,-8.145) (4.429,-8.091) (4.377,-8.037) (4.325,-7.983) (4.274,-7.928) (4.222,-7.874) (4.171,-7.819) (4.120,-7.764) (4.069,-7.709) (4.019,-7.654) (3.968,-7.598) (3.918,-7.542) (3.868,-7.487) (3.818,-7.431) (3.768,-7.374) (3.719,-7.318) (3.669,-7.262) (3.620,-7.205) (3.571,-7.148) (3.523,-7.091) (3.474,-7.034) (3.426,-6.976) (3.378,-6.919) (3.330,-6.861) (3.282,-6.803) (3.235,-6.745) (3.188,-6.687) (3.141,-6.628) (3.094,-6.570) (3.047,-6.511) (3.001,-6.452) (2.955,-6.393) (2.909,-6.333) (2.863,-6.274) (2.818,-6.214) (2.773,-6.154) (2.728,-6.094) (2.683,-6.034) (2.639,-5.974) (2.594,-5.913) (2.550,-5.852) (2.521,-5.812)}
\LC{(-2.798,6.367) (-2.842,6.428) (-2.886,6.488) (-2.931,6.548) (-2.976,6.608) (-3.021,6.668) (-3.067,6.728) (-3.112,6.788) (-3.158,6.847) (-3.204,6.906) (-3.250,6.965) (-3.297,7.024) (-3.343,7.083) (-3.390,7.141) (-3.437,7.200) (-3.484,7.258) (-3.532,7.316) (-3.580,7.374) (-3.627,7.432) (-3.675,7.489) (-3.724,7.547) (-3.772,7.604) (-3.821,7.661) (-3.870,7.718) (-3.919,7.775) (-3.968,7.832) (-4.017,7.888) (-4.067,7.944) (-4.116,8.000) (-4.166,8.056) (-4.216,8.112) (-4.267,8.168) (-4.317,8.223) (-4.368,8.279) (-4.419,8.334) (-4.470,8.389) (-4.521,8.444) (-4.572,8.498) (-4.624,8.553) (-4.676,8.607) (-4.728,8.661) (-4.780,8.715) (-4.832,8.769) (-4.884,8.822) (-4.937,8.876) (-4.990,8.929) (-5.043,8.982) (-5.096,9.035) (-5.149,9.088) (-5.203,9.141) (-5.256,9.193) (-5.310,9.245) (-5.364,9.298) (-5.418,9.349) (-5.472,9.401) (-5.527,9.453) (-5.581,9.504) (-5.636,9.556) (-5.691,9.607) (-5.746,9.658) (-5.801,9.708) (-5.857,9.759) (-5.912,9.809) (-5.968,9.860) (-6.024,9.910) (-6.080,9.960) (-6.136,10.000)}
\LC{(-1.929,0.938) (-1.998,0.911) (-2.066,0.879) (-2.132,0.844) (-2.196,0.804) (-2.257,0.760) (-2.315,0.713) (-2.370,0.661) (-2.421,0.606) (-2.467,0.548) (-2.510,0.486) (-2.547,0.421) (-2.579,0.353) (-2.606,0.283) (-2.627,0.211) (-2.641,0.137) (-2.650,0.063) (-2.652,-0.012) (-2.648,-0.087) (-2.637,-0.161) (-2.621,-0.234) (-2.598,-0.306) (-2.570,-0.375) (-2.536,-0.442) (-2.497,-0.506) (-2.453,-0.567) (-2.405,-0.624) (-2.353,-0.678) (-2.297,-0.728) (-2.238,-0.775) (-2.176,-0.817) (-2.112,-0.855) (-2.045,-0.890) (-1.976,-0.920) (-1.930,-0.938)}
\LC{(-1.501,-10.000) (-1.484,-9.935) (-1.467,-9.862) (-1.451,-9.789) (-1.435,-9.716) (-1.418,-9.643) (-1.402,-9.569) (-1.386,-9.496) (-1.369,-9.423) (-1.353,-9.350) (-1.337,-9.276) (-1.321,-9.203) (-1.305,-9.130) (-1.290,-9.057) (-1.274,-8.983) (-1.258,-8.910) (-1.242,-8.837) (-1.227,-8.763) (-1.211,-8.690) (-1.196,-8.616) (-1.180,-8.543) (-1.165,-8.470) (-1.150,-8.396) (-1.135,-8.323) (-1.120,-8.249) (-1.104,-8.176) (-1.089,-8.102) (-1.075,-8.029) (-1.060,-7.955) (-1.045,-7.882) (-1.030,-7.808) (-1.016,-7.735) (-1.001,-7.661) (-0.986,-7.588) (-0.972,-7.514) (-0.958,-7.440) (-0.943,-7.367) (-0.929,-7.293) (-0.915,-7.219) (-0.901,-7.146) (-0.887,-7.072) (-0.873,-6.998) (-0.859,-6.925) (-0.845,-6.851) (-0.832,-6.777) (-0.818,-6.703) (-0.804,-6.630) (-0.791,-6.556) (-0.777,-6.482) (-0.764,-6.408) (-0.751,-6.334) (-0.738,-6.261) (-0.725,-6.187) (-0.712,-6.113) (-0.699,-6.039) (-0.686,-5.965) (-0.673,-5.891) (-0.660,-5.817) (-0.648,-5.743) (-0.635,-5.669) (-0.623,-5.596) (-0.610,-5.522) (-0.598,-5.448) (-0.586,-5.374) (-0.574,-5.300) (-0.562,-5.226) (-0.550,-5.152) (-0.538,-5.077) (-0.526,-5.003) (-0.522,-4.979)}
\LC{(-5.577,4.285) (-5.652,4.290) (-5.727,4.295) (-5.802,4.299) (-5.877,4.302) (-5.952,4.304) (-6.027,4.305) (-6.102,4.305) (-6.177,4.304) (-6.252,4.303) (-6.327,4.301) (-6.401,4.297) (-6.476,4.293) (-6.551,4.288) (-6.626,4.282) (-6.701,4.276) (-6.775,4.268) (-6.850,4.259) (-6.924,4.250) (-6.998,4.239) (-7.073,4.228) (-7.147,4.216) (-7.220,4.203) (-7.294,4.189) (-7.368,4.174) (-7.441,4.158) (-7.514,4.141) (-7.587,4.123) (-7.659,4.104) (-7.732,4.085) (-7.804,4.064) (-7.876,4.042) (-7.947,4.020) (-8.019,3.996) (-8.089,3.972) (-8.160,3.947) (-8.230,3.920) (-8.300,3.893) (-8.370,3.865) (-8.439,3.836) (-8.507,3.806) (-8.576,3.774) (-8.643,3.742) (-8.711,3.709) (-8.778,3.675) (-8.844,3.641) (-8.910,3.605) (-8.975,3.568) (-9.040,3.530) (-9.104,3.491) (-9.168,3.452) (-9.231,3.411) (-9.293,3.369) (-9.355,3.327) (-9.416,3.283) (-9.477,3.239) (-9.536,3.194) (-9.595,3.147) (-9.654,3.100) (-9.711,3.052) (-9.768,3.003) (-9.824,2.953) (-9.879,2.902) (-9.933,2.850) (-9.986,2.797) (-10.000,2.780)}

    % Dipolo puntual - puntas de flecha
% Generado por scripts/generar_datos.py (no editar a mano).
\FL{1.429}{-0.015}{-91.8}
\FL{1.467}{2.168}{37.2}
\FL{4.190}{3.194}{4.1}
\FL{6.989}{2.463}{-35.4}
\FL{8.330}{-0.035}{-90.7}
\FL{6.953}{-2.489}{-145.3}
\FL{4.146}{-3.191}{175.4}
\FL{1.433}{-2.141}{142.3}
\FL{-1.467}{2.168}{142.8}
\FL{-4.190}{3.194}{175.9}
\FL{-6.989}{2.463}{-144.6}
\FL{-8.330}{-0.035}{-89.3}
\FL{-6.953}{-2.489}{-34.7}
\FL{-4.146}{-3.191}{4.6}
\FL{-1.433}{-2.141}{37.7}
\FL{-1.429}{-0.015}{-88.2}
\FL{3.888}{-8.732}{126.6}
\FL{2.141}{-6.053}{119.5}
\FL{0.773}{-3.164}{110.7}
\FL{-4.014}{-8.749}{52.4}
\FL{-2.211}{-6.078}{59.7}
\FL{-0.799}{-3.182}{68.8}
\FL{8.283}{-6.321}{175.9}
\FL{4.997}{-5.511}{156.6}
\FL{2.162}{-3.665}{137.0}
\FL{-4.602}{2.370}{-163.1}
\FL{-6.550}{-0.011}{-89.7}
\FL{-4.581}{-2.377}{-16.5}
\FL{0.881}{2.999}{65.3}
\FL{2.174}{5.280}{56.0}
\FL{3.771}{7.329}{48.3}
\FL{5.631}{9.179}{41.4}
\FL{-5.151}{3.827}{177.3}
\FL{-7.813}{3.273}{-152.7}
\FL{-9.682}{1.331}{-113.2}
\FL{-9.677}{-1.343}{-66.6}
\FL{-7.802}{-3.279}{-27.1}
\FL{-5.138}{-3.826}{2.8}
\FL{-0.035}{-8.072}{89.6}
\FL{-0.013}{-4.222}{89.7}
\FL{-1.128}{3.214}{119.3}
\FL{-2.813}{5.599}{130.8}
\FL{-4.925}{7.654}{140.6}
\FL{-7.322}{9.325}{149.6}
\FL{-3.065}{4.757}{140.7}
\FL{-5.627}{6.350}{155.4}
\FL{-8.507}{7.248}{170.0}
\FL{-8.690}{-7.903}{13.5}
\FL{-6.169}{-7.023}{25.0}
\FL{-3.895}{-5.671}{36.6}
\FL{-1.952}{-3.877}{49.1}
\FL{5.299}{2.380}{-23.9}
\FL{6.981}{-0.015}{-90.4}
\FL{5.272}{-2.392}{-156.6}
\FL{8.439}{-7.362}{168.7}
\FL{5.486}{-6.366}{154.1}
\FL{4.915}{3.849}{5.5}
\FL{7.491}{3.499}{-21.9}
\FL{9.507}{1.882}{-57.2}
\FL{3.579}{6.089}{43.2}
\FL{5.947}{7.925}{32.5}
\FL{8.605}{9.307}{22.4}
\FL{3.187}{1.458}{-23.0}
\FL{3.151}{-1.473}{-158.1}
\FL{2.160}{-8.592}{111.3}
\FL{1.153}{-5.714}{107.2}
\FL{-4.995}{6.777}{146.6}
\FL{-8.259}{8.389}{160.8}
\FL{9.384}{0.914}{-73.4}
\FL{8.847}{-1.977}{-126.7}
\FL{6.306}{-3.544}{-167.7}
\FL{7.949}{-9.303}{153.6}
\FL{5.005}{-7.467}{142.4}
\FL{2.512}{-5.091}{130.1}
\FL{-8.679}{-5.784}{-3.2}
\FL{-6.027}{-5.529}{13.9}
\FL{-3.557}{-4.524}{30.4}
\FL{8.555}{-5.231}{-172.2}
\FL{5.625}{-5.084}{166.8}
\FL{-0.941}{4.720}{107.0}
\FL{-2.243}{8.288}{112.8}
\FL{-8.254}{-8.950}{22.6}
\FL{-5.066}{-7.187}{35.4}
\FL{-2.308}{-8.603}{67.3}
\FL{-1.253}{-5.795}{71.6}
\FL{3.758}{2.167}{-10.9}
\FL{5.781}{-0.009}{-90.3}
\FL{3.741}{-2.171}{-169.4}
\FL{-0.004}{4.753}{90.1}
\FL{-0.010}{8.253}{90.1}
\FL{6.083}{4.475}{2.1}
\FL{8.815}{3.977}{-23.7}
\FL{-2.961}{3.778}{149.5}
\FL{-5.722}{4.766}{171.2}
\FL{-8.648}{4.626}{-165.1}
\FL{4.561}{6.002}{32.0}
\FL{8.107}{7.526}{14.6}
\FL{-2.725}{1.238}{-156.7}
\FL{-2.701}{-1.247}{-22.5}
\FL{-8.897}{-3.737}{-27.2}
\FL{-6.345}{-4.367}{-1.5}
\FL{-3.980}{1.547}{-149.1}
\FL{-3.966}{-1.555}{-30.5}
\FL{-5.638}{5.666}{161.3}
\FL{-8.531}{6.181}{178.7}
\FL{2.629}{3.737}{35.5}
\FL{5.423}{5.041}{14.6}
\FL{8.497}{5.250}{-7.3}
\FL{2.013}{6.217}{62.9}
\FL{3.508}{8.788}{57.0}
\FL{-2.126}{6.103}{119.1}
\FL{-3.822}{8.756}{125.9}
\FL{0.825}{5.102}{76.2}
\FL{1.763}{8.396}{72.1}
\FL{8.584}{-8.611}{161.4}
\FL{5.940}{-7.430}{150.4}
\FL{5.637}{5.904}{20.8}
\FL{8.529}{6.557}{4.5}
\FL{0.666}{8.106}{83.0}
\FL{-0.752}{7.848}{98.2}
\FL{2.733}{-0.005}{-90.3}
\FL{0.983}{-7.660}{101.0}
\FL{-5.711}{-8.978}{39.9}
\FL{-3.341}{-6.651}{49.2}
\FL{8.932}{-3.512}{-149.6}
\FL{6.364}{-4.257}{-177.0}
\FL{3.714}{-3.846}{159.8}
\FL{-7.737}{-4.911}{-5.8}
\FL{-7.426}{-6.768}{13.6}
\FL{5.415}{-9.049}{137.6}
\FL{3.410}{-6.957}{129.9}
\FL{-4.368}{8.279}{132.6}
\FL{-2.652}{-0.012}{-89.2}
\FL{-0.967}{-7.489}{78.9}
\FL{-7.971}{4.012}{-161.9}
}
  \dipoloRojo{0.85}
  \node[font=\tiny, text=mmarojo] at (0.62,0.05) {$\vp$};
  \marcoMMA{10}{5}{$x/\ell$}{$z/\ell$}
\end{tikzpicture}\\[1pt]{\small\scshape (a)\; dipolo puntual}
\end{minipage}\hfill
\begin{minipage}[t]{0.49\textwidth}\centering
\begin{tikzpicture}[x=0.29cm, y=0.29cm]
  \panelMMA{10}{mmafondo}{1.2}{%
    % Dipolo físico - equipotenciales
% Generado por scripts/generar_datos.py (no editar a mano).
\EQ{(0.250,-0.010) (0.300,-0.012) (0.350,-0.013) (0.400,-0.015) (0.450,-0.018) (0.500,-0.021) (0.550,-0.024) (0.600,-0.028) (0.650,-0.032) (0.700,-0.037) (0.750,-0.043) (0.800,-0.049) (0.850,-0.056) (0.900,-0.064) (0.950,-0.073) (1.000,-0.082) (1.050,-0.093) (1.100,-0.104) (1.150,-0.116) (1.200,-0.130) (1.250,-0.145) (1.300,-0.161) (1.350,-0.178) (1.400,-0.196) (1.450,-0.216) (1.500,-0.238) (1.550,-0.261) (1.600,-0.286) (1.650,-0.313) (1.700,-0.341) (1.750,-0.372) (1.792,-0.400) (1.850,-0.441) (1.894,-0.475) (1.950,-0.521) (2.000,-0.566) (2.050,-0.615) (2.100,-0.668) (2.150,-0.725) (2.189,-0.775) (2.225,-0.824) (2.260,-0.875) (2.292,-0.925) (2.321,-0.975) (2.347,-1.025) (2.372,-1.075) (2.395,-1.125) (2.416,-1.175) (2.435,-1.225) (2.452,-1.275) (2.468,-1.325) (2.482,-1.375) (2.495,-1.425) (2.506,-1.475) (2.515,-1.525) (2.523,-1.575) (2.530,-1.625) (2.535,-1.675) (2.539,-1.725) (2.542,-1.775) (2.543,-1.825) (2.543,-1.875) (2.542,-1.925) (2.540,-1.975) (2.536,-2.025) (2.531,-2.075) (2.525,-2.125) (2.517,-2.175) (2.508,-2.225) (2.498,-2.275) (2.487,-2.325) (2.474,-2.375) (2.460,-2.425) (2.444,-2.475) (2.428,-2.525) (2.410,-2.575) (2.390,-2.625) (2.369,-2.675) (2.347,-2.725) (2.325,-2.771) (2.300,-2.821) (2.275,-2.867) (2.250,-2.912) (2.212,-2.975) (2.180,-3.025) (2.150,-3.069) (2.110,-3.125) (2.075,-3.172) (2.033,-3.225) (2.000,-3.264) (1.950,-3.321) (1.900,-3.374) (1.850,-3.424) (1.800,-3.472) (1.750,-3.517) (1.711,-3.550) (1.650,-3.600) (1.600,-3.638) (1.550,-3.675) (1.500,-3.709) (1.450,-3.742) (1.400,-3.774) (1.356,-3.800) (1.312,-3.825) (1.266,-3.850) (1.218,-3.875) (1.168,-3.900) (1.114,-3.925) (1.057,-3.950) (1.000,-3.974) (0.950,-3.993) (0.900,-4.011) (0.850,-4.028) (0.800,-4.044) (0.750,-4.059) (0.700,-4.073) (0.650,-4.086) (0.600,-4.098) (0.550,-4.108) (0.500,-4.118) (0.450,-4.127) (0.400,-4.135) (0.350,-4.142) (0.300,-4.148) (0.250,-4.153) (0.200,-4.157) (0.150,-4.161) (0.100,-4.163) (0.050,-4.164) (0.000,-4.165) (-0.050,-4.164) (-0.100,-4.163) (-0.150,-4.161) (-0.200,-4.157) (-0.250,-4.153) (-0.300,-4.148) (-0.350,-4.142) (-0.400,-4.135) (-0.450,-4.127) (-0.500,-4.118) (-0.550,-4.108) (-0.600,-4.098) (-0.650,-4.086) (-0.700,-4.073) (-0.750,-4.059) (-0.800,-4.044) (-0.850,-4.028) (-0.900,-4.011) (-0.950,-3.993) (-1.000,-3.974) (-1.050,-3.953) (-1.100,-3.931) (-1.150,-3.909) (-1.200,-3.884) (-1.250,-3.859) (-1.300,-3.832) (-1.350,-3.803) (-1.398,-3.775) (-1.450,-3.742) (-1.500,-3.709) (-1.549,-3.675) (-1.600,-3.638) (-1.649,-3.600) (-1.700,-3.559) (-1.740,-3.525) (-1.796,-3.475) (-1.849,-3.425) (-1.899,-3.375) (-1.946,-3.325) (-1.990,-3.275) (-2.025,-3.234) (-2.072,-3.175) (-2.110,-3.125) (-2.146,-3.075) (-2.175,-3.032) (-2.212,-2.975) (-2.242,-2.925) (-2.271,-2.875) (-2.298,-2.825) (-2.323,-2.775) (-2.347,-2.725) (-2.369,-2.675) (-2.390,-2.625) (-2.410,-2.575) (-2.428,-2.525) (-2.444,-2.475) (-2.460,-2.425) (-2.474,-2.375) (-2.487,-2.325) (-2.498,-2.275) (-2.508,-2.225) (-2.517,-2.175) (-2.525,-2.125) (-2.531,-2.075) (-2.536,-2.025) (-2.540,-1.975) (-2.542,-1.925) (-2.543,-1.875) (-2.543,-1.825) (-2.542,-1.775) (-2.539,-1.725) (-2.535,-1.675) (-2.530,-1.625) (-2.523,-1.575) (-2.515,-1.525) (-2.506,-1.475) (-2.495,-1.425) (-2.482,-1.375) (-2.468,-1.325) (-2.452,-1.275) (-2.435,-1.225) (-2.416,-1.175) (-2.395,-1.125) (-2.372,-1.075) (-2.347,-1.025) (-2.321,-0.975) (-2.292,-0.925) (-2.260,-0.875) (-2.226,-0.825) (-2.189,-0.775) (-2.150,-0.726) (-2.107,-0.675) (-2.060,-0.625) (-2.010,-0.575) (-1.954,-0.525) (-1.900,-0.480) (-1.850,-0.441) (-1.800,-0.406) (-1.754,-0.375) (-1.700,-0.341) (-1.650,-0.313) (-1.600,-0.286) (-1.550,-0.261) (-1.500,-0.238) (-1.450,-0.216) (-1.400,-0.196) (-1.350,-0.178) (-1.300,-0.161) (-1.250,-0.145) (-1.200,-0.130) (-1.150,-0.116) (-1.100,-0.104) (-1.050,-0.093) (-1.000,-0.082) (-0.950,-0.073) (-0.900,-0.064) (-0.850,-0.056) (-0.800,-0.049) (-0.750,-0.043) (-0.700,-0.037) (-0.650,-0.032) (-0.600,-0.028) (-0.550,-0.024) (-0.500,-0.021) (-0.450,-0.018) (-0.400,-0.015) (-0.350,-0.013) (-0.300,-0.012) (-0.250,-0.010)}
\EQ{(0.250,-0.009) (0.300,-0.010) (0.350,-0.012) (0.400,-0.014) (0.450,-0.016) (0.500,-0.019) (0.550,-0.022) (0.600,-0.025) (0.650,-0.029) (0.700,-0.034) (0.750,-0.039) (0.800,-0.044) (0.850,-0.051) (0.900,-0.058) (0.950,-0.065) (1.000,-0.074) (1.050,-0.083) (1.100,-0.094) (1.150,-0.105) (1.200,-0.117) (1.250,-0.130) (1.300,-0.144) (1.350,-0.159) (1.400,-0.176) (1.450,-0.194) (1.500,-0.213) (1.550,-0.233) (1.600,-0.255) (1.650,-0.279) (1.700,-0.304) (1.750,-0.331) (1.800,-0.360) (1.850,-0.391) (1.900,-0.425) (1.950,-0.460) (2.000,-0.498) (2.050,-0.540) (2.091,-0.575) (2.144,-0.625) (2.193,-0.675) (2.238,-0.725) (2.275,-0.769) (2.319,-0.825) (2.350,-0.868) (2.389,-0.925) (2.420,-0.975) (2.449,-1.025) (2.475,-1.074) (2.500,-1.124) (2.524,-1.175) (2.545,-1.225) (2.564,-1.275) (2.582,-1.325) (2.598,-1.375) (2.613,-1.425) (2.626,-1.475) (2.638,-1.525) (2.648,-1.575) (2.657,-1.625) (2.665,-1.675) (2.671,-1.725) (2.676,-1.775) (2.680,-1.825) (2.682,-1.875) (2.683,-1.925) (2.683,-1.975) (2.682,-2.025) (2.680,-2.075) (2.676,-2.125) (2.671,-2.175) (2.665,-2.225) (2.658,-2.275) (2.649,-2.325) (2.639,-2.375) (2.628,-2.425) (2.616,-2.475) (2.603,-2.525) (2.588,-2.575) (2.572,-2.625) (2.555,-2.675) (2.537,-2.725) (2.517,-2.775) (2.496,-2.825) (2.475,-2.871) (2.450,-2.923) (2.425,-2.972) (2.400,-3.019) (2.375,-3.063) (2.338,-3.125) (2.306,-3.175) (2.275,-3.222) (2.238,-3.275) (2.200,-3.325) (2.161,-3.375) (2.125,-3.419) (2.077,-3.475) (2.031,-3.525) (1.983,-3.575) (1.932,-3.625) (1.878,-3.675) (1.825,-3.722) (1.775,-3.763) (1.729,-3.800) (1.675,-3.841) (1.628,-3.875) (1.575,-3.912) (1.525,-3.944) (1.476,-3.975) (1.425,-4.005) (1.375,-4.034) (1.325,-4.061) (1.275,-4.087) (1.225,-4.111) (1.175,-4.134) (1.125,-4.156) (1.075,-4.177) (1.025,-4.197) (0.975,-4.216) (0.925,-4.233) (0.875,-4.250) (0.825,-4.266) (0.775,-4.280) (0.725,-4.294) (0.675,-4.306) (0.625,-4.318) (0.575,-4.329) (0.525,-4.338) (0.475,-4.347) (0.425,-4.355) (0.375,-4.362) (0.325,-4.368) (0.275,-4.374) (0.225,-4.378) (0.175,-4.381) (0.125,-4.384) (0.075,-4.386) (0.025,-4.387) (-0.050,-4.386) (-0.100,-4.385) (-0.150,-4.383) (-0.200,-4.380) (-0.250,-4.376) (-0.300,-4.371) (-0.350,-4.365) (-0.400,-4.359) (-0.450,-4.351) (-0.500,-4.343) (-0.550,-4.334) (-0.600,-4.323) (-0.650,-4.312) (-0.700,-4.300) (-0.750,-4.287) (-0.800,-4.273) (-0.850,-4.258) (-0.900,-4.242) (-0.949,-4.225) (-1.000,-4.207) (-1.050,-4.187) (-1.100,-4.167) (-1.150,-4.145) (-1.200,-4.123) (-1.248,-4.100) (-1.298,-4.075) (-1.345,-4.050) (-1.391,-4.025) (-1.434,-4.000) (-1.500,-3.960) (-1.550,-3.928) (-1.592,-3.900) (-1.650,-3.859) (-1.696,-3.825) (-1.750,-3.784) (-1.791,-3.750) (-1.850,-3.700) (-1.900,-3.655) (-1.950,-3.607) (-2.000,-3.557) (-2.050,-3.504) (-2.099,-3.450) (-2.141,-3.400) (-2.175,-3.358) (-2.219,-3.300) (-2.250,-3.258) (-2.290,-3.200) (-2.322,-3.150) (-2.350,-3.105) (-2.383,-3.050) (-2.410,-3.000) (-2.436,-2.950) (-2.461,-2.900) (-2.484,-2.850) (-2.506,-2.800) (-2.527,-2.750) (-2.546,-2.700) (-2.564,-2.650) (-2.580,-2.600) (-2.596,-2.550) (-2.610,-2.500) (-2.622,-2.450) (-2.634,-2.400) (-2.644,-2.350) (-2.653,-2.300) (-2.661,-2.250) (-2.668,-2.200) (-2.674,-2.150) (-2.678,-2.100) (-2.681,-2.050) (-2.683,-2.000) (-2.683,-1.950) (-2.683,-1.900) (-2.681,-1.850) (-2.678,-1.800) (-2.674,-1.750) (-2.668,-1.700) (-2.661,-1.650) (-2.653,-1.600) (-2.643,-1.550) (-2.632,-1.500) (-2.620,-1.450) (-2.606,-1.400) (-2.590,-1.350) (-2.573,-1.300) (-2.555,-1.250) (-2.534,-1.200) (-2.512,-1.150) (-2.488,-1.100) (-2.463,-1.050) (-2.435,-1.000) (-2.405,-0.950) (-2.375,-0.904) (-2.337,-0.850) (-2.300,-0.801) (-2.259,-0.750) (-2.225,-0.711) (-2.175,-0.657) (-2.125,-0.607) (-2.075,-0.561) (-2.033,-0.525) (-1.975,-0.479) (-1.925,-0.442) (-1.875,-0.408) (-1.825,-0.376) (-1.775,-0.346) (-1.725,-0.318) (-1.675,-0.292) (-1.625,-0.267) (-1.575,-0.244) (-1.525,-0.223) (-1.475,-0.203) (-1.425,-0.185) (-1.375,-0.168) (-1.325,-0.152) (-1.275,-0.137) (-1.225,-0.123) (-1.175,-0.111) (-1.125,-0.099) (-1.075,-0.088) (-1.025,-0.078) (-0.975,-0.070) (-0.925,-0.061) (-0.875,-0.054) (-0.825,-0.047) (-0.775,-0.041) (-0.725,-0.036) (-0.675,-0.031) (-0.625,-0.027) (-0.575,-0.023) (-0.525,-0.020) (-0.475,-0.017) (-0.425,-0.015) (-0.375,-0.013) (-0.325,-0.011) (-0.275,-0.010) (-0.250,-0.009)}
\EQ{(0.250,-0.008) (0.300,-0.009) (0.350,-0.011) (0.400,-0.012) (0.450,-0.014) (0.500,-0.017) (0.550,-0.019) (0.600,-0.022) (0.650,-0.026) (0.700,-0.030) (0.750,-0.034) (0.800,-0.039) (0.850,-0.045) (0.900,-0.051) (0.950,-0.058) (1.000,-0.066) (1.050,-0.074) (1.100,-0.083) (1.150,-0.093) (1.200,-0.104) (1.250,-0.115) (1.300,-0.128) (1.350,-0.141) (1.400,-0.156) (1.450,-0.172) (1.500,-0.188) (1.550,-0.206) (1.599,-0.225) (1.650,-0.246) (1.700,-0.268) (1.750,-0.292) (1.800,-0.317) (1.850,-0.344) (1.900,-0.372) (1.946,-0.400) (2.000,-0.435) (2.050,-0.470) (2.091,-0.500) (2.150,-0.547) (2.200,-0.589) (2.240,-0.625) (2.292,-0.675) (2.340,-0.725) (2.375,-0.764) (2.425,-0.824) (2.465,-0.875) (2.500,-0.924) (2.534,-0.975) (2.566,-1.025) (2.595,-1.075) (2.622,-1.125) (2.648,-1.175) (2.671,-1.225) (2.693,-1.275) (2.713,-1.325) (2.732,-1.375) (2.749,-1.425) (2.765,-1.475) (2.779,-1.525) (2.791,-1.575) (2.803,-1.625) (2.813,-1.675) (2.822,-1.725) (2.829,-1.775) (2.835,-1.825) (2.840,-1.875) (2.844,-1.925) (2.847,-1.975) (2.848,-2.025) (2.849,-2.075) (2.848,-2.125) (2.846,-2.175) (2.843,-2.225) (2.838,-2.275) (2.833,-2.325) (2.826,-2.375) (2.818,-2.425) (2.810,-2.475) (2.800,-2.525) (2.789,-2.575) (2.776,-2.625) (2.763,-2.675) (2.748,-2.725) (2.732,-2.775) (2.715,-2.825) (2.697,-2.875) (2.677,-2.925) (2.657,-2.975) (2.635,-3.025) (2.611,-3.075) (2.586,-3.125) (2.560,-3.175) (2.533,-3.225) (2.504,-3.275) (2.475,-3.322) (2.441,-3.375) (2.407,-3.425) (2.375,-3.470) (2.334,-3.525) (2.300,-3.569) (2.254,-3.625) (2.211,-3.675) (2.175,-3.715) (2.125,-3.767) (2.075,-3.818) (2.025,-3.865) (1.987,-3.900) (1.930,-3.950) (1.875,-3.996) (1.825,-4.035) (1.775,-4.073) (1.725,-4.109) (1.675,-4.143) (1.627,-4.175) (1.575,-4.208) (1.525,-4.238) (1.475,-4.266) (1.425,-4.294) (1.375,-4.320) (1.325,-4.345) (1.275,-4.369) (1.225,-4.392) (1.175,-4.413) (1.125,-4.434) (1.075,-4.453) (1.025,-4.472) (0.975,-4.489) (0.925,-4.506) (0.875,-4.521) (0.825,-4.536) (0.775,-4.549) (0.725,-4.562) (0.675,-4.574) (0.625,-4.585) (0.575,-4.595) (0.525,-4.604) (0.475,-4.612) (0.425,-4.620) (0.375,-4.626) (0.325,-4.632) (0.275,-4.637) (0.225,-4.641) (0.175,-4.644) (0.125,-4.647) (0.075,-4.649) (0.025,-4.649) (-0.050,-4.649) (-0.100,-4.648) (-0.150,-4.646) (-0.200,-4.643) (-0.250,-4.639) (-0.300,-4.635) (-0.350,-4.629) (-0.400,-4.623) (-0.450,-4.616) (-0.500,-4.608) (-0.550,-4.599) (-0.600,-4.590) (-0.650,-4.579) (-0.700,-4.568) (-0.750,-4.556) (-0.800,-4.543) (-0.850,-4.528) (-0.900,-4.513) (-0.950,-4.497) (-1.000,-4.480) (-1.050,-4.463) (-1.100,-4.444) (-1.150,-4.424) (-1.200,-4.402) (-1.250,-4.380) (-1.300,-4.357) (-1.350,-4.333) (-1.400,-4.307) (-1.450,-4.280) (-1.500,-4.252) (-1.546,-4.225) (-1.600,-4.192) (-1.650,-4.160) (-1.700,-4.126) (-1.750,-4.091) (-1.800,-4.054) (-1.850,-4.016) (-1.900,-3.975) (-1.950,-3.933) (-2.000,-3.888) (-2.041,-3.850) (-2.093,-3.800) (-2.142,-3.750) (-2.188,-3.700) (-2.225,-3.659) (-2.275,-3.600) (-2.315,-3.550) (-2.350,-3.504) (-2.389,-3.450) (-2.424,-3.400) (-2.457,-3.350) (-2.488,-3.300) (-2.518,-3.250) (-2.547,-3.200) (-2.573,-3.150) (-2.599,-3.100) (-2.623,-3.050) (-2.646,-3.000) (-2.667,-2.950) (-2.687,-2.900) (-2.706,-2.850) (-2.724,-2.800) (-2.740,-2.750) (-2.756,-2.700) (-2.770,-2.650) (-2.783,-2.600) (-2.794,-2.550) (-2.805,-2.500) (-2.814,-2.450) (-2.822,-2.400) (-2.830,-2.350) (-2.836,-2.300) (-2.841,-2.250) (-2.844,-2.200) (-2.847,-2.150) (-2.848,-2.100) (-2.849,-2.050) (-2.848,-2.000) (-2.846,-1.950) (-2.842,-1.900) (-2.838,-1.850) (-2.832,-1.800) (-2.826,-1.750) (-2.817,-1.700) (-2.808,-1.650) (-2.797,-1.600) (-2.785,-1.550) (-2.772,-1.500) (-2.757,-1.450) (-2.741,-1.400) (-2.723,-1.350) (-2.703,-1.300) (-2.683,-1.250) (-2.660,-1.200) (-2.635,-1.150) (-2.609,-1.100) (-2.581,-1.050) (-2.550,-1.000) (-2.518,-0.950) (-2.483,-0.900) (-2.450,-0.856) (-2.405,-0.800) (-2.362,-0.750) (-2.325,-0.709) (-2.275,-0.659) (-2.225,-0.611) (-2.184,-0.575) (-2.125,-0.526) (-2.075,-0.488) (-2.025,-0.452) (-1.975,-0.419) (-1.925,-0.387) (-1.875,-0.358) (-1.825,-0.330) (-1.775,-0.304) (-1.725,-0.280) (-1.675,-0.257) (-1.625,-0.236) (-1.575,-0.216) (-1.525,-0.197) (-1.475,-0.180) (-1.425,-0.164) (-1.375,-0.148) (-1.325,-0.134) (-1.275,-0.121) (-1.225,-0.109) (-1.175,-0.098) (-1.125,-0.088) (-1.075,-0.078) (-1.025,-0.070) (-0.975,-0.062) (-0.925,-0.055) (-0.875,-0.048) (-0.825,-0.042) (-0.775,-0.037) (-0.725,-0.032) (-0.675,-0.028) (-0.625,-0.024) (-0.575,-0.021) (-0.525,-0.018) (-0.475,-0.015) (-0.425,-0.013) (-0.375,-0.011) (-0.325,-0.010) (-0.275,-0.009) (-0.250,-0.008)}
\EQ{(0.250,-0.007) (0.300,-0.008) (0.350,-0.009) (0.400,-0.011) (0.450,-0.012) (0.500,-0.014) (0.550,-0.017) (0.600,-0.020) (0.650,-0.023) (0.700,-0.026) (0.750,-0.030) (0.800,-0.034) (0.850,-0.039) (0.900,-0.045) (0.950,-0.051) (1.000,-0.057) (1.050,-0.065) (1.100,-0.073) (1.150,-0.081) (1.200,-0.090) (1.250,-0.101) (1.300,-0.112) (1.350,-0.123) (1.400,-0.136) (1.450,-0.149) (1.500,-0.164) (1.550,-0.180) (1.600,-0.196) (1.650,-0.214) (1.700,-0.233) (1.750,-0.253) (1.800,-0.275) (1.850,-0.298) (1.900,-0.322) (1.950,-0.348) (2.000,-0.375) (2.050,-0.404) (2.100,-0.435) (2.150,-0.468) (2.196,-0.500) (2.250,-0.540) (2.294,-0.575) (2.350,-0.622) (2.400,-0.667) (2.450,-0.715) (2.500,-0.767) (2.550,-0.823) (2.594,-0.875) (2.625,-0.915) (2.669,-0.975) (2.700,-1.020) (2.735,-1.075) (2.765,-1.125) (2.793,-1.175) (2.820,-1.225) (2.844,-1.275) (2.867,-1.325) (2.888,-1.375) (2.908,-1.425) (2.926,-1.475) (2.943,-1.525) (2.959,-1.575) (2.973,-1.625) (2.986,-1.675) (2.997,-1.725) (3.008,-1.775) (3.017,-1.825) (3.024,-1.875) (3.031,-1.925) (3.037,-1.975) (3.041,-2.025) (3.044,-2.075) (3.047,-2.125) (3.048,-2.175) (3.048,-2.225) (3.047,-2.275) (3.044,-2.325) (3.041,-2.375) (3.037,-2.425) (3.032,-2.475) (3.025,-2.525) (3.018,-2.575) (3.009,-2.625) (3.000,-2.675) (2.989,-2.725) (2.977,-2.775) (2.964,-2.825) (2.950,-2.875) (2.935,-2.925) (2.919,-2.975) (2.902,-3.025) (2.883,-3.075) (2.864,-3.125) (2.843,-3.175) (2.821,-3.225) (2.798,-3.275) (2.775,-3.321) (2.750,-3.370) (2.725,-3.416) (2.700,-3.460) (2.662,-3.525) (2.630,-3.575) (2.600,-3.621) (2.563,-3.675) (2.527,-3.725) (2.489,-3.775) (2.450,-3.825) (2.408,-3.875) (2.375,-3.914) (2.325,-3.969) (2.275,-4.022) (2.225,-4.072) (2.175,-4.120) (2.125,-4.166) (2.087,-4.200) (2.028,-4.250) (1.975,-4.292) (1.933,-4.325) (1.875,-4.368) (1.829,-4.400) (1.775,-4.437) (1.725,-4.469) (1.676,-4.500) (1.625,-4.530) (1.575,-4.559) (1.525,-4.586) (1.475,-4.613) (1.425,-4.638) (1.375,-4.662) (1.325,-4.685) (1.275,-4.707) (1.225,-4.727) (1.175,-4.747) (1.125,-4.766) (1.075,-4.784) (1.025,-4.801) (0.975,-4.818) (0.925,-4.833) (0.875,-4.847) (0.825,-4.861) (0.775,-4.873) (0.725,-4.885) (0.675,-4.896) (0.625,-4.907) (0.575,-4.916) (0.525,-4.924) (0.475,-4.932) (0.425,-4.939) (0.375,-4.945) (0.325,-4.951) (0.275,-4.955) (0.225,-4.959) (0.175,-4.962) (0.125,-4.965) (0.075,-4.966) (0.025,-4.967) (-0.050,-4.967) (-0.100,-4.965) (-0.150,-4.963) (-0.200,-4.961) (-0.250,-4.957) (-0.300,-4.953) (-0.350,-4.948) (-0.400,-4.942) (-0.450,-4.936) (-0.500,-4.928) (-0.550,-4.920) (-0.600,-4.911) (-0.650,-4.902) (-0.700,-4.891) (-0.750,-4.880) (-0.800,-4.867) (-0.850,-4.854) (-0.900,-4.840) (-0.950,-4.825) (-1.000,-4.810) (-1.050,-4.793) (-1.100,-4.775) (-1.150,-4.757) (-1.200,-4.738) (-1.250,-4.717) (-1.300,-4.696) (-1.346,-4.675) (-1.400,-4.650) (-1.450,-4.625) (-1.499,-4.600) (-1.546,-4.575) (-1.591,-4.550) (-1.650,-4.516) (-1.700,-4.485) (-1.750,-4.453) (-1.793,-4.425) (-1.850,-4.385) (-1.899,-4.350) (-1.950,-4.312) (-1.997,-4.275) (-2.050,-4.231) (-2.100,-4.189) (-2.143,-4.150) (-2.197,-4.100) (-2.248,-4.050) (-2.296,-4.000) (-2.343,-3.950) (-2.387,-3.900) (-2.425,-3.855) (-2.470,-3.800) (-2.508,-3.750) (-2.545,-3.700) (-2.575,-3.658) (-2.614,-3.600) (-2.646,-3.550) (-2.675,-3.503) (-2.706,-3.450) (-2.734,-3.400) (-2.760,-3.350) (-2.786,-3.300) (-2.809,-3.250) (-2.832,-3.200) (-2.853,-3.150) (-2.874,-3.100) (-2.893,-3.050) (-2.911,-3.000) (-2.927,-2.950) (-2.943,-2.900) (-2.957,-2.850) (-2.971,-2.800) (-2.983,-2.750) (-2.994,-2.700) (-3.005,-2.650) (-3.014,-2.600) (-3.022,-2.550) (-3.029,-2.500) (-3.034,-2.450) (-3.039,-2.400) (-3.043,-2.350) (-3.046,-2.300) (-3.047,-2.250) (-3.048,-2.200) (-3.047,-2.150) (-3.046,-2.100) (-3.043,-2.050) (-3.039,-2.000) (-3.034,-1.950) (-3.028,-1.900) (-3.021,-1.850) (-3.012,-1.800) (-3.003,-1.750) (-2.992,-1.700) (-2.979,-1.650) (-2.966,-1.600) (-2.951,-1.550) (-2.935,-1.500) (-2.918,-1.450) (-2.898,-1.400) (-2.878,-1.350) (-2.856,-1.300) (-2.832,-1.250) (-2.807,-1.200) (-2.780,-1.150) (-2.751,-1.100) (-2.720,-1.050) (-2.686,-1.000) (-2.651,-0.950) (-2.613,-0.900) (-2.575,-0.852) (-2.530,-0.800) (-2.484,-0.750) (-2.435,-0.700) (-2.382,-0.650) (-2.325,-0.601) (-2.275,-0.560) (-2.230,-0.525) (-2.175,-0.485) (-2.125,-0.451) (-2.075,-0.419) (-2.025,-0.389) (-1.975,-0.361) (-1.925,-0.335) (-1.875,-0.310) (-1.825,-0.286) (-1.775,-0.264) (-1.725,-0.243) (-1.675,-0.223) (-1.625,-0.205) (-1.575,-0.188) (-1.525,-0.172) (-1.475,-0.157) (-1.425,-0.143) (-1.375,-0.129) (-1.325,-0.117) (-1.275,-0.106) (-1.225,-0.095) (-1.175,-0.086) (-1.125,-0.077) (-1.075,-0.069) (-1.025,-0.061) (-0.975,-0.054) (-0.925,-0.048) (-0.875,-0.042) (-0.825,-0.037) (-0.775,-0.032) (-0.725,-0.028) (-0.675,-0.024) (-0.625,-0.021) (-0.575,-0.018) (-0.525,-0.016) (-0.475,-0.013) (-0.425,-0.012) (-0.375,-0.010) (-0.325,-0.009) (-0.275,-0.008) (-0.250,-0.007)}
\EQ{(0.250,-0.006) (0.300,-0.007) (0.350,-0.008) (0.400,-0.009) (0.450,-0.011) (0.500,-0.012) (0.550,-0.014) (0.600,-0.017) (0.650,-0.019) (0.700,-0.022) (0.750,-0.026) (0.800,-0.030) (0.850,-0.034) (0.900,-0.038) (0.950,-0.044) (1.000,-0.049) (1.050,-0.055) (1.100,-0.062) (1.150,-0.069) (1.200,-0.077) (1.250,-0.086) (1.300,-0.095) (1.350,-0.105) (1.400,-0.116) (1.450,-0.128) (1.500,-0.140) (1.550,-0.153) (1.600,-0.167) (1.650,-0.183) (1.700,-0.199) (1.750,-0.216) (1.800,-0.234) (1.850,-0.253) (1.900,-0.273) (1.950,-0.295) (2.000,-0.318) (2.050,-0.342) (2.100,-0.367) (2.150,-0.394) (2.200,-0.423) (2.245,-0.450) (2.300,-0.485) (2.350,-0.519) (2.394,-0.550) (2.450,-0.593) (2.491,-0.625) (2.549,-0.675) (2.600,-0.721) (2.650,-0.770) (2.700,-0.821) (2.749,-0.875) (2.791,-0.925) (2.825,-0.967) (2.869,-1.025) (2.900,-1.069) (2.937,-1.125) (2.968,-1.175) (2.998,-1.225) (3.025,-1.274) (3.050,-1.322) (3.075,-1.373) (3.099,-1.425) (3.120,-1.475) (3.140,-1.525) (3.159,-1.575) (3.176,-1.625) (3.192,-1.675) (3.207,-1.725) (3.220,-1.775) (3.233,-1.825) (3.244,-1.875) (3.254,-1.925) (3.263,-1.975) (3.270,-2.025) (3.277,-2.075) (3.282,-2.125) (3.287,-2.175) (3.290,-2.225) (3.293,-2.275) (3.294,-2.325) (3.295,-2.375) (3.294,-2.425) (3.293,-2.475) (3.290,-2.525) (3.287,-2.575) (3.282,-2.625) (3.276,-2.675) (3.270,-2.725) (3.262,-2.775) (3.254,-2.825) (3.245,-2.875) (3.234,-2.925) (3.223,-2.975) (3.210,-3.025) (3.197,-3.075) (3.182,-3.125) (3.167,-3.175) (3.150,-3.225) (3.133,-3.275) (3.114,-3.325) (3.094,-3.375) (3.073,-3.425) (3.051,-3.475) (3.028,-3.525) (3.004,-3.575) (2.978,-3.625) (2.951,-3.675) (2.925,-3.722) (2.894,-3.775) (2.864,-3.825) (2.832,-3.875) (2.800,-3.923) (2.764,-3.975) (2.727,-4.025) (2.689,-4.075) (2.650,-4.125) (2.609,-4.175) (2.575,-4.214) (2.525,-4.270) (2.475,-4.324) (2.425,-4.375) (2.375,-4.424) (2.325,-4.471) (2.275,-4.516) (2.235,-4.550) (2.175,-4.600) (2.125,-4.640) (2.079,-4.675) (2.025,-4.715) (1.975,-4.750) (1.925,-4.784) (1.875,-4.817) (1.825,-4.848) (1.781,-4.875) (1.725,-4.907) (1.675,-4.935) (1.625,-4.962) (1.575,-4.988) (1.525,-5.013) (1.475,-5.037) (1.425,-5.059) (1.375,-5.081) (1.325,-5.102) (1.275,-5.122) (1.225,-5.141) (1.175,-5.159) (1.125,-5.177) (1.075,-5.193) (1.025,-5.209) (0.975,-5.223) (0.925,-5.238) (0.875,-5.251) (0.825,-5.263) (0.775,-5.275) (0.725,-5.286) (0.675,-5.296) (0.625,-5.305) (0.575,-5.314) (0.525,-5.322) (0.475,-5.329) (0.425,-5.335) (0.375,-5.341) (0.325,-5.346) (0.275,-5.350) (0.225,-5.354) (0.175,-5.357) (0.125,-5.359) (0.075,-5.360) (0.025,-5.361) (-0.050,-5.361) (-0.100,-5.360) (-0.150,-5.358) (-0.200,-5.355) (-0.250,-5.352) (-0.300,-5.348) (-0.350,-5.344) (-0.400,-5.338) (-0.450,-5.332) (-0.500,-5.325) (-0.550,-5.318) (-0.600,-5.310) (-0.650,-5.301) (-0.700,-5.291) (-0.750,-5.280) (-0.800,-5.269) (-0.850,-5.257) (-0.900,-5.244) (-0.950,-5.231) (-1.000,-5.216) (-1.050,-5.201) (-1.100,-5.185) (-1.150,-5.168) (-1.200,-5.150) (-1.250,-5.132) (-1.300,-5.112) (-1.350,-5.092) (-1.400,-5.070) (-1.450,-5.048) (-1.499,-5.025) (-1.550,-5.001) (-1.600,-4.975) (-1.648,-4.950) (-1.694,-4.925) (-1.738,-4.900) (-1.800,-4.863) (-1.850,-4.833) (-1.900,-4.801) (-1.950,-4.767) (-2.000,-4.733) (-2.045,-4.700) (-2.100,-4.659) (-2.144,-4.625) (-2.200,-4.580) (-2.250,-4.537) (-2.293,-4.500) (-2.347,-4.450) (-2.400,-4.400) (-2.450,-4.350) (-2.497,-4.300) (-2.543,-4.250) (-2.587,-4.200) (-2.625,-4.155) (-2.670,-4.100) (-2.708,-4.050) (-2.746,-4.000) (-2.775,-3.959) (-2.815,-3.900) (-2.848,-3.850) (-2.875,-3.807) (-2.909,-3.750) (-2.938,-3.700) (-2.965,-3.650) (-2.991,-3.600) (-3.016,-3.550) (-3.040,-3.500) (-3.062,-3.450) (-3.084,-3.400) (-3.104,-3.350) (-3.123,-3.300) (-3.141,-3.250) (-3.159,-3.200) (-3.175,-3.150) (-3.190,-3.100) (-3.204,-3.050) (-3.217,-3.000) (-3.228,-2.950) (-3.239,-2.900) (-3.249,-2.850) (-3.258,-2.800) (-3.266,-2.750) (-3.273,-2.700) (-3.279,-2.650) (-3.284,-2.600) (-3.288,-2.550) (-3.292,-2.500) (-3.294,-2.450) (-3.295,-2.400) (-3.295,-2.350) (-3.294,-2.300) (-3.292,-2.250) (-3.289,-2.200) (-3.285,-2.150) (-3.280,-2.100) (-3.274,-2.050) (-3.267,-2.000) (-3.258,-1.950) (-3.249,-1.900) (-3.238,-1.850) (-3.227,-1.800) (-3.214,-1.750) (-3.200,-1.701) (-3.184,-1.650) (-3.168,-1.600) (-3.150,-1.551) (-3.131,-1.500) (-3.110,-1.450) (-3.088,-1.400) (-3.064,-1.350) (-3.039,-1.300) (-3.012,-1.250) (-2.983,-1.200) (-2.953,-1.150) (-2.925,-1.107) (-2.887,-1.050) (-2.850,-1.000) (-2.811,-0.950) (-2.775,-0.906) (-2.727,-0.850) (-2.680,-0.800) (-2.630,-0.750) (-2.577,-0.700) (-2.525,-0.654) (-2.475,-0.613) (-2.427,-0.575) (-2.375,-0.537) (-2.325,-0.502) (-2.275,-0.469) (-2.225,-0.438) (-2.175,-0.408) (-2.125,-0.381) (-2.075,-0.354) (-2.025,-0.329) (-1.975,-0.306) (-1.925,-0.284) (-1.875,-0.263) (-1.825,-0.243) (-1.777,-0.225) (-1.725,-0.207) (-1.675,-0.190) (-1.625,-0.175) (-1.575,-0.160) (-1.525,-0.147) (-1.475,-0.134) (-1.425,-0.122) (-1.375,-0.111) (-1.325,-0.100) (-1.275,-0.091) (-1.225,-0.082) (-1.175,-0.073) (-1.125,-0.066) (-1.075,-0.059) (-1.025,-0.052) (-0.975,-0.046) (-0.925,-0.041) (-0.875,-0.036) (-0.825,-0.032) (-0.775,-0.028) (-0.725,-0.024) (-0.675,-0.021) (-0.625,-0.018) (-0.575,-0.016) (-0.525,-0.013) (-0.475,-0.012) (-0.425,-0.010) (-0.375,-0.009) (-0.325,-0.007) (-0.275,-0.007) (-0.250,-0.006)}
\EQ{(0.250,-0.005) (0.300,-0.006) (0.350,-0.007) (0.400,-0.008) (0.450,-0.009) (0.500,-0.010) (0.550,-0.012) (0.600,-0.014) (0.650,-0.016) (0.700,-0.019) (0.750,-0.021) (0.800,-0.025) (0.850,-0.028) (0.900,-0.032) (0.950,-0.036) (1.000,-0.041) (1.050,-0.046) (1.100,-0.052) (1.150,-0.058) (1.200,-0.064) (1.250,-0.072) (1.300,-0.079) (1.350,-0.088) (1.400,-0.097) (1.450,-0.106) (1.500,-0.116) (1.550,-0.127) (1.600,-0.139) (1.650,-0.151) (1.700,-0.165) (1.750,-0.179) (1.800,-0.193) (1.850,-0.209) (1.898,-0.225) (1.950,-0.243) (2.000,-0.262) (2.050,-0.282) (2.100,-0.302) (2.150,-0.324) (2.200,-0.347) (2.250,-0.371) (2.300,-0.397) (2.350,-0.423) (2.397,-0.450) (2.450,-0.481) (2.500,-0.513) (2.550,-0.546) (2.593,-0.575) (2.650,-0.617) (2.693,-0.650) (2.750,-0.696) (2.800,-0.739) (2.840,-0.775) (2.892,-0.825) (2.942,-0.875) (2.988,-0.925) (3.025,-0.967) (3.074,-1.025) (3.113,-1.075) (3.150,-1.125) (3.186,-1.175) (3.219,-1.225) (3.250,-1.274) (3.280,-1.325) (3.308,-1.375) (3.335,-1.425) (3.360,-1.475) (3.384,-1.525) (3.406,-1.575) (3.427,-1.625) (3.447,-1.675) (3.465,-1.725) (3.482,-1.775) (3.498,-1.825) (3.513,-1.875) (3.527,-1.925) (3.540,-1.975) (3.551,-2.025) (3.562,-2.075) (3.571,-2.125) (3.580,-2.175) (3.587,-2.225) (3.594,-2.275) (3.599,-2.325) (3.604,-2.375) (3.607,-2.425) (3.610,-2.475) (3.611,-2.525) (3.612,-2.575) (3.612,-2.625) (3.611,-2.675) (3.609,-2.725) (3.606,-2.775) (3.602,-2.825) (3.598,-2.875) (3.592,-2.925) (3.586,-2.975) (3.578,-3.025) (3.570,-3.075) (3.561,-3.125) (3.551,-3.175) (3.540,-3.225) (3.528,-3.275) (3.515,-3.325) (3.501,-3.375) (3.487,-3.425) (3.471,-3.475) (3.455,-3.525) (3.437,-3.575) (3.419,-3.625) (3.400,-3.673) (3.379,-3.725) (3.357,-3.775) (3.334,-3.825) (3.311,-3.875) (3.286,-3.925) (3.260,-3.975) (3.234,-4.025) (3.205,-4.075) (3.176,-4.125) (3.150,-4.168) (3.114,-4.225) (3.081,-4.275) (3.050,-4.321) (3.011,-4.375) (2.975,-4.424) (2.936,-4.475) (2.900,-4.520) (2.855,-4.575) (2.811,-4.625) (2.775,-4.666) (2.725,-4.720) (2.675,-4.771) (2.625,-4.821) (2.575,-4.869) (2.525,-4.915) (2.485,-4.950) (2.426,-5.000) (2.375,-5.042) (2.334,-5.075) (2.275,-5.120) (2.234,-5.150) (2.175,-5.192) (2.127,-5.225) (2.075,-5.260) (2.025,-5.291) (1.975,-5.322) (1.928,-5.350) (1.885,-5.375) (1.840,-5.400) (1.794,-5.425) (1.746,-5.450) (1.700,-5.473) (1.650,-5.497) (1.600,-5.521) (1.550,-5.543) (1.500,-5.565) (1.450,-5.586) (1.400,-5.606) (1.350,-5.625) (1.300,-5.643) (1.250,-5.661) (1.200,-5.677) (1.150,-5.693) (1.100,-5.709) (1.050,-5.723) (1.000,-5.737) (0.950,-5.750) (0.900,-5.762) (0.850,-5.774) (0.800,-5.785) (0.750,-5.795) (0.700,-5.805) (0.650,-5.813) (0.600,-5.822) (0.550,-5.829) (0.500,-5.836) (0.450,-5.842) (0.400,-5.848) (0.350,-5.853) (0.300,-5.857) (0.250,-5.860) (0.200,-5.863) (0.150,-5.866) (0.100,-5.867) (0.050,-5.868) (0.000,-5.868) (-0.050,-5.868) (-0.100,-5.867) (-0.150,-5.866) (-0.200,-5.863) (-0.250,-5.860) (-0.300,-5.857) (-0.350,-5.853) (-0.400,-5.848) (-0.450,-5.842) (-0.500,-5.836) (-0.550,-5.829) (-0.600,-5.822) (-0.650,-5.813) (-0.700,-5.805) (-0.750,-5.795) (-0.800,-5.785) (-0.850,-5.774) (-0.900,-5.762) (-0.950,-5.750) (-1.000,-5.737) (-1.050,-5.723) (-1.100,-5.709) (-1.150,-5.693) (-1.200,-5.677) (-1.250,-5.661) (-1.300,-5.643) (-1.349,-5.625) (-1.400,-5.606) (-1.450,-5.586) (-1.500,-5.565) (-1.550,-5.543) (-1.600,-5.521) (-1.650,-5.497) (-1.696,-5.475) (-1.746,-5.450) (-1.794,-5.425) (-1.840,-5.400) (-1.885,-5.375) (-1.928,-5.350) (-1.975,-5.322) (-2.025,-5.291) (-2.075,-5.260) (-2.125,-5.226) (-2.175,-5.192) (-2.225,-5.157) (-2.268,-5.125) (-2.325,-5.082) (-2.365,-5.050) (-2.425,-5.001) (-2.475,-4.959) (-2.513,-4.925) (-2.568,-4.875) (-2.621,-4.825) (-2.671,-4.775) (-2.720,-4.725) (-2.767,-4.675) (-2.811,-4.625) (-2.850,-4.580) (-2.896,-4.525) (-2.936,-4.475) (-2.974,-4.425) (-3.011,-4.375) (-3.047,-4.325) (-3.081,-4.275) (-3.114,-4.225) (-3.146,-4.175) (-3.175,-4.127) (-3.205,-4.075) (-3.234,-4.025) (-3.260,-3.975) (-3.286,-3.925) (-3.311,-3.875) (-3.334,-3.825) (-3.357,-3.775) (-3.379,-3.725) (-3.399,-3.675) (-3.419,-3.625) (-3.437,-3.575) (-3.455,-3.525) (-3.471,-3.475) (-3.487,-3.425) (-3.501,-3.375) (-3.515,-3.325) (-3.528,-3.275) (-3.540,-3.225) (-3.551,-3.175) (-3.561,-3.125) (-3.570,-3.075) (-3.578,-3.025) (-3.586,-2.975) (-3.592,-2.925) (-3.598,-2.875) (-3.602,-2.825) (-3.606,-2.775) (-3.609,-2.725) (-3.611,-2.675) (-3.612,-2.625) (-3.612,-2.575) (-3.611,-2.525) (-3.610,-2.475) (-3.607,-2.425) (-3.604,-2.375) (-3.599,-2.325) (-3.594,-2.275) (-3.587,-2.225) (-3.580,-2.175) (-3.571,-2.125) (-3.562,-2.075) (-3.551,-2.025) (-3.540,-1.975) (-3.527,-1.925) (-3.513,-1.875) (-3.498,-1.825) (-3.482,-1.775) (-3.465,-1.725) (-3.447,-1.675) (-3.427,-1.625) (-3.406,-1.575) (-3.384,-1.525) (-3.360,-1.475) (-3.335,-1.425) (-3.308,-1.375) (-3.280,-1.325) (-3.250,-1.275) (-3.219,-1.225) (-3.186,-1.175) (-3.150,-1.125) (-3.113,-1.075) (-3.075,-1.026) (-3.032,-0.975) (-2.988,-0.925) (-2.950,-0.884) (-2.900,-0.833) (-2.850,-0.784) (-2.813,-0.750) (-2.755,-0.700) (-2.700,-0.655) (-2.650,-0.617) (-2.600,-0.580) (-2.557,-0.550) (-2.500,-0.513) (-2.450,-0.481) (-2.400,-0.452) (-2.353,-0.425) (-2.307,-0.400) (-2.258,-0.375) (-2.207,-0.350) (-2.152,-0.325) (-2.100,-0.302) (-2.050,-0.282) (-2.000,-0.262) (-1.950,-0.243) (-1.900,-0.226) (-1.850,-0.209) (-1.800,-0.193) (-1.750,-0.179) (-1.700,-0.165) (-1.650,-0.151) (-1.600,-0.139) (-1.550,-0.127) (-1.500,-0.116) (-1.450,-0.106) (-1.400,-0.097) (-1.350,-0.088) (-1.300,-0.079) (-1.250,-0.072) (-1.200,-0.064) (-1.150,-0.058) (-1.100,-0.052) (-1.050,-0.046) (-1.000,-0.041) (-0.950,-0.036) (-0.900,-0.032) (-0.850,-0.028) (-0.800,-0.025) (-0.750,-0.021) (-0.700,-0.019) (-0.650,-0.016) (-0.600,-0.014) (-0.550,-0.012) (-0.500,-0.010) (-0.450,-0.009) (-0.400,-0.008) (-0.350,-0.007) (-0.300,-0.006) (-0.250,-0.005)}
\EQ{(0.250,-0.004) (0.300,-0.005) (0.350,-0.005) (0.400,-0.006) (0.450,-0.007) (0.500,-0.008) (0.550,-0.010) (0.600,-0.011) (0.650,-0.013) (0.700,-0.015) (0.750,-0.017) (0.800,-0.020) (0.850,-0.022) (0.900,-0.026) (0.950,-0.029) (1.000,-0.033) (1.050,-0.037) (1.100,-0.041) (1.150,-0.046) (1.200,-0.052) (1.250,-0.057) (1.300,-0.063) (1.350,-0.070) (1.400,-0.077) (1.450,-0.085) (1.500,-0.093) (1.550,-0.102) (1.600,-0.111) (1.650,-0.121) (1.700,-0.131) (1.750,-0.142) (1.800,-0.154) (1.850,-0.166) (1.900,-0.179) (1.950,-0.193) (2.000,-0.208) (2.050,-0.223) (2.100,-0.239) (2.150,-0.256) (2.200,-0.274) (2.250,-0.293) (2.300,-0.313) (2.350,-0.333) (2.400,-0.355) (2.450,-0.378) (2.497,-0.400) (2.548,-0.425) (2.596,-0.450) (2.642,-0.475) (2.686,-0.500) (2.750,-0.538) (2.800,-0.570) (2.847,-0.600) (2.900,-0.637) (2.950,-0.673) (3.000,-0.711) (3.050,-0.750) (3.100,-0.792) (3.138,-0.825) (3.193,-0.875) (3.245,-0.925) (3.294,-0.975) (3.341,-1.025) (3.375,-1.063) (3.425,-1.122) (3.467,-1.175) (3.500,-1.218) (3.542,-1.275) (3.575,-1.323) (3.609,-1.375) (3.640,-1.425) (3.670,-1.475) (3.698,-1.525) (3.725,-1.574) (3.750,-1.623) (3.775,-1.675) (3.798,-1.725) (3.820,-1.775) (3.841,-1.825) (3.860,-1.875) (3.878,-1.925) (3.896,-1.975) (3.912,-2.025) (3.927,-2.075) (3.941,-2.125) (3.954,-2.175) (3.966,-2.225) (3.977,-2.275) (3.988,-2.325) (3.997,-2.375) (4.005,-2.425) (4.013,-2.475) (4.020,-2.525) (4.025,-2.575) (4.030,-2.625) (4.034,-2.675) (4.037,-2.725) (4.040,-2.775) (4.041,-2.825) (4.042,-2.875) (4.042,-2.925) (4.041,-2.975) (4.039,-3.025) (4.037,-3.075) (4.033,-3.125) (4.029,-3.175) (4.024,-3.225) (4.018,-3.275) (4.012,-3.325) (4.004,-3.375) (3.996,-3.425) (3.987,-3.475) (3.977,-3.525) (3.967,-3.575) (3.955,-3.625) (3.943,-3.675) (3.930,-3.725) (3.916,-3.775) (3.901,-3.825) (3.886,-3.875) (3.869,-3.925) (3.852,-3.975) (3.834,-4.025) (3.815,-4.075) (3.795,-4.125) (3.775,-4.172) (3.752,-4.225) (3.729,-4.275) (3.706,-4.325) (3.681,-4.375) (3.655,-4.425) (3.629,-4.475) (3.601,-4.525) (3.575,-4.571) (3.543,-4.625) (3.512,-4.675) (3.480,-4.725) (3.450,-4.771) (3.413,-4.825) (3.378,-4.875) (3.341,-4.925) (3.303,-4.975) (3.264,-5.025) (3.225,-5.073) (3.181,-5.125) (3.138,-5.175) (3.100,-5.217) (3.050,-5.271) (3.000,-5.322) (2.950,-5.372) (2.900,-5.420) (2.850,-5.467) (2.800,-5.512) (2.756,-5.550) (2.700,-5.598) (2.650,-5.638) (2.603,-5.675) (2.550,-5.716) (2.504,-5.750) (2.450,-5.789) (2.400,-5.823) (2.350,-5.857) (2.300,-5.889) (2.250,-5.921) (2.202,-5.950) (2.150,-5.981) (2.100,-6.010) (2.050,-6.037) (2.000,-6.064) (1.950,-6.090) (1.900,-6.115) (1.850,-6.139) (1.800,-6.163) (1.750,-6.186) (1.700,-6.208) (1.650,-6.229) (1.600,-6.249) (1.550,-6.269) (1.500,-6.288) (1.450,-6.306) (1.400,-6.323) (1.350,-6.340) (1.300,-6.356) (1.250,-6.372) (1.200,-6.387) (1.150,-6.401) (1.100,-6.414) (1.050,-6.427) (1.000,-6.439) (0.950,-6.451) (0.900,-6.462) (0.850,-6.472) (0.800,-6.482) (0.750,-6.491) (0.700,-6.499) (0.650,-6.507) (0.600,-6.515) (0.550,-6.521) (0.500,-6.527) (0.450,-6.533) (0.400,-6.538) (0.350,-6.542) (0.300,-6.546) (0.250,-6.549) (0.200,-6.552) (0.150,-6.554) (0.100,-6.555) (0.050,-6.556) (0.000,-6.556) (-0.050,-6.556) (-0.100,-6.555) (-0.150,-6.554) (-0.200,-6.552) (-0.250,-6.549) (-0.300,-6.546) (-0.350,-6.542) (-0.400,-6.538) (-0.450,-6.533) (-0.500,-6.527) (-0.550,-6.521) (-0.600,-6.515) (-0.650,-6.507) (-0.700,-6.499) (-0.750,-6.491) (-0.800,-6.482) (-0.850,-6.472) (-0.900,-6.462) (-0.950,-6.451) (-1.000,-6.439) (-1.050,-6.427) (-1.100,-6.414) (-1.150,-6.401) (-1.200,-6.387) (-1.250,-6.372) (-1.300,-6.356) (-1.350,-6.340) (-1.400,-6.323) (-1.450,-6.306) (-1.500,-6.288) (-1.550,-6.269) (-1.598,-6.250) (-1.650,-6.229) (-1.700,-6.208) (-1.750,-6.186) (-1.800,-6.163) (-1.850,-6.139) (-1.900,-6.115) (-1.950,-6.090) (-2.000,-6.064) (-2.050,-6.037) (-2.100,-6.010) (-2.150,-5.981) (-2.200,-5.951) (-2.243,-5.925) (-2.300,-5.889) (-2.350,-5.857) (-2.397,-5.825) (-2.450,-5.789) (-2.500,-5.753) (-2.550,-5.716) (-2.600,-5.678) (-2.650,-5.638) (-2.697,-5.600) (-2.750,-5.556) (-2.800,-5.512) (-2.841,-5.475) (-2.895,-5.425) (-2.947,-5.375) (-2.997,-5.325) (-3.046,-5.275) (-3.092,-5.225) (-3.138,-5.175) (-3.175,-5.132) (-3.223,-5.075) (-3.264,-5.025) (-3.300,-4.979) (-3.341,-4.925) (-3.375,-4.879) (-3.413,-4.825) (-3.447,-4.775) (-3.475,-4.733) (-3.512,-4.675) (-3.543,-4.625) (-3.573,-4.575) (-3.600,-4.527) (-3.625,-4.482) (-3.650,-4.435) (-3.675,-4.387) (-3.700,-4.337) (-3.725,-4.284) (-3.750,-4.229) (-3.774,-4.175) (-3.795,-4.125) (-3.815,-4.075) (-3.834,-4.025) (-3.852,-3.975) (-3.869,-3.925) (-3.886,-3.875) (-3.901,-3.825) (-3.916,-3.775) (-3.930,-3.725) (-3.943,-3.675) (-3.955,-3.625) (-3.967,-3.575) (-3.977,-3.525) (-3.987,-3.475) (-3.996,-3.425) (-4.004,-3.375) (-4.012,-3.325) (-4.018,-3.275) (-4.024,-3.225) (-4.029,-3.175) (-4.033,-3.125) (-4.037,-3.075) (-4.039,-3.025) (-4.041,-2.975) (-4.042,-2.925) (-4.042,-2.875) (-4.041,-2.825) (-4.040,-2.775) (-4.037,-2.725) (-4.034,-2.675) (-4.030,-2.625) (-4.025,-2.575) (-4.020,-2.525) (-4.013,-2.475) (-4.005,-2.425) (-3.997,-2.375) (-3.988,-2.325) (-3.977,-2.275) (-3.966,-2.225) (-3.954,-2.175) (-3.941,-2.125) (-3.927,-2.075) (-3.912,-2.025) (-3.896,-1.975) (-3.878,-1.925) (-3.860,-1.875) (-3.841,-1.825) (-3.820,-1.775) (-3.798,-1.725) (-3.775,-1.675) (-3.751,-1.625) (-3.725,-1.575) (-3.700,-1.528) (-3.675,-1.483) (-3.640,-1.425) (-3.609,-1.375) (-3.576,-1.325) (-3.542,-1.275) (-3.505,-1.225) (-3.475,-1.185) (-3.427,-1.125) (-3.385,-1.075) (-3.350,-1.035) (-3.300,-0.981) (-3.250,-0.930) (-3.200,-0.882) (-3.150,-0.836) (-3.110,-0.800) (-3.050,-0.750) (-3.000,-0.711) (-2.953,-0.675) (-2.900,-0.637) (-2.850,-0.602) (-2.800,-0.570) (-2.750,-0.538) (-2.700,-0.508) (-2.650,-0.480) (-2.600,-0.452) (-2.550,-0.426) (-2.500,-0.401) (-2.450,-0.378) (-2.400,-0.355) (-2.350,-0.333) (-2.300,-0.313) (-2.250,-0.293) (-2.202,-0.275) (-2.150,-0.256) (-2.100,-0.239) (-2.050,-0.223) (-2.000,-0.208) (-1.950,-0.193) (-1.900,-0.179) (-1.850,-0.166) (-1.800,-0.154) (-1.750,-0.142) (-1.700,-0.131) (-1.650,-0.121) (-1.600,-0.111) (-1.550,-0.102) (-1.500,-0.093) (-1.450,-0.085) (-1.400,-0.077) (-1.350,-0.070) (-1.300,-0.063) (-1.250,-0.057) (-1.200,-0.052) (-1.150,-0.046) (-1.100,-0.041) (-1.050,-0.037) (-1.000,-0.033) (-0.950,-0.029) (-0.900,-0.026) (-0.850,-0.022) (-0.800,-0.020) (-0.750,-0.017) (-0.700,-0.015) (-0.650,-0.013) (-0.600,-0.011) (-0.550,-0.010) (-0.500,-0.008) (-0.450,-0.007) (-0.400,-0.006) (-0.350,-0.005) (-0.300,-0.005) (-0.250,-0.004)}
\EQ{(0.250,-0.003) (0.300,-0.003) (0.350,-0.004) (0.400,-0.005) (0.450,-0.005) (0.500,-0.006) (0.550,-0.007) (0.600,-0.008) (0.650,-0.010) (0.700,-0.011) (0.750,-0.013) (0.800,-0.015) (0.850,-0.017) (0.900,-0.019) (0.950,-0.022) (1.000,-0.025) (1.050,-0.028) (1.100,-0.031) (1.150,-0.035) (1.200,-0.039) (1.250,-0.043) (1.300,-0.047) (1.350,-0.052) (1.400,-0.058) (1.450,-0.063) (1.500,-0.070) (1.550,-0.076) (1.600,-0.083) (1.650,-0.090) (1.700,-0.098) (1.750,-0.106) (1.800,-0.115) (1.850,-0.124) (1.900,-0.134) (1.950,-0.144) (2.000,-0.155) (2.050,-0.166) (2.100,-0.178) (2.150,-0.191) (2.200,-0.204) (2.250,-0.218) (2.300,-0.232) (2.350,-0.247) (2.400,-0.263) (2.450,-0.279) (2.500,-0.296) (2.550,-0.314) (2.600,-0.333) (2.650,-0.353) (2.700,-0.373) (2.750,-0.394) (2.800,-0.416) (2.850,-0.439) (2.900,-0.463) (2.950,-0.488) (3.000,-0.514) (3.050,-0.541) (3.100,-0.569) (3.150,-0.599) (3.194,-0.625) (3.250,-0.661) (3.300,-0.694) (3.346,-0.725) (3.400,-0.764) (3.448,-0.800) (3.500,-0.840) (3.543,-0.875) (3.600,-0.923) (3.650,-0.967) (3.700,-1.013) (3.738,-1.050) (3.775,-1.087) (3.812,-1.125) (3.850,-1.165) (3.900,-1.221) (3.946,-1.275) (3.987,-1.325) (4.025,-1.373) (4.064,-1.425) (4.100,-1.475) (4.134,-1.525) (4.167,-1.575) (4.199,-1.625) (4.225,-1.668) (4.258,-1.725) (4.286,-1.775) (4.312,-1.825) (4.338,-1.875) (4.362,-1.925) (4.385,-1.975) (4.407,-2.025) (4.428,-2.075) (4.448,-2.125) (4.467,-2.175) (4.485,-2.225) (4.502,-2.275) (4.519,-2.325) (4.534,-2.375) (4.548,-2.425) (4.562,-2.475) (4.575,-2.525) (4.587,-2.575) (4.598,-2.625) (4.608,-2.675) (4.617,-2.725) (4.626,-2.775) (4.634,-2.825) (4.641,-2.875) (4.647,-2.925) (4.653,-2.975) (4.658,-3.025) (4.662,-3.075) (4.665,-3.125) (4.668,-3.175) (4.670,-3.225) (4.671,-3.275) (4.671,-3.325) (4.671,-3.375) (4.670,-3.425) (4.668,-3.475) (4.666,-3.525) (4.663,-3.575) (4.659,-3.625) (4.655,-3.675) (4.650,-3.725) (4.644,-3.775) (4.637,-3.825) (4.630,-3.875) (4.622,-3.925) (4.613,-3.975) (4.604,-4.025) (4.594,-4.075) (4.583,-4.125) (4.572,-4.175) (4.560,-4.225) (4.547,-4.275) (4.533,-4.325) (4.519,-4.375) (4.504,-4.425) (4.488,-4.475) (4.472,-4.525) (4.454,-4.575) (4.436,-4.625) (4.418,-4.675) (4.398,-4.725) (4.378,-4.775) (4.357,-4.825) (4.335,-4.875) (4.312,-4.925) (4.288,-4.975) (4.264,-5.025) (4.239,-5.075) (4.213,-5.125) (4.186,-5.175) (4.158,-5.225) (4.130,-5.275) (4.100,-5.325) (4.069,-5.375) (4.038,-5.425) (4.005,-5.475) (3.975,-5.520) (3.937,-5.575) (3.902,-5.625) (3.865,-5.675) (3.827,-5.725) (3.789,-5.775) (3.750,-5.823) (3.707,-5.875) (3.665,-5.925) (3.625,-5.970) (3.576,-6.025) (3.529,-6.075) (3.481,-6.125) (3.432,-6.175) (3.381,-6.225) (3.328,-6.275) (3.275,-6.324) (3.225,-6.368) (3.175,-6.412) (3.129,-6.450) (3.075,-6.495) (3.025,-6.534) (2.975,-6.573) (2.925,-6.611) (2.875,-6.647) (2.825,-6.683) (2.775,-6.717) (2.726,-6.750) (2.675,-6.784) (2.625,-6.816) (2.575,-6.846) (2.528,-6.875) (2.475,-6.906) (2.425,-6.935) (2.375,-6.962) (2.325,-6.989) (2.275,-7.015) (2.225,-7.041) (2.175,-7.066) (2.125,-7.090) (2.075,-7.113) (2.025,-7.136) (1.975,-7.158) (1.925,-7.179) (1.875,-7.200) (1.825,-7.220) (1.775,-7.239) (1.725,-7.258) (1.678,-7.275) (1.625,-7.294) (1.575,-7.311) (1.525,-7.327) (1.475,-7.343) (1.425,-7.358) (1.375,-7.372) (1.325,-7.386) (1.275,-7.400) (1.225,-7.413) (1.176,-7.425) (1.125,-7.437) (1.075,-7.448) (1.025,-7.459) (0.975,-7.469) (0.925,-7.479) (0.875,-7.488) (0.825,-7.497) (0.775,-7.505) (0.725,-7.512) (0.675,-7.519) (0.625,-7.526) (0.575,-7.532) (0.525,-7.538) (0.475,-7.543) (0.425,-7.547) (0.375,-7.551) (0.325,-7.555) (0.275,-7.558) (0.225,-7.560) (0.175,-7.562) (0.125,-7.564) (0.075,-7.565) (0.025,-7.565) (-0.050,-7.565) (-0.100,-7.564) (-0.150,-7.563) (-0.200,-7.561) (-0.250,-7.559) (-0.300,-7.556) (-0.350,-7.553) (-0.400,-7.549) (-0.450,-7.545) (-0.500,-7.540) (-0.550,-7.535) (-0.600,-7.529) (-0.650,-7.523) (-0.700,-7.516) (-0.750,-7.509) (-0.800,-7.501) (-0.850,-7.492) (-0.900,-7.484) (-0.950,-7.474) (-1.000,-7.464) (-1.050,-7.454) (-1.100,-7.443) (-1.150,-7.431) (-1.200,-7.419) (-1.250,-7.406) (-1.300,-7.393) (-1.350,-7.380) (-1.400,-7.365) (-1.450,-7.350) (-1.500,-7.335) (-1.550,-7.319) (-1.600,-7.302) (-1.650,-7.285) (-1.700,-7.267) (-1.750,-7.249) (-1.800,-7.229) (-1.850,-7.210) (-1.900,-7.189) (-1.950,-7.168) (-2.000,-7.147) (-2.049,-7.125) (-2.100,-7.101) (-2.150,-7.078) (-2.200,-7.053) (-2.250,-7.028) (-2.300,-7.002) (-2.350,-6.976) (-2.397,-6.950) (-2.442,-6.925) (-2.485,-6.900) (-2.550,-6.862) (-2.600,-6.831) (-2.649,-6.800) (-2.700,-6.767) (-2.750,-6.734) (-2.800,-6.700) (-2.850,-6.665) (-2.900,-6.629) (-2.950,-6.592) (-3.000,-6.554) (-3.050,-6.515) (-3.099,-6.475) (-3.150,-6.433) (-3.188,-6.400) (-3.246,-6.350) (-3.300,-6.301) (-3.350,-6.254) (-3.400,-6.206) (-3.450,-6.157) (-3.500,-6.106) (-3.550,-6.053) (-3.599,-6.000) (-3.643,-5.950) (-3.686,-5.900) (-3.725,-5.854) (-3.769,-5.800) (-3.800,-5.761) (-3.847,-5.700) (-3.884,-5.650) (-3.920,-5.600) (-3.950,-5.557) (-3.989,-5.500) (-4.022,-5.450) (-4.050,-5.406) (-4.085,-5.350) (-4.115,-5.300) (-4.144,-5.250) (-4.172,-5.200) (-4.200,-5.150) (-4.225,-5.102) (-4.250,-5.053) (-4.275,-5.003) (-4.300,-4.951) (-4.323,-4.900) (-4.346,-4.850) (-4.367,-4.800) (-4.388,-4.750) (-4.408,-4.700) (-4.427,-4.650) (-4.445,-4.600) (-4.463,-4.550) (-4.480,-4.500) (-4.496,-4.450) (-4.511,-4.400) (-4.526,-4.350) (-4.540,-4.300) (-4.553,-4.250) (-4.566,-4.200) (-4.578,-4.150) (-4.589,-4.100) (-4.599,-4.050) (-4.609,-4.000) (-4.618,-3.950) (-4.626,-3.900) (-4.634,-3.850) (-4.641,-3.800) (-4.647,-3.750) (-4.652,-3.700) (-4.657,-3.650) (-4.661,-3.600) (-4.665,-3.550) (-4.667,-3.500) (-4.669,-3.450) (-4.671,-3.400) (-4.671,-3.350) (-4.671,-3.300) (-4.670,-3.250) (-4.669,-3.200) (-4.667,-3.150) (-4.664,-3.100) (-4.660,-3.050) (-4.655,-3.000) (-4.650,-2.950) (-4.644,-2.900) (-4.638,-2.850) (-4.630,-2.800) (-4.622,-2.750) (-4.613,-2.700) (-4.603,-2.650) (-4.592,-2.600) (-4.581,-2.550) (-4.569,-2.500) (-4.555,-2.450) (-4.541,-2.400) (-4.526,-2.350) (-4.511,-2.300) (-4.494,-2.250) (-4.476,-2.200) (-4.458,-2.150) (-4.438,-2.100) (-4.418,-2.050) (-4.396,-2.000) (-4.375,-1.953) (-4.350,-1.900) (-4.325,-1.850) (-4.300,-1.802) (-4.275,-1.755) (-4.250,-1.711) (-4.214,-1.650) (-4.183,-1.600) (-4.151,-1.550) (-4.117,-1.500) (-4.082,-1.450) (-4.050,-1.406) (-4.007,-1.350) (-3.975,-1.310) (-3.925,-1.250) (-3.882,-1.200) (-3.836,-1.150) (-3.800,-1.112) (-3.764,-1.075) (-3.725,-1.037) (-3.686,-1.000) (-3.631,-0.950) (-3.575,-0.902) (-3.525,-0.860) (-3.481,-0.825) (-3.425,-0.782) (-3.381,-0.750) (-3.325,-0.711) (-3.275,-0.677) (-3.225,-0.645) (-3.175,-0.614) (-3.125,-0.584) (-3.075,-0.555) (-3.025,-0.527) (-2.975,-0.501) (-2.925,-0.476) (-2.875,-0.451) (-2.825,-0.428) (-2.775,-0.405) (-2.725,-0.383) (-2.675,-0.363) (-2.625,-0.343) (-2.575,-0.324) (-2.525,-0.305) (-2.475,-0.288) (-2.425,-0.271) (-2.375,-0.255) (-2.325,-0.239) (-2.276,-0.225) (-2.225,-0.211) (-2.175,-0.197) (-2.125,-0.184) (-2.075,-0.172) (-2.025,-0.161) (-1.975,-0.149) (-1.925,-0.139) (-1.875,-0.129) (-1.825,-0.120) (-1.775,-0.111) (-1.725,-0.102) (-1.675,-0.094) (-1.625,-0.087) (-1.575,-0.079) (-1.525,-0.073) (-1.475,-0.066) (-1.425,-0.061) (-1.375,-0.055) (-1.325,-0.050) (-1.275,-0.045) (-1.225,-0.041) (-1.175,-0.037) (-1.125,-0.033) (-1.075,-0.029) (-1.025,-0.026) (-0.975,-0.023) (-0.925,-0.020) (-0.875,-0.018) (-0.825,-0.016) (-0.775,-0.014) (-0.725,-0.012) (-0.675,-0.010) (-0.625,-0.009) (-0.575,-0.008) (-0.525,-0.007) (-0.475,-0.006) (-0.425,-0.005) (-0.375,-0.004) (-0.325,-0.004) (-0.275,-0.003) (-0.250,-0.003)}
\EQ{(0.250,-0.002) (0.300,-0.002) (0.350,-0.003) (0.400,-0.003) (0.450,-0.004) (0.500,-0.004) (0.550,-0.005) (0.600,-0.006) (0.650,-0.006) (0.700,-0.007) (0.750,-0.009) (0.800,-0.010) (0.850,-0.011) (0.900,-0.013) (0.950,-0.014) (1.000,-0.016) (1.050,-0.018) (1.100,-0.021) (1.150,-0.023) (1.200,-0.026) (1.250,-0.029) (1.300,-0.032) (1.350,-0.035) (1.400,-0.038) (1.450,-0.042) (1.500,-0.046) (1.550,-0.051) (1.600,-0.055) (1.650,-0.060) (1.700,-0.065) (1.750,-0.071) (1.800,-0.076) (1.850,-0.083) (1.900,-0.089) (1.950,-0.096) (2.000,-0.103) (2.050,-0.110) (2.100,-0.118) (2.150,-0.126) (2.200,-0.135) (2.250,-0.144) (2.300,-0.153) (2.350,-0.163) (2.400,-0.174) (2.450,-0.184) (2.500,-0.196) (2.550,-0.207) (2.600,-0.219) (2.650,-0.232) (2.700,-0.245) (2.750,-0.259) (2.800,-0.273) (2.850,-0.287) (2.900,-0.303) (2.950,-0.319) (3.000,-0.335) (3.050,-0.352) (3.100,-0.369) (3.150,-0.388) (3.200,-0.407) (3.247,-0.425) (3.300,-0.446) (3.350,-0.467) (3.400,-0.489) (3.450,-0.511) (3.500,-0.534) (3.550,-0.558) (3.600,-0.583) (3.650,-0.608) (3.700,-0.635) (3.750,-0.662) (3.800,-0.690) (3.850,-0.719) (3.900,-0.749) (3.950,-0.780) (4.000,-0.812) (4.050,-0.846) (4.093,-0.875) (4.150,-0.915) (4.197,-0.950) (4.250,-0.990) (4.295,-1.025) (4.350,-1.070) (4.400,-1.112) (4.444,-1.150) (4.500,-1.200) (4.550,-1.247) (4.600,-1.295) (4.650,-1.346) (4.700,-1.398) (4.748,-1.450) (4.792,-1.500) (4.825,-1.539) (4.875,-1.599) (4.916,-1.650) (4.950,-1.695) (4.991,-1.750) (5.025,-1.797) (5.061,-1.850) (5.095,-1.900) (5.125,-1.947) (5.158,-2.000) (5.188,-2.050) (5.217,-2.100) (5.245,-2.150) (5.272,-2.200) (5.298,-2.250) (5.323,-2.300) (5.347,-2.350) (5.370,-2.400) (5.393,-2.450) (5.414,-2.500) (5.435,-2.550) (5.455,-2.600) (5.474,-2.650) (5.492,-2.700) (5.510,-2.750) (5.526,-2.800) (5.542,-2.850) (5.558,-2.900) (5.572,-2.950) (5.586,-3.000) (5.599,-3.050) (5.612,-3.100) (5.623,-3.150) (5.634,-3.200) (5.645,-3.250) (5.655,-3.300) (5.664,-3.350) (5.672,-3.400) (5.680,-3.450) (5.687,-3.500) (5.694,-3.550) (5.700,-3.600) (5.705,-3.650) (5.710,-3.700) (5.714,-3.750) (5.717,-3.800) (5.720,-3.850) (5.722,-3.900) (5.724,-3.950) (5.725,-4.000) (5.726,-4.050) (5.726,-4.100) (5.725,-4.150) (5.724,-4.200) (5.722,-4.250) (5.720,-4.300) (5.717,-4.350) (5.714,-4.400) (5.710,-4.450) (5.705,-4.500) (5.700,-4.550) (5.694,-4.600) (5.688,-4.650) (5.681,-4.700) (5.673,-4.750) (5.665,-4.800) (5.657,-4.850) (5.648,-4.900) (5.638,-4.950) (5.628,-5.000) (5.617,-5.050) (5.605,-5.100) (5.593,-5.150) (5.581,-5.200) (5.567,-5.250) (5.553,-5.300) (5.539,-5.350) (5.524,-5.400) (5.509,-5.450) (5.492,-5.500) (5.476,-5.550) (5.458,-5.600) (5.440,-5.650) (5.422,-5.700) (5.402,-5.750) (5.382,-5.800) (5.362,-5.850) (5.341,-5.900) (5.319,-5.950) (5.296,-6.000) (5.275,-6.046) (5.250,-6.099) (5.225,-6.150) (5.200,-6.200) (5.175,-6.248) (5.150,-6.295) (5.125,-6.341) (5.100,-6.386) (5.075,-6.430) (5.049,-6.475) (5.019,-6.525) (4.988,-6.575) (4.957,-6.625) (4.925,-6.674) (4.892,-6.725) (4.858,-6.775) (4.825,-6.822) (4.787,-6.875) (4.751,-6.925) (4.714,-6.975) (4.675,-7.025) (4.636,-7.075) (4.600,-7.120) (4.555,-7.175) (4.513,-7.225) (4.475,-7.269) (4.425,-7.325) (4.380,-7.375) (4.333,-7.425) (4.286,-7.475) (4.250,-7.512) (4.212,-7.550) (4.175,-7.586) (4.135,-7.625) (4.082,-7.675) (4.028,-7.725) (3.975,-7.773) (3.925,-7.817) (3.886,-7.850) (3.827,-7.900) (3.775,-7.942) (3.734,-7.975) (3.675,-8.021) (3.625,-8.059) (3.575,-8.097) (3.525,-8.133) (3.475,-8.169) (3.430,-8.200) (3.375,-8.238) (3.325,-8.271) (3.281,-8.300) (3.225,-8.335) (3.175,-8.367) (3.125,-8.397) (3.078,-8.425) (3.025,-8.456) (2.975,-8.484) (2.925,-8.512) (2.875,-8.540) (2.825,-8.566) (2.775,-8.592) (2.725,-8.617) (2.675,-8.642) (2.625,-8.666) (2.575,-8.690) (2.525,-8.713) (2.475,-8.736) (2.425,-8.758) (2.375,-8.779) (2.325,-8.800) (2.275,-8.820) (2.225,-8.840) (2.175,-8.859) (2.125,-8.878) (2.075,-8.897) (2.025,-8.914) (1.975,-8.932) (1.925,-8.948) (1.875,-8.965) (1.825,-8.981) (1.775,-8.996) (1.725,-9.011) (1.677,-9.025) (1.625,-9.039) (1.575,-9.053) (1.525,-9.066) (1.475,-9.079) (1.425,-9.091) (1.375,-9.103) (1.325,-9.114) (1.275,-9.125) (1.225,-9.135) (1.175,-9.145) (1.125,-9.155) (1.075,-9.164) (1.025,-9.172) (0.975,-9.181) (0.925,-9.189) (0.875,-9.196) (0.825,-9.203) (0.775,-9.209) (0.725,-9.216) (0.675,-9.221) (0.625,-9.227) (0.575,-9.232) (0.525,-9.236) (0.475,-9.240) (0.425,-9.244) (0.375,-9.247) (0.325,-9.250) (0.275,-9.252) (0.225,-9.254) (0.175,-9.256) (0.125,-9.257) (0.075,-9.258) (0.025,-9.258) (-0.050,-9.258) (-0.100,-9.258) (-0.150,-9.257) (-0.200,-9.255) (-0.250,-9.253) (-0.300,-9.251) (-0.350,-9.249) (-0.400,-9.246) (-0.450,-9.242) (-0.500,-9.238) (-0.550,-9.234) (-0.600,-9.229) (-0.650,-9.224) (-0.700,-9.219) (-0.750,-9.213) (-0.800,-9.206) (-0.850,-9.199) (-0.900,-9.192) (-0.950,-9.185) (-1.000,-9.177) (-1.050,-9.168) (-1.100,-9.159) (-1.149,-9.150) (-1.200,-9.140) (-1.250,-9.130) (-1.300,-9.119) (-1.350,-9.108) (-1.400,-9.097) (-1.450,-9.085) (-1.500,-9.072) (-1.550,-9.060) (-1.600,-9.046) (-1.650,-9.033) (-1.700,-9.018) (-1.750,-9.004) (-1.800,-8.988) (-1.850,-8.973) (-1.900,-8.957) (-1.950,-8.940) (-2.000,-8.923) (-2.050,-8.906) (-2.100,-8.887) (-2.150,-8.869) (-2.199,-8.850) (-2.250,-8.830) (-2.300,-8.810) (-2.350,-8.790) (-2.400,-8.768) (-2.450,-8.747) (-2.499,-8.725) (-2.550,-8.702) (-2.600,-8.678) (-2.650,-8.654) (-2.700,-8.630) (-2.750,-8.605) (-2.800,-8.579) (-2.850,-8.553) (-2.900,-8.526) (-2.947,-8.500) (-2.992,-8.475) (-3.035,-8.450) (-3.100,-8.412) (-3.150,-8.382) (-3.200,-8.351) (-3.250,-8.320) (-3.300,-8.287) (-3.350,-8.254) (-3.394,-8.225) (-3.450,-8.186) (-3.500,-8.151) (-3.550,-8.115) (-3.600,-8.078) (-3.650,-8.040) (-3.700,-8.002) (-3.750,-7.962) (-3.796,-7.925) (-3.850,-7.881) (-3.900,-7.838) (-3.944,-7.800) (-4.000,-7.750) (-4.050,-7.705) (-4.100,-7.658) (-4.150,-7.611) (-4.187,-7.575) (-4.225,-7.537) (-4.261,-7.500) (-4.300,-7.460) (-4.350,-7.407) (-4.400,-7.353) (-4.448,-7.300) (-4.491,-7.250) (-4.525,-7.211) (-4.575,-7.151) (-4.616,-7.100) (-4.650,-7.058) (-4.695,-7.000) (-4.725,-6.960) (-4.769,-6.900) (-4.800,-6.857) (-4.840,-6.800) (-4.875,-6.750) (-4.908,-6.700) (-4.941,-6.650) (-4.973,-6.600) (-5.000,-6.556) (-5.034,-6.500) (-5.063,-6.450) (-5.092,-6.400) (-5.120,-6.350) (-5.147,-6.300) (-5.174,-6.250) (-5.200,-6.200) (-5.225,-6.150) (-5.250,-6.100) (-5.273,-6.050) (-5.296,-6.000) (-5.319,-5.950) (-5.341,-5.900) (-5.362,-5.850) (-5.382,-5.800) (-5.402,-5.750) (-5.422,-5.700) (-5.440,-5.650) (-5.458,-5.600) (-5.475,-5.552) (-5.492,-5.500) (-5.509,-5.450) (-5.524,-5.400) (-5.539,-5.350) (-5.553,-5.300) (-5.567,-5.250) (-5.581,-5.200) (-5.593,-5.150) (-5.605,-5.100) (-5.617,-5.050) (-5.628,-5.000) (-5.638,-4.950) (-5.648,-4.900) (-5.657,-4.850) (-5.665,-4.800) (-5.673,-4.750) (-5.681,-4.700) (-5.688,-4.650) (-5.694,-4.600) (-5.700,-4.550) (-5.705,-4.500) (-5.710,-4.450) (-5.714,-4.400) (-5.717,-4.350) (-5.720,-4.300) (-5.722,-4.250) (-5.724,-4.200) (-5.725,-4.150) (-5.726,-4.100) (-5.726,-4.050) (-5.725,-4.000) (-5.724,-3.950) (-5.722,-3.900) (-5.720,-3.850) (-5.717,-3.800) (-5.714,-3.750) (-5.710,-3.700) (-5.705,-3.650) (-5.700,-3.600) (-5.694,-3.550) (-5.687,-3.500) (-5.680,-3.450) (-5.672,-3.400) (-5.664,-3.350) (-5.655,-3.300) (-5.645,-3.250) (-5.634,-3.200) (-5.623,-3.150) (-5.612,-3.100) (-5.599,-3.050) (-5.586,-3.000) (-5.572,-2.950) (-5.558,-2.900) (-5.542,-2.850) (-5.526,-2.800) (-5.510,-2.750) (-5.492,-2.700) (-5.474,-2.650) (-5.455,-2.600) (-5.435,-2.550) (-5.414,-2.500) (-5.393,-2.450) (-5.370,-2.400) (-5.347,-2.350) (-5.325,-2.304) (-5.300,-2.254) (-5.275,-2.206) (-5.250,-2.159) (-5.225,-2.114) (-5.200,-2.071) (-5.173,-2.025) (-5.143,-1.975) (-5.111,-1.925) (-5.078,-1.875) (-5.050,-1.833) (-5.009,-1.775) (-4.975,-1.728) (-4.935,-1.675) (-4.900,-1.630) (-4.855,-1.575) (-4.814,-1.525) (-4.775,-1.481) (-4.725,-1.425) (-4.678,-1.375) (-4.630,-1.325) (-4.579,-1.275) (-4.527,-1.225) (-4.475,-1.177) (-4.425,-1.133) (-4.386,-1.100) (-4.326,-1.050) (-4.275,-1.009) (-4.231,-0.975) (-4.175,-0.933) (-4.129,-0.900) (-4.075,-0.863) (-4.025,-0.829) (-3.981,-0.800) (-3.925,-0.765) (-3.875,-0.734) (-3.825,-0.705) (-3.775,-0.676) (-3.728,-0.650) (-3.682,-0.625) (-3.634,-0.600) (-3.584,-0.575) (-3.533,-0.550) (-3.480,-0.525) (-3.425,-0.500) (-3.375,-0.478) (-3.325,-0.457) (-3.275,-0.436) (-3.225,-0.416) (-3.175,-0.397) (-3.125,-0.379) (-3.075,-0.361) (-3.025,-0.343) (-2.975,-0.327) (-2.925,-0.311) (-2.875,-0.295) (-2.825,-0.280) (-2.775,-0.266) (-2.725,-0.252) (-2.675,-0.238) (-2.625,-0.226) (-2.575,-0.213) (-2.525,-0.201) (-2.475,-0.190) (-2.425,-0.179) (-2.375,-0.168) (-2.325,-0.158) (-2.275,-0.149) (-2.225,-0.139) (-2.175,-0.131) (-2.125,-0.122) (-2.075,-0.114) (-2.025,-0.107) (-1.975,-0.099) (-1.925,-0.092) (-1.875,-0.086) (-1.825,-0.079) (-1.775,-0.074) (-1.725,-0.068) (-1.675,-0.063) (-1.625,-0.058) (-1.575,-0.053) (-1.525,-0.048) (-1.475,-0.044) (-1.425,-0.040) (-1.375,-0.037) (-1.325,-0.033) (-1.275,-0.030) (-1.225,-0.027) (-1.175,-0.024) (-1.125,-0.022) (-1.075,-0.020) (-1.025,-0.017) (-0.975,-0.015) (-0.925,-0.014) (-0.875,-0.012) (-0.825,-0.011) (-0.775,-0.009) (-0.725,-0.008) (-0.675,-0.007) (-0.625,-0.006) (-0.575,-0.005) (-0.525,-0.004) (-0.475,-0.004) (-0.425,-0.003) (-0.375,-0.003) (-0.325,-0.002) (-0.275,-0.002) (-0.250,-0.002)}
\EQ{(-6.558,-10.000) (-6.597,-9.950) (-6.636,-9.900) (-6.674,-9.850) (-6.712,-9.800) (-6.748,-9.750) (-6.785,-9.700) (-6.820,-9.650) (-6.850,-9.607) (-6.889,-9.550) (-6.923,-9.500) (-6.956,-9.450) (-6.989,-9.400) (-7.021,-9.350) (-7.050,-9.304) (-7.083,-9.250) (-7.114,-9.200) (-7.144,-9.150) (-7.173,-9.100) (-7.200,-9.053) (-7.225,-9.009) (-7.250,-8.964) (-7.275,-8.918) (-7.300,-8.872) (-7.325,-8.825) (-7.350,-8.776) (-7.375,-8.727) (-7.400,-8.677) (-7.425,-8.626) (-7.449,-8.575) (-7.473,-8.525) (-7.496,-8.475) (-7.519,-8.425) (-7.541,-8.375) (-7.562,-8.325) (-7.584,-8.275) (-7.604,-8.225) (-7.625,-8.175) (-7.645,-8.125) (-7.664,-8.075) (-7.683,-8.025) (-7.701,-7.975) (-7.719,-7.925) (-7.737,-7.875) (-7.754,-7.825) (-7.771,-7.775) (-7.787,-7.725) (-7.803,-7.675) (-7.819,-7.625) (-7.834,-7.575) (-7.849,-7.525) (-7.863,-7.475) (-7.877,-7.425) (-7.890,-7.375) (-7.903,-7.325) (-7.916,-7.275) (-7.928,-7.225) (-7.940,-7.175) (-7.951,-7.125) (-7.962,-7.075) (-7.972,-7.025) (-7.983,-6.975) (-7.992,-6.925) (-8.002,-6.875) (-8.010,-6.825) (-8.019,-6.775) (-8.027,-6.725) (-8.035,-6.675) (-8.042,-6.625) (-8.049,-6.575) (-8.055,-6.525) (-8.062,-6.475) (-8.067,-6.425) (-8.073,-6.375) (-8.077,-6.325) (-8.082,-6.275) (-8.086,-6.225) (-8.090,-6.175) (-8.093,-6.125) (-8.096,-6.075) (-8.098,-6.025) (-8.100,-5.975) (-8.102,-5.925) (-8.103,-5.875) (-8.104,-5.825) (-8.105,-5.775) (-8.105,-5.725) (-8.104,-5.675) (-8.103,-5.625) (-8.102,-5.575) (-8.100,-5.525) (-8.098,-5.475) (-8.096,-5.425) (-8.093,-5.375) (-8.090,-5.325) (-8.086,-5.275) (-8.082,-5.225) (-8.077,-5.175) (-8.072,-5.125) (-8.067,-5.075) (-8.061,-5.025) (-8.055,-4.975) (-8.048,-4.925) (-8.041,-4.875) (-8.033,-4.825) (-8.025,-4.775) (-8.016,-4.725) (-8.007,-4.675) (-7.998,-4.625) (-7.988,-4.575) (-7.977,-4.525) (-7.967,-4.475) (-7.955,-4.425) (-7.943,-4.375) (-7.931,-4.325) (-7.918,-4.275) (-7.905,-4.225) (-7.891,-4.175) (-7.877,-4.125) (-7.862,-4.075) (-7.847,-4.025) (-7.831,-3.975) (-7.815,-3.925) (-7.798,-3.875) (-7.780,-3.825) (-7.762,-3.775) (-7.744,-3.725) (-7.725,-3.676) (-7.705,-3.625) (-7.685,-3.575) (-7.664,-3.525) (-7.643,-3.475) (-7.621,-3.425) (-7.600,-3.379) (-7.575,-3.325) (-7.551,-3.275) (-7.526,-3.225) (-7.501,-3.175) (-7.475,-3.125) (-7.450,-3.077) (-7.425,-3.031) (-7.400,-2.986) (-7.375,-2.942) (-7.336,-2.875) (-7.306,-2.825) (-7.275,-2.775) (-7.243,-2.725) (-7.211,-2.675) (-7.178,-2.625) (-7.144,-2.575) (-7.109,-2.525) (-7.075,-2.478) (-7.036,-2.425) (-7.000,-2.377) (-6.960,-2.325) (-6.925,-2.281) (-6.880,-2.225) (-6.838,-2.175) (-6.800,-2.130) (-6.752,-2.075) (-6.707,-2.025) (-6.661,-1.975) (-6.625,-1.937) (-6.590,-1.900) (-6.550,-1.860) (-6.500,-1.810) (-6.464,-1.775) (-6.425,-1.738) (-6.384,-1.700) (-6.329,-1.650) (-6.275,-1.602) (-6.225,-1.559) (-6.185,-1.525) (-6.125,-1.476) (-6.075,-1.435) (-6.030,-1.400) (-5.975,-1.358) (-5.931,-1.325) (-5.875,-1.285) (-5.826,-1.250) (-5.775,-1.214) (-5.725,-1.180) (-5.679,-1.150) (-5.625,-1.115) (-5.575,-1.083) (-5.525,-1.052) (-5.480,-1.025) (-5.425,-0.993) (-5.375,-0.964) (-5.325,-0.936) (-5.275,-0.908) (-5.225,-0.881) (-5.175,-0.855) (-5.125,-0.830) (-5.075,-0.805) (-5.025,-0.780) (-4.975,-0.756) (-4.925,-0.733) (-4.875,-0.710) (-4.825,-0.688) (-4.775,-0.666) (-4.725,-0.645) (-4.678,-0.625) (-4.625,-0.604) (-4.575,-0.584) (-4.525,-0.565) (-4.475,-0.546) (-4.425,-0.527) (-4.375,-0.509) (-4.325,-0.492) (-4.276,-0.475) (-4.225,-0.458) (-4.175,-0.442) (-4.125,-0.426) (-4.075,-0.411) (-4.025,-0.396) (-3.975,-0.381) (-3.925,-0.367) (-3.875,-0.353) (-3.825,-0.340) (-3.775,-0.326) (-3.725,-0.314) (-3.675,-0.301) (-3.625,-0.289) (-3.575,-0.278) (-3.525,-0.266) (-3.475,-0.255) (-3.425,-0.244) (-3.375,-0.234) (-3.325,-0.224) (-3.275,-0.214) (-3.225,-0.204) (-3.175,-0.195) (-3.125,-0.186) (-3.075,-0.178) (-3.025,-0.169) (-2.975,-0.161) (-2.925,-0.153) (-2.875,-0.146) (-2.825,-0.139) (-2.775,-0.132) (-2.725,-0.125) (-2.675,-0.118) (-2.625,-0.112) (-2.575,-0.106) (-2.525,-0.100) (-2.475,-0.094) (-2.425,-0.089) (-2.375,-0.084) (-2.325,-0.079) (-2.275,-0.074) (-2.225,-0.069) (-2.175,-0.065) (-2.125,-0.061) (-2.075,-0.057) (-2.025,-0.053) (-1.975,-0.050) (-1.925,-0.046) (-1.875,-0.043) (-1.825,-0.040) (-1.775,-0.037) (-1.725,-0.034) (-1.675,-0.031) (-1.625,-0.029) (-1.575,-0.026) (-1.525,-0.024) (-1.475,-0.022) (-1.425,-0.020) (-1.375,-0.018) (-1.325,-0.017) (-1.275,-0.015) (-1.225,-0.014) (-1.175,-0.012) (-1.125,-0.011) (-1.075,-0.010) (-1.025,-0.009) (-0.975,-0.008) (-0.925,-0.007) (-0.875,-0.006) (-0.825,-0.005) (-0.775,-0.005) (-0.725,-0.004) (-0.675,-0.003) (-0.625,-0.003) (-0.575,-0.003) (-0.525,-0.002) (-0.475,-0.002) (-0.425,-0.002) (-0.375,-0.001) (-0.325,-0.001) (-0.275,-0.001) (-0.250,-0.001)}
\EQ{(0.250,-0.001) (0.300,-0.001) (0.350,-0.001) (0.400,-0.002) (0.450,-0.002) (0.500,-0.002) (0.550,-0.002) (0.600,-0.003) (0.650,-0.003) (0.700,-0.004) (0.750,-0.004) (0.800,-0.005) (0.850,-0.006) (0.900,-0.006) (0.950,-0.007) (1.000,-0.008) (1.050,-0.009) (1.100,-0.010) (1.150,-0.012) (1.200,-0.013) (1.250,-0.014) (1.300,-0.016) (1.350,-0.017) (1.400,-0.019) (1.450,-0.021) (1.500,-0.023) (1.550,-0.025) (1.600,-0.028) (1.650,-0.030) (1.700,-0.033) (1.750,-0.035) (1.800,-0.038) (1.850,-0.041) (1.900,-0.044) (1.950,-0.048) (2.000,-0.051) (2.050,-0.055) (2.100,-0.059) (2.150,-0.063) (2.200,-0.067) (2.250,-0.072) (2.300,-0.076) (2.350,-0.081) (2.400,-0.086) (2.450,-0.092) (2.500,-0.097) (2.550,-0.103) (2.600,-0.109) (2.650,-0.115) (2.700,-0.121) (2.750,-0.128) (2.800,-0.135) (2.850,-0.142) (2.900,-0.150) (2.950,-0.157) (3.000,-0.165) (3.050,-0.173) (3.100,-0.182) (3.150,-0.191) (3.200,-0.200) (3.250,-0.209) (3.300,-0.219) (3.350,-0.229) (3.400,-0.239) (3.450,-0.250) (3.500,-0.261) (3.550,-0.272) (3.600,-0.283) (3.650,-0.295) (3.700,-0.307) (3.750,-0.320) (3.800,-0.333) (3.850,-0.346) (3.900,-0.360) (3.950,-0.374) (4.000,-0.388) (4.050,-0.403) (4.100,-0.418) (4.150,-0.434) (4.200,-0.450) (4.250,-0.466) (4.300,-0.483) (4.348,-0.500) (4.400,-0.518) (4.450,-0.536) (4.500,-0.555) (4.550,-0.574) (4.600,-0.594) (4.650,-0.614) (4.700,-0.634) (4.750,-0.655) (4.796,-0.675) (4.850,-0.699) (4.900,-0.721) (4.950,-0.744) (5.000,-0.768) (5.050,-0.792) (5.100,-0.817) (5.150,-0.842) (5.200,-0.868) (5.250,-0.895) (5.300,-0.922) (5.350,-0.950) (5.394,-0.975) (5.450,-1.007) (5.500,-1.037) (5.550,-1.068) (5.600,-1.099) (5.650,-1.131) (5.700,-1.164) (5.750,-1.197) (5.800,-1.232) (5.850,-1.267) (5.896,-1.300) (5.950,-1.339) (5.997,-1.375) (6.050,-1.416) (6.093,-1.450) (6.150,-1.496) (6.200,-1.537) (6.244,-1.575) (6.300,-1.624) (6.350,-1.669) (6.400,-1.715) (6.438,-1.750) (6.475,-1.786) (6.515,-1.825) (6.565,-1.875) (6.600,-1.911) (6.637,-1.950) (6.675,-1.990) (6.725,-2.045) (6.774,-2.100) (6.817,-2.150) (6.850,-2.189) (6.900,-2.250) (6.940,-2.300) (6.975,-2.344) (7.018,-2.400) (7.050,-2.444) (7.091,-2.500) (7.125,-2.548) (7.161,-2.600) (7.194,-2.650) (7.225,-2.697) (7.259,-2.750) (7.291,-2.800) (7.321,-2.850) (7.350,-2.899) (7.380,-2.950) (7.408,-3.000) (7.435,-3.050) (7.462,-3.100) (7.488,-3.150) (7.514,-3.200) (7.539,-3.250) (7.563,-3.300) (7.587,-3.350) (7.609,-3.400) (7.632,-3.450) (7.653,-3.500) (7.675,-3.550) (7.695,-3.600) (7.715,-3.650) (7.734,-3.700) (7.753,-3.750) (7.771,-3.800) (7.789,-3.850) (7.806,-3.900) (7.823,-3.950) (7.839,-4.000) (7.855,-4.050) (7.870,-4.100) (7.884,-4.150) (7.898,-4.200) (7.912,-4.250) (7.925,-4.300) (7.937,-4.350) (7.949,-4.400) (7.961,-4.450) (7.972,-4.500) (7.983,-4.550) (7.993,-4.600) (8.003,-4.650) (8.012,-4.700) (8.021,-4.750) (8.029,-4.800) (8.037,-4.850) (8.044,-4.900) (8.051,-4.950) (8.058,-5.000) (8.064,-5.050) (8.070,-5.100) (8.075,-5.150) (8.080,-5.200) (8.084,-5.250) (8.088,-5.300) (8.091,-5.350) (8.095,-5.400) (8.097,-5.450) (8.099,-5.500) (8.101,-5.550) (8.103,-5.600) (8.104,-5.650) (8.104,-5.700) (8.105,-5.750) (8.104,-5.800) (8.104,-5.850) (8.103,-5.900) (8.101,-5.950) (8.099,-6.000) (8.097,-6.050) (8.094,-6.100) (8.091,-6.150) (8.088,-6.200) (8.084,-6.250) (8.080,-6.300) (8.075,-6.350) (8.070,-6.400) (8.064,-6.450) (8.059,-6.500) (8.052,-6.550) (8.046,-6.600) (8.039,-6.650) (8.031,-6.700) (8.023,-6.750) (8.015,-6.800) (8.006,-6.850) (7.997,-6.900) (7.987,-6.950) (7.978,-7.000) (7.967,-7.050) (7.956,-7.100) (7.945,-7.150) (7.934,-7.200) (7.922,-7.250) (7.909,-7.300) (7.897,-7.350) (7.883,-7.400) (7.870,-7.450) (7.856,-7.500) (7.841,-7.550) (7.826,-7.600) (7.811,-7.650) (7.795,-7.700) (7.779,-7.750) (7.763,-7.800) (7.746,-7.850) (7.728,-7.900) (7.710,-7.950) (7.692,-8.000) (7.673,-8.050) (7.654,-8.100) (7.635,-8.150) (7.615,-8.200) (7.594,-8.250) (7.573,-8.300) (7.552,-8.350) (7.530,-8.400) (7.507,-8.450) (7.485,-8.500) (7.461,-8.550) (7.438,-8.600) (7.413,-8.650) (7.389,-8.700) (7.363,-8.750) (7.338,-8.800) (7.311,-8.850) (7.285,-8.900) (7.258,-8.950) (7.230,-9.000) (7.202,-9.050) (7.175,-9.096) (7.144,-9.150) (7.114,-9.200) (7.083,-9.250) (7.052,-9.300) (7.025,-9.344) (6.989,-9.400) (6.956,-9.450) (6.925,-9.497) (6.889,-9.550) (6.855,-9.600) (6.825,-9.643) (6.785,-9.700) (6.750,-9.748) (6.712,-9.800) (6.675,-9.849) (6.636,-9.900) (6.600,-9.947) (6.558,-10.000)}
\EQ{(0.250,0.001) (0.300,0.001) (0.350,0.001) (0.400,0.002) (0.450,0.002) (0.500,0.002) (0.550,0.002) (0.600,0.003) (0.650,0.003) (0.700,0.004) (0.750,0.004) (0.800,0.005) (0.850,0.006) (0.900,0.006) (0.950,0.007) (1.000,0.008) (1.050,0.009) (1.100,0.010) (1.150,0.012) (1.200,0.013) (1.250,0.014) (1.300,0.016) (1.350,0.017) (1.400,0.019) (1.450,0.021) (1.500,0.023) (1.550,0.025) (1.600,0.028) (1.650,0.030) (1.700,0.033) (1.750,0.035) (1.800,0.038) (1.850,0.041) (1.900,0.044) (1.950,0.048) (2.000,0.051) (2.050,0.055) (2.100,0.059) (2.150,0.063) (2.200,0.067) (2.250,0.072) (2.300,0.076) (2.350,0.081) (2.400,0.086) (2.450,0.092) (2.500,0.097) (2.550,0.103) (2.600,0.109) (2.650,0.115) (2.700,0.121) (2.750,0.128) (2.800,0.135) (2.850,0.142) (2.900,0.150) (2.950,0.157) (3.000,0.165) (3.050,0.173) (3.100,0.182) (3.150,0.191) (3.200,0.200) (3.250,0.209) (3.300,0.219) (3.350,0.229) (3.400,0.239) (3.450,0.250) (3.500,0.261) (3.550,0.272) (3.600,0.283) (3.650,0.295) (3.700,0.307) (3.750,0.320) (3.800,0.333) (3.850,0.346) (3.900,0.360) (3.950,0.374) (4.000,0.388) (4.050,0.403) (4.100,0.418) (4.150,0.434) (4.200,0.450) (4.250,0.466) (4.300,0.483) (4.348,0.500) (4.400,0.518) (4.450,0.536) (4.500,0.555) (4.550,0.574) (4.600,0.594) (4.650,0.614) (4.700,0.634) (4.750,0.655) (4.796,0.675) (4.850,0.699) (4.900,0.721) (4.950,0.744) (5.000,0.768) (5.050,0.792) (5.100,0.817) (5.150,0.842) (5.200,0.868) (5.250,0.895) (5.300,0.922) (5.350,0.950) (5.394,0.975) (5.450,1.007) (5.500,1.037) (5.550,1.068) (5.600,1.099) (5.650,1.131) (5.700,1.164) (5.750,1.197) (5.800,1.232) (5.850,1.267) (5.896,1.300) (5.950,1.339) (5.997,1.375) (6.050,1.416) (6.093,1.450) (6.150,1.496) (6.200,1.537) (6.244,1.575) (6.300,1.624) (6.350,1.669) (6.400,1.715) (6.438,1.750) (6.475,1.786) (6.515,1.825) (6.565,1.875) (6.600,1.911) (6.637,1.950) (6.675,1.990) (6.725,2.045) (6.774,2.100) (6.817,2.150) (6.850,2.189) (6.900,2.250) (6.940,2.300) (6.975,2.344) (7.018,2.400) (7.050,2.444) (7.091,2.500) (7.125,2.548) (7.161,2.600) (7.194,2.650) (7.225,2.697) (7.259,2.750) (7.291,2.800) (7.321,2.850) (7.350,2.899) (7.380,2.950) (7.408,3.000) (7.435,3.050) (7.462,3.100) (7.488,3.150) (7.514,3.200) (7.539,3.250) (7.563,3.300) (7.587,3.350) (7.609,3.400) (7.632,3.450) (7.653,3.500) (7.675,3.550) (7.695,3.600) (7.715,3.650) (7.734,3.700) (7.753,3.750) (7.771,3.800) (7.789,3.850) (7.806,3.900) (7.823,3.950) (7.839,4.000) (7.855,4.050) (7.870,4.100) (7.884,4.150) (7.898,4.200) (7.912,4.250) (7.925,4.300) (7.937,4.350) (7.949,4.400) (7.961,4.450) (7.972,4.500) (7.983,4.550) (7.993,4.600) (8.003,4.650) (8.012,4.700) (8.021,4.750) (8.029,4.800) (8.037,4.850) (8.044,4.900) (8.051,4.950) (8.058,5.000) (8.064,5.050) (8.070,5.100) (8.075,5.150) (8.080,5.200) (8.084,5.250) (8.088,5.300) (8.091,5.350) (8.095,5.400) (8.097,5.450) (8.099,5.500) (8.101,5.550) (8.103,5.600) (8.104,5.650) (8.104,5.700) (8.105,5.750) (8.104,5.800) (8.104,5.850) (8.103,5.900) (8.101,5.950) (8.099,6.000) (8.097,6.050) (8.094,6.100) (8.091,6.150) (8.088,6.200) (8.084,6.250) (8.080,6.300) (8.075,6.350) (8.070,6.400) (8.064,6.450) (8.059,6.500) (8.052,6.550) (8.046,6.600) (8.039,6.650) (8.031,6.700) (8.023,6.750) (8.015,6.800) (8.006,6.850) (7.997,6.900) (7.987,6.950) (7.978,7.000) (7.967,7.050) (7.956,7.100) (7.945,7.150) (7.934,7.200) (7.922,7.250) (7.909,7.300) (7.897,7.350) (7.883,7.400) (7.870,7.450) (7.856,7.500) (7.841,7.550) (7.826,7.600) (7.811,7.650) (7.795,7.700) (7.779,7.750) (7.763,7.800) (7.746,7.850) (7.728,7.900) (7.710,7.950) (7.692,8.000) (7.673,8.050) (7.654,8.100) (7.635,8.150) (7.615,8.200) (7.594,8.250) (7.573,8.300) (7.552,8.350) (7.530,8.400) (7.507,8.450) (7.485,8.500) (7.461,8.550) (7.438,8.600) (7.413,8.650) (7.389,8.700) (7.363,8.750) (7.338,8.800) (7.311,8.850) (7.285,8.900) (7.258,8.950) (7.230,9.000) (7.202,9.050) (7.175,9.096) (7.144,9.150) (7.114,9.200) (7.083,9.250) (7.052,9.300) (7.025,9.344) (6.989,9.400) (6.956,9.450) (6.925,9.497) (6.889,9.550) (6.855,9.600) (6.825,9.643) (6.785,9.700) (6.750,9.748) (6.712,9.800) (6.675,9.849) (6.636,9.900) (6.600,9.947) (6.558,10.000)}
\EQ{(-6.558,10.000) (-6.597,9.950) (-6.636,9.900) (-6.674,9.850) (-6.712,9.800) (-6.748,9.750) (-6.785,9.700) (-6.820,9.650) (-6.850,9.607) (-6.889,9.550) (-6.923,9.500) (-6.956,9.450) (-6.989,9.400) (-7.021,9.350) (-7.050,9.304) (-7.083,9.250) (-7.114,9.200) (-7.144,9.150) (-7.173,9.100) (-7.200,9.053) (-7.225,9.009) (-7.250,8.964) (-7.275,8.918) (-7.300,8.872) (-7.325,8.825) (-7.350,8.776) (-7.375,8.727) (-7.400,8.677) (-7.425,8.626) (-7.449,8.575) (-7.473,8.525) (-7.496,8.475) (-7.519,8.425) (-7.541,8.375) (-7.562,8.325) (-7.584,8.275) (-7.604,8.225) (-7.625,8.175) (-7.645,8.125) (-7.664,8.075) (-7.683,8.025) (-7.701,7.975) (-7.719,7.925) (-7.737,7.875) (-7.754,7.825) (-7.771,7.775) (-7.787,7.725) (-7.803,7.675) (-7.819,7.625) (-7.834,7.575) (-7.849,7.525) (-7.863,7.475) (-7.877,7.425) (-7.890,7.375) (-7.903,7.325) (-7.916,7.275) (-7.928,7.225) (-7.940,7.175) (-7.951,7.125) (-7.962,7.075) (-7.972,7.025) (-7.983,6.975) (-7.992,6.925) (-8.002,6.875) (-8.010,6.825) (-8.019,6.775) (-8.027,6.725) (-8.035,6.675) (-8.042,6.625) (-8.049,6.575) (-8.055,6.525) (-8.062,6.475) (-8.067,6.425) (-8.073,6.375) (-8.077,6.325) (-8.082,6.275) (-8.086,6.225) (-8.090,6.175) (-8.093,6.125) (-8.096,6.075) (-8.098,6.025) (-8.100,5.975) (-8.102,5.925) (-8.103,5.875) (-8.104,5.825) (-8.105,5.775) (-8.105,5.725) (-8.104,5.675) (-8.103,5.625) (-8.102,5.575) (-8.100,5.525) (-8.098,5.475) (-8.096,5.425) (-8.093,5.375) (-8.090,5.325) (-8.086,5.275) (-8.082,5.225) (-8.077,5.175) (-8.072,5.125) (-8.067,5.075) (-8.061,5.025) (-8.055,4.975) (-8.048,4.925) (-8.041,4.875) (-8.033,4.825) (-8.025,4.775) (-8.016,4.725) (-8.007,4.675) (-7.998,4.625) (-7.988,4.575) (-7.977,4.525) (-7.967,4.475) (-7.955,4.425) (-7.943,4.375) (-7.931,4.325) (-7.918,4.275) (-7.905,4.225) (-7.891,4.175) (-7.877,4.125) (-7.862,4.075) (-7.847,4.025) (-7.831,3.975) (-7.815,3.925) (-7.798,3.875) (-7.780,3.825) (-7.762,3.775) (-7.744,3.725) (-7.725,3.676) (-7.705,3.625) (-7.685,3.575) (-7.664,3.525) (-7.643,3.475) (-7.621,3.425) (-7.600,3.379) (-7.575,3.325) (-7.551,3.275) (-7.526,3.225) (-7.501,3.175) (-7.475,3.125) (-7.450,3.077) (-7.425,3.031) (-7.400,2.986) (-7.375,2.942) (-7.336,2.875) (-7.306,2.825) (-7.275,2.775) (-7.243,2.725) (-7.211,2.675) (-7.178,2.625) (-7.144,2.575) (-7.109,2.525) (-7.075,2.478) (-7.036,2.425) (-7.000,2.377) (-6.960,2.325) (-6.925,2.281) (-6.880,2.225) (-6.838,2.175) (-6.800,2.130) (-6.752,2.075) (-6.707,2.025) (-6.661,1.975) (-6.625,1.937) (-6.590,1.900) (-6.550,1.860) (-6.500,1.810) (-6.464,1.775) (-6.425,1.738) (-6.384,1.700) (-6.329,1.650) (-6.275,1.602) (-6.225,1.559) (-6.185,1.525) (-6.125,1.476) (-6.075,1.435) (-6.030,1.400) (-5.975,1.358) (-5.931,1.325) (-5.875,1.285) (-5.826,1.250) (-5.775,1.214) (-5.725,1.180) (-5.679,1.150) (-5.625,1.115) (-5.575,1.083) (-5.525,1.052) (-5.480,1.025) (-5.425,0.993) (-5.375,0.964) (-5.325,0.936) (-5.275,0.908) (-5.225,0.881) (-5.175,0.855) (-5.125,0.830) (-5.075,0.805) (-5.025,0.780) (-4.975,0.756) (-4.925,0.733) (-4.875,0.710) (-4.825,0.688) (-4.775,0.666) (-4.725,0.645) (-4.678,0.625) (-4.625,0.604) (-4.575,0.584) (-4.525,0.565) (-4.475,0.546) (-4.425,0.527) (-4.375,0.509) (-4.325,0.492) (-4.276,0.475) (-4.225,0.458) (-4.175,0.442) (-4.125,0.426) (-4.075,0.411) (-4.025,0.396) (-3.975,0.381) (-3.925,0.367) (-3.875,0.353) (-3.825,0.340) (-3.775,0.326) (-3.725,0.314) (-3.675,0.301) (-3.625,0.289) (-3.575,0.278) (-3.525,0.266) (-3.475,0.255) (-3.425,0.244) (-3.375,0.234) (-3.325,0.224) (-3.275,0.214) (-3.225,0.204) (-3.175,0.195) (-3.125,0.186) (-3.075,0.178) (-3.025,0.169) (-2.975,0.161) (-2.925,0.153) (-2.875,0.146) (-2.825,0.139) (-2.775,0.132) (-2.725,0.125) (-2.675,0.118) (-2.625,0.112) (-2.575,0.106) (-2.525,0.100) (-2.475,0.094) (-2.425,0.089) (-2.375,0.084) (-2.325,0.079) (-2.275,0.074) (-2.225,0.069) (-2.175,0.065) (-2.125,0.061) (-2.075,0.057) (-2.025,0.053) (-1.975,0.050) (-1.925,0.046) (-1.875,0.043) (-1.825,0.040) (-1.775,0.037) (-1.725,0.034) (-1.675,0.031) (-1.625,0.029) (-1.575,0.026) (-1.525,0.024) (-1.475,0.022) (-1.425,0.020) (-1.375,0.018) (-1.325,0.017) (-1.275,0.015) (-1.225,0.014) (-1.175,0.012) (-1.125,0.011) (-1.075,0.010) (-1.025,0.009) (-0.975,0.008) (-0.925,0.007) (-0.875,0.006) (-0.825,0.005) (-0.775,0.005) (-0.725,0.004) (-0.675,0.003) (-0.625,0.003) (-0.575,0.003) (-0.525,0.002) (-0.475,0.002) (-0.425,0.002) (-0.375,0.001) (-0.325,0.001) (-0.275,0.001) (-0.250,0.001)}
\EQ{(0.250,0.002) (0.300,0.002) (0.350,0.003) (0.400,0.003) (0.450,0.004) (0.500,0.004) (0.550,0.005) (0.600,0.006) (0.650,0.006) (0.700,0.007) (0.750,0.009) (0.800,0.010) (0.850,0.011) (0.900,0.013) (0.950,0.014) (1.000,0.016) (1.050,0.018) (1.100,0.021) (1.150,0.023) (1.200,0.026) (1.250,0.029) (1.300,0.032) (1.350,0.035) (1.400,0.038) (1.450,0.042) (1.500,0.046) (1.550,0.051) (1.600,0.055) (1.650,0.060) (1.700,0.065) (1.750,0.071) (1.800,0.076) (1.850,0.083) (1.900,0.089) (1.950,0.096) (2.000,0.103) (2.050,0.110) (2.100,0.118) (2.150,0.126) (2.200,0.135) (2.250,0.144) (2.300,0.153) (2.350,0.163) (2.400,0.174) (2.450,0.184) (2.500,0.196) (2.550,0.207) (2.600,0.219) (2.650,0.232) (2.700,0.245) (2.750,0.259) (2.800,0.273) (2.850,0.287) (2.900,0.303) (2.950,0.319) (3.000,0.335) (3.050,0.352) (3.100,0.369) (3.150,0.388) (3.200,0.407) (3.247,0.425) (3.300,0.446) (3.350,0.467) (3.400,0.489) (3.450,0.511) (3.500,0.534) (3.550,0.558) (3.600,0.583) (3.650,0.608) (3.700,0.635) (3.750,0.662) (3.800,0.690) (3.850,0.719) (3.900,0.749) (3.950,0.780) (4.000,0.812) (4.050,0.846) (4.093,0.875) (4.150,0.915) (4.197,0.950) (4.250,0.990) (4.295,1.025) (4.350,1.070) (4.400,1.112) (4.444,1.150) (4.500,1.200) (4.550,1.247) (4.600,1.295) (4.650,1.346) (4.700,1.398) (4.748,1.450) (4.792,1.500) (4.825,1.539) (4.875,1.599) (4.916,1.650) (4.950,1.695) (4.991,1.750) (5.025,1.797) (5.061,1.850) (5.095,1.900) (5.125,1.947) (5.158,2.000) (5.188,2.050) (5.217,2.100) (5.245,2.150) (5.272,2.200) (5.298,2.250) (5.323,2.300) (5.347,2.350) (5.370,2.400) (5.393,2.450) (5.414,2.500) (5.435,2.550) (5.455,2.600) (5.474,2.650) (5.492,2.700) (5.510,2.750) (5.526,2.800) (5.542,2.850) (5.558,2.900) (5.572,2.950) (5.586,3.000) (5.599,3.050) (5.612,3.100) (5.623,3.150) (5.634,3.200) (5.645,3.250) (5.655,3.300) (5.664,3.350) (5.672,3.400) (5.680,3.450) (5.687,3.500) (5.694,3.550) (5.700,3.600) (5.705,3.650) (5.710,3.700) (5.714,3.750) (5.717,3.800) (5.720,3.850) (5.722,3.900) (5.724,3.950) (5.725,4.000) (5.726,4.050) (5.726,4.100) (5.725,4.150) (5.724,4.200) (5.722,4.250) (5.720,4.300) (5.717,4.350) (5.714,4.400) (5.710,4.450) (5.705,4.500) (5.700,4.550) (5.694,4.600) (5.688,4.650) (5.681,4.700) (5.673,4.750) (5.665,4.800) (5.657,4.850) (5.648,4.900) (5.638,4.950) (5.628,5.000) (5.617,5.050) (5.605,5.100) (5.593,5.150) (5.581,5.200) (5.567,5.250) (5.553,5.300) (5.539,5.350) (5.524,5.400) (5.509,5.450) (5.492,5.500) (5.476,5.550) (5.458,5.600) (5.440,5.650) (5.422,5.700) (5.402,5.750) (5.382,5.800) (5.362,5.850) (5.341,5.900) (5.319,5.950) (5.296,6.000) (5.275,6.046) (5.250,6.099) (5.225,6.150) (5.200,6.200) (5.175,6.248) (5.150,6.295) (5.125,6.341) (5.100,6.386) (5.075,6.430) (5.049,6.475) (5.019,6.525) (4.988,6.575) (4.957,6.625) (4.925,6.674) (4.892,6.725) (4.858,6.775) (4.825,6.822) (4.787,6.875) (4.751,6.925) (4.714,6.975) (4.675,7.025) (4.636,7.075) (4.600,7.120) (4.555,7.175) (4.513,7.225) (4.475,7.269) (4.425,7.325) (4.380,7.375) (4.333,7.425) (4.286,7.475) (4.250,7.512) (4.212,7.550) (4.175,7.586) (4.135,7.625) (4.082,7.675) (4.028,7.725) (3.975,7.773) (3.925,7.817) (3.886,7.850) (3.827,7.900) (3.775,7.942) (3.734,7.975) (3.675,8.021) (3.625,8.059) (3.575,8.097) (3.525,8.133) (3.475,8.169) (3.430,8.200) (3.375,8.238) (3.325,8.271) (3.281,8.300) (3.225,8.335) (3.175,8.367) (3.125,8.397) (3.078,8.425) (3.025,8.456) (2.975,8.484) (2.925,8.512) (2.875,8.540) (2.825,8.566) (2.775,8.592) (2.725,8.617) (2.675,8.642) (2.625,8.666) (2.575,8.690) (2.525,8.713) (2.475,8.736) (2.425,8.758) (2.375,8.779) (2.325,8.800) (2.275,8.820) (2.225,8.840) (2.175,8.859) (2.125,8.878) (2.075,8.897) (2.025,8.914) (1.975,8.932) (1.925,8.948) (1.875,8.965) (1.825,8.981) (1.775,8.996) (1.725,9.011) (1.677,9.025) (1.625,9.039) (1.575,9.053) (1.525,9.066) (1.475,9.079) (1.425,9.091) (1.375,9.103) (1.325,9.114) (1.275,9.125) (1.225,9.135) (1.175,9.145) (1.125,9.155) (1.075,9.164) (1.025,9.172) (0.975,9.181) (0.925,9.189) (0.875,9.196) (0.825,9.203) (0.775,9.209) (0.725,9.216) (0.675,9.221) (0.625,9.227) (0.575,9.232) (0.525,9.236) (0.475,9.240) (0.425,9.244) (0.375,9.247) (0.325,9.250) (0.275,9.252) (0.225,9.254) (0.175,9.256) (0.125,9.257) (0.075,9.258) (0.025,9.258) (-0.050,9.258) (-0.100,9.258) (-0.150,9.257) (-0.200,9.255) (-0.250,9.253) (-0.300,9.251) (-0.350,9.249) (-0.400,9.246) (-0.450,9.242) (-0.500,9.238) (-0.550,9.234) (-0.600,9.229) (-0.650,9.224) (-0.700,9.219) (-0.750,9.213) (-0.800,9.206) (-0.850,9.199) (-0.900,9.192) (-0.950,9.185) (-1.000,9.177) (-1.050,9.168) (-1.100,9.159) (-1.149,9.150) (-1.200,9.140) (-1.250,9.130) (-1.300,9.119) (-1.350,9.108) (-1.400,9.097) (-1.450,9.085) (-1.500,9.072) (-1.550,9.060) (-1.600,9.046) (-1.650,9.033) (-1.700,9.018) (-1.750,9.004) (-1.800,8.988) (-1.850,8.973) (-1.900,8.957) (-1.950,8.940) (-2.000,8.923) (-2.050,8.906) (-2.100,8.887) (-2.150,8.869) (-2.199,8.850) (-2.250,8.830) (-2.300,8.810) (-2.350,8.790) (-2.400,8.768) (-2.450,8.747) (-2.499,8.725) (-2.550,8.702) (-2.600,8.678) (-2.650,8.654) (-2.700,8.630) (-2.750,8.605) (-2.800,8.579) (-2.850,8.553) (-2.900,8.526) (-2.947,8.500) (-2.992,8.475) (-3.035,8.450) (-3.100,8.412) (-3.150,8.382) (-3.200,8.351) (-3.250,8.320) (-3.300,8.287) (-3.350,8.254) (-3.394,8.225) (-3.450,8.186) (-3.500,8.151) (-3.550,8.115) (-3.600,8.078) (-3.650,8.040) (-3.700,8.002) (-3.750,7.962) (-3.796,7.925) (-3.850,7.881) (-3.900,7.838) (-3.944,7.800) (-4.000,7.750) (-4.050,7.705) (-4.100,7.658) (-4.150,7.611) (-4.187,7.575) (-4.225,7.537) (-4.261,7.500) (-4.300,7.460) (-4.350,7.407) (-4.400,7.353) (-4.448,7.300) (-4.491,7.250) (-4.525,7.211) (-4.575,7.151) (-4.616,7.100) (-4.650,7.058) (-4.695,7.000) (-4.725,6.960) (-4.769,6.900) (-4.800,6.857) (-4.840,6.800) (-4.875,6.750) (-4.908,6.700) (-4.941,6.650) (-4.973,6.600) (-5.000,6.556) (-5.034,6.500) (-5.063,6.450) (-5.092,6.400) (-5.120,6.350) (-5.147,6.300) (-5.174,6.250) (-5.200,6.200) (-5.225,6.150) (-5.250,6.100) (-5.273,6.050) (-5.296,6.000) (-5.319,5.950) (-5.341,5.900) (-5.362,5.850) (-5.382,5.800) (-5.402,5.750) (-5.422,5.700) (-5.440,5.650) (-5.458,5.600) (-5.475,5.552) (-5.492,5.500) (-5.509,5.450) (-5.524,5.400) (-5.539,5.350) (-5.553,5.300) (-5.567,5.250) (-5.581,5.200) (-5.593,5.150) (-5.605,5.100) (-5.617,5.050) (-5.628,5.000) (-5.638,4.950) (-5.648,4.900) (-5.657,4.850) (-5.665,4.800) (-5.673,4.750) (-5.681,4.700) (-5.688,4.650) (-5.694,4.600) (-5.700,4.550) (-5.705,4.500) (-5.710,4.450) (-5.714,4.400) (-5.717,4.350) (-5.720,4.300) (-5.722,4.250) (-5.724,4.200) (-5.725,4.150) (-5.726,4.100) (-5.726,4.050) (-5.725,4.000) (-5.724,3.950) (-5.722,3.900) (-5.720,3.850) (-5.717,3.800) (-5.714,3.750) (-5.710,3.700) (-5.705,3.650) (-5.700,3.600) (-5.694,3.550) (-5.687,3.500) (-5.680,3.450) (-5.672,3.400) (-5.664,3.350) (-5.655,3.300) (-5.645,3.250) (-5.634,3.200) (-5.623,3.150) (-5.612,3.100) (-5.599,3.050) (-5.586,3.000) (-5.572,2.950) (-5.558,2.900) (-5.542,2.850) (-5.526,2.800) (-5.510,2.750) (-5.492,2.700) (-5.474,2.650) (-5.455,2.600) (-5.435,2.550) (-5.414,2.500) (-5.393,2.450) (-5.370,2.400) (-5.347,2.350) (-5.325,2.304) (-5.300,2.254) (-5.275,2.206) (-5.250,2.159) (-5.225,2.114) (-5.200,2.071) (-5.173,2.025) (-5.143,1.975) (-5.111,1.925) (-5.078,1.875) (-5.050,1.833) (-5.009,1.775) (-4.975,1.728) (-4.935,1.675) (-4.900,1.630) (-4.855,1.575) (-4.814,1.525) (-4.775,1.481) (-4.725,1.425) (-4.678,1.375) (-4.630,1.325) (-4.579,1.275) (-4.527,1.225) (-4.475,1.177) (-4.425,1.133) (-4.386,1.100) (-4.326,1.050) (-4.275,1.009) (-4.231,0.975) (-4.175,0.933) (-4.129,0.900) (-4.075,0.863) (-4.025,0.829) (-3.981,0.800) (-3.925,0.765) (-3.875,0.734) (-3.825,0.705) (-3.775,0.676) (-3.728,0.650) (-3.682,0.625) (-3.634,0.600) (-3.584,0.575) (-3.533,0.550) (-3.480,0.525) (-3.425,0.500) (-3.375,0.478) (-3.325,0.457) (-3.275,0.436) (-3.225,0.416) (-3.175,0.397) (-3.125,0.379) (-3.075,0.361) (-3.025,0.343) (-2.975,0.327) (-2.925,0.311) (-2.875,0.295) (-2.825,0.280) (-2.775,0.266) (-2.725,0.252) (-2.675,0.238) (-2.625,0.226) (-2.575,0.213) (-2.525,0.201) (-2.475,0.190) (-2.425,0.179) (-2.375,0.168) (-2.325,0.158) (-2.275,0.149) (-2.225,0.139) (-2.175,0.131) (-2.125,0.122) (-2.075,0.114) (-2.025,0.107) (-1.975,0.099) (-1.925,0.092) (-1.875,0.086) (-1.825,0.079) (-1.775,0.074) (-1.725,0.068) (-1.675,0.063) (-1.625,0.058) (-1.575,0.053) (-1.525,0.048) (-1.475,0.044) (-1.425,0.040) (-1.375,0.037) (-1.325,0.033) (-1.275,0.030) (-1.225,0.027) (-1.175,0.024) (-1.125,0.022) (-1.075,0.020) (-1.025,0.017) (-0.975,0.015) (-0.925,0.014) (-0.875,0.012) (-0.825,0.011) (-0.775,0.009) (-0.725,0.008) (-0.675,0.007) (-0.625,0.006) (-0.575,0.005) (-0.525,0.004) (-0.475,0.004) (-0.425,0.003) (-0.375,0.003) (-0.325,0.002) (-0.275,0.002) (-0.250,0.002)}
\EQ{(0.250,0.003) (0.300,0.003) (0.350,0.004) (0.400,0.005) (0.450,0.005) (0.500,0.006) (0.550,0.007) (0.600,0.008) (0.650,0.010) (0.700,0.011) (0.750,0.013) (0.800,0.015) (0.850,0.017) (0.900,0.019) (0.950,0.022) (1.000,0.025) (1.050,0.028) (1.100,0.031) (1.150,0.035) (1.200,0.039) (1.250,0.043) (1.300,0.047) (1.350,0.052) (1.400,0.058) (1.450,0.063) (1.500,0.070) (1.550,0.076) (1.600,0.083) (1.650,0.090) (1.700,0.098) (1.750,0.106) (1.800,0.115) (1.850,0.124) (1.900,0.134) (1.950,0.144) (2.000,0.155) (2.050,0.166) (2.100,0.178) (2.150,0.191) (2.200,0.204) (2.250,0.218) (2.300,0.232) (2.350,0.247) (2.400,0.263) (2.450,0.279) (2.500,0.296) (2.550,0.314) (2.600,0.333) (2.650,0.353) (2.700,0.373) (2.750,0.394) (2.800,0.416) (2.850,0.439) (2.900,0.463) (2.950,0.488) (3.000,0.514) (3.050,0.541) (3.100,0.569) (3.150,0.599) (3.194,0.625) (3.250,0.661) (3.300,0.694) (3.346,0.725) (3.400,0.764) (3.448,0.800) (3.500,0.840) (3.543,0.875) (3.600,0.923) (3.650,0.967) (3.700,1.013) (3.738,1.050) (3.775,1.087) (3.812,1.125) (3.850,1.165) (3.900,1.221) (3.946,1.275) (3.987,1.325) (4.025,1.373) (4.064,1.425) (4.100,1.475) (4.134,1.525) (4.167,1.575) (4.199,1.625) (4.225,1.668) (4.258,1.725) (4.286,1.775) (4.312,1.825) (4.338,1.875) (4.362,1.925) (4.385,1.975) (4.407,2.025) (4.428,2.075) (4.448,2.125) (4.467,2.175) (4.485,2.225) (4.502,2.275) (4.519,2.325) (4.534,2.375) (4.548,2.425) (4.562,2.475) (4.575,2.525) (4.587,2.575) (4.598,2.625) (4.608,2.675) (4.617,2.725) (4.626,2.775) (4.634,2.825) (4.641,2.875) (4.647,2.925) (4.653,2.975) (4.658,3.025) (4.662,3.075) (4.665,3.125) (4.668,3.175) (4.670,3.225) (4.671,3.275) (4.671,3.325) (4.671,3.375) (4.670,3.425) (4.668,3.475) (4.666,3.525) (4.663,3.575) (4.659,3.625) (4.655,3.675) (4.650,3.725) (4.644,3.775) (4.637,3.825) (4.630,3.875) (4.622,3.925) (4.613,3.975) (4.604,4.025) (4.594,4.075) (4.583,4.125) (4.572,4.175) (4.560,4.225) (4.547,4.275) (4.533,4.325) (4.519,4.375) (4.504,4.425) (4.488,4.475) (4.472,4.525) (4.454,4.575) (4.436,4.625) (4.418,4.675) (4.398,4.725) (4.378,4.775) (4.357,4.825) (4.335,4.875) (4.312,4.925) (4.288,4.975) (4.264,5.025) (4.239,5.075) (4.213,5.125) (4.186,5.175) (4.158,5.225) (4.130,5.275) (4.100,5.325) (4.069,5.375) (4.038,5.425) (4.005,5.475) (3.975,5.520) (3.937,5.575) (3.902,5.625) (3.865,5.675) (3.827,5.725) (3.789,5.775) (3.750,5.823) (3.707,5.875) (3.665,5.925) (3.625,5.970) (3.576,6.025) (3.529,6.075) (3.481,6.125) (3.432,6.175) (3.381,6.225) (3.328,6.275) (3.275,6.324) (3.225,6.368) (3.175,6.412) (3.129,6.450) (3.075,6.495) (3.025,6.534) (2.975,6.573) (2.925,6.611) (2.875,6.647) (2.825,6.683) (2.775,6.717) (2.726,6.750) (2.675,6.784) (2.625,6.816) (2.575,6.846) (2.528,6.875) (2.475,6.906) (2.425,6.935) (2.375,6.962) (2.325,6.989) (2.275,7.015) (2.225,7.041) (2.175,7.066) (2.125,7.090) (2.075,7.113) (2.025,7.136) (1.975,7.158) (1.925,7.179) (1.875,7.200) (1.825,7.220) (1.775,7.239) (1.725,7.258) (1.678,7.275) (1.625,7.294) (1.575,7.311) (1.525,7.327) (1.475,7.343) (1.425,7.358) (1.375,7.372) (1.325,7.386) (1.275,7.400) (1.225,7.413) (1.176,7.425) (1.125,7.437) (1.075,7.448) (1.025,7.459) (0.975,7.469) (0.925,7.479) (0.875,7.488) (0.825,7.497) (0.775,7.505) (0.725,7.512) (0.675,7.519) (0.625,7.526) (0.575,7.532) (0.525,7.538) (0.475,7.543) (0.425,7.547) (0.375,7.551) (0.325,7.555) (0.275,7.558) (0.225,7.560) (0.175,7.562) (0.125,7.564) (0.075,7.565) (0.025,7.565) (-0.050,7.565) (-0.100,7.564) (-0.150,7.563) (-0.200,7.561) (-0.250,7.559) (-0.300,7.556) (-0.350,7.553) (-0.400,7.549) (-0.450,7.545) (-0.500,7.540) (-0.550,7.535) (-0.600,7.529) (-0.650,7.523) (-0.700,7.516) (-0.750,7.509) (-0.800,7.501) (-0.850,7.492) (-0.900,7.484) (-0.950,7.474) (-1.000,7.464) (-1.050,7.454) (-1.100,7.443) (-1.150,7.431) (-1.200,7.419) (-1.250,7.406) (-1.300,7.393) (-1.350,7.380) (-1.400,7.365) (-1.450,7.350) (-1.500,7.335) (-1.550,7.319) (-1.600,7.302) (-1.650,7.285) (-1.700,7.267) (-1.750,7.249) (-1.800,7.229) (-1.850,7.210) (-1.900,7.189) (-1.950,7.168) (-2.000,7.147) (-2.049,7.125) (-2.100,7.101) (-2.150,7.078) (-2.200,7.053) (-2.250,7.028) (-2.300,7.002) (-2.350,6.976) (-2.397,6.950) (-2.442,6.925) (-2.485,6.900) (-2.550,6.862) (-2.600,6.831) (-2.649,6.800) (-2.700,6.767) (-2.750,6.734) (-2.800,6.700) (-2.850,6.665) (-2.900,6.629) (-2.950,6.592) (-3.000,6.554) (-3.050,6.515) (-3.099,6.475) (-3.150,6.433) (-3.188,6.400) (-3.246,6.350) (-3.300,6.301) (-3.350,6.254) (-3.400,6.206) (-3.450,6.157) (-3.500,6.106) (-3.550,6.053) (-3.599,6.000) (-3.643,5.950) (-3.686,5.900) (-3.725,5.854) (-3.769,5.800) (-3.800,5.761) (-3.847,5.700) (-3.884,5.650) (-3.920,5.600) (-3.950,5.557) (-3.989,5.500) (-4.022,5.450) (-4.050,5.406) (-4.085,5.350) (-4.115,5.300) (-4.144,5.250) (-4.172,5.200) (-4.200,5.150) (-4.225,5.102) (-4.250,5.053) (-4.275,5.003) (-4.300,4.951) (-4.323,4.900) (-4.346,4.850) (-4.367,4.800) (-4.388,4.750) (-4.408,4.700) (-4.427,4.650) (-4.445,4.600) (-4.463,4.550) (-4.480,4.500) (-4.496,4.450) (-4.511,4.400) (-4.526,4.350) (-4.540,4.300) (-4.553,4.250) (-4.566,4.200) (-4.578,4.150) (-4.589,4.100) (-4.599,4.050) (-4.609,4.000) (-4.618,3.950) (-4.626,3.900) (-4.634,3.850) (-4.641,3.800) (-4.647,3.750) (-4.652,3.700) (-4.657,3.650) (-4.661,3.600) (-4.665,3.550) (-4.667,3.500) (-4.669,3.450) (-4.671,3.400) (-4.671,3.350) (-4.671,3.300) (-4.670,3.250) (-4.669,3.200) (-4.667,3.150) (-4.664,3.100) (-4.660,3.050) (-4.655,3.000) (-4.650,2.950) (-4.644,2.900) (-4.638,2.850) (-4.630,2.800) (-4.622,2.750) (-4.613,2.700) (-4.603,2.650) (-4.592,2.600) (-4.581,2.550) (-4.569,2.500) (-4.555,2.450) (-4.541,2.400) (-4.526,2.350) (-4.511,2.300) (-4.494,2.250) (-4.476,2.200) (-4.458,2.150) (-4.438,2.100) (-4.418,2.050) (-4.396,2.000) (-4.375,1.953) (-4.350,1.900) (-4.325,1.850) (-4.300,1.802) (-4.275,1.755) (-4.250,1.711) (-4.214,1.650) (-4.183,1.600) (-4.151,1.550) (-4.117,1.500) (-4.082,1.450) (-4.050,1.406) (-4.007,1.350) (-3.975,1.310) (-3.925,1.250) (-3.882,1.200) (-3.836,1.150) (-3.800,1.112) (-3.764,1.075) (-3.725,1.037) (-3.686,1.000) (-3.631,0.950) (-3.575,0.902) (-3.525,0.860) (-3.481,0.825) (-3.425,0.782) (-3.381,0.750) (-3.325,0.711) (-3.275,0.677) (-3.225,0.645) (-3.175,0.614) (-3.125,0.584) (-3.075,0.555) (-3.025,0.527) (-2.975,0.501) (-2.925,0.476) (-2.875,0.451) (-2.825,0.428) (-2.775,0.405) (-2.725,0.383) (-2.675,0.363) (-2.625,0.343) (-2.575,0.324) (-2.525,0.305) (-2.475,0.288) (-2.425,0.271) (-2.375,0.255) (-2.325,0.239) (-2.276,0.225) (-2.225,0.211) (-2.175,0.197) (-2.125,0.184) (-2.075,0.172) (-2.025,0.161) (-1.975,0.149) (-1.925,0.139) (-1.875,0.129) (-1.825,0.120) (-1.775,0.111) (-1.725,0.102) (-1.675,0.094) (-1.625,0.087) (-1.575,0.079) (-1.525,0.073) (-1.475,0.066) (-1.425,0.061) (-1.375,0.055) (-1.325,0.050) (-1.275,0.045) (-1.225,0.041) (-1.175,0.037) (-1.125,0.033) (-1.075,0.029) (-1.025,0.026) (-0.975,0.023) (-0.925,0.020) (-0.875,0.018) (-0.825,0.016) (-0.775,0.014) (-0.725,0.012) (-0.675,0.010) (-0.625,0.009) (-0.575,0.008) (-0.525,0.007) (-0.475,0.006) (-0.425,0.005) (-0.375,0.004) (-0.325,0.004) (-0.275,0.003) (-0.250,0.003)}
\EQ{(0.250,0.004) (0.300,0.005) (0.350,0.005) (0.400,0.006) (0.450,0.007) (0.500,0.008) (0.550,0.010) (0.600,0.011) (0.650,0.013) (0.700,0.015) (0.750,0.017) (0.800,0.020) (0.850,0.022) (0.900,0.026) (0.950,0.029) (1.000,0.033) (1.050,0.037) (1.100,0.041) (1.150,0.046) (1.200,0.052) (1.250,0.057) (1.300,0.063) (1.350,0.070) (1.400,0.077) (1.450,0.085) (1.500,0.093) (1.550,0.102) (1.600,0.111) (1.650,0.121) (1.700,0.131) (1.750,0.142) (1.800,0.154) (1.850,0.166) (1.900,0.179) (1.950,0.193) (2.000,0.208) (2.050,0.223) (2.100,0.239) (2.150,0.256) (2.200,0.274) (2.250,0.293) (2.300,0.313) (2.350,0.333) (2.400,0.355) (2.450,0.378) (2.497,0.400) (2.548,0.425) (2.596,0.450) (2.642,0.475) (2.686,0.500) (2.750,0.538) (2.800,0.570) (2.847,0.600) (2.900,0.637) (2.950,0.673) (3.000,0.711) (3.050,0.750) (3.100,0.792) (3.138,0.825) (3.193,0.875) (3.245,0.925) (3.294,0.975) (3.341,1.025) (3.375,1.063) (3.425,1.122) (3.467,1.175) (3.500,1.218) (3.542,1.275) (3.575,1.323) (3.609,1.375) (3.640,1.425) (3.670,1.475) (3.698,1.525) (3.725,1.574) (3.750,1.623) (3.775,1.675) (3.798,1.725) (3.820,1.775) (3.841,1.825) (3.860,1.875) (3.878,1.925) (3.896,1.975) (3.912,2.025) (3.927,2.075) (3.941,2.125) (3.954,2.175) (3.966,2.225) (3.977,2.275) (3.988,2.325) (3.997,2.375) (4.005,2.425) (4.013,2.475) (4.020,2.525) (4.025,2.575) (4.030,2.625) (4.034,2.675) (4.037,2.725) (4.040,2.775) (4.041,2.825) (4.042,2.875) (4.042,2.925) (4.041,2.975) (4.039,3.025) (4.037,3.075) (4.033,3.125) (4.029,3.175) (4.024,3.225) (4.018,3.275) (4.012,3.325) (4.004,3.375) (3.996,3.425) (3.987,3.475) (3.977,3.525) (3.967,3.575) (3.955,3.625) (3.943,3.675) (3.930,3.725) (3.916,3.775) (3.901,3.825) (3.886,3.875) (3.869,3.925) (3.852,3.975) (3.834,4.025) (3.815,4.075) (3.795,4.125) (3.775,4.172) (3.752,4.225) (3.729,4.275) (3.706,4.325) (3.681,4.375) (3.655,4.425) (3.629,4.475) (3.601,4.525) (3.575,4.571) (3.543,4.625) (3.512,4.675) (3.480,4.725) (3.450,4.771) (3.413,4.825) (3.378,4.875) (3.341,4.925) (3.303,4.975) (3.264,5.025) (3.225,5.073) (3.181,5.125) (3.138,5.175) (3.100,5.217) (3.050,5.271) (3.000,5.322) (2.950,5.372) (2.900,5.420) (2.850,5.467) (2.800,5.512) (2.756,5.550) (2.700,5.598) (2.650,5.638) (2.603,5.675) (2.550,5.716) (2.504,5.750) (2.450,5.789) (2.400,5.823) (2.350,5.857) (2.300,5.889) (2.250,5.921) (2.202,5.950) (2.150,5.981) (2.100,6.010) (2.050,6.037) (2.000,6.064) (1.950,6.090) (1.900,6.115) (1.850,6.139) (1.800,6.163) (1.750,6.186) (1.700,6.208) (1.650,6.229) (1.600,6.249) (1.550,6.269) (1.500,6.288) (1.450,6.306) (1.400,6.323) (1.350,6.340) (1.300,6.356) (1.250,6.372) (1.200,6.387) (1.150,6.401) (1.100,6.414) (1.050,6.427) (1.000,6.439) (0.950,6.451) (0.900,6.462) (0.850,6.472) (0.800,6.482) (0.750,6.491) (0.700,6.499) (0.650,6.507) (0.600,6.515) (0.550,6.521) (0.500,6.527) (0.450,6.533) (0.400,6.538) (0.350,6.542) (0.300,6.546) (0.250,6.549) (0.200,6.552) (0.150,6.554) (0.100,6.555) (0.050,6.556) (0.000,6.556) (-0.050,6.556) (-0.100,6.555) (-0.150,6.554) (-0.200,6.552) (-0.250,6.549) (-0.300,6.546) (-0.350,6.542) (-0.400,6.538) (-0.450,6.533) (-0.500,6.527) (-0.550,6.521) (-0.600,6.515) (-0.650,6.507) (-0.700,6.499) (-0.750,6.491) (-0.800,6.482) (-0.850,6.472) (-0.900,6.462) (-0.950,6.451) (-1.000,6.439) (-1.050,6.427) (-1.100,6.414) (-1.150,6.401) (-1.200,6.387) (-1.250,6.372) (-1.300,6.356) (-1.350,6.340) (-1.400,6.323) (-1.450,6.306) (-1.500,6.288) (-1.550,6.269) (-1.598,6.250) (-1.650,6.229) (-1.700,6.208) (-1.750,6.186) (-1.800,6.163) (-1.850,6.139) (-1.900,6.115) (-1.950,6.090) (-2.000,6.064) (-2.050,6.037) (-2.100,6.010) (-2.150,5.981) (-2.200,5.951) (-2.243,5.925) (-2.300,5.889) (-2.350,5.857) (-2.397,5.825) (-2.450,5.789) (-2.500,5.753) (-2.550,5.716) (-2.600,5.678) (-2.650,5.638) (-2.697,5.600) (-2.750,5.556) (-2.800,5.512) (-2.841,5.475) (-2.895,5.425) (-2.947,5.375) (-2.997,5.325) (-3.046,5.275) (-3.092,5.225) (-3.138,5.175) (-3.175,5.132) (-3.223,5.075) (-3.264,5.025) (-3.300,4.979) (-3.341,4.925) (-3.375,4.879) (-3.413,4.825) (-3.447,4.775) (-3.475,4.733) (-3.512,4.675) (-3.543,4.625) (-3.573,4.575) (-3.600,4.527) (-3.625,4.482) (-3.650,4.435) (-3.675,4.387) (-3.700,4.337) (-3.725,4.284) (-3.750,4.229) (-3.774,4.175) (-3.795,4.125) (-3.815,4.075) (-3.834,4.025) (-3.852,3.975) (-3.869,3.925) (-3.886,3.875) (-3.901,3.825) (-3.916,3.775) (-3.930,3.725) (-3.943,3.675) (-3.955,3.625) (-3.967,3.575) (-3.977,3.525) (-3.987,3.475) (-3.996,3.425) (-4.004,3.375) (-4.012,3.325) (-4.018,3.275) (-4.024,3.225) (-4.029,3.175) (-4.033,3.125) (-4.037,3.075) (-4.039,3.025) (-4.041,2.975) (-4.042,2.925) (-4.042,2.875) (-4.041,2.825) (-4.040,2.775) (-4.037,2.725) (-4.034,2.675) (-4.030,2.625) (-4.025,2.575) (-4.020,2.525) (-4.013,2.475) (-4.005,2.425) (-3.997,2.375) (-3.988,2.325) (-3.977,2.275) (-3.966,2.225) (-3.954,2.175) (-3.941,2.125) (-3.927,2.075) (-3.912,2.025) (-3.896,1.975) (-3.878,1.925) (-3.860,1.875) (-3.841,1.825) (-3.820,1.775) (-3.798,1.725) (-3.775,1.675) (-3.751,1.625) (-3.725,1.575) (-3.700,1.528) (-3.675,1.483) (-3.640,1.425) (-3.609,1.375) (-3.576,1.325) (-3.542,1.275) (-3.505,1.225) (-3.475,1.185) (-3.427,1.125) (-3.385,1.075) (-3.350,1.035) (-3.300,0.981) (-3.250,0.930) (-3.200,0.882) (-3.150,0.836) (-3.110,0.800) (-3.050,0.750) (-3.000,0.711) (-2.953,0.675) (-2.900,0.637) (-2.850,0.602) (-2.800,0.570) (-2.750,0.538) (-2.700,0.508) (-2.650,0.480) (-2.600,0.452) (-2.550,0.426) (-2.500,0.401) (-2.450,0.378) (-2.400,0.355) (-2.350,0.333) (-2.300,0.313) (-2.250,0.293) (-2.202,0.275) (-2.150,0.256) (-2.100,0.239) (-2.050,0.223) (-2.000,0.208) (-1.950,0.193) (-1.900,0.179) (-1.850,0.166) (-1.800,0.154) (-1.750,0.142) (-1.700,0.131) (-1.650,0.121) (-1.600,0.111) (-1.550,0.102) (-1.500,0.093) (-1.450,0.085) (-1.400,0.077) (-1.350,0.070) (-1.300,0.063) (-1.250,0.057) (-1.200,0.052) (-1.150,0.046) (-1.100,0.041) (-1.050,0.037) (-1.000,0.033) (-0.950,0.029) (-0.900,0.026) (-0.850,0.022) (-0.800,0.020) (-0.750,0.017) (-0.700,0.015) (-0.650,0.013) (-0.600,0.011) (-0.550,0.010) (-0.500,0.008) (-0.450,0.007) (-0.400,0.006) (-0.350,0.005) (-0.300,0.005) (-0.250,0.004)}
\EQ{(0.250,0.005) (0.300,0.006) (0.350,0.007) (0.400,0.008) (0.450,0.009) (0.500,0.010) (0.550,0.012) (0.600,0.014) (0.650,0.016) (0.700,0.019) (0.750,0.021) (0.800,0.025) (0.850,0.028) (0.900,0.032) (0.950,0.036) (1.000,0.041) (1.050,0.046) (1.100,0.052) (1.150,0.058) (1.200,0.064) (1.250,0.072) (1.300,0.079) (1.350,0.088) (1.400,0.097) (1.450,0.106) (1.500,0.116) (1.550,0.127) (1.600,0.139) (1.650,0.151) (1.700,0.165) (1.750,0.179) (1.800,0.193) (1.850,0.209) (1.898,0.225) (1.950,0.243) (2.000,0.262) (2.050,0.282) (2.100,0.302) (2.150,0.324) (2.200,0.347) (2.250,0.371) (2.300,0.397) (2.350,0.423) (2.397,0.450) (2.450,0.481) (2.500,0.513) (2.550,0.546) (2.593,0.575) (2.650,0.617) (2.693,0.650) (2.750,0.696) (2.800,0.739) (2.840,0.775) (2.892,0.825) (2.942,0.875) (2.988,0.925) (3.025,0.967) (3.074,1.025) (3.113,1.075) (3.150,1.125) (3.186,1.175) (3.219,1.225) (3.250,1.274) (3.280,1.325) (3.308,1.375) (3.335,1.425) (3.360,1.475) (3.384,1.525) (3.406,1.575) (3.427,1.625) (3.447,1.675) (3.465,1.725) (3.482,1.775) (3.498,1.825) (3.513,1.875) (3.527,1.925) (3.540,1.975) (3.551,2.025) (3.562,2.075) (3.571,2.125) (3.580,2.175) (3.587,2.225) (3.594,2.275) (3.599,2.325) (3.604,2.375) (3.607,2.425) (3.610,2.475) (3.611,2.525) (3.612,2.575) (3.612,2.625) (3.611,2.675) (3.609,2.725) (3.606,2.775) (3.602,2.825) (3.598,2.875) (3.592,2.925) (3.586,2.975) (3.578,3.025) (3.570,3.075) (3.561,3.125) (3.551,3.175) (3.540,3.225) (3.528,3.275) (3.515,3.325) (3.501,3.375) (3.487,3.425) (3.471,3.475) (3.455,3.525) (3.437,3.575) (3.419,3.625) (3.400,3.673) (3.379,3.725) (3.357,3.775) (3.334,3.825) (3.311,3.875) (3.286,3.925) (3.260,3.975) (3.234,4.025) (3.205,4.075) (3.176,4.125) (3.150,4.168) (3.114,4.225) (3.081,4.275) (3.050,4.321) (3.011,4.375) (2.975,4.424) (2.936,4.475) (2.900,4.520) (2.855,4.575) (2.811,4.625) (2.775,4.666) (2.725,4.720) (2.675,4.771) (2.625,4.821) (2.575,4.869) (2.525,4.915) (2.485,4.950) (2.426,5.000) (2.375,5.042) (2.334,5.075) (2.275,5.120) (2.234,5.150) (2.175,5.192) (2.127,5.225) (2.075,5.260) (2.025,5.291) (1.975,5.322) (1.928,5.350) (1.885,5.375) (1.840,5.400) (1.794,5.425) (1.746,5.450) (1.700,5.473) (1.650,5.497) (1.600,5.521) (1.550,5.543) (1.500,5.565) (1.450,5.586) (1.400,5.606) (1.350,5.625) (1.300,5.643) (1.250,5.661) (1.200,5.677) (1.150,5.693) (1.100,5.709) (1.050,5.723) (1.000,5.737) (0.950,5.750) (0.900,5.762) (0.850,5.774) (0.800,5.785) (0.750,5.795) (0.700,5.805) (0.650,5.813) (0.600,5.822) (0.550,5.829) (0.500,5.836) (0.450,5.842) (0.400,5.848) (0.350,5.853) (0.300,5.857) (0.250,5.860) (0.200,5.863) (0.150,5.866) (0.100,5.867) (0.050,5.868) (0.000,5.868) (-0.050,5.868) (-0.100,5.867) (-0.150,5.866) (-0.200,5.863) (-0.250,5.860) (-0.300,5.857) (-0.350,5.853) (-0.400,5.848) (-0.450,5.842) (-0.500,5.836) (-0.550,5.829) (-0.600,5.822) (-0.650,5.813) (-0.700,5.805) (-0.750,5.795) (-0.800,5.785) (-0.850,5.774) (-0.900,5.762) (-0.950,5.750) (-1.000,5.737) (-1.050,5.723) (-1.100,5.709) (-1.150,5.693) (-1.200,5.677) (-1.250,5.661) (-1.300,5.643) (-1.349,5.625) (-1.400,5.606) (-1.450,5.586) (-1.500,5.565) (-1.550,5.543) (-1.600,5.521) (-1.650,5.497) (-1.696,5.475) (-1.746,5.450) (-1.794,5.425) (-1.840,5.400) (-1.885,5.375) (-1.928,5.350) (-1.975,5.322) (-2.025,5.291) (-2.075,5.260) (-2.125,5.226) (-2.175,5.192) (-2.225,5.157) (-2.268,5.125) (-2.325,5.082) (-2.365,5.050) (-2.425,5.001) (-2.475,4.959) (-2.513,4.925) (-2.568,4.875) (-2.621,4.825) (-2.671,4.775) (-2.720,4.725) (-2.767,4.675) (-2.811,4.625) (-2.850,4.580) (-2.896,4.525) (-2.936,4.475) (-2.974,4.425) (-3.011,4.375) (-3.047,4.325) (-3.081,4.275) (-3.114,4.225) (-3.146,4.175) (-3.175,4.127) (-3.205,4.075) (-3.234,4.025) (-3.260,3.975) (-3.286,3.925) (-3.311,3.875) (-3.334,3.825) (-3.357,3.775) (-3.379,3.725) (-3.399,3.675) (-3.419,3.625) (-3.437,3.575) (-3.455,3.525) (-3.471,3.475) (-3.487,3.425) (-3.501,3.375) (-3.515,3.325) (-3.528,3.275) (-3.540,3.225) (-3.551,3.175) (-3.561,3.125) (-3.570,3.075) (-3.578,3.025) (-3.586,2.975) (-3.592,2.925) (-3.598,2.875) (-3.602,2.825) (-3.606,2.775) (-3.609,2.725) (-3.611,2.675) (-3.612,2.625) (-3.612,2.575) (-3.611,2.525) (-3.610,2.475) (-3.607,2.425) (-3.604,2.375) (-3.599,2.325) (-3.594,2.275) (-3.587,2.225) (-3.580,2.175) (-3.571,2.125) (-3.562,2.075) (-3.551,2.025) (-3.540,1.975) (-3.527,1.925) (-3.513,1.875) (-3.498,1.825) (-3.482,1.775) (-3.465,1.725) (-3.447,1.675) (-3.427,1.625) (-3.406,1.575) (-3.384,1.525) (-3.360,1.475) (-3.335,1.425) (-3.308,1.375) (-3.280,1.325) (-3.250,1.275) (-3.219,1.225) (-3.186,1.175) (-3.150,1.125) (-3.113,1.075) (-3.075,1.026) (-3.032,0.975) (-2.988,0.925) (-2.950,0.884) (-2.900,0.833) (-2.850,0.784) (-2.813,0.750) (-2.755,0.700) (-2.700,0.655) (-2.650,0.617) (-2.600,0.580) (-2.557,0.550) (-2.500,0.513) (-2.450,0.481) (-2.400,0.452) (-2.353,0.425) (-2.307,0.400) (-2.258,0.375) (-2.207,0.350) (-2.152,0.325) (-2.100,0.302) (-2.050,0.282) (-2.000,0.262) (-1.950,0.243) (-1.900,0.226) (-1.850,0.209) (-1.800,0.193) (-1.750,0.179) (-1.700,0.165) (-1.650,0.151) (-1.600,0.139) (-1.550,0.127) (-1.500,0.116) (-1.450,0.106) (-1.400,0.097) (-1.350,0.088) (-1.300,0.079) (-1.250,0.072) (-1.200,0.064) (-1.150,0.058) (-1.100,0.052) (-1.050,0.046) (-1.000,0.041) (-0.950,0.036) (-0.900,0.032) (-0.850,0.028) (-0.800,0.025) (-0.750,0.021) (-0.700,0.019) (-0.650,0.016) (-0.600,0.014) (-0.550,0.012) (-0.500,0.010) (-0.450,0.009) (-0.400,0.008) (-0.350,0.007) (-0.300,0.006) (-0.250,0.005)}
\EQ{(0.250,0.006) (0.300,0.007) (0.350,0.008) (0.400,0.009) (0.450,0.011) (0.500,0.012) (0.550,0.014) (0.600,0.017) (0.650,0.019) (0.700,0.022) (0.750,0.026) (0.800,0.030) (0.850,0.034) (0.900,0.038) (0.950,0.044) (1.000,0.049) (1.050,0.055) (1.100,0.062) (1.150,0.069) (1.200,0.077) (1.250,0.086) (1.300,0.095) (1.350,0.105) (1.400,0.116) (1.450,0.128) (1.500,0.140) (1.550,0.153) (1.600,0.167) (1.650,0.183) (1.700,0.199) (1.750,0.216) (1.800,0.234) (1.850,0.253) (1.900,0.273) (1.950,0.295) (2.000,0.318) (2.050,0.342) (2.100,0.367) (2.150,0.394) (2.200,0.423) (2.245,0.450) (2.300,0.485) (2.350,0.519) (2.394,0.550) (2.450,0.593) (2.491,0.625) (2.549,0.675) (2.600,0.721) (2.650,0.770) (2.700,0.821) (2.749,0.875) (2.791,0.925) (2.825,0.967) (2.869,1.025) (2.900,1.069) (2.937,1.125) (2.968,1.175) (2.998,1.225) (3.025,1.274) (3.050,1.322) (3.075,1.373) (3.099,1.425) (3.120,1.475) (3.140,1.525) (3.159,1.575) (3.176,1.625) (3.192,1.675) (3.207,1.725) (3.220,1.775) (3.233,1.825) (3.244,1.875) (3.254,1.925) (3.263,1.975) (3.270,2.025) (3.277,2.075) (3.282,2.125) (3.287,2.175) (3.290,2.225) (3.293,2.275) (3.294,2.325) (3.295,2.375) (3.294,2.425) (3.293,2.475) (3.290,2.525) (3.287,2.575) (3.282,2.625) (3.276,2.675) (3.270,2.725) (3.262,2.775) (3.254,2.825) (3.245,2.875) (3.234,2.925) (3.223,2.975) (3.210,3.025) (3.197,3.075) (3.182,3.125) (3.167,3.175) (3.150,3.225) (3.133,3.275) (3.114,3.325) (3.094,3.375) (3.073,3.425) (3.051,3.475) (3.028,3.525) (3.004,3.575) (2.978,3.625) (2.951,3.675) (2.925,3.722) (2.894,3.775) (2.864,3.825) (2.832,3.875) (2.800,3.923) (2.764,3.975) (2.727,4.025) (2.689,4.075) (2.650,4.125) (2.609,4.175) (2.575,4.214) (2.525,4.270) (2.475,4.324) (2.425,4.375) (2.375,4.424) (2.325,4.471) (2.275,4.516) (2.235,4.550) (2.175,4.600) (2.125,4.640) (2.079,4.675) (2.025,4.715) (1.975,4.750) (1.925,4.784) (1.875,4.817) (1.825,4.848) (1.781,4.875) (1.725,4.907) (1.675,4.935) (1.625,4.962) (1.575,4.988) (1.525,5.013) (1.475,5.037) (1.425,5.059) (1.375,5.081) (1.325,5.102) (1.275,5.122) (1.225,5.141) (1.175,5.159) (1.125,5.177) (1.075,5.193) (1.025,5.209) (0.975,5.223) (0.925,5.238) (0.875,5.251) (0.825,5.263) (0.775,5.275) (0.725,5.286) (0.675,5.296) (0.625,5.305) (0.575,5.314) (0.525,5.322) (0.475,5.329) (0.425,5.335) (0.375,5.341) (0.325,5.346) (0.275,5.350) (0.225,5.354) (0.175,5.357) (0.125,5.359) (0.075,5.360) (0.025,5.361) (-0.050,5.361) (-0.100,5.360) (-0.150,5.358) (-0.200,5.355) (-0.250,5.352) (-0.300,5.348) (-0.350,5.344) (-0.400,5.338) (-0.450,5.332) (-0.500,5.325) (-0.550,5.318) (-0.600,5.310) (-0.650,5.301) (-0.700,5.291) (-0.750,5.280) (-0.800,5.269) (-0.850,5.257) (-0.900,5.244) (-0.950,5.231) (-1.000,5.216) (-1.050,5.201) (-1.100,5.185) (-1.150,5.168) (-1.200,5.150) (-1.250,5.132) (-1.300,5.112) (-1.350,5.092) (-1.400,5.070) (-1.450,5.048) (-1.499,5.025) (-1.550,5.001) (-1.600,4.975) (-1.648,4.950) (-1.694,4.925) (-1.738,4.900) (-1.800,4.863) (-1.850,4.833) (-1.900,4.801) (-1.950,4.767) (-2.000,4.733) (-2.045,4.700) (-2.100,4.659) (-2.144,4.625) (-2.200,4.580) (-2.250,4.537) (-2.293,4.500) (-2.347,4.450) (-2.400,4.400) (-2.450,4.350) (-2.497,4.300) (-2.543,4.250) (-2.587,4.200) (-2.625,4.155) (-2.670,4.100) (-2.708,4.050) (-2.746,4.000) (-2.775,3.959) (-2.815,3.900) (-2.848,3.850) (-2.875,3.807) (-2.909,3.750) (-2.938,3.700) (-2.965,3.650) (-2.991,3.600) (-3.016,3.550) (-3.040,3.500) (-3.062,3.450) (-3.084,3.400) (-3.104,3.350) (-3.123,3.300) (-3.141,3.250) (-3.159,3.200) (-3.175,3.150) (-3.190,3.100) (-3.204,3.050) (-3.217,3.000) (-3.228,2.950) (-3.239,2.900) (-3.249,2.850) (-3.258,2.800) (-3.266,2.750) (-3.273,2.700) (-3.279,2.650) (-3.284,2.600) (-3.288,2.550) (-3.292,2.500) (-3.294,2.450) (-3.295,2.400) (-3.295,2.350) (-3.294,2.300) (-3.292,2.250) (-3.289,2.200) (-3.285,2.150) (-3.280,2.100) (-3.274,2.050) (-3.267,2.000) (-3.258,1.950) (-3.249,1.900) (-3.238,1.850) (-3.227,1.800) (-3.214,1.750) (-3.200,1.701) (-3.184,1.650) (-3.168,1.600) (-3.150,1.551) (-3.131,1.500) (-3.110,1.450) (-3.088,1.400) (-3.064,1.350) (-3.039,1.300) (-3.012,1.250) (-2.983,1.200) (-2.953,1.150) (-2.925,1.107) (-2.887,1.050) (-2.850,1.000) (-2.811,0.950) (-2.775,0.906) (-2.727,0.850) (-2.680,0.800) (-2.630,0.750) (-2.577,0.700) (-2.525,0.654) (-2.475,0.613) (-2.427,0.575) (-2.375,0.537) (-2.325,0.502) (-2.275,0.469) (-2.225,0.438) (-2.175,0.408) (-2.125,0.381) (-2.075,0.354) (-2.025,0.329) (-1.975,0.306) (-1.925,0.284) (-1.875,0.263) (-1.825,0.243) (-1.777,0.225) (-1.725,0.207) (-1.675,0.190) (-1.625,0.175) (-1.575,0.160) (-1.525,0.147) (-1.475,0.134) (-1.425,0.122) (-1.375,0.111) (-1.325,0.100) (-1.275,0.091) (-1.225,0.082) (-1.175,0.073) (-1.125,0.066) (-1.075,0.059) (-1.025,0.052) (-0.975,0.046) (-0.925,0.041) (-0.875,0.036) (-0.825,0.032) (-0.775,0.028) (-0.725,0.024) (-0.675,0.021) (-0.625,0.018) (-0.575,0.016) (-0.525,0.013) (-0.475,0.012) (-0.425,0.010) (-0.375,0.009) (-0.325,0.007) (-0.275,0.007) (-0.250,0.006)}
\EQ{(0.250,0.007) (0.300,0.008) (0.350,0.009) (0.400,0.011) (0.450,0.012) (0.500,0.014) (0.550,0.017) (0.600,0.020) (0.650,0.023) (0.700,0.026) (0.750,0.030) (0.800,0.034) (0.850,0.039) (0.900,0.045) (0.950,0.051) (1.000,0.057) (1.050,0.065) (1.100,0.073) (1.150,0.081) (1.200,0.090) (1.250,0.101) (1.300,0.112) (1.350,0.123) (1.400,0.136) (1.450,0.149) (1.500,0.164) (1.550,0.180) (1.600,0.196) (1.650,0.214) (1.700,0.233) (1.750,0.253) (1.800,0.275) (1.850,0.298) (1.900,0.322) (1.950,0.348) (2.000,0.375) (2.050,0.404) (2.100,0.435) (2.150,0.468) (2.196,0.500) (2.250,0.540) (2.294,0.575) (2.350,0.622) (2.400,0.667) (2.450,0.715) (2.500,0.767) (2.550,0.823) (2.594,0.875) (2.625,0.915) (2.669,0.975) (2.700,1.020) (2.735,1.075) (2.765,1.125) (2.793,1.175) (2.820,1.225) (2.844,1.275) (2.867,1.325) (2.888,1.375) (2.908,1.425) (2.926,1.475) (2.943,1.525) (2.959,1.575) (2.973,1.625) (2.986,1.675) (2.997,1.725) (3.008,1.775) (3.017,1.825) (3.024,1.875) (3.031,1.925) (3.037,1.975) (3.041,2.025) (3.044,2.075) (3.047,2.125) (3.048,2.175) (3.048,2.225) (3.047,2.275) (3.044,2.325) (3.041,2.375) (3.037,2.425) (3.032,2.475) (3.025,2.525) (3.018,2.575) (3.009,2.625) (3.000,2.675) (2.989,2.725) (2.977,2.775) (2.964,2.825) (2.950,2.875) (2.935,2.925) (2.919,2.975) (2.902,3.025) (2.883,3.075) (2.864,3.125) (2.843,3.175) (2.821,3.225) (2.798,3.275) (2.775,3.321) (2.750,3.370) (2.725,3.416) (2.700,3.460) (2.662,3.525) (2.630,3.575) (2.600,3.621) (2.563,3.675) (2.527,3.725) (2.489,3.775) (2.450,3.825) (2.408,3.875) (2.375,3.914) (2.325,3.969) (2.275,4.022) (2.225,4.072) (2.175,4.120) (2.125,4.166) (2.087,4.200) (2.028,4.250) (1.975,4.292) (1.933,4.325) (1.875,4.368) (1.829,4.400) (1.775,4.437) (1.725,4.469) (1.676,4.500) (1.625,4.530) (1.575,4.559) (1.525,4.586) (1.475,4.613) (1.425,4.638) (1.375,4.662) (1.325,4.685) (1.275,4.707) (1.225,4.727) (1.175,4.747) (1.125,4.766) (1.075,4.784) (1.025,4.801) (0.975,4.818) (0.925,4.833) (0.875,4.847) (0.825,4.861) (0.775,4.873) (0.725,4.885) (0.675,4.896) (0.625,4.907) (0.575,4.916) (0.525,4.924) (0.475,4.932) (0.425,4.939) (0.375,4.945) (0.325,4.951) (0.275,4.955) (0.225,4.959) (0.175,4.962) (0.125,4.965) (0.075,4.966) (0.025,4.967) (-0.050,4.967) (-0.100,4.965) (-0.150,4.963) (-0.200,4.961) (-0.250,4.957) (-0.300,4.953) (-0.350,4.948) (-0.400,4.942) (-0.450,4.936) (-0.500,4.928) (-0.550,4.920) (-0.600,4.911) (-0.650,4.902) (-0.700,4.891) (-0.750,4.880) (-0.800,4.867) (-0.850,4.854) (-0.900,4.840) (-0.950,4.825) (-1.000,4.810) (-1.050,4.793) (-1.100,4.775) (-1.150,4.757) (-1.200,4.738) (-1.250,4.717) (-1.300,4.696) (-1.346,4.675) (-1.400,4.650) (-1.450,4.625) (-1.499,4.600) (-1.546,4.575) (-1.591,4.550) (-1.650,4.516) (-1.700,4.485) (-1.750,4.453) (-1.793,4.425) (-1.850,4.385) (-1.899,4.350) (-1.950,4.312) (-1.997,4.275) (-2.050,4.231) (-2.100,4.189) (-2.143,4.150) (-2.197,4.100) (-2.248,4.050) (-2.296,4.000) (-2.343,3.950) (-2.387,3.900) (-2.425,3.855) (-2.470,3.800) (-2.508,3.750) (-2.545,3.700) (-2.575,3.658) (-2.614,3.600) (-2.646,3.550) (-2.675,3.503) (-2.706,3.450) (-2.734,3.400) (-2.760,3.350) (-2.786,3.300) (-2.809,3.250) (-2.832,3.200) (-2.853,3.150) (-2.874,3.100) (-2.893,3.050) (-2.911,3.000) (-2.927,2.950) (-2.943,2.900) (-2.957,2.850) (-2.971,2.800) (-2.983,2.750) (-2.994,2.700) (-3.005,2.650) (-3.014,2.600) (-3.022,2.550) (-3.029,2.500) (-3.034,2.450) (-3.039,2.400) (-3.043,2.350) (-3.046,2.300) (-3.047,2.250) (-3.048,2.200) (-3.047,2.150) (-3.046,2.100) (-3.043,2.050) (-3.039,2.000) (-3.034,1.950) (-3.028,1.900) (-3.021,1.850) (-3.012,1.800) (-3.003,1.750) (-2.992,1.700) (-2.979,1.650) (-2.966,1.600) (-2.951,1.550) (-2.935,1.500) (-2.918,1.450) (-2.898,1.400) (-2.878,1.350) (-2.856,1.300) (-2.832,1.250) (-2.807,1.200) (-2.780,1.150) (-2.751,1.100) (-2.720,1.050) (-2.686,1.000) (-2.651,0.950) (-2.613,0.900) (-2.575,0.852) (-2.530,0.800) (-2.484,0.750) (-2.435,0.700) (-2.382,0.650) (-2.325,0.601) (-2.275,0.560) (-2.230,0.525) (-2.175,0.485) (-2.125,0.451) (-2.075,0.419) (-2.025,0.389) (-1.975,0.361) (-1.925,0.335) (-1.875,0.310) (-1.825,0.286) (-1.775,0.264) (-1.725,0.243) (-1.675,0.223) (-1.625,0.205) (-1.575,0.188) (-1.525,0.172) (-1.475,0.157) (-1.425,0.143) (-1.375,0.129) (-1.325,0.117) (-1.275,0.106) (-1.225,0.095) (-1.175,0.086) (-1.125,0.077) (-1.075,0.069) (-1.025,0.061) (-0.975,0.054) (-0.925,0.048) (-0.875,0.042) (-0.825,0.037) (-0.775,0.032) (-0.725,0.028) (-0.675,0.024) (-0.625,0.021) (-0.575,0.018) (-0.525,0.016) (-0.475,0.013) (-0.425,0.012) (-0.375,0.010) (-0.325,0.009) (-0.275,0.008) (-0.250,0.007)}
\EQ{(0.250,0.008) (0.300,0.009) (0.350,0.011) (0.400,0.012) (0.450,0.014) (0.500,0.017) (0.550,0.019) (0.600,0.022) (0.650,0.026) (0.700,0.030) (0.750,0.034) (0.800,0.039) (0.850,0.045) (0.900,0.051) (0.950,0.058) (1.000,0.066) (1.050,0.074) (1.100,0.083) (1.150,0.093) (1.200,0.104) (1.250,0.115) (1.300,0.128) (1.350,0.141) (1.400,0.156) (1.450,0.172) (1.500,0.188) (1.550,0.206) (1.599,0.225) (1.650,0.246) (1.700,0.268) (1.750,0.292) (1.800,0.317) (1.850,0.344) (1.900,0.372) (1.946,0.400) (2.000,0.435) (2.050,0.470) (2.091,0.500) (2.150,0.547) (2.200,0.589) (2.240,0.625) (2.292,0.675) (2.340,0.725) (2.375,0.764) (2.425,0.824) (2.465,0.875) (2.500,0.924) (2.534,0.975) (2.566,1.025) (2.595,1.075) (2.622,1.125) (2.648,1.175) (2.671,1.225) (2.693,1.275) (2.713,1.325) (2.732,1.375) (2.749,1.425) (2.765,1.475) (2.779,1.525) (2.791,1.575) (2.803,1.625) (2.813,1.675) (2.822,1.725) (2.829,1.775) (2.835,1.825) (2.840,1.875) (2.844,1.925) (2.847,1.975) (2.848,2.025) (2.849,2.075) (2.848,2.125) (2.846,2.175) (2.843,2.225) (2.838,2.275) (2.833,2.325) (2.826,2.375) (2.818,2.425) (2.810,2.475) (2.800,2.525) (2.789,2.575) (2.776,2.625) (2.763,2.675) (2.748,2.725) (2.732,2.775) (2.715,2.825) (2.697,2.875) (2.677,2.925) (2.657,2.975) (2.635,3.025) (2.611,3.075) (2.586,3.125) (2.560,3.175) (2.533,3.225) (2.504,3.275) (2.475,3.322) (2.441,3.375) (2.407,3.425) (2.375,3.470) (2.334,3.525) (2.300,3.569) (2.254,3.625) (2.211,3.675) (2.175,3.715) (2.125,3.767) (2.075,3.818) (2.025,3.865) (1.987,3.900) (1.930,3.950) (1.875,3.996) (1.825,4.035) (1.775,4.073) (1.725,4.109) (1.675,4.143) (1.627,4.175) (1.575,4.208) (1.525,4.238) (1.475,4.266) (1.425,4.294) (1.375,4.320) (1.325,4.345) (1.275,4.369) (1.225,4.392) (1.175,4.413) (1.125,4.434) (1.075,4.453) (1.025,4.472) (0.975,4.489) (0.925,4.506) (0.875,4.521) (0.825,4.536) (0.775,4.549) (0.725,4.562) (0.675,4.574) (0.625,4.585) (0.575,4.595) (0.525,4.604) (0.475,4.612) (0.425,4.620) (0.375,4.626) (0.325,4.632) (0.275,4.637) (0.225,4.641) (0.175,4.644) (0.125,4.647) (0.075,4.649) (0.025,4.649) (-0.050,4.649) (-0.100,4.648) (-0.150,4.646) (-0.200,4.643) (-0.250,4.639) (-0.300,4.635) (-0.350,4.629) (-0.400,4.623) (-0.450,4.616) (-0.500,4.608) (-0.550,4.599) (-0.600,4.590) (-0.650,4.579) (-0.700,4.568) (-0.750,4.556) (-0.800,4.543) (-0.850,4.528) (-0.900,4.513) (-0.950,4.497) (-1.000,4.480) (-1.050,4.463) (-1.100,4.444) (-1.150,4.424) (-1.200,4.402) (-1.250,4.380) (-1.300,4.357) (-1.350,4.333) (-1.400,4.307) (-1.450,4.280) (-1.500,4.252) (-1.546,4.225) (-1.600,4.192) (-1.650,4.160) (-1.700,4.126) (-1.750,4.091) (-1.800,4.054) (-1.850,4.016) (-1.900,3.975) (-1.950,3.933) (-2.000,3.888) (-2.041,3.850) (-2.093,3.800) (-2.142,3.750) (-2.188,3.700) (-2.225,3.659) (-2.275,3.600) (-2.315,3.550) (-2.350,3.504) (-2.389,3.450) (-2.424,3.400) (-2.457,3.350) (-2.488,3.300) (-2.518,3.250) (-2.547,3.200) (-2.573,3.150) (-2.599,3.100) (-2.623,3.050) (-2.646,3.000) (-2.667,2.950) (-2.687,2.900) (-2.706,2.850) (-2.724,2.800) (-2.740,2.750) (-2.756,2.700) (-2.770,2.650) (-2.783,2.600) (-2.794,2.550) (-2.805,2.500) (-2.814,2.450) (-2.822,2.400) (-2.830,2.350) (-2.836,2.300) (-2.841,2.250) (-2.844,2.200) (-2.847,2.150) (-2.848,2.100) (-2.849,2.050) (-2.848,2.000) (-2.846,1.950) (-2.842,1.900) (-2.838,1.850) (-2.832,1.800) (-2.826,1.750) (-2.817,1.700) (-2.808,1.650) (-2.797,1.600) (-2.785,1.550) (-2.772,1.500) (-2.757,1.450) (-2.741,1.400) (-2.723,1.350) (-2.703,1.300) (-2.683,1.250) (-2.660,1.200) (-2.635,1.150) (-2.609,1.100) (-2.581,1.050) (-2.550,1.000) (-2.518,0.950) (-2.483,0.900) (-2.450,0.856) (-2.405,0.800) (-2.362,0.750) (-2.325,0.709) (-2.275,0.659) (-2.225,0.611) (-2.184,0.575) (-2.125,0.526) (-2.075,0.488) (-2.025,0.452) (-1.975,0.419) (-1.925,0.387) (-1.875,0.358) (-1.825,0.330) (-1.775,0.304) (-1.725,0.280) (-1.675,0.257) (-1.625,0.236) (-1.575,0.216) (-1.525,0.197) (-1.475,0.180) (-1.425,0.164) (-1.375,0.148) (-1.325,0.134) (-1.275,0.121) (-1.225,0.109) (-1.175,0.098) (-1.125,0.088) (-1.075,0.078) (-1.025,0.070) (-0.975,0.062) (-0.925,0.055) (-0.875,0.048) (-0.825,0.042) (-0.775,0.037) (-0.725,0.032) (-0.675,0.028) (-0.625,0.024) (-0.575,0.021) (-0.525,0.018) (-0.475,0.015) (-0.425,0.013) (-0.375,0.011) (-0.325,0.010) (-0.275,0.009) (-0.250,0.008)}
\EQ{(0.250,0.009) (0.300,0.010) (0.350,0.012) (0.400,0.014) (0.450,0.016) (0.500,0.019) (0.550,0.022) (0.600,0.025) (0.650,0.029) (0.700,0.034) (0.750,0.039) (0.800,0.044) (0.850,0.051) (0.900,0.058) (0.950,0.065) (1.000,0.074) (1.050,0.083) (1.100,0.094) (1.150,0.105) (1.200,0.117) (1.250,0.130) (1.300,0.144) (1.350,0.159) (1.400,0.176) (1.450,0.194) (1.500,0.213) (1.550,0.233) (1.600,0.255) (1.650,0.279) (1.700,0.304) (1.750,0.331) (1.800,0.360) (1.850,0.391) (1.900,0.425) (1.950,0.460) (2.000,0.498) (2.050,0.540) (2.091,0.575) (2.144,0.625) (2.193,0.675) (2.238,0.725) (2.275,0.769) (2.319,0.825) (2.350,0.868) (2.389,0.925) (2.420,0.975) (2.449,1.025) (2.475,1.074) (2.500,1.124) (2.524,1.175) (2.545,1.225) (2.564,1.275) (2.582,1.325) (2.598,1.375) (2.613,1.425) (2.626,1.475) (2.638,1.525) (2.648,1.575) (2.657,1.625) (2.665,1.675) (2.671,1.725) (2.676,1.775) (2.680,1.825) (2.682,1.875) (2.683,1.925) (2.683,1.975) (2.682,2.025) (2.680,2.075) (2.676,2.125) (2.671,2.175) (2.665,2.225) (2.658,2.275) (2.649,2.325) (2.639,2.375) (2.628,2.425) (2.616,2.475) (2.603,2.525) (2.588,2.575) (2.572,2.625) (2.555,2.675) (2.537,2.725) (2.517,2.775) (2.496,2.825) (2.475,2.871) (2.450,2.923) (2.425,2.972) (2.400,3.019) (2.375,3.063) (2.338,3.125) (2.306,3.175) (2.275,3.222) (2.238,3.275) (2.200,3.325) (2.161,3.375) (2.125,3.419) (2.077,3.475) (2.031,3.525) (1.983,3.575) (1.932,3.625) (1.878,3.675) (1.825,3.722) (1.775,3.763) (1.729,3.800) (1.675,3.841) (1.628,3.875) (1.575,3.912) (1.525,3.944) (1.476,3.975) (1.425,4.005) (1.375,4.034) (1.325,4.061) (1.275,4.087) (1.225,4.111) (1.175,4.134) (1.125,4.156) (1.075,4.177) (1.025,4.197) (0.975,4.216) (0.925,4.233) (0.875,4.250) (0.825,4.266) (0.775,4.280) (0.725,4.294) (0.675,4.306) (0.625,4.318) (0.575,4.329) (0.525,4.338) (0.475,4.347) (0.425,4.355) (0.375,4.362) (0.325,4.368) (0.275,4.374) (0.225,4.378) (0.175,4.381) (0.125,4.384) (0.075,4.386) (0.025,4.387) (-0.050,4.386) (-0.100,4.385) (-0.150,4.383) (-0.200,4.380) (-0.250,4.376) (-0.300,4.371) (-0.350,4.365) (-0.400,4.359) (-0.450,4.351) (-0.500,4.343) (-0.550,4.334) (-0.600,4.323) (-0.650,4.312) (-0.700,4.300) (-0.750,4.287) (-0.800,4.273) (-0.850,4.258) (-0.900,4.242) (-0.949,4.225) (-1.000,4.207) (-1.050,4.187) (-1.100,4.167) (-1.150,4.145) (-1.200,4.123) (-1.248,4.100) (-1.298,4.075) (-1.345,4.050) (-1.391,4.025) (-1.434,4.000) (-1.500,3.960) (-1.550,3.928) (-1.592,3.900) (-1.650,3.859) (-1.696,3.825) (-1.750,3.784) (-1.791,3.750) (-1.850,3.700) (-1.900,3.655) (-1.950,3.607) (-2.000,3.557) (-2.050,3.504) (-2.099,3.450) (-2.141,3.400) (-2.175,3.358) (-2.219,3.300) (-2.250,3.258) (-2.290,3.200) (-2.322,3.150) (-2.350,3.105) (-2.383,3.050) (-2.410,3.000) (-2.436,2.950) (-2.461,2.900) (-2.484,2.850) (-2.506,2.800) (-2.527,2.750) (-2.546,2.700) (-2.564,2.650) (-2.580,2.600) (-2.596,2.550) (-2.610,2.500) (-2.622,2.450) (-2.634,2.400) (-2.644,2.350) (-2.653,2.300) (-2.661,2.250) (-2.668,2.200) (-2.674,2.150) (-2.678,2.100) (-2.681,2.050) (-2.683,2.000) (-2.683,1.950) (-2.683,1.900) (-2.681,1.850) (-2.678,1.800) (-2.674,1.750) (-2.668,1.700) (-2.661,1.650) (-2.653,1.600) (-2.643,1.550) (-2.632,1.500) (-2.620,1.450) (-2.606,1.400) (-2.590,1.350) (-2.573,1.300) (-2.555,1.250) (-2.534,1.200) (-2.512,1.150) (-2.488,1.100) (-2.463,1.050) (-2.435,1.000) (-2.405,0.950) (-2.375,0.904) (-2.337,0.850) (-2.300,0.801) (-2.259,0.750) (-2.225,0.711) (-2.175,0.657) (-2.125,0.607) (-2.075,0.561) (-2.033,0.525) (-1.975,0.479) (-1.925,0.442) (-1.875,0.408) (-1.825,0.376) (-1.775,0.346) (-1.725,0.318) (-1.675,0.292) (-1.625,0.267) (-1.575,0.244) (-1.525,0.223) (-1.475,0.203) (-1.425,0.185) (-1.375,0.168) (-1.325,0.152) (-1.275,0.137) (-1.225,0.123) (-1.175,0.111) (-1.125,0.099) (-1.075,0.088) (-1.025,0.078) (-0.975,0.070) (-0.925,0.061) (-0.875,0.054) (-0.825,0.047) (-0.775,0.041) (-0.725,0.036) (-0.675,0.031) (-0.625,0.027) (-0.575,0.023) (-0.525,0.020) (-0.475,0.017) (-0.425,0.015) (-0.375,0.013) (-0.325,0.011) (-0.275,0.010) (-0.250,0.009)}
\EQ{(0.250,0.010) (0.300,0.012) (0.350,0.013) (0.400,0.015) (0.450,0.018) (0.500,0.021) (0.550,0.024) (0.600,0.028) (0.650,0.032) (0.700,0.037) (0.750,0.043) (0.800,0.049) (0.850,0.056) (0.900,0.064) (0.950,0.073) (1.000,0.082) (1.050,0.093) (1.100,0.104) (1.150,0.116) (1.200,0.130) (1.250,0.145) (1.300,0.161) (1.350,0.178) (1.400,0.196) (1.450,0.216) (1.500,0.238) (1.550,0.261) (1.600,0.286) (1.650,0.313) (1.700,0.341) (1.750,0.372) (1.792,0.400) (1.850,0.441) (1.894,0.475) (1.950,0.521) (2.000,0.566) (2.050,0.615) (2.100,0.668) (2.150,0.725) (2.189,0.775) (2.225,0.824) (2.260,0.875) (2.292,0.925) (2.321,0.975) (2.347,1.025) (2.372,1.075) (2.395,1.125) (2.416,1.175) (2.435,1.225) (2.452,1.275) (2.468,1.325) (2.482,1.375) (2.495,1.425) (2.506,1.475) (2.515,1.525) (2.523,1.575) (2.530,1.625) (2.535,1.675) (2.539,1.725) (2.542,1.775) (2.543,1.825) (2.543,1.875) (2.542,1.925) (2.540,1.975) (2.536,2.025) (2.531,2.075) (2.525,2.125) (2.517,2.175) (2.508,2.225) (2.498,2.275) (2.487,2.325) (2.474,2.375) (2.460,2.425) (2.444,2.475) (2.428,2.525) (2.410,2.575) (2.390,2.625) (2.369,2.675) (2.347,2.725) (2.325,2.771) (2.300,2.821) (2.275,2.867) (2.250,2.912) (2.212,2.975) (2.180,3.025) (2.150,3.069) (2.110,3.125) (2.075,3.172) (2.033,3.225) (2.000,3.264) (1.950,3.321) (1.900,3.374) (1.850,3.424) (1.800,3.472) (1.750,3.517) (1.711,3.550) (1.650,3.600) (1.600,3.638) (1.550,3.675) (1.500,3.709) (1.450,3.742) (1.400,3.774) (1.356,3.800) (1.312,3.825) (1.266,3.850) (1.218,3.875) (1.168,3.900) (1.114,3.925) (1.057,3.950) (1.000,3.974) (0.950,3.993) (0.900,4.011) (0.850,4.028) (0.800,4.044) (0.750,4.059) (0.700,4.073) (0.650,4.086) (0.600,4.098) (0.550,4.108) (0.500,4.118) (0.450,4.127) (0.400,4.135) (0.350,4.142) (0.300,4.148) (0.250,4.153) (0.200,4.157) (0.150,4.161) (0.100,4.163) (0.050,4.164) (0.000,4.165) (-0.050,4.164) (-0.100,4.163) (-0.150,4.161) (-0.200,4.157) (-0.250,4.153) (-0.300,4.148) (-0.350,4.142) (-0.400,4.135) (-0.450,4.127) (-0.500,4.118) (-0.550,4.108) (-0.600,4.098) (-0.650,4.086) (-0.700,4.073) (-0.750,4.059) (-0.800,4.044) (-0.850,4.028) (-0.900,4.011) (-0.950,3.993) (-1.000,3.974) (-1.050,3.953) (-1.100,3.931) (-1.150,3.909) (-1.200,3.884) (-1.250,3.859) (-1.300,3.832) (-1.350,3.803) (-1.398,3.775) (-1.450,3.742) (-1.500,3.709) (-1.549,3.675) (-1.600,3.638) (-1.649,3.600) (-1.700,3.559) (-1.740,3.525) (-1.796,3.475) (-1.849,3.425) (-1.899,3.375) (-1.946,3.325) (-1.990,3.275) (-2.025,3.234) (-2.072,3.175) (-2.110,3.125) (-2.146,3.075) (-2.175,3.032) (-2.212,2.975) (-2.242,2.925) (-2.271,2.875) (-2.298,2.825) (-2.323,2.775) (-2.347,2.725) (-2.369,2.675) (-2.390,2.625) (-2.410,2.575) (-2.428,2.525) (-2.444,2.475) (-2.460,2.425) (-2.474,2.375) (-2.487,2.325) (-2.498,2.275) (-2.508,2.225) (-2.517,2.175) (-2.525,2.125) (-2.531,2.075) (-2.536,2.025) (-2.540,1.975) (-2.542,1.925) (-2.543,1.875) (-2.543,1.825) (-2.542,1.775) (-2.539,1.725) (-2.535,1.675) (-2.530,1.625) (-2.523,1.575) (-2.515,1.525) (-2.506,1.475) (-2.495,1.425) (-2.482,1.375) (-2.468,1.325) (-2.452,1.275) (-2.435,1.225) (-2.416,1.175) (-2.395,1.125) (-2.372,1.075) (-2.347,1.025) (-2.321,0.975) (-2.292,0.925) (-2.260,0.875) (-2.226,0.825) (-2.189,0.775) (-2.150,0.726) (-2.107,0.675) (-2.060,0.625) (-2.010,0.575) (-1.954,0.525) (-1.900,0.480) (-1.850,0.441) (-1.800,0.406) (-1.754,0.375) (-1.700,0.341) (-1.650,0.313) (-1.600,0.286) (-1.550,0.261) (-1.500,0.238) (-1.450,0.216) (-1.400,0.196) (-1.350,0.178) (-1.300,0.161) (-1.250,0.145) (-1.200,0.130) (-1.150,0.116) (-1.100,0.104) (-1.050,0.093) (-1.000,0.082) (-0.950,0.073) (-0.900,0.064) (-0.850,0.056) (-0.800,0.049) (-0.750,0.043) (-0.700,0.037) (-0.650,0.032) (-0.600,0.028) (-0.550,0.024) (-0.500,0.021) (-0.450,0.018) (-0.400,0.015) (-0.350,0.013) (-0.300,0.012) (-0.250,0.010)}

    \draw[mmaEQ] (-10,0) -- (10,0);
    % Dipolo físico - líneas de campo
% Generado por scripts/generar_datos.py (no editar a mano).
\LC{(1.086,0.495) (1.147,0.452) (1.204,0.403) (1.254,0.348) (1.298,0.287) (1.334,0.221) (1.360,0.150) (1.376,0.077) (1.382,0.003) (1.377,-0.072) (1.361,-0.146) (1.336,-0.216) (1.301,-0.282) (1.258,-0.344) (1.207,-0.399) (1.151,-0.449) (1.090,-0.492)}
\LC{(0.450,1.109) (0.501,1.164) (0.552,1.219) (0.605,1.272) (0.658,1.325) (0.713,1.376) (0.768,1.426) (0.825,1.476) (0.882,1.524) (0.941,1.571) (1.000,1.617) (1.060,1.662) (1.121,1.706) (1.183,1.748) (1.245,1.790) (1.308,1.830) (1.372,1.869) (1.437,1.907) (1.503,1.943) (1.569,1.979) (1.636,2.013) (1.703,2.046) (1.771,2.077) (1.840,2.108) (1.909,2.137) (1.978,2.165) (2.049,2.191) (2.119,2.216) (2.190,2.240) (2.262,2.262) (2.334,2.284) (2.406,2.303) (2.479,2.322) (2.552,2.339) (2.625,2.354) (2.699,2.369) (2.773,2.381) (2.847,2.393) (2.921,2.403) (2.996,2.411) (3.071,2.418) (3.145,2.424) (3.220,2.428) (3.295,2.430) (3.370,2.431) (3.445,2.430) (3.520,2.428) (3.595,2.425) (3.670,2.420) (3.745,2.413) (3.819,2.404) (3.893,2.394) (3.968,2.383) (4.041,2.370) (4.115,2.355) (4.188,2.338) (4.261,2.320) (4.333,2.301) (4.405,2.279) (4.476,2.256) (4.547,2.231) (4.617,2.205) (4.687,2.177) (4.756,2.147) (4.824,2.115) (4.891,2.082) (4.957,2.047) (5.023,2.011) (5.087,1.973) (5.151,1.933) (5.213,1.891) (5.274,1.848) (5.335,1.803) (5.393,1.756) (5.451,1.708) (5.507,1.658) (5.561,1.606) (5.614,1.553) (5.666,1.499) (5.715,1.443) (5.763,1.385) (5.810,1.326) (5.854,1.265) (5.896,1.204) (5.937,1.140) (5.975,1.076) (6.011,1.010) (6.045,0.943) (6.077,0.875) (6.107,0.807) (6.134,0.737) (6.158,0.666) (6.181,0.594) (6.200,0.522) (6.218,0.449) (6.232,0.375) (6.244,0.301) (6.254,0.227) (6.260,0.152) (6.265,0.077) (6.266,0.002) (6.265,-0.073) (6.261,-0.148) (6.254,-0.222) (6.245,-0.297) (6.233,-0.371) (6.219,-0.444) (6.202,-0.517) (6.182,-0.590) (6.160,-0.661) (6.135,-0.732) (6.108,-0.802) (6.079,-0.871) (6.047,-0.939) (6.013,-1.006) (5.977,-1.072) (5.939,-1.136) (5.899,-1.200) (5.857,-1.262) (5.812,-1.322) (5.766,-1.381) (5.718,-1.439) (5.669,-1.495) (5.617,-1.550) (5.564,-1.603) (5.510,-1.655) (5.454,-1.705) (5.397,-1.753) (5.338,-1.800) (5.278,-1.845) (5.217,-1.888) (5.155,-1.930) (5.091,-1.970) (5.027,-2.008) (4.962,-2.045) (4.895,-2.080) (4.828,-2.113) (4.760,-2.145) (4.691,-2.175) (4.622,-2.203) (4.552,-2.230) (4.481,-2.255) (4.409,-2.278) (4.338,-2.299) (4.265,-2.319) (4.193,-2.337) (4.119,-2.354) (4.046,-2.369) (3.972,-2.382) (3.898,-2.394) (3.824,-2.404) (3.749,-2.412) (3.674,-2.419) (3.600,-2.424) (3.525,-2.428) (3.450,-2.430) (3.375,-2.431) (3.300,-2.430) (3.225,-2.428) (3.150,-2.424) (3.075,-2.418) (3.000,-2.412) (2.926,-2.403) (2.852,-2.393) (2.777,-2.382) (2.704,-2.369) (2.630,-2.355) (2.556,-2.340) (2.483,-2.323) (2.411,-2.305) (2.338,-2.285) (2.266,-2.264) (2.195,-2.241) (2.124,-2.218) (2.053,-2.193) (1.983,-2.166) (1.913,-2.138) (1.844,-2.110) (1.775,-2.079) (1.707,-2.048) (1.640,-2.015) (1.573,-1.981) (1.507,-1.946) (1.441,-1.909) (1.376,-1.871) (1.312,-1.832) (1.249,-1.792) (1.186,-1.751) (1.125,-1.708) (1.064,-1.665) (1.003,-1.620) (0.944,-1.574) (0.886,-1.527) (0.828,-1.479) (0.772,-1.430) (0.716,-1.379) (0.661,-1.328) (0.608,-1.275) (0.555,-1.222) (0.504,-1.167) (0.453,-1.112) (0.437,-1.093)}
\LC{(-0.450,1.109) (-0.501,1.164) (-0.552,1.219) (-0.605,1.272) (-0.658,1.325) (-0.713,1.376) (-0.768,1.426) (-0.825,1.476) (-0.882,1.524) (-0.941,1.571) (-1.000,1.617) (-1.060,1.662) (-1.121,1.706) (-1.183,1.748) (-1.245,1.790) (-1.308,1.830) (-1.372,1.869) (-1.437,1.907) (-1.503,1.943) (-1.569,1.979) (-1.636,2.013) (-1.703,2.046) (-1.771,2.077) (-1.840,2.108) (-1.909,2.137) (-1.978,2.165) (-2.049,2.191) (-2.119,2.216) (-2.190,2.240) (-2.262,2.262) (-2.334,2.284) (-2.406,2.303) (-2.479,2.322) (-2.552,2.339) (-2.625,2.354) (-2.699,2.369) (-2.773,2.381) (-2.847,2.393) (-2.921,2.403) (-2.996,2.411) (-3.071,2.418) (-3.145,2.424) (-3.220,2.428) (-3.295,2.430) (-3.370,2.431) (-3.445,2.430) (-3.520,2.428) (-3.595,2.425) (-3.670,2.420) (-3.745,2.413) (-3.819,2.404) (-3.893,2.394) (-3.968,2.383) (-4.041,2.370) (-4.115,2.355) (-4.188,2.338) (-4.261,2.320) (-4.333,2.301) (-4.405,2.279) (-4.476,2.256) (-4.547,2.231) (-4.617,2.205) (-4.687,2.177) (-4.756,2.147) (-4.824,2.115) (-4.891,2.082) (-4.957,2.047) (-5.023,2.011) (-5.087,1.973) (-5.151,1.933) (-5.213,1.891) (-5.274,1.848) (-5.335,1.803) (-5.393,1.756) (-5.451,1.708) (-5.507,1.658) (-5.561,1.606) (-5.614,1.553) (-5.666,1.499) (-5.715,1.443) (-5.763,1.385) (-5.810,1.326) (-5.854,1.265) (-5.896,1.204) (-5.937,1.140) (-5.975,1.076) (-6.011,1.010) (-6.045,0.943) (-6.077,0.875) (-6.107,0.807) (-6.134,0.737) (-6.158,0.666) (-6.181,0.594) (-6.200,0.522) (-6.218,0.449) (-6.232,0.375) (-6.244,0.301) (-6.254,0.227) (-6.260,0.152) (-6.265,0.077) (-6.266,0.002) (-6.265,-0.073) (-6.261,-0.148) (-6.254,-0.222) (-6.245,-0.297) (-6.233,-0.371) (-6.219,-0.444) (-6.202,-0.517) (-6.182,-0.590) (-6.160,-0.661) (-6.135,-0.732) (-6.108,-0.802) (-6.079,-0.871) (-6.047,-0.939) (-6.013,-1.006) (-5.977,-1.072) (-5.939,-1.136) (-5.899,-1.200) (-5.857,-1.262) (-5.812,-1.322) (-5.766,-1.381) (-5.718,-1.439) (-5.669,-1.495) (-5.617,-1.550) (-5.564,-1.603) (-5.510,-1.655) (-5.454,-1.705) (-5.397,-1.753) (-5.338,-1.800) (-5.278,-1.845) (-5.217,-1.888) (-5.155,-1.930) (-5.091,-1.970) (-5.027,-2.008) (-4.962,-2.045) (-4.895,-2.080) (-4.828,-2.113) (-4.760,-2.145) (-4.691,-2.175) (-4.622,-2.203) (-4.552,-2.230) (-4.481,-2.255) (-4.409,-2.278) (-4.338,-2.299) (-4.265,-2.319) (-4.193,-2.337) (-4.119,-2.354) (-4.046,-2.369) (-3.972,-2.382) (-3.898,-2.394) (-3.824,-2.404) (-3.749,-2.412) (-3.674,-2.419) (-3.600,-2.424) (-3.525,-2.428) (-3.450,-2.430) (-3.375,-2.431) (-3.300,-2.430) (-3.225,-2.428) (-3.150,-2.424) (-3.075,-2.418) (-3.000,-2.412) (-2.926,-2.403) (-2.852,-2.393) (-2.777,-2.382) (-2.704,-2.369) (-2.630,-2.355) (-2.556,-2.340) (-2.483,-2.323) (-2.411,-2.305) (-2.338,-2.285) (-2.266,-2.264) (-2.195,-2.241) (-2.124,-2.218) (-2.053,-2.193) (-1.983,-2.166) (-1.913,-2.138) (-1.844,-2.110) (-1.775,-2.079) (-1.707,-2.048) (-1.640,-2.015) (-1.573,-1.981) (-1.507,-1.946) (-1.441,-1.909) (-1.376,-1.871) (-1.312,-1.832) (-1.249,-1.792) (-1.186,-1.751) (-1.125,-1.708) (-1.064,-1.665) (-1.003,-1.620) (-0.944,-1.574) (-0.886,-1.527) (-0.828,-1.479) (-0.772,-1.430) (-0.716,-1.379) (-0.661,-1.328) (-0.608,-1.275) (-0.555,-1.222) (-0.504,-1.167) (-0.453,-1.112) (-0.437,-1.093)}
\LC{(-1.086,0.495) (-1.147,0.452) (-1.204,0.403) (-1.254,0.348) (-1.298,0.287) (-1.334,0.221) (-1.360,0.150) (-1.376,0.077) (-1.382,0.003) (-1.377,-0.072) (-1.361,-0.146) (-1.336,-0.216) (-1.301,-0.282) (-1.258,-0.344) (-1.207,-0.399) (-1.151,-0.449) (-1.090,-0.492)}
\LC{(4.895,-10.000) (4.847,-9.949) (4.799,-9.892) (4.752,-9.834) (4.704,-9.776) (4.656,-9.718) (4.609,-9.660) (4.562,-9.601) (4.515,-9.543) (4.468,-9.484) (4.421,-9.426) (4.375,-9.367) (4.329,-9.308) (4.282,-9.249) (4.236,-9.190) (4.190,-9.130) (4.145,-9.071) (4.099,-9.011) (4.053,-8.952) (4.008,-8.892) (3.963,-8.832) (3.918,-8.772) (3.873,-8.712) (3.829,-8.652) (3.784,-8.591) (3.740,-8.531) (3.696,-8.470) (3.652,-8.409) (3.608,-8.349) (3.564,-8.288) (3.521,-8.226) (3.477,-8.165) (3.434,-8.104) (3.391,-8.042) (3.348,-7.981) (3.306,-7.919) (3.263,-7.857) (3.221,-7.795) (3.179,-7.733) (3.137,-7.671) (3.095,-7.609) (3.053,-7.547) (3.012,-7.484) (2.971,-7.421) (2.929,-7.359) (2.889,-7.296) (2.848,-7.233) (2.807,-7.170) (2.767,-7.107) (2.727,-7.043) (2.687,-6.980) (2.647,-6.916) (2.607,-6.852) (2.568,-6.789) (2.529,-6.725) (2.490,-6.661) (2.451,-6.597) (2.412,-6.532) (2.374,-6.468) (2.336,-6.403) (2.298,-6.339) (2.260,-6.274) (2.222,-6.209) (2.185,-6.144) (2.147,-6.079) (2.110,-6.014) (2.073,-5.948) (2.037,-5.883) (2.000,-5.817) (1.964,-5.752) (1.928,-5.686) (1.892,-5.620) (1.857,-5.554) (1.821,-5.488) (1.786,-5.422) (1.751,-5.355) (1.717,-5.289) (1.682,-5.222) (1.648,-5.156) (1.614,-5.089) (1.580,-5.022) (1.546,-4.955) (1.513,-4.888) (1.480,-4.820) (1.447,-4.753) (1.414,-4.685) (1.382,-4.618) (1.350,-4.550) (1.318,-4.482) (1.286,-4.414) (1.255,-4.346) (1.223,-4.278) (1.192,-4.210) (1.162,-4.141) (1.131,-4.073) (1.101,-4.004) (1.071,-3.935) (1.041,-3.866) (1.012,-3.797) (0.983,-3.728) (0.954,-3.659) (0.925,-3.590) (0.897,-3.520) (0.869,-3.451) (0.841,-3.381) (0.814,-3.311) (0.786,-3.241) (0.760,-3.171) (0.733,-3.101) (0.707,-3.031) (0.680,-2.961) (0.655,-2.890) (0.629,-2.820) (0.604,-2.749) (0.579,-2.678) (0.555,-2.607) (0.531,-2.536) (0.507,-2.465) (0.483,-2.394) (0.460,-2.323) (0.437,-2.251) (0.415,-2.180) (0.393,-2.108) (0.371,-2.036) (0.349,-1.964) (0.328,-1.893) (0.307,-1.820) (0.287,-1.748) (0.267,-1.676) (0.247,-1.604) (0.228,-1.531)}
\LC{(-5.083,-10.000) (-5.034,-9.952) (-4.985,-9.896) (-4.935,-9.839) (-4.886,-9.782) (-4.838,-9.725) (-4.789,-9.668) (-4.741,-9.611) (-4.692,-9.553) (-4.644,-9.496) (-4.596,-9.438) (-4.548,-9.381) (-4.500,-9.323) (-4.453,-9.265) (-4.405,-9.207) (-4.358,-9.148) (-4.311,-9.090) (-4.264,-9.032) (-4.217,-8.973) (-4.171,-8.914) (-4.124,-8.855) (-4.078,-8.796) (-4.032,-8.737) (-3.986,-8.678) (-3.940,-8.619) (-3.894,-8.559) (-3.849,-8.499) (-3.804,-8.440) (-3.758,-8.380) (-3.713,-8.320) (-3.669,-8.260) (-3.624,-8.199) (-3.580,-8.139) (-3.535,-8.078) (-3.491,-8.018) (-3.447,-7.957) (-3.403,-7.896) (-3.360,-7.835) (-3.316,-7.774) (-3.273,-7.713) (-3.230,-7.651) (-3.187,-7.590) (-3.144,-7.528) (-3.102,-7.466) (-3.060,-7.404) (-3.017,-7.342) (-2.976,-7.280) (-2.934,-7.218) (-2.892,-7.155) (-2.851,-7.093) (-2.809,-7.030) (-2.768,-6.968) (-2.728,-6.905) (-2.687,-6.842) (-2.646,-6.778) (-2.606,-6.715) (-2.566,-6.652) (-2.526,-6.588) (-2.487,-6.525) (-2.447,-6.461) (-2.408,-6.397) (-2.369,-6.333) (-2.330,-6.269) (-2.291,-6.205) (-2.253,-6.140) (-2.214,-6.076) (-2.176,-6.011) (-2.139,-5.946) (-2.101,-5.881) (-2.064,-5.816) (-2.026,-5.751) (-1.989,-5.686) (-1.953,-5.621) (-1.916,-5.555) (-1.880,-5.490) (-1.844,-5.424) (-1.808,-5.358) (-1.772,-5.292) (-1.737,-5.226) (-1.701,-5.160) (-1.666,-5.093) (-1.632,-5.027) (-1.597,-4.960) (-1.563,-4.894) (-1.529,-4.827) (-1.495,-4.760) (-1.461,-4.693) (-1.428,-4.626) (-1.395,-4.558) (-1.362,-4.491) (-1.329,-4.423) (-1.297,-4.356) (-1.265,-4.288) (-1.233,-4.220) (-1.202,-4.152) (-1.170,-4.084) (-1.139,-4.015) (-1.108,-3.947) (-1.078,-3.879) (-1.048,-3.810) (-1.018,-3.741) (-0.988,-3.672) (-0.959,-3.603) (-0.929,-3.534) (-0.900,-3.465) (-0.872,-3.396) (-0.844,-3.326) (-0.816,-3.257) (-0.788,-3.187) (-0.760,-3.117) (-0.733,-3.047) (-0.706,-2.977) (-0.680,-2.907) (-0.654,-2.837) (-0.628,-2.766) (-0.602,-2.696) (-0.577,-2.625) (-0.552,-2.555) (-0.528,-2.484) (-0.503,-2.413) (-0.479,-2.342) (-0.456,-2.270) (-0.433,-2.199) (-0.410,-2.128) (-0.387,-2.056) (-0.365,-1.984) (-0.343,-1.913) (-0.322,-1.841) (-0.301,-1.769) (-0.280,-1.697) (-0.260,-1.624) (-0.240,-1.552)}
\LC{(10.000,-6.214) (9.948,-6.223) (9.874,-6.231) (9.799,-6.239) (9.724,-6.246) (9.650,-6.253) (9.575,-6.259) (9.500,-6.264) (9.425,-6.269) (9.350,-6.273) (9.275,-6.277) (9.200,-6.280) (9.125,-6.282) (9.050,-6.284) (8.976,-6.285) (8.901,-6.286) (8.826,-6.286) (8.751,-6.285) (8.676,-6.284) (8.601,-6.282) (8.526,-6.280) (8.451,-6.277) (8.376,-6.273) (8.301,-6.269) (8.226,-6.264) (8.151,-6.259) (8.076,-6.253) (8.002,-6.246) (7.927,-6.239) (7.852,-6.232) (7.778,-6.223) (7.703,-6.215) (7.629,-6.205) (7.555,-6.195) (7.480,-6.185) (7.406,-6.174) (7.332,-6.162) (7.258,-6.150) (7.184,-6.137) (7.110,-6.124) (7.037,-6.110) (6.963,-6.096) (6.890,-6.081) (6.816,-6.065) (6.743,-6.049) (6.670,-6.033) (6.597,-6.016) (6.524,-5.998) (6.451,-5.980) (6.379,-5.961) (6.306,-5.942) (6.234,-5.922) (6.162,-5.901) (6.089,-5.880) (6.018,-5.859) (5.946,-5.837) (5.874,-5.815) (5.803,-5.791) (5.732,-5.768) (5.661,-5.744) (5.590,-5.719) (5.519,-5.694) (5.449,-5.668) (5.378,-5.642) (5.308,-5.616) (5.239,-5.588) (5.169,-5.561) (5.099,-5.532) (5.030,-5.504) (4.961,-5.475) (4.892,-5.445) (4.824,-5.415) (4.755,-5.384) (4.687,-5.352) (4.619,-5.321) (4.551,-5.288) (4.484,-5.256) (4.417,-5.222) (4.350,-5.189) (4.283,-5.154) (4.216,-5.120) (4.150,-5.084) (4.084,-5.049) (4.019,-5.013) (3.953,-4.976) (3.888,-4.939) (3.823,-4.901) (3.759,-4.863) (3.694,-4.824) (3.630,-4.785) (3.567,-4.746) (3.503,-4.706) (3.440,-4.665) (3.377,-4.624) (3.315,-4.583) (3.253,-4.541) (3.191,-4.498) (3.129,-4.455) (3.068,-4.412) (3.007,-4.368) (2.947,-4.324) (2.886,-4.279) (2.827,-4.234) (2.767,-4.188) (2.708,-4.142) (2.649,-4.096) (2.591,-4.049) (2.533,-4.001) (2.475,-3.953) (2.418,-3.905) (2.361,-3.856) (2.304,-3.807) (2.248,-3.757) (2.192,-3.707) (2.137,-3.656) (2.082,-3.605) (2.027,-3.554) (1.973,-3.502) (1.920,-3.449) (1.866,-3.396) (1.814,-3.343) (1.761,-3.290) (1.709,-3.235) (1.658,-3.181) (1.607,-3.126) (1.556,-3.070) (1.506,-3.015) (1.457,-2.958) (1.408,-2.902) (1.359,-2.845) (1.311,-2.787) (1.263,-2.729) (1.216,-2.671) (1.170,-2.612) (1.124,-2.553) (1.078,-2.493) (1.033,-2.433) (0.989,-2.372) (0.945,-2.312) (0.931,-2.291)}
\LC{(0.308,1.617) (0.332,1.688) (0.356,1.759) (0.381,1.830) (0.406,1.901) (0.432,1.971) (0.458,2.041) (0.484,2.112) (0.511,2.182) (0.539,2.251) (0.566,2.321) (0.595,2.391) (0.623,2.460) (0.652,2.529) (0.681,2.598) (0.711,2.667) (0.741,2.736) (0.772,2.804) (0.802,2.873) (0.834,2.941) (0.865,3.009) (0.897,3.077) (0.929,3.145) (0.962,3.212) (0.995,3.280) (1.028,3.347) (1.062,3.414) (1.096,3.481) (1.130,3.547) (1.164,3.614) (1.199,3.680) (1.235,3.746) (1.270,3.813) (1.306,3.878) (1.342,3.944) (1.379,4.010) (1.415,4.075) (1.453,4.140) (1.490,4.205) (1.528,4.270) (1.566,4.335) (1.604,4.399) (1.642,4.464) (1.681,4.528) (1.720,4.592) (1.760,4.656) (1.799,4.719) (1.839,4.783) (1.879,4.846) (1.920,4.909) (1.961,4.972) (2.002,5.035) (2.043,5.098) (2.084,5.160) (2.126,5.222) (2.168,5.284) (2.211,5.346) (2.253,5.408) (2.296,5.470) (2.339,5.531) (2.382,5.592) (2.426,5.653) (2.470,5.714) (2.514,5.775) (2.558,5.836) (2.603,5.896) (2.647,5.956) (2.692,6.016) (2.738,6.076) (2.783,6.136) (2.829,6.195) (2.875,6.254) (2.921,6.313) (2.967,6.372) (3.014,6.431) (3.061,6.490) (3.108,6.548) (3.155,6.606) (3.202,6.664) (3.250,6.722) (3.298,6.780) (3.346,6.838) (3.394,6.895) (3.443,6.952) (3.492,7.009) (3.541,7.066) (3.590,7.122) (3.639,7.179) (3.689,7.235) (3.739,7.291) (3.789,7.347) (3.839,7.403) (3.889,7.458) (3.940,7.514) (3.991,7.569) (4.042,7.624) (4.093,7.679) (4.144,7.733) (4.196,7.788) (4.248,7.842) (4.299,7.896) (4.352,7.950) (4.404,8.004) (4.457,8.057) (4.509,8.111) (4.562,8.164) (4.615,8.217) (4.668,8.270) (4.722,8.322) (4.776,8.375) (4.829,8.427) (4.883,8.479) (4.937,8.531) (4.992,8.583) (5.046,8.634) (5.101,8.685) (5.156,8.737) (5.211,8.788) (5.266,8.838) (5.322,8.889) (5.377,8.939) (5.433,8.989) (5.489,9.039) (5.545,9.089) (5.601,9.139) (5.657,9.188) (5.714,9.238) (5.771,9.287) (5.828,9.336) (5.885,9.384) (5.942,9.433) (5.999,9.481) (6.057,9.529) (6.114,9.577) (6.172,9.625) (6.230,9.672) (6.288,9.720) (6.347,9.767) (6.405,9.814) (6.464,9.861) (6.522,9.907) (6.581,9.954) (6.640,10.000)}
\LC{(3.378,2.882) (3.452,2.893) (3.526,2.904) (3.601,2.913) (3.675,2.921) (3.750,2.927) (3.825,2.933) (3.900,2.937) (3.975,2.940) (4.050,2.942) (4.125,2.942) (4.200,2.941) (4.275,2.939) (4.349,2.936) (4.424,2.931) (4.499,2.925) (4.574,2.918) (4.648,2.909) (4.723,2.899) (4.797,2.888) (4.871,2.876) (4.944,2.862) (5.018,2.846) (5.091,2.830) (5.164,2.812) (5.236,2.792) (5.308,2.772) (5.380,2.750) (5.451,2.726) (5.522,2.701) (5.592,2.675) (5.662,2.647) (5.731,2.618) (5.800,2.588) (5.867,2.556) (5.935,2.523) (6.001,2.488) (6.067,2.452) (6.132,2.414) (6.196,2.376) (6.259,2.335) (6.322,2.294) (6.383,2.251) (6.444,2.206) (6.503,2.160) (6.561,2.113) (6.619,2.065) (6.675,2.015) (6.729,1.964) (6.783,1.912) (6.835,1.858) (6.886,1.803) (6.936,1.747) (6.984,1.689) (7.031,1.630) (7.076,1.570) (7.120,1.509) (7.162,1.447) (7.202,1.384) (7.241,1.320) (7.277,1.254) (7.312,1.188) (7.346,1.121) (7.377,1.053) (7.407,0.984) (7.434,0.914) (7.460,0.844) (7.483,0.772) (7.505,0.701) (7.524,0.628) (7.542,0.555) (7.557,0.482) (7.570,0.408) (7.581,0.334) (7.589,0.259) (7.596,0.184) (7.600,0.110) (7.602,0.035) (7.602,-0.040) (7.600,-0.115) (7.596,-0.190) (7.589,-0.265) (7.580,-0.339) (7.569,-0.414) (7.556,-0.487) (7.540,-0.561) (7.523,-0.634) (7.503,-0.706) (7.482,-0.778) (7.458,-0.849) (7.432,-0.920) (7.404,-0.989) (7.375,-1.058) (7.343,-1.126) (7.310,-1.193) (7.275,-1.260) (7.238,-1.325) (7.199,-1.389) (7.158,-1.452) (7.116,-1.514) (7.073,-1.575) (7.027,-1.635) (6.981,-1.693) (6.932,-1.751) (6.883,-1.807) (6.831,-1.862) (6.779,-1.916) (6.725,-1.968) (6.670,-2.019) (6.614,-2.069) (6.557,-2.117) (6.498,-2.164) (6.439,-2.210) (6.378,-2.254) (6.317,-2.297) (6.254,-2.338) (6.191,-2.379) (6.127,-2.417) (6.062,-2.455) (5.996,-2.491) (5.929,-2.525) (5.862,-2.558) (5.794,-2.590) (5.726,-2.620) (5.656,-2.649) (5.587,-2.677) (5.516,-2.703) (5.446,-2.728) (5.374,-2.751) (5.303,-2.773) (5.231,-2.794) (5.158,-2.813) (5.085,-2.831) (5.012,-2.848) (4.939,-2.863) (4.865,-2.877) (4.791,-2.889) (4.717,-2.900) (4.642,-2.910) (4.568,-2.918) (4.493,-2.926) (4.419,-2.932) (4.344,-2.936) (4.269,-2.939) (4.194,-2.941) (4.119,-2.942) (4.044,-2.942) (3.969,-2.940) (3.894,-2.937) (3.819,-2.932) (3.744,-2.927) (3.670,-2.920) (3.595,-2.912) (3.521,-2.903) (3.446,-2.892) (3.397,-2.885)}
\LC{(-1.547,2.471) (-1.606,2.518) (-1.665,2.564) (-1.724,2.609) (-1.785,2.654) (-1.846,2.698) (-1.907,2.741) (-1.969,2.783) (-2.031,2.825) (-2.094,2.866) (-2.158,2.906) (-2.222,2.945) (-2.286,2.983) (-2.351,3.021) (-2.416,3.058) (-2.482,3.094) (-2.548,3.130) (-2.615,3.164) (-2.682,3.198) (-2.749,3.231) (-2.817,3.263) (-2.885,3.294) (-2.953,3.325) (-3.022,3.354) (-3.091,3.383) (-3.161,3.411) (-3.231,3.438) (-3.301,3.465) (-3.372,3.490) (-3.442,3.515) (-3.514,3.539) (-3.585,3.562) (-3.657,3.584) (-3.729,3.605) (-3.801,3.626) (-3.873,3.645) (-3.946,3.664) (-4.019,3.681) (-4.092,3.698) (-4.165,3.714) (-4.239,3.729) (-4.312,3.743) (-4.386,3.757) (-4.460,3.769) (-4.534,3.780) (-4.608,3.791) (-4.683,3.800) (-4.757,3.809) (-4.832,3.817) (-4.907,3.824) (-4.981,3.830) (-5.056,3.834) (-5.131,3.838) (-5.206,3.841) (-5.281,3.843) (-5.356,3.844) (-5.431,3.845) (-5.506,3.844) (-5.581,3.842) (-5.656,3.839) (-5.731,3.835) (-5.806,3.830) (-5.880,3.825) (-5.955,3.818) (-6.030,3.810) (-6.104,3.801) (-6.179,3.791) (-6.253,3.780) (-6.327,3.769) (-6.401,3.756) (-6.474,3.742) (-6.548,3.727) (-6.621,3.711) (-6.694,3.694) (-6.767,3.676) (-6.839,3.656) (-6.912,3.636) (-6.984,3.615) (-7.055,3.593) (-7.126,3.569) (-7.197,3.545) (-7.268,3.519) (-7.338,3.493) (-7.408,3.465) (-7.477,3.437) (-7.546,3.407) (-7.614,3.376) (-7.682,3.344) (-7.750,3.311) (-7.816,3.277) (-7.883,3.242) (-7.948,3.206) (-8.013,3.169) (-8.078,3.130) (-8.142,3.091) (-8.205,3.050) (-8.267,3.009) (-8.329,2.966) (-8.390,2.923) (-8.450,2.878) (-8.510,2.832) (-8.568,2.785) (-8.626,2.738) (-8.683,2.689) (-8.739,2.639) (-8.794,2.588) (-8.848,2.536) (-8.901,2.483) (-8.954,2.429) (-9.005,2.374) (-9.055,2.319) (-9.104,2.262) (-9.152,2.204) (-9.199,2.146) (-9.244,2.086) (-9.289,2.026) (-9.332,1.964) (-9.374,1.902) (-9.415,1.839) (-9.454,1.775) (-9.492,1.711) (-9.529,1.646) (-9.564,1.579) (-9.598,1.513) (-9.631,1.445) (-9.662,1.377) (-9.692,1.308) (-9.720,1.239) (-9.747,1.168) (-9.772,1.098) (-9.796,1.027) (-9.818,0.955) (-9.839,0.883) (-9.857,0.810) (-9.875,0.737) (-9.891,0.664) (-9.905,0.590) (-9.917,0.516) (-9.928,0.442) (-9.937,0.368) (-9.944,0.293) (-9.950,0.218) (-9.954,0.143) (-9.957,0.068) (-9.957,-0.007) (-9.956,-0.082) (-9.954,-0.156) (-9.949,-0.231) (-9.943,-0.306) (-9.935,-0.381) (-9.926,-0.455) (-9.915,-0.529) (-9.902,-0.603) (-9.888,-0.677) (-9.872,-0.750) (-9.854,-0.823) (-9.835,-0.895) (-9.814,-0.968) (-9.792,-1.039) (-9.768,-1.110) (-9.742,-1.181) (-9.715,-1.251) (-9.687,-1.320) (-9.657,-1.389) (-9.626,-1.457) (-9.593,-1.524) (-9.558,-1.591) (-9.523,-1.657) (-9.486,-1.722) (-9.447,-1.787) (-9.408,-1.850) (-9.367,-1.913) (-9.325,-1.975) (-9.281,-2.036) (-9.236,-2.097) (-9.191,-2.156) (-9.144,-2.214) (-9.095,-2.272) (-9.046,-2.328) (-8.996,-2.384) (-8.945,-2.439) (-8.892,-2.492) (-8.839,-2.545) (-8.785,-2.597) (-8.729,-2.648) (-8.673,-2.697) (-8.616,-2.746) (-8.558,-2.794) (-8.499,-2.840) (-8.440,-2.886) (-8.379,-2.930) (-8.318,-2.974) (-8.256,-3.016) (-8.194,-3.057) (-8.131,-3.098) (-8.067,-3.137) (-8.002,-3.175) (-7.937,-3.212) (-7.871,-3.248) (-7.805,-3.283) (-7.738,-3.317) (-7.670,-3.350) (-7.602,-3.381) (-7.534,-3.412) (-7.465,-3.442) (-7.396,-3.470) (-7.326,-3.498) (-7.256,-3.524) (-7.185,-3.549) (-7.114,-3.574) (-7.043,-3.597) (-6.971,-3.619) (-6.899,-3.640) (-6.827,-3.660) (-6.754,-3.679) (-6.681,-3.697) (-6.608,-3.714) (-6.535,-3.729) (-6.462,-3.744) (-6.388,-3.758) (-6.314,-3.771) (-6.240,-3.782) (-6.166,-3.793) (-6.091,-3.803) (-6.017,-3.811) (-5.942,-3.819) (-5.867,-3.826) (-5.793,-3.831) (-5.718,-3.836) (-5.643,-3.840) (-5.568,-3.842) (-5.493,-3.844) (-5.418,-3.845) (-5.343,-3.844) (-5.268,-3.843) (-5.193,-3.841) (-5.118,-3.838) (-5.043,-3.834) (-4.968,-3.829) (-4.894,-3.823) (-4.819,-3.816) (-4.744,-3.808) (-4.670,-3.799) (-4.595,-3.789) (-4.521,-3.778) (-4.447,-3.767) (-4.373,-3.754) (-4.299,-3.741) (-4.226,-3.727) (-4.152,-3.711) (-4.079,-3.695) (-4.006,-3.678) (-3.933,-3.660) (-3.861,-3.642) (-3.788,-3.622) (-3.716,-3.601) (-3.644,-3.580) (-3.573,-3.558) (-3.501,-3.535) (-3.430,-3.511) (-3.359,-3.486) (-3.289,-3.460) (-3.219,-3.434) (-3.149,-3.406) (-3.079,-3.378) (-3.010,-3.349) (-2.941,-3.319) (-2.873,-3.289) (-2.805,-3.257) (-2.737,-3.225) (-2.670,-3.192) (-2.603,-3.158) (-2.536,-3.123) (-2.470,-3.088) (-2.405,-3.052) (-2.339,-3.015) (-2.275,-2.977) (-2.210,-2.938) (-2.147,-2.899) (-2.083,-2.859) (-2.020,-2.818) (-1.958,-2.776) (-1.896,-2.733) (-1.835,-2.690) (-1.774,-2.646) (-1.714,-2.602) (-1.654,-2.556) (-1.595,-2.510) (-1.556,-2.479)}
\LC{(-0.049,-10.000) (-0.048,-9.947) (-0.048,-9.872) (-0.047,-9.797) (-0.047,-9.722) (-0.046,-9.647) (-0.045,-9.572) (-0.045,-9.497) (-0.044,-9.422) (-0.044,-9.347) (-0.043,-9.272) (-0.043,-9.197) (-0.042,-9.122) (-0.042,-9.047) (-0.041,-8.972) (-0.041,-8.897) (-0.040,-8.822) (-0.040,-8.747) (-0.039,-8.672) (-0.039,-8.597) (-0.038,-8.522) (-0.038,-8.447) (-0.037,-8.372) (-0.037,-8.297) (-0.036,-8.222) (-0.036,-8.147) (-0.035,-8.072) (-0.035,-7.997) (-0.034,-7.922) (-0.034,-7.847) (-0.033,-7.772) (-0.033,-7.697) (-0.032,-7.622) (-0.032,-7.547) (-0.031,-7.472) (-0.031,-7.397) (-0.030,-7.322) (-0.030,-7.247) (-0.029,-7.172) (-0.029,-7.097) (-0.029,-7.022) (-0.028,-6.947) (-0.028,-6.872) (-0.027,-6.797) (-0.027,-6.722) (-0.026,-6.647) (-0.026,-6.572) (-0.025,-6.497) (-0.025,-6.422) (-0.024,-6.347) (-0.024,-6.272) (-0.024,-6.197) (-0.023,-6.122) (-0.023,-6.047) (-0.022,-5.972) (-0.022,-5.897) (-0.021,-5.822) (-0.021,-5.747) (-0.021,-5.672) (-0.020,-5.597) (-0.020,-5.522) (-0.019,-5.447) (-0.019,-5.372) (-0.019,-5.297) (-0.018,-5.222) (-0.018,-5.147) (-0.017,-5.072) (-0.017,-4.997) (-0.017,-4.922) (-0.016,-4.847) (-0.016,-4.772) (-0.016,-4.697) (-0.015,-4.622) (-0.015,-4.547) (-0.014,-4.472) (-0.014,-4.397) (-0.014,-4.322) (-0.013,-4.247) (-0.013,-4.172) (-0.013,-4.097) (-0.012,-4.022) (-0.012,-3.947) (-0.012,-3.872) (-0.011,-3.797) (-0.011,-3.722) (-0.011,-3.647) (-0.010,-3.572) (-0.010,-3.497) (-0.010,-3.422) (-0.009,-3.347) (-0.009,-3.272) (-0.009,-3.197) (-0.008,-3.122) (-0.008,-3.047) (-0.008,-2.972) (-0.007,-2.897) (-0.007,-2.822) (-0.007,-2.747) (-0.006,-2.672) (-0.006,-2.597) (-0.006,-2.522) (-0.006,-2.447) (-0.005,-2.372) (-0.005,-2.347)}
\LC{(-0.371,1.674) (-0.398,1.743) (-0.426,1.813) (-0.454,1.883) (-0.483,1.952) (-0.512,2.021) (-0.541,2.090) (-0.572,2.159) (-0.602,2.227) (-0.633,2.295) (-0.665,2.363) (-0.697,2.431) (-0.729,2.499) (-0.762,2.566) (-0.795,2.634) (-0.829,2.701) (-0.863,2.768) (-0.897,2.834) (-0.932,2.901) (-0.967,2.967) (-1.003,3.033) (-1.039,3.098) (-1.075,3.164) (-1.112,3.229) (-1.149,3.295) (-1.187,3.359) (-1.225,3.424) (-1.263,3.489) (-1.302,3.553) (-1.341,3.617) (-1.380,3.681) (-1.420,3.745) (-1.460,3.808) (-1.500,3.871) (-1.541,3.934) (-1.582,3.997) (-1.623,4.059) (-1.665,4.122) (-1.707,4.184) (-1.749,4.246) (-1.792,4.307) (-1.835,4.369) (-1.878,4.430) (-1.922,4.491) (-1.966,4.552) (-2.010,4.613) (-2.054,4.673) (-2.099,4.733) (-2.144,4.793) (-2.190,4.853) (-2.235,4.912) (-2.281,4.971) (-2.328,5.030) (-2.374,5.089) (-2.421,5.148) (-2.468,5.206) (-2.516,5.264) (-2.563,5.322) (-2.611,5.380) (-2.659,5.437) (-2.708,5.495) (-2.757,5.552) (-2.806,5.608) (-2.855,5.665) (-2.904,5.721) (-2.954,5.777) (-3.004,5.833) (-3.054,5.889) (-3.105,5.944) (-3.156,6.000) (-3.207,6.055) (-3.258,6.109) (-3.310,6.164) (-3.361,6.218) (-3.413,6.272) (-3.466,6.326) (-3.518,6.380) (-3.571,6.433) (-3.624,6.486) (-3.677,6.539) (-3.730,6.592) (-3.784,6.644) (-3.838,6.696) (-3.892,6.748) (-3.946,6.800) (-4.001,6.851) (-4.055,6.903) (-4.110,6.954) (-4.166,7.004) (-4.221,7.055) (-4.277,7.105) (-4.332,7.155) (-4.388,7.205) (-4.445,7.255) (-4.501,7.304) (-4.558,7.353) (-4.615,7.402) (-4.672,7.451) (-4.729,7.499) (-4.786,7.548) (-4.844,7.595) (-4.902,7.643) (-4.960,7.691) (-5.018,7.738) (-5.077,7.785) (-5.135,7.832) (-5.194,7.878) (-5.253,7.924) (-5.313,7.970) (-5.372,8.016) (-5.432,8.062) (-5.491,8.107) (-5.551,8.152) (-5.611,8.197) (-5.672,8.241) (-5.732,8.286) (-5.793,8.330) (-5.854,8.374) (-5.915,8.417) (-5.976,8.461) (-6.038,8.504) (-6.099,8.546) (-6.161,8.589) (-6.223,8.631) (-6.285,8.673) (-6.347,8.715) (-6.409,8.757) (-6.472,8.798) (-6.535,8.839) (-6.598,8.880) (-6.661,8.921) (-6.724,8.961) (-6.787,9.001) (-6.851,9.041) (-6.914,9.081) (-6.978,9.120) (-7.042,9.159) (-7.107,9.198) (-7.171,9.237) (-7.235,9.275) (-7.300,9.313) (-7.365,9.351) (-7.430,9.388) (-7.495,9.426) (-7.560,9.463) (-7.625,9.500) (-7.691,9.536) (-7.756,9.572) (-7.822,9.609) (-7.888,9.644) (-7.954,9.680) (-8.020,9.715) (-8.087,9.750) (-8.153,9.785) (-8.220,9.819) (-8.286,9.854) (-8.353,9.888) (-8.420,9.921) (-8.487,9.955) (-8.555,9.988) (-8.600,10.000)}
\LC{(-1.946,3.725) (-1.997,3.779) (-2.049,3.834) (-2.101,3.888) (-2.154,3.941) (-2.206,3.995) (-2.260,4.047) (-2.313,4.100) (-2.367,4.152) (-2.421,4.204) (-2.476,4.255) (-2.531,4.306) (-2.586,4.357) (-2.642,4.407) (-2.698,4.457) (-2.754,4.506) (-2.811,4.556) (-2.868,4.604) (-2.925,4.653) (-2.983,4.701) (-3.041,4.748) (-3.099,4.795) (-3.158,4.842) (-3.217,4.889) (-3.276,4.935) (-3.336,4.980) (-3.395,5.025) (-3.456,5.070) (-3.516,5.115) (-3.577,5.159) (-3.638,5.202) (-3.699,5.246) (-3.760,5.289) (-3.822,5.331) (-3.884,5.373) (-3.947,5.415) (-4.009,5.456) (-4.072,5.497) (-4.135,5.537) (-4.199,5.577) (-4.262,5.617) (-4.326,5.656) (-4.391,5.695) (-4.455,5.734) (-4.520,5.772) (-4.585,5.809) (-4.650,5.846) (-4.715,5.883) (-4.781,5.920) (-4.846,5.956) (-4.912,5.991) (-4.979,6.026) (-5.045,6.061) (-5.112,6.095) (-5.179,6.129) (-5.246,6.162) (-5.313,6.196) (-5.381,6.228) (-5.449,6.260) (-5.517,6.292) (-5.585,6.323) (-5.653,6.354) (-5.722,6.385) (-5.790,6.415) (-5.859,6.444) (-5.928,6.474) (-5.998,6.502) (-6.067,6.531) (-6.137,6.558) (-6.206,6.586) (-6.276,6.613) (-6.347,6.639) (-6.417,6.666) (-6.487,6.691) (-6.558,6.717) (-6.629,6.741) (-6.700,6.766) (-6.771,6.790) (-6.842,6.813) (-6.913,6.836) (-6.985,6.859) (-7.057,6.881) (-7.128,6.902) (-7.200,6.924) (-7.272,6.944) (-7.345,6.965) (-7.417,6.984) (-7.489,7.004) (-7.562,7.023) (-7.635,7.041) (-7.707,7.059) (-7.780,7.077) (-7.853,7.094) (-7.926,7.111) (-8.000,7.127) (-8.073,7.142) (-8.147,7.158) (-8.220,7.172) (-8.294,7.187) (-8.367,7.200) (-8.441,7.214) (-8.515,7.227) (-8.589,7.239) (-8.663,7.251) (-8.737,7.262) (-8.811,7.273) (-8.886,7.284) (-8.960,7.294) (-9.034,7.303) (-9.109,7.312) (-9.183,7.321) (-9.258,7.329) (-9.333,7.336) (-9.407,7.343) (-9.482,7.350) (-9.557,7.356) (-9.631,7.362) (-9.706,7.367) (-9.781,7.371) (-9.856,7.376) (-9.931,7.379) (-10.000,7.382)}
\LC{(-10.000,-8.148) (-9.943,-8.138) (-9.869,-8.127) (-9.795,-8.117) (-9.720,-8.105) (-9.646,-8.094) (-9.572,-8.082) (-9.498,-8.069) (-9.424,-8.056) (-9.351,-8.043) (-9.277,-8.029) (-9.203,-8.015) (-9.130,-8.000) (-9.056,-7.985) (-8.983,-7.970) (-8.910,-7.954) (-8.836,-7.938) (-8.763,-7.921) (-8.690,-7.904) (-8.617,-7.886) (-8.545,-7.868) (-8.472,-7.850) (-8.399,-7.831) (-8.327,-7.812) (-8.254,-7.792) (-8.182,-7.772) (-8.110,-7.752) (-8.038,-7.731) (-7.966,-7.710) (-7.894,-7.688) (-7.822,-7.666) (-7.751,-7.644) (-7.679,-7.621) (-7.608,-7.597) (-7.537,-7.574) (-7.466,-7.550) (-7.395,-7.525) (-7.324,-7.500) (-7.254,-7.475) (-7.183,-7.449) (-7.113,-7.423) (-7.043,-7.397) (-6.973,-7.370) (-6.903,-7.343) (-6.833,-7.315) (-6.763,-7.287) (-6.694,-7.259) (-6.625,-7.230) (-6.556,-7.201) (-6.487,-7.171) (-6.418,-7.141) (-6.349,-7.111) (-6.281,-7.080) (-6.213,-7.049) (-6.145,-7.018) (-6.077,-6.986) (-6.009,-6.954) (-5.941,-6.921) (-5.874,-6.888) (-5.807,-6.855) (-5.740,-6.821) (-5.673,-6.787) (-5.607,-6.752) (-5.540,-6.718) (-5.474,-6.682) (-5.408,-6.647) (-5.342,-6.611) (-5.277,-6.574) (-5.211,-6.538) (-5.146,-6.500) (-5.081,-6.463) (-5.016,-6.425) (-4.952,-6.387) (-4.887,-6.348) (-4.823,-6.309) (-4.760,-6.270) (-4.696,-6.230) (-4.632,-6.190) (-4.569,-6.150) (-4.506,-6.109) (-4.444,-6.068) (-4.381,-6.026) (-4.319,-5.985) (-4.257,-5.942) (-4.195,-5.900) (-4.134,-5.857) (-4.072,-5.814) (-4.011,-5.770) (-3.951,-5.726) (-3.890,-5.682) (-3.830,-5.637) (-3.770,-5.592) (-3.710,-5.547) (-3.651,-5.501) (-3.592,-5.455) (-3.533,-5.408) (-3.474,-5.362) (-3.416,-5.314) (-3.358,-5.267) (-3.300,-5.219) (-3.242,-5.171) (-3.185,-5.122) (-3.128,-5.074) (-3.072,-5.024) (-3.015,-4.975) (-2.959,-4.925) (-2.904,-4.875) (-2.848,-4.824) (-2.793,-4.773) (-2.739,-4.722) (-2.684,-4.670) (-2.630,-4.618) (-2.576,-4.566) (-2.523,-4.514) (-2.470,-4.461) (-2.417,-4.407) (-2.364,-4.354) (-2.312,-4.300) (-2.260,-4.246) (-2.209,-4.191) (-2.158,-4.136) (-2.107,-4.081) (-2.057,-4.025) (-2.007,-3.969) (-1.957,-3.913) (-1.908,-3.857) (-1.859,-3.800) (-1.811,-3.742) (-1.763,-3.685) (-1.715,-3.627) (-1.668,-3.569) (-1.621,-3.510) (-1.574,-3.451) (-1.528,-3.392) (-1.482,-3.333) (-1.437,-3.273) (-1.392,-3.213) (-1.348,-3.153) (-1.304,-3.092) (-1.260,-3.031) (-1.217,-2.969) (-1.174,-2.908) (-1.146,-2.866)}
\LC{(-5.708,3.384) (-5.781,3.371) (-5.855,3.357) (-5.928,3.341) (-6.001,3.324) (-6.074,3.306) (-6.147,3.287) (-6.219,3.267) (-6.291,3.245) (-6.362,3.223) (-6.433,3.199) (-6.504,3.174) (-6.574,3.148) (-6.644,3.120) (-6.714,3.092) (-6.782,3.062) (-6.851,3.031) (-6.918,2.999) (-6.986,2.966) (-7.052,2.931) (-7.118,2.895) (-7.183,2.858) (-7.248,2.820) (-7.312,2.781) (-7.375,2.740) (-7.437,2.699) (-7.499,2.656) (-7.560,2.612) (-7.620,2.567) (-7.679,2.521) (-7.737,2.473) (-7.794,2.424) (-7.850,2.375) (-7.905,2.324) (-7.959,2.272) (-8.012,2.219) (-8.064,2.164) (-8.114,2.109) (-8.164,2.053) (-8.212,1.995) (-8.259,1.937) (-8.305,1.877) (-8.349,1.817) (-8.392,1.756) (-8.434,1.693) (-8.474,1.630) (-8.513,1.566) (-8.551,1.501) (-8.586,1.435) (-8.621,1.368) (-8.653,1.301) (-8.685,1.233) (-8.714,1.164) (-8.742,1.094) (-8.768,1.024) (-8.793,0.953) (-8.815,0.881) (-8.836,0.809) (-8.856,0.737) (-8.873,0.664) (-8.889,0.591) (-8.902,0.517) (-8.914,0.443) (-8.925,0.368) (-8.933,0.294) (-8.939,0.219) (-8.944,0.144) (-8.947,0.069) (-8.947,-0.006) (-8.946,-0.081) (-8.943,-0.156) (-8.938,-0.230) (-8.932,-0.305) (-8.923,-0.380) (-8.913,-0.454) (-8.901,-0.528) (-8.886,-0.601) (-8.870,-0.675) (-8.853,-0.748) (-8.833,-0.820) (-8.812,-0.892) (-8.789,-0.963) (-8.764,-1.034) (-8.738,-1.104) (-8.710,-1.174) (-8.680,-1.243) (-8.649,-1.311) (-8.616,-1.378) (-8.581,-1.445) (-8.545,-1.511) (-8.507,-1.575) (-8.468,-1.640) (-8.428,-1.703) (-8.386,-1.765) (-8.343,-1.826) (-8.298,-1.886) (-8.252,-1.946) (-8.205,-2.004) (-8.157,-2.061) (-8.107,-2.117) (-8.056,-2.173) (-8.004,-2.227) (-7.951,-2.280) (-7.897,-2.331) (-7.842,-2.382) (-7.785,-2.432) (-7.728,-2.480) (-7.670,-2.528) (-7.611,-2.574) (-7.551,-2.619) (-7.490,-2.662) (-7.428,-2.705) (-7.366,-2.747) (-7.302,-2.787) (-7.238,-2.826) (-7.174,-2.864) (-7.108,-2.901) (-7.042,-2.936) (-6.976,-2.971) (-6.908,-3.004) (-6.841,-3.036) (-6.772,-3.066) (-6.703,-3.096) (-6.634,-3.124) (-6.564,-3.152) (-6.494,-3.178) (-6.423,-3.202) (-6.352,-3.226) (-6.280,-3.249) (-6.208,-3.270) (-6.136,-3.290) (-6.063,-3.309) (-5.990,-3.327) (-5.917,-3.343) (-5.844,-3.359) (-5.770,-3.373) (-5.697,-3.386)}
\LC{(10.000,-7.567) (9.949,-7.562) (9.875,-7.557) (9.800,-7.551) (9.725,-7.545) (9.650,-7.539) (9.576,-7.532) (9.501,-7.524) (9.427,-7.516) (9.352,-7.507) (9.278,-7.498) (9.203,-7.489) (9.129,-7.479) (9.055,-7.468) (8.980,-7.457) (8.906,-7.446) (8.832,-7.434) (8.758,-7.422) (8.684,-7.409) (8.610,-7.396) (8.537,-7.382) (8.463,-7.368) (8.390,-7.354) (8.316,-7.339) (8.243,-7.323) (8.169,-7.307) (8.096,-7.291) (8.023,-7.274) (7.950,-7.256) (7.877,-7.239) (7.805,-7.220) (7.732,-7.202) (7.659,-7.183) (7.587,-7.163) (7.515,-7.143) (7.443,-7.123) (7.371,-7.102) (7.299,-7.080) (7.227,-7.059) (7.155,-7.036) (7.084,-7.014) (7.012,-6.991) (6.941,-6.967) (6.870,-6.943) (6.799,-6.919) (6.728,-6.894) (6.658,-6.869) (6.587,-6.843) (6.517,-6.817) (6.447,-6.790) (6.377,-6.763) (6.307,-6.736) (6.237,-6.708) (6.168,-6.680) (6.099,-6.651) (6.030,-6.622) (5.961,-6.592) (5.892,-6.562) (5.823,-6.532) (5.755,-6.501) (5.687,-6.470) (5.619,-6.438) (5.551,-6.406) (5.483,-6.374) (5.416,-6.341) (5.349,-6.308) (5.282,-6.274) (5.215,-6.240) (5.148,-6.206) (5.082,-6.171) (5.016,-6.135) (4.950,-6.100) (4.884,-6.064) (4.818,-6.027) (4.753,-5.990) (4.688,-5.953) (4.623,-5.915) (4.559,-5.877) (4.494,-5.838) (4.430,-5.799) (4.366,-5.760) (4.303,-5.720) (4.239,-5.680) (4.176,-5.640) (4.113,-5.599) (4.051,-5.558) (3.988,-5.516) (3.947,-5.488)}
\LC{(1.522,2.462) (1.580,2.510) (1.639,2.556) (1.698,2.602) (1.758,2.647) (1.818,2.692) (1.879,2.735) (1.941,2.778) (2.003,2.821) (2.065,2.862) (2.128,2.903) (2.192,2.943) (2.256,2.982) (2.320,3.020) (2.385,3.058) (2.451,3.094) (2.516,3.130) (2.583,3.166) (2.649,3.200) (2.716,3.234) (2.784,3.267) (2.851,3.299) (2.920,3.330) (2.988,3.361) (3.057,3.390) (3.126,3.419) (3.196,3.447) (3.266,3.474) (3.336,3.501) (3.406,3.526) (3.477,3.551) (3.548,3.575) (3.620,3.598) (3.691,3.620) (3.763,3.641) (3.836,3.662) (3.908,3.681) (3.981,3.700) (4.053,3.718) (4.126,3.735) (4.200,3.751) (4.273,3.766) (4.347,3.780) (4.421,3.793) (4.495,3.806) (4.569,3.818) (4.643,3.828) (4.717,3.838) (4.792,3.847) (4.866,3.855) (4.941,3.862) (5.016,3.868) (5.091,3.873) (5.165,3.877) (5.240,3.880) (5.315,3.883) (5.390,3.884) (5.465,3.884) (5.540,3.884) (5.615,3.882) (5.690,3.880) (5.765,3.876) (5.840,3.872) (5.915,3.866) (5.990,3.860) (6.064,3.852) (6.139,3.844) (6.213,3.835) (6.287,3.824) (6.362,3.813) (6.435,3.800) (6.509,3.787) (6.583,3.772) (6.656,3.757) (6.729,3.740) (6.802,3.722) (6.875,3.704) (6.947,3.684) (7.019,3.663) (7.091,3.642) (7.163,3.619) (7.234,3.595) (7.304,3.570) (7.375,3.544) (7.445,3.517) (7.514,3.489) (7.583,3.460) (7.652,3.430) (7.720,3.399) (7.788,3.366) (7.855,3.333) (7.922,3.299) (7.988,3.263) (8.053,3.226) (8.118,3.189) (8.182,3.150) (8.246,3.110) (8.309,3.070) (8.371,3.028) (8.433,2.985) (8.494,2.941) (8.554,2.896) (8.613,2.850) (8.671,2.803) (8.729,2.755) (8.786,2.706) (8.841,2.656) (8.896,2.605) (8.950,2.553) (9.003,2.500) (9.055,2.446) (9.106,2.391) (9.156,2.335) (9.205,2.278) (9.253,2.220) (9.300,2.161) (9.345,2.102) (9.389,2.041) (9.433,1.980) (9.474,1.918) (9.515,1.855) (9.554,1.791) (9.592,1.726) (9.629,1.661) (9.665,1.595) (9.699,1.528) (9.731,1.460) (9.762,1.392) (9.792,1.323) (9.820,1.254) (9.847,1.184) (9.873,1.113) (9.896,1.042) (9.919,0.970) (9.939,0.898) (9.958,0.826) (9.976,0.753) (9.992,0.679) (10.000,0.630)}
\LC{(-3.266,2.966) (-3.339,2.982) (-3.413,2.996) (-3.486,3.010) (-3.560,3.022) (-3.634,3.034) (-3.709,3.044) (-3.783,3.053) (-3.858,3.061) (-3.932,3.067) (-4.007,3.073) (-4.082,3.077) (-4.157,3.080) (-4.232,3.082) (-4.307,3.083) (-4.382,3.082) (-4.457,3.081) (-4.532,3.078) (-4.607,3.074) (-4.682,3.068) (-4.756,3.062) (-4.831,3.054) (-4.905,3.045) (-4.980,3.034) (-5.054,3.023) (-5.128,3.010) (-5.201,2.995) (-5.275,2.980) (-5.348,2.963) (-5.420,2.945) (-5.493,2.926) (-5.565,2.905) (-5.637,2.883) (-5.708,2.859) (-5.779,2.835) (-5.849,2.809) (-5.919,2.781) (-5.988,2.753) (-6.057,2.723) (-6.125,2.691) (-6.193,2.659) (-6.259,2.625) (-6.326,2.589) (-6.391,2.552) (-6.456,2.514) (-6.519,2.475) (-6.582,2.434) (-6.645,2.392) (-6.706,2.349) (-6.766,2.304) (-6.825,2.258) (-6.884,2.211) (-6.941,2.163) (-6.997,2.113) (-7.052,2.062) (-7.105,2.009) (-7.158,1.956) (-7.209,1.901) (-7.259,1.845) (-7.308,1.788) (-7.355,1.730) (-7.401,1.670) (-7.445,1.610) (-7.487,1.548) (-7.529,1.485) (-7.568,1.422) (-7.606,1.357) (-7.642,1.291) (-7.677,1.225) (-7.710,1.157) (-7.741,1.089) (-7.770,1.020) (-7.797,0.950) (-7.822,0.879) (-7.846,0.808) (-7.867,0.736) (-7.887,0.664) (-7.904,0.591) (-7.920,0.518) (-7.933,0.444) (-7.945,0.370) (-7.954,0.295) (-7.961,0.221) (-7.967,0.146) (-7.970,0.071) (-7.971,-0.004) (-7.969,-0.079) (-7.966,-0.154) (-7.961,-0.229) (-7.953,-0.304) (-7.944,-0.378) (-7.932,-0.452) (-7.918,-0.526) (-7.903,-0.599) (-7.885,-0.672) (-7.865,-0.744) (-7.843,-0.816) (-7.820,-0.887) (-7.794,-0.958) (-7.767,-1.027) (-7.737,-1.097) (-7.706,-1.165) (-7.673,-1.232) (-7.638,-1.299) (-7.602,-1.364) (-7.564,-1.429) (-7.524,-1.492) (-7.483,-1.555) (-7.440,-1.616) (-7.395,-1.677) (-7.350,-1.736) (-7.302,-1.794) (-7.254,-1.851) (-7.204,-1.907) (-7.152,-1.962) (-7.099,-2.015) (-7.046,-2.068) (-6.991,-2.118) (-6.934,-2.168) (-6.877,-2.217) (-6.819,-2.264) (-6.759,-2.309) (-6.699,-2.354) (-6.638,-2.397) (-6.575,-2.439) (-6.512,-2.480) (-6.448,-2.519) (-6.384,-2.557) (-6.318,-2.593) (-6.252,-2.628) (-6.185,-2.662) (-6.118,-2.695) (-6.049,-2.726) (-5.981,-2.756) (-5.911,-2.784) (-5.841,-2.812) (-5.771,-2.837) (-5.700,-2.862) (-5.629,-2.885) (-5.557,-2.907) (-5.485,-2.928) (-5.412,-2.947) (-5.340,-2.965) (-5.266,-2.982) (-5.193,-2.997) (-5.119,-3.011) (-5.045,-3.024) (-4.971,-3.036) (-4.897,-3.046) (-4.823,-3.055) (-4.748,-3.063) (-4.673,-3.069) (-4.598,-3.074) (-4.524,-3.078) (-4.449,-3.081) (-4.374,-3.083) (-4.299,-3.083) (-4.224,-3.082) (-4.149,-3.080) (-4.074,-3.077) (-3.999,-3.072) (-3.924,-3.067) (-3.849,-3.060) (-3.775,-3.052) (-3.700,-3.043) (-3.626,-3.032) (-3.552,-3.021) (-3.478,-3.009) (-3.404,-2.995) (-3.331,-2.980) (-3.257,-2.964)}
\LC{(2.562,5.059) (2.612,5.115) (2.662,5.171) (2.713,5.226) (2.763,5.282) (2.814,5.337) (2.865,5.391) (2.917,5.446) (2.969,5.500) (3.021,5.554) (3.073,5.608) (3.125,5.662) (3.178,5.715) (3.231,5.768) (3.285,5.820) (3.338,5.873) (3.392,5.925) (3.446,5.977) (3.501,6.029) (3.555,6.080) (3.610,6.131) (3.665,6.182) (3.721,6.233) (3.776,6.283) (3.832,6.333) (3.888,6.383) (3.944,6.432) (4.001,6.482) (4.058,6.531) (4.115,6.579) (4.172,6.628) (4.229,6.676) (4.287,6.724) (4.345,6.772) (4.403,6.819) (4.462,6.866) (4.520,6.913) (4.579,6.960) (4.638,7.006) (4.697,7.052) (4.757,7.098) (4.816,7.143) (4.876,7.188) (4.936,7.233) (4.996,7.278) (5.057,7.322) (5.118,7.366) (5.178,7.410) (5.239,7.454) (5.301,7.497) (5.362,7.540) (5.424,7.583) (5.486,7.625) (5.548,7.667) (5.610,7.709) (5.672,7.751) (5.735,7.792) (5.798,7.833) (5.861,7.874) (5.924,7.914) (5.987,7.954) (6.051,7.994) (6.114,8.034) (6.178,8.073) (6.242,8.112) (6.307,8.151) (6.371,8.190) (6.435,8.228) (6.500,8.266) (6.565,8.303) (6.630,8.340) (6.695,8.378) (6.761,8.414) (6.826,8.451) (6.892,8.487) (6.958,8.523) (7.024,8.558) (7.090,8.594) (7.156,8.629) (7.223,8.663) (7.290,8.698) (7.356,8.732) (7.423,8.766) (7.490,8.799) (7.558,8.832) (7.625,8.865) (7.693,8.898) (7.760,8.930) (7.828,8.962) (7.896,8.994) (7.964,9.026) (8.032,9.057) (8.101,9.087) (8.169,9.118) (8.238,9.148) (8.307,9.178) (8.375,9.208) (8.444,9.237) (8.514,9.266) (8.583,9.295) (8.652,9.323) (8.722,9.352) (8.791,9.379) (8.861,9.407) (8.931,9.434) (9.001,9.461) (9.071,9.488) (9.141,9.514) (9.212,9.540) (9.282,9.566) (9.353,9.591) (9.424,9.616) (9.494,9.641) (9.565,9.665) (9.636,9.689) (9.707,9.713) (9.779,9.737) (9.850,9.760) (9.921,9.783) (9.993,9.805) (10.000,9.813)}
\LC{(1.896,1.632) (1.970,1.642) (2.045,1.649) (2.120,1.654) (2.195,1.657) (2.270,1.658) (2.345,1.656) (2.420,1.653) (2.494,1.647) (2.569,1.639) (2.643,1.628) (2.717,1.615) (2.790,1.600) (2.863,1.582) (2.936,1.562) (3.007,1.539) (3.078,1.514) (3.148,1.487) (3.216,1.457) (3.284,1.425) (3.350,1.390) (3.415,1.353) (3.479,1.313) (3.541,1.270) (3.601,1.226) (3.660,1.179) (3.716,1.129) (3.770,1.077) (3.822,1.023) (3.871,0.967) (3.918,0.908) (3.962,0.848) (4.004,0.785) (4.042,0.720) (4.077,0.654) (4.109,0.586) (4.137,0.517) (4.162,0.446) (4.184,0.374) (4.201,0.301) (4.215,0.228) (4.225,0.153) (4.231,0.079) (4.233,0.004) (4.231,-0.071) (4.226,-0.146) (4.216,-0.220) (4.203,-0.294) (4.185,-0.367) (4.165,-0.439) (4.140,-0.510) (4.112,-0.580) (4.080,-0.648) (4.045,-0.714) (4.008,-0.779) (3.967,-0.842) (3.923,-0.902) (3.876,-0.961) (3.827,-1.018) (3.775,-1.072) (3.721,-1.124) (3.665,-1.174) (3.607,-1.221) (3.547,-1.266) (3.485,-1.309) (3.422,-1.349) (3.357,-1.386) (3.291,-1.421) (3.223,-1.454) (3.154,-1.484) (3.085,-1.512) (3.014,-1.537) (2.943,-1.560) (2.870,-1.580) (2.798,-1.598) (2.724,-1.614) (2.650,-1.627) (2.576,-1.638) (2.502,-1.646) (2.427,-1.652) (2.352,-1.656) (2.277,-1.658) (2.202,-1.657) (2.127,-1.654) (2.052,-1.650) (1.978,-1.642) (1.903,-1.633) (1.878,-1.630)}
\LC{(2.741,-10.000) (2.712,-9.933) (2.682,-9.864) (2.653,-9.795) (2.624,-9.726) (2.595,-9.657) (2.566,-9.588) (2.537,-9.519) (2.508,-9.449) (2.479,-9.380) (2.451,-9.311) (2.422,-9.241) (2.394,-9.172) (2.366,-9.102) (2.338,-9.033) (2.310,-8.963) (2.282,-8.894) (2.254,-8.824) (2.226,-8.754) (2.199,-8.685) (2.171,-8.615) (2.144,-8.545) (2.117,-8.475) (2.090,-8.405) (2.063,-8.335) (2.036,-8.265) (2.009,-8.195) (1.983,-8.125) (1.956,-8.055) (1.930,-7.985) (1.903,-7.914) (1.877,-7.844) (1.851,-7.774) (1.825,-7.703) (1.799,-7.633) (1.774,-7.563) (1.748,-7.492) (1.723,-7.421) (1.697,-7.351) (1.672,-7.280) (1.647,-7.210) (1.622,-7.139) (1.597,-7.068) (1.572,-6.997) (1.548,-6.926) (1.523,-6.856) (1.499,-6.785) (1.475,-6.714) (1.451,-6.643) (1.427,-6.571) (1.403,-6.500) (1.379,-6.429) (1.356,-6.358) (1.332,-6.287) (1.309,-6.215) (1.286,-6.144) (1.263,-6.073) (1.240,-6.001) (1.217,-5.930) (1.194,-5.858) (1.172,-5.787) (1.150,-5.715) (1.127,-5.644) (1.105,-5.572) (1.083,-5.500) (1.062,-5.428) (1.040,-5.357) (1.018,-5.285) (0.997,-5.213) (0.976,-5.141) (0.955,-5.069) (0.934,-4.997) (0.913,-4.925) (0.893,-4.853) (0.872,-4.781) (0.852,-4.708) (0.832,-4.636) (0.812,-4.564) (0.792,-4.492) (0.772,-4.419) (0.753,-4.347) (0.746,-4.323)}
\LC{(-3.495,5.655) (-3.552,5.704) (-3.609,5.752) (-3.667,5.801) (-3.725,5.849) (-3.783,5.896) (-3.841,5.944) (-3.899,5.991) (-3.958,6.037) (-4.017,6.084) (-4.076,6.130) (-4.135,6.175) (-4.195,6.221) (-4.255,6.266) (-4.315,6.311) (-4.376,6.355) (-4.436,6.399) (-4.497,6.443) (-4.558,6.487) (-4.619,6.530) (-4.681,6.573) (-4.743,6.615) (-4.805,6.658) (-4.867,6.699) (-4.929,6.741) (-4.992,6.782) (-5.055,6.823) (-5.118,6.864) (-5.181,6.904) (-5.245,6.944) (-5.308,6.984) (-5.372,7.023) (-5.436,7.062) (-5.500,7.101) (-5.565,7.139) (-5.630,7.177) (-5.694,7.215) (-5.759,7.252) (-5.825,7.289) (-5.890,7.326) (-5.956,7.362) (-6.022,7.398) (-6.087,7.434) (-6.154,7.469) (-6.220,7.504) (-6.286,7.539) (-6.353,7.573) (-6.420,7.607) (-6.487,7.641) (-6.554,7.674) (-6.622,7.707) (-6.689,7.740) (-6.757,7.772) (-6.825,7.804) (-6.893,7.836) (-6.961,7.867) (-7.029,7.898) (-7.098,7.928) (-7.166,7.959) (-7.235,7.988) (-7.304,8.018) (-7.373,8.047) (-7.442,8.076) (-7.512,8.105) (-7.581,8.133) (-7.651,8.160) (-7.721,8.188) (-7.791,8.215) (-7.861,8.242) (-7.931,8.268) (-8.001,8.294) (-8.072,8.320) (-8.142,8.345) (-8.213,8.370) (-8.284,8.395) (-8.355,8.419) (-8.426,8.443) (-8.497,8.466) (-8.569,8.489) (-8.640,8.512) (-8.712,8.535) (-8.783,8.557) (-8.855,8.579) (-8.927,8.600) (-8.999,8.621) (-9.071,8.642) (-9.143,8.662) (-9.216,8.682) (-9.288,8.701) (-9.360,8.720) (-9.433,8.739) (-9.506,8.758) (-9.579,8.776) (-9.651,8.793) (-9.724,8.811) (-9.797,8.828) (-9.871,8.844) (-9.944,8.860) (-10.000,8.876)}
\LC{(8.690,2.175) (8.738,2.118) (8.785,2.059) (8.831,2.000) (8.876,1.940) (8.919,1.878) (8.961,1.816) (9.002,1.753) (9.041,1.689) (9.079,1.625) (9.115,1.559) (9.150,1.493) (9.184,1.426) (9.216,1.358) (9.247,1.289) (9.276,1.220) (9.303,1.151) (9.329,1.080) (9.354,1.009) (9.376,0.938) (9.397,0.866) (9.417,0.793) (9.435,0.720) (9.451,0.647) (9.465,0.574) (9.478,0.500) (9.489,0.425) (9.498,0.351) (9.505,0.276) (9.511,0.202) (9.515,0.127) (9.517,0.052) (9.517,-0.023) (9.516,-0.098) (9.512,-0.173) (9.507,-0.248) (9.501,-0.323) (9.492,-0.397) (9.482,-0.472) (9.470,-0.546) (9.456,-0.619) (9.441,-0.693) (9.424,-0.766) (9.405,-0.838) (9.385,-0.910) (9.363,-0.982) (9.339,-1.053) (9.313,-1.124) (9.287,-1.194) (9.258,-1.263) (9.228,-1.332) (9.196,-1.400) (9.163,-1.467) (9.129,-1.534) (9.093,-1.600) (9.056,-1.665) (9.017,-1.729) (8.977,-1.792) (8.935,-1.855) (8.892,-1.916) (8.848,-1.977) (8.803,-2.037) (8.756,-2.096) (8.709,-2.154) (8.660,-2.210) (8.610,-2.266) (8.559,-2.321) (8.507,-2.375) (8.453,-2.428) (8.399,-2.480) (8.344,-2.530) (8.288,-2.580) (8.231,-2.629) (8.173,-2.676) (8.114,-2.723) (8.054,-2.768) (7.993,-2.812) (7.932,-2.856) (7.870,-2.898) (7.807,-2.939) (7.744,-2.978) (7.679,-3.017) (7.614,-3.055) (7.549,-3.091) (7.483,-3.127) (7.416,-3.161) (7.349,-3.194) (7.281,-3.226) (7.213,-3.257) (7.144,-3.286) (7.074,-3.315) (7.005,-3.343) (6.934,-3.369) (6.864,-3.394) (6.793,-3.418) (6.721,-3.441) (6.650,-3.463) (6.578,-3.484) (6.505,-3.504) (6.433,-3.523) (6.360,-3.540) (6.287,-3.557) (6.213,-3.572) (6.140,-3.586) (6.066,-3.600) (5.992,-3.612) (5.918,-3.623) (5.843,-3.633) (5.769,-3.642) (5.694,-3.650) (5.620,-3.657) (5.545,-3.662) (5.470,-3.667) (5.395,-3.671) (5.320,-3.673) (5.245,-3.675) (5.170,-3.676) (5.095,-3.675) (5.020,-3.674) (4.945,-3.672) (4.870,-3.668) (4.795,-3.664) (4.721,-3.658) (4.646,-3.652) (4.571,-3.645) (4.497,-3.636) (4.422,-3.627) (4.348,-3.617) (4.274,-3.605) (4.200,-3.593) (4.126,-3.580) (4.052,-3.566) (3.979,-3.551) (3.906,-3.535) (3.832,-3.518) (3.760,-3.500) (3.687,-3.482) (3.615,-3.462) (3.542,-3.441) (3.471,-3.420) (3.399,-3.397) (3.328,-3.374) (3.257,-3.350) (3.186,-3.325) (3.116,-3.299) (3.046,-3.272) (2.976,-3.244) (2.907,-3.216) (2.838,-3.186) (2.769,-3.156) (2.701,-3.125) (2.633,-3.093) (2.566,-3.060) (2.499,-3.026) (2.432,-2.992) (2.366,-2.957) (2.300,-2.920) (2.235,-2.883) (2.170,-2.846) (2.106,-2.807) (2.042,-2.768) (1.979,-2.727) (1.916,-2.686) (1.854,-2.644) (1.792,-2.602) (1.731,-2.558) (1.711,-2.544)}
\LC{(9.553,-10.000) (9.483,-9.975) (9.413,-9.948) (9.343,-9.920) (9.273,-9.893) (9.203,-9.865) (9.134,-9.837) (9.064,-9.808) (8.995,-9.780) (8.926,-9.751) (8.857,-9.722) (8.788,-9.692) (8.719,-9.662) (8.651,-9.632) (8.582,-9.602) (8.514,-9.571) (8.445,-9.540) (8.377,-9.509) (8.309,-9.477) (8.241,-9.445) (8.173,-9.413) (8.106,-9.381) (8.038,-9.348) (7.971,-9.315) (7.904,-9.282) (7.836,-9.248) (7.769,-9.215) (7.703,-9.181) (7.636,-9.146) (7.569,-9.112) (7.503,-9.077) (7.437,-9.042) (7.371,-9.006) (7.305,-8.970) (7.239,-8.934) (7.173,-8.898) (7.108,-8.861) (7.043,-8.825) (6.977,-8.788) (6.912,-8.750) (6.848,-8.712) (6.783,-8.675) (6.718,-8.636) (6.654,-8.598) (6.590,-8.559) (6.526,-8.520) (6.462,-8.481) (6.398,-8.441) (6.334,-8.401) (6.271,-8.361) (6.208,-8.321) (6.145,-8.280) (6.082,-8.239) (6.019,-8.198) (5.957,-8.157) (5.894,-8.115) (5.832,-8.073) (5.770,-8.031) (5.709,-7.988) (5.647,-7.946) (5.585,-7.903) (5.524,-7.859) (5.463,-7.816) (5.402,-7.772) (5.342,-7.728) (5.281,-7.684) (5.221,-7.639) (5.161,-7.594) (5.101,-7.549) (5.041,-7.504) (4.982,-7.458) (4.922,-7.412) (4.863,-7.366) (4.804,-7.319) (4.746,-7.273) (4.687,-7.226) (4.629,-7.179) (4.571,-7.131) (4.513,-7.083) (4.455,-7.035) (4.398,-6.987) (4.341,-6.939) (4.284,-6.890) (4.227,-6.841) (4.170,-6.792) (4.114,-6.742) (4.058,-6.692) (4.002,-6.642) (3.946,-6.592) (3.891,-6.542) (3.836,-6.491) (3.781,-6.440) (3.726,-6.389) (3.671,-6.337) (3.617,-6.285) (3.563,-6.233) (3.509,-6.181) (3.456,-6.129) (3.402,-6.076) (3.349,-6.023) (3.297,-5.970) (3.244,-5.916) (3.192,-5.862) (3.140,-5.808) (3.088,-5.754) (3.036,-5.699) (2.985,-5.645) (2.934,-5.590) (2.883,-5.535) (2.833,-5.479) (2.783,-5.423) (2.733,-5.367) (2.683,-5.311) (2.634,-5.255) (2.585,-5.198) (2.536,-5.141) (2.487,-5.084) (2.439,-5.026) (2.391,-4.969) (2.343,-4.911) (2.296,-4.853) (2.249,-4.794) (2.202,-4.736) (2.156,-4.677) (2.110,-4.618) (2.064,-4.558) (2.018,-4.499) (1.973,-4.439) (1.928,-4.379) (1.883,-4.319) (1.839,-4.258) (1.795,-4.197) (1.751,-4.136) (1.708,-4.075) (1.665,-4.014) (1.622,-3.952) (1.580,-3.890) (1.538,-3.828) (1.497,-3.766) (1.483,-3.745)}
\LC{(-10.000,-5.598) (-9.942,-5.614) (-9.868,-5.630) (-9.795,-5.644) (-9.721,-5.658) (-9.647,-5.672) (-9.574,-5.685) (-9.499,-5.697) (-9.425,-5.708) (-9.351,-5.719) (-9.277,-5.729) (-9.202,-5.738) (-9.128,-5.747) (-9.053,-5.755) (-8.979,-5.762) (-8.904,-5.769) (-8.829,-5.775) (-8.754,-5.780) (-8.680,-5.785) (-8.605,-5.789) (-8.530,-5.792) (-8.455,-5.795) (-8.380,-5.797) (-8.305,-5.799) (-8.230,-5.799) (-8.155,-5.799) (-8.080,-5.799) (-8.005,-5.798) (-7.930,-5.796) (-7.855,-5.794) (-7.780,-5.790) (-7.705,-5.787) (-7.630,-5.782) (-7.555,-5.777) (-7.481,-5.772) (-7.406,-5.766) (-7.331,-5.759) (-7.257,-5.751) (-7.182,-5.743) (-7.107,-5.734) (-7.033,-5.725) (-6.959,-5.715) (-6.884,-5.705) (-6.810,-5.693) (-6.736,-5.682) (-6.662,-5.669) (-6.588,-5.656) (-6.515,-5.643) (-6.441,-5.628) (-6.367,-5.614) (-6.294,-5.598) (-6.221,-5.582) (-6.148,-5.566) (-6.075,-5.548) (-6.002,-5.531) (-5.929,-5.512) (-5.857,-5.493) (-5.784,-5.474) (-5.712,-5.454) (-5.640,-5.433) (-5.568,-5.412) (-5.496,-5.390) (-5.425,-5.367) (-5.353,-5.344) (-5.282,-5.321) (-5.211,-5.297) (-5.140,-5.272) (-5.069,-5.247) (-4.999,-5.221) (-4.929,-5.195) (-4.859,-5.168) (-4.789,-5.140) (-4.720,-5.112) (-4.650,-5.084) (-4.581,-5.054) (-4.512,-5.025) (-4.444,-4.994) (-4.375,-4.964) (-4.307,-4.932) (-4.239,-4.900) (-4.172,-4.868) (-4.104,-4.835) (-4.037,-4.802) (-3.970,-4.768) (-3.904,-4.733) (-3.837,-4.698) (-3.771,-4.662) (-3.706,-4.626) (-3.640,-4.589) (-3.575,-4.552) (-3.510,-4.514) (-3.446,-4.476) (-3.382,-4.437) (-3.318,-4.398) (-3.254,-4.358) (-3.191,-4.318) (-3.128,-4.277) (-3.065,-4.236) (-3.003,-4.194) (-2.941,-4.152) (-2.880,-4.109) (-2.819,-4.065) (-2.758,-4.022) (-2.697,-3.977) (-2.637,-3.932) (-2.577,-3.887) (-2.518,-3.841) (-2.459,-3.795)}
\LC{(10.000,-4.870) (9.938,-4.895) (9.867,-4.920) (9.796,-4.943) (9.724,-4.966) (9.653,-4.988) (9.581,-5.010) (9.509,-5.030) (9.436,-5.050) (9.364,-5.069) (9.291,-5.087) (9.218,-5.105) (9.145,-5.122) (9.072,-5.138) (8.998,-5.153) (8.925,-5.168) (8.851,-5.181) (8.777,-5.194) (8.703,-5.207) (8.629,-5.218) (8.555,-5.229) (8.480,-5.239) (8.406,-5.248) (8.332,-5.257) (8.257,-5.265) (8.182,-5.272) (8.108,-5.278) (8.033,-5.284) (7.958,-5.289) (7.883,-5.293) (7.808,-5.297) (7.733,-5.300) (7.658,-5.302) (7.583,-5.303) (7.508,-5.304) (7.433,-5.304) (7.358,-5.303) (7.283,-5.302) (7.208,-5.299) (7.133,-5.297) (7.058,-5.293) (6.984,-5.289) (6.909,-5.284) (6.834,-5.278) (6.759,-5.272) (6.684,-5.265) (6.610,-5.257) (6.535,-5.249) (6.461,-5.240) (6.387,-5.230) (6.312,-5.220) (6.238,-5.209) (6.164,-5.197) (6.090,-5.185) (6.016,-5.172) (5.942,-5.158) (5.869,-5.144) (5.795,-5.129) (5.722,-5.113) (5.649,-5.097) (5.576,-5.080) (5.503,-5.062) (5.430,-5.044) (5.357,-5.025) (5.285,-5.006) (5.213,-4.986) (5.141,-4.965) (5.069,-4.944) (4.997,-4.922) (4.926,-4.899) (4.854,-4.876) (4.783,-4.852) (4.712,-4.827) (4.642,-4.802) (4.571,-4.776) (4.501,-4.750) (4.431,-4.723)}
\LC{(-0.451,2.963) (-0.469,3.036) (-0.487,3.109) (-0.505,3.181) (-0.523,3.254) (-0.541,3.327) (-0.560,3.400) (-0.579,3.472) (-0.598,3.545) (-0.617,3.617) (-0.636,3.690) (-0.656,3.762) (-0.676,3.834) (-0.696,3.907) (-0.716,3.979) (-0.736,4.051) (-0.757,4.123) (-0.777,4.195) (-0.798,4.267) (-0.819,4.339) (-0.840,4.411) (-0.862,4.483) (-0.883,4.555) (-0.905,4.627) (-0.927,4.699) (-0.949,4.770) (-0.971,4.842) (-0.993,4.914) (-1.016,4.985) (-1.039,5.057) (-1.062,5.128) (-1.085,5.199) (-1.108,5.271) (-1.131,5.342) (-1.155,5.413) (-1.178,5.484) (-1.202,5.555) (-1.226,5.627) (-1.250,5.698) (-1.274,5.769) (-1.299,5.839) (-1.323,5.910) (-1.348,5.981) (-1.373,6.052) (-1.398,6.123) (-1.423,6.193) (-1.448,6.264) (-1.474,6.334) (-1.499,6.405) (-1.525,6.475) (-1.551,6.546) (-1.577,6.616) (-1.603,6.686) (-1.629,6.757) (-1.656,6.827) (-1.682,6.897) (-1.709,6.967) (-1.736,7.037) (-1.763,7.107) (-1.790,7.177) (-1.817,7.247) (-1.845,7.317) (-1.872,7.387) (-1.900,7.456) (-1.928,7.526) (-1.955,7.596) (-1.983,7.665) (-2.012,7.735) (-2.040,7.804) (-2.068,7.873) (-2.097,7.943) (-2.125,8.012) (-2.154,8.081) (-2.183,8.151) (-2.212,8.220) (-2.241,8.289) (-2.271,8.358) (-2.300,8.427) (-2.330,8.496) (-2.359,8.565) (-2.389,8.634) (-2.419,8.702) (-2.449,8.771) (-2.479,8.840) (-2.509,8.908) (-2.540,8.977) (-2.570,9.046) (-2.601,9.114) (-2.631,9.182) (-2.662,9.251) (-2.693,9.319) (-2.724,9.387) (-2.755,9.456) (-2.787,9.524) (-2.818,9.592) (-2.850,9.660) (-2.881,9.728) (-2.913,9.796) (-2.945,9.864) (-2.977,9.932) (-3.009,9.999) (-3.020,10.000)}
\LC{(-10.000,-9.565) (-9.952,-9.544) (-9.880,-9.523) (-9.808,-9.502) (-9.736,-9.480) (-9.664,-9.458) (-9.593,-9.436) (-9.521,-9.413) (-9.450,-9.390) (-9.379,-9.367) (-9.308,-9.343) (-9.237,-9.319) (-9.166,-9.295) (-9.095,-9.270) (-9.024,-9.245) (-8.953,-9.220) (-8.883,-9.194) (-8.813,-9.168) (-8.742,-9.142) (-8.672,-9.115) (-8.602,-9.088) (-8.532,-9.061) (-8.463,-9.034) (-8.393,-9.006) (-8.323,-8.978) (-8.254,-8.949) (-8.185,-8.920) (-8.116,-8.891) (-8.047,-8.862) (-7.978,-8.832) (-7.909,-8.802) (-7.841,-8.772) (-7.772,-8.741) (-7.704,-8.710) (-7.636,-8.679) (-7.568,-8.647) (-7.500,-8.615) (-7.432,-8.583) (-7.364,-8.550) (-7.297,-8.517) (-7.230,-8.484) (-7.163,-8.451) (-7.096,-8.417) (-7.029,-8.383) (-6.962,-8.349) (-6.896,-8.314) (-6.829,-8.279) (-6.763,-8.244) (-6.697,-8.208) (-6.631,-8.172) (-6.566,-8.136) (-6.500,-8.100) (-6.435,-8.063) (-6.370,-8.026) (-6.305,-7.988) (-6.240,-7.951) (-6.175,-7.913) (-6.111,-7.874) (-6.046,-7.836) (-5.982,-7.797) (-5.918,-7.758) (-5.854,-7.718) (-5.791,-7.678) (-5.727,-7.638) (-5.664,-7.598) (-5.601,-7.557) (-5.538,-7.516) (-5.476,-7.475) (-5.413,-7.434) (-5.351,-7.392) (-5.289,-7.350) (-5.227,-7.307) (-5.165,-7.265) (-5.104,-7.222) (-5.043,-7.178) (-4.982,-7.135) (-4.921,-7.091) (-4.860,-7.047) (-4.800,-7.002) (-4.740,-6.958) (-4.680,-6.913) (-4.620,-6.867) (-4.560,-6.822) (-4.501,-6.776) (-4.442,-6.730) (-4.383,-6.683) (-4.324,-6.637) (-4.266,-6.590) (-4.207,-6.542) (-4.149,-6.495) (-4.092,-6.447) (-4.034,-6.399) (-3.977,-6.350) (-3.920,-6.302) (-3.863,-6.253) (-3.806,-6.204) (-3.750,-6.154) (-3.694,-6.104) (-3.657,-6.071)}
\LC{(-2.932,-10.000) (-2.901,-9.953) (-2.869,-9.885) (-2.838,-9.817) (-2.807,-9.749) (-2.777,-9.680) (-2.746,-9.612) (-2.715,-9.543) (-2.685,-9.475) (-2.654,-9.406) (-2.624,-9.338) (-2.594,-9.269) (-2.564,-9.200) (-2.534,-9.132) (-2.504,-9.063) (-2.474,-8.994) (-2.445,-8.925) (-2.415,-8.856) (-2.386,-8.787) (-2.357,-8.718) (-2.327,-8.649) (-2.298,-8.580) (-2.269,-8.511) (-2.241,-8.441) (-2.212,-8.372) (-2.183,-8.303) (-2.155,-8.233) (-2.127,-8.164) (-2.099,-8.094) (-2.070,-8.025) (-2.043,-7.955) (-2.015,-7.885) (-1.987,-7.816) (-1.959,-7.746) (-1.932,-7.676) (-1.905,-7.606) (-1.878,-7.536) (-1.850,-7.467) (-1.824,-7.397) (-1.797,-7.326) (-1.770,-7.256) (-1.744,-7.186) (-1.717,-7.116) (-1.691,-7.046) (-1.665,-6.976) (-1.639,-6.905) (-1.613,-6.835) (-1.587,-6.764) (-1.561,-6.694) (-1.536,-6.623) (-1.511,-6.553) (-1.485,-6.482) (-1.460,-6.411) (-1.435,-6.341) (-1.411,-6.270) (-1.386,-6.199) (-1.361,-6.128) (-1.337,-6.057) (-1.313,-5.986) (-1.289,-5.915) (-1.265,-5.844) (-1.241,-5.773) (-1.217,-5.702) (-1.194,-5.631) (-1.170,-5.559) (-1.147,-5.488) (-1.124,-5.417) (-1.101,-5.345) (-1.079,-5.274) (-1.056,-5.202) (-1.034,-5.131) (-1.011,-5.059) (-0.989,-4.987) (-0.967,-4.916) (-0.945,-4.844) (-0.924,-4.772) (-0.902,-4.700) (-0.881,-4.628) (-0.860,-4.556) (-0.839,-4.484) (-0.818,-4.412) (-0.811,-4.388)}
\LC{(-0.002,3.003) (-0.002,3.078) (-0.002,3.153) (-0.002,3.228) (-0.003,3.303) (-0.003,3.378) (-0.003,3.453) (-0.003,3.528) (-0.003,3.603) (-0.003,3.678) (-0.003,3.753) (-0.003,3.828) (-0.003,3.903) (-0.003,3.978) (-0.003,4.053) (-0.004,4.128) (-0.004,4.203) (-0.004,4.278) (-0.004,4.353) (-0.004,4.428) (-0.004,4.503) (-0.004,4.578) (-0.004,4.653) (-0.004,4.728) (-0.004,4.803) (-0.005,4.878) (-0.005,4.953) (-0.005,5.028) (-0.005,5.103) (-0.005,5.178) (-0.005,5.253) (-0.005,5.328) (-0.005,5.403) (-0.005,5.478) (-0.006,5.553) (-0.006,5.628) (-0.006,5.703) (-0.006,5.778) (-0.006,5.853) (-0.006,5.928) (-0.006,6.003) (-0.006,6.078) (-0.007,6.153) (-0.007,6.228) (-0.007,6.303) (-0.007,6.378) (-0.007,6.453) (-0.007,6.528) (-0.007,6.603) (-0.007,6.678) (-0.007,6.753) (-0.008,6.828) (-0.008,6.903) (-0.008,6.978) (-0.008,7.053) (-0.008,7.128) (-0.008,7.203) (-0.008,7.278) (-0.009,7.353) (-0.009,7.428) (-0.009,7.503) (-0.009,7.578) (-0.009,7.653) (-0.009,7.728) (-0.009,7.803) (-0.009,7.878) (-0.010,7.953) (-0.010,8.028) (-0.010,8.103) (-0.010,8.178) (-0.010,8.253) (-0.010,8.328) (-0.010,8.403) (-0.011,8.478) (-0.011,8.553) (-0.011,8.628) (-0.011,8.703) (-0.011,8.778) (-0.011,8.853) (-0.011,8.928) (-0.012,9.003) (-0.012,9.078) (-0.012,9.153) (-0.012,9.228) (-0.012,9.303) (-0.012,9.378) (-0.012,9.453) (-0.013,9.528) (-0.013,9.603) (-0.013,9.678) (-0.013,9.753) (-0.013,9.828) (-0.013,9.903) (-0.013,9.978) (-0.014,10.000)}
\LC{(4.652,4.285) (4.725,4.303) (4.798,4.319) (4.871,4.335) (4.945,4.350) (5.018,4.365) (5.092,4.378) (5.166,4.391) (5.240,4.403) (5.314,4.415) (5.388,4.425) (5.463,4.435) (5.537,4.444) (5.612,4.452) (5.686,4.459) (5.761,4.466) (5.836,4.471) (5.911,4.476) (5.986,4.480) (6.061,4.483) (6.135,4.486) (6.210,4.487) (6.285,4.488) (6.360,4.488) (6.435,4.487) (6.510,4.485) (6.585,4.483) (6.660,4.479) (6.735,4.475) (6.810,4.470) (6.885,4.464) (6.959,4.457) (7.034,4.449) (7.109,4.441) (7.183,4.431) (7.257,4.421) (7.331,4.410) (7.405,4.398) (7.479,4.385) (7.553,4.371) (7.627,4.356) (7.700,4.341) (7.773,4.324) (7.846,4.307) (7.919,4.289) (7.991,4.269) (8.064,4.249) (8.136,4.228) (8.207,4.207) (8.279,4.184) (8.350,4.160) (8.421,4.135) (8.491,4.110) (8.561,4.084) (8.631,4.056) (8.701,4.028) (8.770,3.999) (8.839,3.969) (8.907,3.938) (8.975,3.906) (9.042,3.873) (9.109,3.839) (9.176,3.804) (9.242,3.769) (9.307,3.732) (9.372,3.695) (9.437,3.656) (9.500,3.617) (9.564,3.577) (9.627,3.536) (9.689,3.494) (9.750,3.451) (9.811,3.407) (9.871,3.362) (9.931,3.317) (9.990,3.270) (10.000,3.255)}
\LC{(-3.065,3.859) (-3.131,3.895) (-3.197,3.931) (-3.263,3.966) (-3.330,4.001) (-3.396,4.034) (-3.464,4.068) (-3.531,4.100) (-3.599,4.132) (-3.667,4.163) (-3.736,4.194) (-3.805,4.224) (-3.874,4.253) (-3.943,4.282) (-4.013,4.310) (-4.082,4.337) (-4.152,4.364) (-4.223,4.390) (-4.293,4.415) (-4.364,4.440) (-4.435,4.464) (-4.507,4.487) (-4.578,4.509) (-4.650,4.531) (-4.722,4.553) (-4.794,4.573) (-4.866,4.593) (-4.939,4.612) (-5.011,4.631) (-5.084,4.648) (-5.157,4.665) (-5.231,4.682) (-5.304,4.697) (-5.377,4.712) (-5.451,4.727) (-5.525,4.740) (-5.599,4.753) (-5.673,4.765) (-5.747,4.776) (-5.821,4.787) (-5.895,4.797) (-5.970,4.806) (-6.044,4.814) (-6.119,4.822) (-6.194,4.829) (-6.268,4.835) (-6.343,4.841) (-6.418,4.845) (-6.493,4.849) (-6.568,4.852) (-6.643,4.855) (-6.718,4.856) (-6.793,4.857) (-6.868,4.857) (-6.943,4.857) (-7.018,4.855) (-7.093,4.853) (-7.168,4.850) (-7.243,4.846) (-7.318,4.842) (-7.392,4.836) (-7.467,4.830) (-7.542,4.823) (-7.616,4.816) (-7.691,4.807) (-7.765,4.798) (-7.840,4.788) (-7.914,4.777) (-7.988,4.765) (-8.062,4.752) (-8.136,4.739) (-8.209,4.725) (-8.283,4.710) (-8.356,4.694) (-8.429,4.677) (-8.502,4.660) (-8.575,4.641) (-8.647,4.622) (-8.720,4.602) (-8.792,4.582) (-8.864,4.560) (-8.935,4.537) (-9.006,4.514) (-9.077,4.490) (-9.148,4.465) (-9.218,4.439) (-9.289,4.412) (-9.358,4.385) (-9.428,4.356) (-9.497,4.327) (-9.566,4.297) (-9.634,4.266) (-9.702,4.234) (-9.769,4.202) (-9.836,4.168) (-9.903,4.134) (-9.969,4.099) (-10.000,4.075)}
\LC{(2.988,4.855) (3.044,4.904) (3.102,4.953) (3.159,5.001) (3.217,5.049) (3.275,5.096) (3.333,5.143) (3.392,5.190) (3.451,5.237) (3.510,5.283) (3.570,5.328) (3.629,5.374) (3.689,5.418) (3.750,5.463) (3.810,5.507) (3.871,5.551) (3.932,5.594) (3.994,5.638) (4.055,5.680) (4.117,5.723) (4.179,5.765) (4.242,5.806) (4.305,5.847) (4.367,5.888) (4.431,5.929) (4.494,5.969) (4.558,6.008) (4.622,6.048) (4.686,6.087) (4.750,6.125) (4.814,6.163) (4.879,6.201) (4.944,6.239) (5.010,6.276) (5.075,6.312) (5.141,6.349) (5.206,6.384) (5.273,6.420) (5.339,6.455) (5.405,6.490) (5.472,6.524) (5.539,6.558) (5.606,6.591) (5.673,6.624) (5.741,6.657) (5.808,6.690) (5.876,6.721) (5.944,6.753) (6.013,6.784) (6.081,6.815) (6.150,6.845) (6.218,6.875) (6.287,6.905) (6.356,6.934) (6.426,6.963) (6.495,6.991) (6.565,7.019) (6.635,7.046) (6.704,7.073) (6.775,7.100) (6.845,7.126) (6.915,7.152) (6.986,7.178) (7.056,7.203) (7.127,7.227) (7.198,7.252) (7.269,7.275) (7.341,7.299) (7.412,7.322) (7.484,7.344) (7.555,7.366) (7.627,7.388) (7.699,7.409) (7.771,7.430) (7.843,7.451) (7.915,7.471) (7.988,7.490) (8.060,7.509) (8.133,7.528) (8.206,7.546) (8.279,7.564) (8.352,7.582) (8.425,7.599) (8.498,7.615) (8.571,7.631) (8.644,7.647) (8.718,7.662) (8.791,7.677) (8.865,7.691) (8.939,7.705) (9.012,7.719) (9.086,7.732) (9.160,7.744) (9.234,7.757) (9.308,7.768) (9.382,7.780) (9.457,7.790) (9.531,7.801) (9.605,7.811) (9.680,7.820) (9.754,7.829) (9.829,7.838) (9.903,7.846) (9.978,7.853) (10.000,7.856)}
\LC{(-1.438,1.372) (-1.512,1.387) (-1.586,1.399) (-1.660,1.408) (-1.735,1.415) (-1.810,1.420) (-1.885,1.422) (-1.960,1.422) (-2.035,1.419) (-2.110,1.413) (-2.184,1.405) (-2.258,1.394) (-2.332,1.380) (-2.405,1.364) (-2.478,1.345) (-2.549,1.323) (-2.620,1.298) (-2.690,1.270) (-2.758,1.239) (-2.825,1.206) (-2.891,1.169) (-2.955,1.130) (-3.017,1.088) (-3.077,1.043) (-3.134,0.995) (-3.190,0.945) (-3.243,0.891) (-3.293,0.835) (-3.340,0.777) (-3.383,0.716) (-3.424,0.653) (-3.461,0.588) (-3.494,0.520) (-3.523,0.451) (-3.548,0.381) (-3.569,0.309) (-3.586,0.236) (-3.598,0.162) (-3.605,0.087) (-3.608,0.012) (-3.607,-0.063) (-3.601,-0.138) (-3.590,-0.212) (-3.575,-0.285) (-3.555,-0.358) (-3.532,-0.429) (-3.504,-0.498) (-3.472,-0.566) (-3.436,-0.632) (-3.397,-0.696) (-3.354,-0.758) (-3.308,-0.817) (-3.259,-0.874) (-3.207,-0.928) (-3.153,-0.979) (-3.095,-1.028) (-3.036,-1.074) (-2.975,-1.117) (-2.911,-1.157) (-2.846,-1.194) (-2.780,-1.229) (-2.712,-1.260) (-2.643,-1.289) (-2.572,-1.315) (-2.501,-1.338) (-2.429,-1.358) (-2.356,-1.375) (-2.282,-1.390) (-2.208,-1.402) (-2.133,-1.411) (-2.059,-1.417) (-1.984,-1.421) (-1.909,-1.422) (-1.834,-1.421) (-1.759,-1.417) (-1.684,-1.411) (-1.610,-1.402) (-1.536,-1.391) (-1.462,-1.377) (-1.437,-1.372)}
\LC{(-2.727,1.915) (-2.802,1.912) (-2.876,1.906) (-2.951,1.899) (-3.025,1.890) (-3.100,1.878) (-3.173,1.865) (-3.247,1.850) (-3.320,1.832) (-3.392,1.813) (-3.464,1.791) (-3.535,1.767) (-3.605,1.741) (-3.675,1.713) (-3.744,1.683) (-3.811,1.650) (-3.878,1.616) (-3.943,1.579) (-4.008,1.541) (-4.070,1.500) (-4.132,1.457) (-4.192,1.411) (-4.250,1.364) (-4.306,1.315) (-4.361,1.264) (-4.414,1.210) (-4.465,1.155) (-4.513,1.098) (-4.559,1.039) (-4.603,0.978) (-4.644,0.915) (-4.683,0.851) (-4.719,0.785) (-4.753,0.718) (-4.783,0.650) (-4.811,0.580) (-4.835,0.509) (-4.856,0.437) (-4.874,0.364) (-4.889,0.291) (-4.900,0.217) (-4.909,0.142) (-4.913,0.067) (-4.915,-0.008) (-4.913,-0.083) (-4.907,-0.158) (-4.898,-0.232) (-4.886,-0.306) (-4.871,-0.379) (-4.852,-0.452) (-4.830,-0.524) (-4.805,-0.594) (-4.777,-0.664) (-4.746,-0.732) (-4.712,-0.799) (-4.675,-0.865) (-4.636,-0.928) (-4.594,-0.991) (-4.550,-1.051) (-4.503,-1.110) (-4.454,-1.167) (-4.403,-1.221) (-4.350,-1.274) (-4.295,-1.325) (-4.238,-1.374) (-4.179,-1.421) (-4.119,-1.466) (-4.057,-1.508) (-3.994,-1.549) (-3.930,-1.587) (-3.864,-1.623) (-3.797,-1.657) (-3.730,-1.689) (-3.661,-1.719) (-3.591,-1.747) (-3.520,-1.772) (-3.449,-1.795) (-3.377,-1.817) (-3.305,-1.836) (-3.232,-1.853) (-3.158,-1.868) (-3.084,-1.881) (-3.010,-1.892) (-2.936,-1.901) (-2.861,-1.907) (-2.786,-1.912) (-2.736,-1.915)}
\LC{(-4.344,5.136) (-4.411,5.170) (-4.479,5.202) (-4.547,5.235) (-4.615,5.266) (-4.683,5.298) (-4.751,5.328) (-4.820,5.358) (-4.889,5.388) (-4.958,5.417) (-5.027,5.446) (-5.097,5.474) (-5.166,5.502) (-5.236,5.529) (-5.307,5.555) (-5.377,5.581) (-5.447,5.607) (-5.518,5.632) (-5.589,5.656) (-5.660,5.680) (-5.731,5.703) (-5.803,5.726) (-5.874,5.749) (-5.946,5.770) (-6.018,5.792) (-6.090,5.812) (-6.162,5.832) (-6.235,5.852) (-6.307,5.871) (-6.380,5.890) (-6.453,5.907) (-6.526,5.925) (-6.599,5.942) (-6.672,5.958) (-6.745,5.974) (-6.819,5.989) (-6.893,6.003) (-6.966,6.018) (-7.040,6.031) (-7.114,6.044) (-7.188,6.056) (-7.262,6.068) (-7.336,6.079) (-7.410,6.090) (-7.485,6.100) (-7.559,6.109) (-7.633,6.118) (-7.708,6.127) (-7.783,6.134) (-7.857,6.141) (-7.932,6.148) (-8.007,6.154) (-8.082,6.159) (-8.156,6.164) (-8.231,6.168) (-8.306,6.172) (-8.381,6.175) (-8.456,6.177) (-8.531,6.179) (-8.606,6.180) (-8.681,6.181) (-8.756,6.181) (-8.831,6.180) (-8.906,6.179) (-8.981,6.177) (-9.056,6.175) (-9.131,6.172) (-9.206,6.168) (-9.281,6.164) (-9.356,6.159) (-9.430,6.154) (-9.505,6.147) (-9.580,6.141) (-9.654,6.133) (-9.729,6.125) (-9.804,6.117) (-9.878,6.107) (-9.952,6.097) (-10.000,6.090)}
\LC{(2.341,3.550) (2.401,3.596) (2.461,3.641) (2.521,3.685) (2.582,3.729) (2.643,3.772) (2.705,3.815) (2.767,3.857) (2.829,3.899) (2.892,3.940) (2.955,3.980) (3.019,4.020) (3.082,4.060) (3.146,4.098) (3.211,4.137) (3.276,4.174) (3.341,4.212) (3.406,4.248) (3.472,4.284) (3.538,4.320) (3.605,4.354) (3.672,4.389) (3.739,4.422) (3.806,4.455) (3.873,4.488) (3.941,4.520) (4.009,4.551) (4.078,4.582) (4.147,4.612) (4.216,4.641) (4.285,4.670) (4.354,4.699) (4.424,4.726) (4.494,4.753) (4.564,4.780) (4.634,4.806) (4.705,4.831) (4.776,4.856) (4.847,4.880) (4.918,4.903) (4.990,4.926) (5.061,4.948) (5.133,4.970) (5.205,4.991) (5.277,5.011) (5.350,5.031) (5.422,5.050) (5.495,5.068) (5.568,5.086) (5.641,5.103) (5.714,5.119) (5.787,5.135) (5.861,5.150) (5.934,5.165) (6.008,5.179) (6.082,5.192) (6.156,5.205) (6.230,5.216) (6.304,5.228) (6.378,5.238) (6.453,5.248) (6.527,5.257) (6.602,5.266) (6.676,5.274) (6.751,5.281) (6.825,5.287) (6.900,5.293) (6.975,5.298) (7.050,5.303) (7.125,5.306) (7.200,5.310) (7.275,5.312) (7.350,5.314) (7.425,5.315) (7.500,5.315) (7.575,5.314) (7.650,5.313) (7.725,5.311) (7.800,5.309) (7.875,5.306) (7.949,5.302) (8.024,5.297) (8.099,5.291) (8.174,5.285) (8.249,5.278) (8.323,5.271) (8.398,5.262) (8.472,5.253) (8.546,5.243) (8.621,5.233) (8.695,5.222) (8.769,5.210) (8.843,5.197) (8.917,5.183) (8.990,5.169) (9.064,5.154) (9.137,5.138) (9.210,5.122) (9.283,5.105) (9.356,5.086) (9.429,5.068) (9.501,5.048) (9.573,5.028) (9.645,5.007) (9.717,4.985) (9.788,4.962) (9.860,4.939) (9.931,4.915) (10.000,4.890)}
\LC{(1.360,4.868) (1.391,4.936) (1.421,5.005) (1.452,5.073) (1.483,5.141) (1.515,5.209) (1.546,5.277) (1.578,5.345) (1.610,5.413) (1.642,5.481) (1.674,5.549) (1.707,5.616) (1.740,5.684) (1.772,5.751) (1.806,5.818) (1.839,5.886) (1.872,5.953) (1.906,6.020) (1.940,6.087) (1.974,6.153) (2.008,6.220) (2.043,6.287) (2.077,6.353) (2.112,6.420) (2.147,6.486) (2.182,6.552) (2.218,6.618) (2.253,6.684) (2.289,6.750) (2.325,6.816) (2.361,6.882) (2.397,6.948) (2.434,7.013) (2.470,7.079) (2.507,7.144) (2.544,7.209) (2.581,7.274) (2.619,7.340) (2.656,7.404) (2.694,7.469) (2.731,7.534) (2.769,7.599) (2.808,7.663) (2.846,7.728) (2.884,7.792) (2.923,7.857) (2.962,7.921) (3.001,7.985) (3.040,8.049) (3.079,8.113) (3.119,8.176) (3.158,8.240) (3.198,8.304) (3.238,8.367) (3.278,8.431) (3.318,8.494) (3.359,8.557) (3.399,8.620) (3.440,8.683) (3.481,8.746) (3.522,8.809) (3.563,8.872) (3.604,8.934) (3.646,8.997) (3.687,9.059) (3.729,9.121) (3.771,9.184) (3.813,9.246) (3.855,9.308) (3.897,9.370) (3.940,9.431) (3.983,9.493) (4.025,9.555) (4.068,9.616) (4.111,9.678) (4.154,9.739) (4.198,9.800) (4.241,9.861) (4.285,9.922) (4.329,9.983) (4.343,10.000)}
\LC{(-1.409,4.730) (-1.442,4.797) (-1.474,4.865) (-1.507,4.932) (-1.540,5.000) (-1.573,5.067) (-1.606,5.134) (-1.640,5.201) (-1.674,5.268) (-1.708,5.335) (-1.742,5.402) (-1.777,5.468) (-1.812,5.535) (-1.846,5.601) (-1.882,5.667) (-1.917,5.733) (-1.953,5.799) (-1.988,5.865) (-2.024,5.931) (-2.060,5.997) (-2.097,6.062) (-2.133,6.128) (-2.170,6.193) (-2.207,6.258) (-2.244,6.324) (-2.282,6.389) (-2.319,6.454) (-2.357,6.518) (-2.395,6.583) (-2.433,6.648) (-2.472,6.712) (-2.510,6.776) (-2.549,6.841) (-2.588,6.905) (-2.627,6.969) (-2.666,7.033) (-2.706,7.096) (-2.745,7.160) (-2.785,7.224) (-2.825,7.287) (-2.865,7.350) (-2.906,7.414) (-2.946,7.477) (-2.987,7.540) (-3.028,7.603) (-3.069,7.665) (-3.110,7.728) (-3.152,7.790) (-3.193,7.853) (-3.235,7.915) (-3.277,7.977) (-3.319,8.039) (-3.361,8.101) (-3.404,8.163) (-3.446,8.225) (-3.489,8.287) (-3.532,8.348) (-3.575,8.409) (-3.618,8.471) (-3.662,8.532) (-3.705,8.593) (-3.749,8.654) (-3.793,8.715) (-3.837,8.775) (-3.881,8.836) (-3.926,8.896) (-3.970,8.957) (-4.015,9.017) (-4.060,9.077) (-4.105,9.137) (-4.150,9.197) (-4.195,9.256) (-4.241,9.316) (-4.287,9.376) (-4.332,9.435) (-4.378,9.494) (-4.424,9.553) (-4.471,9.613) (-4.517,9.671) (-4.564,9.730) (-4.610,9.789) (-4.657,9.848) (-4.704,9.906) (-4.751,9.964) (-4.783,10.000)}
\LC{(-10.000,-3.295) (-9.955,-3.342) (-9.896,-3.387) (-9.835,-3.432) (-9.774,-3.475) (-9.713,-3.518) (-9.650,-3.560) (-9.588,-3.601) (-9.524,-3.641) (-9.460,-3.680) (-9.396,-3.718) (-9.330,-3.755) (-9.265,-3.791) (-9.199,-3.827) (-9.132,-3.861) (-9.065,-3.895) (-8.998,-3.928) (-8.930,-3.959) (-8.861,-3.990) (-8.792,-4.020) (-8.723,-4.049) (-8.654,-4.077) (-8.584,-4.104) (-8.514,-4.131) (-8.443,-4.156) (-8.372,-4.181) (-8.301,-4.204) (-8.229,-4.227) (-8.158,-4.248) (-8.086,-4.269) (-8.013,-4.289) (-7.941,-4.308) (-7.868,-4.326) (-7.795,-4.344) (-7.722,-4.360) (-7.648,-4.376) (-7.575,-4.390) (-7.501,-4.404) (-7.427,-4.417) (-7.353,-4.429) (-7.279,-4.440) (-7.205,-4.450) (-7.130,-4.459) (-7.056,-4.468) (-6.981,-4.476) (-6.906,-4.482) (-6.832,-4.488) (-6.757,-4.493) (-6.682,-4.498) (-6.607,-4.501) (-6.532,-4.503) (-6.457,-4.505) (-6.382,-4.506) (-6.307,-4.506) (-6.232,-4.505) (-6.157,-4.504) (-6.082,-4.501) (-6.007,-4.498) (-5.932,-4.494) (-5.858,-4.489) (-5.783,-4.483) (-5.708,-4.477) (-5.633,-4.469) (-5.559,-4.461) (-5.484,-4.452) (-5.410,-4.443) (-5.336,-4.432) (-5.262,-4.421) (-5.188,-4.409) (-5.114,-4.396) (-5.040,-4.382) (-4.966,-4.368) (-4.893,-4.353) (-4.820,-4.337) (-4.746,-4.320) (-4.674,-4.303) (-4.601,-4.285) (-4.528,-4.266) (-4.456,-4.246) (-4.384,-4.226) (-4.312,-4.204) (-4.288,-4.197)}
\LC{(0.452,3.470) (0.467,3.543) (0.482,3.617) (0.497,3.690) (0.512,3.764) (0.528,3.837) (0.544,3.910) (0.560,3.984) (0.576,4.057) (0.592,4.130) (0.608,4.203) (0.624,4.277) (0.641,4.350) (0.658,4.423) (0.674,4.496) (0.691,4.569) (0.709,4.642) (0.726,4.715) (0.743,4.788) (0.761,4.861) (0.778,4.934) (0.796,5.007) (0.814,5.080) (0.832,5.152) (0.850,5.225) (0.868,5.298) (0.887,5.371) (0.905,5.443) (0.924,5.516) (0.943,5.588) (0.962,5.661) (0.981,5.734) (1.000,5.806) (1.019,5.879) (1.038,5.951) (1.058,6.023) (1.078,6.096) (1.097,6.168) (1.117,6.241) (1.137,6.313) (1.157,6.385) (1.177,6.457) (1.198,6.530) (1.218,6.602) (1.239,6.674) (1.259,6.746) (1.280,6.818) (1.301,6.890) (1.322,6.962) (1.343,7.034) (1.364,7.106) (1.385,7.178) (1.407,7.250) (1.428,7.322) (1.450,7.393) (1.472,7.465) (1.493,7.537) (1.515,7.609) (1.537,7.680) (1.560,7.752) (1.582,7.824) (1.604,7.895) (1.627,7.967) (1.649,8.038) (1.672,8.110) (1.695,8.181) (1.717,8.253) (1.740,8.324) (1.763,8.396) (1.787,8.467) (1.810,8.538) (1.833,8.609) (1.857,8.681) (1.880,8.752) (1.904,8.823) (1.927,8.894) (1.951,8.965) (1.975,9.036) (1.999,9.108) (2.023,9.179) (2.048,9.250) (2.072,9.320) (2.096,9.391) (2.121,9.462) (2.145,9.533) (2.170,9.604) (2.195,9.675) (2.220,9.746) (2.245,9.816) (2.270,9.887) (2.295,9.958) (2.311,10.000)}
\LC{(10.000,-9.011) (9.930,-8.994) (9.857,-8.977) (9.784,-8.959) (9.711,-8.941) (9.638,-8.923) (9.566,-8.904) (9.493,-8.885) (9.421,-8.866) (9.348,-8.846) (9.276,-8.826) (9.204,-8.805) (9.132,-8.784) (9.060,-8.763) (8.988,-8.741) (8.916,-8.719) (8.845,-8.697) (8.773,-8.674) (8.702,-8.651) (8.631,-8.628) (8.560,-8.604) (8.489,-8.580) (8.418,-8.556) (8.347,-8.531) (8.276,-8.506) (8.206,-8.480) (8.135,-8.454) (8.065,-8.428) (7.995,-8.402) (7.925,-8.375) (7.855,-8.348) (7.785,-8.320) (7.716,-8.292) (7.646,-8.264) (7.577,-8.235) (7.508,-8.206) (7.439,-8.177) (7.370,-8.147) (7.301,-8.117) (7.232,-8.087) (7.164,-8.056) (7.096,-8.025) (7.027,-7.994) (6.959,-7.963) (6.892,-7.931) (6.824,-7.898) (6.756,-7.866) (6.689,-7.833) (6.622,-7.799) (6.555,-7.766) (6.488,-7.732) (6.421,-7.697) (6.355,-7.663) (6.288,-7.628) (6.222,-7.592) (6.156,-7.557) (6.090,-7.521) (6.025,-7.484) (5.959,-7.448) (5.894,-7.411) (5.829,-7.374) (5.764,-7.336) (5.699,-7.298) (5.635,-7.260) (5.571,-7.221) (5.506,-7.182) (5.442,-7.143) (5.379,-7.104) (5.315,-7.064) (5.252,-7.024) (5.189,-6.983) (5.126,-6.942) (5.063,-6.901) (5.001,-6.860) (4.938,-6.818) (4.876,-6.776) (4.814,-6.734) (4.753,-6.691) (4.712,-6.662)}
\LC{(4.295,5.298) (4.361,5.334) (4.427,5.369) (4.494,5.404) (4.561,5.438) (4.628,5.472) (4.695,5.505) (4.762,5.538) (4.830,5.570) (4.898,5.602) (4.966,5.634) (5.034,5.664) (5.103,5.695) (5.172,5.725) (5.240,5.754) (5.310,5.783) (5.379,5.812) (5.449,5.840) (5.518,5.867) (5.588,5.894) (5.658,5.921) (5.729,5.947) (5.799,5.972) (5.870,5.998) (5.941,6.022) (6.012,6.046) (6.083,6.070) (6.154,6.093) (6.226,6.115) (6.298,6.137) (6.370,6.159) (6.442,6.180) (6.514,6.201) (6.586,6.221) (6.658,6.240) (6.731,6.259) (6.804,6.278) (6.876,6.295) (6.949,6.313) (7.022,6.330) (7.096,6.346) (7.169,6.362) (7.242,6.377) (7.316,6.392) (7.389,6.407) (7.463,6.420) (7.537,6.434) (7.611,6.446) (7.685,6.458) (7.759,6.470) (7.833,6.481) (7.907,6.492) (7.982,6.502) (8.056,6.511) (8.131,6.520) (8.205,6.529) (8.280,6.536) (8.354,6.544) (8.429,6.551) (8.504,6.557) (8.579,6.562) (8.653,6.567) (8.728,6.572) (8.803,6.576) (8.878,6.579) (8.953,6.582) (9.028,6.585) (9.103,6.586) (9.178,6.588) (9.253,6.588) (9.328,6.588) (9.403,6.588) (9.478,6.587) (9.553,6.585) (9.628,6.583) (9.703,6.580) (9.778,6.576) (9.853,6.572) (9.928,6.568) (10.000,6.563)}
\LC{(0.445,6.219) (0.453,6.293) (0.461,6.368) (0.470,6.443) (0.478,6.517) (0.486,6.592) (0.495,6.666) (0.503,6.741) (0.511,6.815) (0.520,6.890) (0.528,6.964) (0.537,7.039) (0.546,7.113) (0.554,7.188) (0.563,7.262) (0.572,7.337) (0.581,7.411) (0.590,7.486) (0.598,7.560) (0.607,7.635) (0.616,7.709) (0.625,7.784) (0.634,7.858) (0.644,7.932) (0.653,8.007) (0.662,8.081) (0.671,8.156) (0.680,8.230) (0.690,8.305) (0.699,8.379) (0.709,8.453) (0.718,8.528) (0.728,8.602) (0.737,8.677) (0.747,8.751) (0.756,8.825) (0.766,8.900) (0.776,8.974) (0.785,9.048) (0.795,9.123) (0.805,9.197) (0.815,9.272) (0.825,9.346) (0.835,9.420) (0.845,9.495) (0.855,9.569) (0.865,9.643) (0.875,9.718) (0.885,9.792) (0.895,9.866) (0.905,9.940) (0.915,10.000)}
\LC{(-0.456,5.643) (-0.465,5.718) (-0.474,5.792) (-0.483,5.867) (-0.493,5.941) (-0.502,6.015) (-0.512,6.090) (-0.521,6.164) (-0.531,6.239) (-0.540,6.313) (-0.550,6.387) (-0.560,6.462) (-0.570,6.536) (-0.579,6.610) (-0.589,6.685) (-0.599,6.759) (-0.609,6.833) (-0.619,6.908) (-0.629,6.982) (-0.640,7.056) (-0.650,7.131) (-0.660,7.205) (-0.671,7.279) (-0.681,7.353) (-0.691,7.428) (-0.702,7.502) (-0.712,7.576) (-0.723,7.650) (-0.734,7.725) (-0.744,7.799) (-0.755,7.873) (-0.766,7.947) (-0.777,8.022) (-0.788,8.096) (-0.799,8.170) (-0.810,8.244) (-0.821,8.318) (-0.832,8.393) (-0.843,8.467) (-0.854,8.541) (-0.866,8.615) (-0.877,8.689) (-0.888,8.763) (-0.900,8.837) (-0.911,8.912) (-0.923,8.986) (-0.934,9.060) (-0.946,9.134) (-0.957,9.208) (-0.969,9.282) (-0.981,9.356) (-0.993,9.430) (-1.004,9.504) (-1.016,9.578) (-1.028,9.652) (-1.040,9.726) (-1.052,9.800) (-1.064,9.874) (-1.076,9.948) (-1.088,10.000)}
\LC{(1.026,1.052) (1.099,1.067) (1.173,1.078) (1.248,1.087) (1.322,1.093) (1.397,1.095) (1.472,1.094) (1.547,1.090) (1.622,1.082) (1.696,1.071) (1.770,1.057) (1.842,1.039) (1.914,1.017) (1.985,0.992) (2.054,0.963) (2.121,0.930) (2.187,0.893) (2.250,0.853) (2.311,0.809) (2.369,0.761) (2.423,0.710) (2.475,0.655) (2.522,0.597) (2.566,0.536) (2.605,0.472) (2.639,0.405) (2.668,0.336) (2.692,0.265) (2.710,0.192) (2.722,0.119) (2.729,0.044) (2.729,-0.031) (2.724,-0.106) (2.712,-0.180) (2.695,-0.253) (2.672,-0.324) (2.644,-0.394) (2.611,-0.461) (2.573,-0.526) (2.530,-0.587) (2.483,-0.646) (2.432,-0.701) (2.378,-0.753) (2.321,-0.801) (2.260,-0.846) (2.198,-0.887) (2.132,-0.924) (2.065,-0.957) (1.996,-0.987) (1.926,-1.013) (1.855,-1.035) (1.782,-1.054) (1.709,-1.069) (1.634,-1.081) (1.560,-1.089) (1.485,-1.094) (1.410,-1.095) (1.335,-1.093) (1.260,-1.088) (1.186,-1.080) (1.112,-1.069) (1.038,-1.055)}
\LC{(1.474,-10.000) (1.458,-9.936) (1.442,-9.863) (1.425,-9.789) (1.409,-9.716) (1.393,-9.643) (1.377,-9.570) (1.361,-9.496) (1.345,-9.423) (1.329,-9.350) (1.313,-9.276) (1.298,-9.203) (1.282,-9.130) (1.266,-9.056) (1.251,-8.983) (1.235,-8.910) (1.220,-8.836) (1.204,-8.763) (1.189,-8.689) (1.174,-8.616) (1.159,-8.543) (1.143,-8.469) (1.128,-8.396) (1.113,-8.322) (1.098,-8.249) (1.084,-8.175) (1.069,-8.102) (1.054,-8.028) (1.039,-7.955) (1.025,-7.881) (1.010,-7.807) (0.996,-7.734) (0.982,-7.660) (0.967,-7.587) (0.953,-7.513) (0.939,-7.439) (0.925,-7.366) (0.911,-7.292) (0.897,-7.218) (0.883,-7.145) (0.869,-7.071) (0.855,-6.997) (0.842,-6.923) (0.828,-6.850) (0.814,-6.776) (0.801,-6.702) (0.788,-6.628) (0.774,-6.554) (0.761,-6.481) (0.748,-6.407) (0.735,-6.333) (0.722,-6.259) (0.709,-6.185) (0.696,-6.111) (0.683,-6.037) (0.670,-5.963) (0.658,-5.890) (0.645,-5.816) (0.633,-5.742) (0.620,-5.668) (0.608,-5.594) (0.596,-5.520) (0.584,-5.446) (0.572,-5.372)}
\LC{(3.662,1.966) (3.733,1.944) (3.805,1.921) (3.876,1.896) (3.946,1.869) (4.015,1.840) (4.083,1.809) (4.150,1.776) (4.217,1.741) (4.282,1.705) (4.346,1.666) (4.409,1.625) (4.471,1.582) (4.531,1.537) (4.589,1.490) (4.646,1.442) (4.702,1.391) (4.755,1.339) (4.807,1.285) (4.857,1.229) (4.905,1.171) (4.951,1.111) (4.994,1.050) (5.035,0.987) (5.074,0.923) (5.110,0.858) (5.144,0.791) (5.175,0.722) (5.203,0.653) (5.229,0.582) (5.252,0.511) (5.271,0.439) (5.288,0.366) (5.302,0.292) (5.313,0.218) (5.320,0.143) (5.324,0.068) (5.326,-0.007) (5.324,-0.082) (5.319,-0.157) (5.311,-0.231) (5.300,-0.305) (5.285,-0.379) (5.268,-0.452) (5.248,-0.524) (5.225,-0.596) (5.198,-0.666) (5.170,-0.735) (5.138,-0.803) (5.104,-0.870) (5.067,-0.935) (5.028,-0.999) (4.986,-1.062) (4.942,-1.122) (4.896,-1.182) (4.848,-1.239) (4.798,-1.295) (4.746,-1.349) (4.692,-1.401) (4.636,-1.451) (4.579,-1.499) (4.520,-1.545) (4.459,-1.590) (4.398,-1.632) (4.334,-1.673) (4.270,-1.711) (4.205,-1.748) (4.138,-1.783) (4.070,-1.815) (4.002,-1.846) (3.933,-1.874) (3.863,-1.901) (3.792,-1.926) (3.720,-1.948) (3.648,-1.969)}
\LC{(-7.036,-10.000) (-6.975,-9.958) (-6.914,-9.914) (-6.854,-9.870) (-6.793,-9.826) (-6.732,-9.781) (-6.672,-9.737) (-6.612,-9.692) (-6.552,-9.647) (-6.492,-9.602) (-6.432,-9.556) (-6.373,-9.511) (-6.314,-9.465) (-6.254,-9.419) (-6.195,-9.373) (-6.136,-9.326) (-6.078,-9.280) (-6.019,-9.233) (-5.961,-9.186) (-5.902,-9.139) (-5.844,-9.091) (-5.786,-9.044) (-5.728,-8.996) (-5.671,-8.948) (-5.613,-8.899) (-5.556,-8.851) (-5.499,-8.802) (-5.442,-8.754) (-5.385,-8.705) (-5.329,-8.655) (-5.273,-8.606) (-5.216,-8.556) (-5.160,-8.506) (-5.104,-8.456) (-5.049,-8.406) (-4.993,-8.356) (-4.938,-8.305) (-4.883,-8.254) (-4.828,-8.203) (-4.773,-8.152) (-4.718,-8.100) (-4.664,-8.049) (-4.610,-7.997) (-4.556,-7.945) (-4.502,-7.893) (-4.448,-7.840) (-4.395,-7.788) (-4.342,-7.735) (-4.288,-7.682) (-4.236,-7.629) (-4.183,-7.575) (-4.130,-7.522) (-4.078,-7.468) (-4.026,-7.414) (-3.974,-7.360) (-3.923,-7.306) (-3.871,-7.251) (-3.820,-7.196) (-3.769,-7.141) (-3.718,-7.086) (-3.667,-7.031) (-3.617,-6.975) (-3.567,-6.920) (-3.517,-6.864) (-3.467,-6.808) (-3.417,-6.751) (-3.368,-6.695) (-3.319,-6.638) (-3.270,-6.581) (-3.221,-6.524) (-3.173,-6.467) (-3.125,-6.410) (-3.077,-6.352) (-3.029,-6.294) (-2.981,-6.236) (-2.934,-6.178) (-2.887,-6.120) (-2.840,-6.061) (-2.793,-6.002) (-2.747,-5.943) (-2.701,-5.884) (-2.655,-5.825) (-2.609,-5.765) (-2.564,-5.706) (-2.519,-5.646) (-2.474,-5.586) (-2.429,-5.525) (-2.385,-5.465)}
\LC{(10.000,-2.634) (9.966,-2.689) (9.914,-2.743) (9.861,-2.796) (9.807,-2.848) (9.753,-2.900) (9.697,-2.950) (9.641,-3.000) (9.584,-3.048) (9.526,-3.096) (9.468,-3.143) (9.408,-3.189) (9.348,-3.234) (9.287,-3.278) (9.226,-3.321) (9.164,-3.363) (9.101,-3.404) (9.038,-3.444) (8.974,-3.483) (8.909,-3.522) (8.844,-3.559) (8.779,-3.595) (8.713,-3.631) (8.646,-3.665) (8.579,-3.698) (8.511,-3.731) (8.443,-3.762) (8.375,-3.793) (8.306,-3.823) (8.236,-3.851) (8.167,-3.879) (8.097,-3.906) (8.026,-3.931) (7.955,-3.956) (7.884,-3.980) (7.813,-4.003) (7.741,-4.025) (7.669,-4.046) (7.597,-4.066) (7.524,-4.085) (7.451,-4.103) (7.378,-4.120) (7.305,-4.136) (7.232,-4.151) (7.158,-4.166) (7.084,-4.179) (7.011,-4.192) (6.936,-4.203) (6.862,-4.214) (6.788,-4.224) (6.713,-4.233) (6.639,-4.241) (6.564,-4.248) (6.489,-4.254) (6.415,-4.259) (6.340,-4.263) (6.265,-4.267) (6.190,-4.269) (6.115,-4.271) (6.040,-4.272) (5.965,-4.272) (5.890,-4.271) (5.815,-4.269) (5.740,-4.266) (5.665,-4.263) (5.590,-4.258) (5.515,-4.253) (5.441,-4.247) (5.366,-4.240) (5.291,-4.232) (5.217,-4.224) (5.142,-4.214) (5.068,-4.204) (4.994,-4.193) (4.920,-4.181) (4.846,-4.168) (4.772,-4.154) (4.699,-4.140) (4.625,-4.125) (4.552,-4.109) (4.479,-4.092) (4.455,-4.086)}
\LC{(-10.000,-7.077) (-9.926,-7.077) (-9.851,-7.077) (-9.776,-7.075) (-9.701,-7.074) (-9.626,-7.071) (-9.552,-7.068) (-9.477,-7.065) (-9.402,-7.061) (-9.327,-7.057) (-9.252,-7.052) (-9.177,-7.047) (-9.102,-7.041) (-9.028,-7.034) (-8.953,-7.027) (-8.878,-7.020) (-8.804,-7.012) (-8.729,-7.003) (-8.655,-6.994) (-8.580,-6.985) (-8.506,-6.975) (-8.432,-6.964) (-8.358,-6.953) (-8.284,-6.942) (-8.210,-6.930) (-8.136,-6.918) (-8.062,-6.905) (-7.988,-6.891) (-7.914,-6.877) (-7.841,-6.863) (-7.767,-6.848) (-7.694,-6.832) (-7.620,-6.816) (-7.547,-6.800) (-7.474,-6.783) (-7.401,-6.766) (-7.328,-6.748) (-7.256,-6.729) (-7.183,-6.711) (-7.111,-6.691) (-7.038,-6.671) (-6.966,-6.651) (-6.894,-6.630) (-6.822,-6.609) (-6.750,-6.588) (-6.679,-6.565) (-6.607,-6.543) (-6.536,-6.520) (-6.465,-6.496) (-6.393,-6.472) (-6.323,-6.448) (-6.252,-6.423) (-6.181,-6.397) (-6.111,-6.371) (-6.041,-6.345) (-5.971,-6.318) (-5.901,-6.291) (-5.831,-6.263) (-5.762,-6.235) (-5.692,-6.206) (-5.623,-6.177) (-5.554,-6.148) (-5.485,-6.118) (-5.417,-6.087) (-5.349,-6.057) (-5.280,-6.025) (-5.212,-5.993) (-5.145,-5.961) (-5.077,-5.929) (-5.010,-5.896)}
\LC{(-1.487,0.918) (-1.560,0.900) (-1.632,0.878) (-1.702,0.852) (-1.771,0.822) (-1.837,0.787) (-1.901,0.748) (-1.962,0.705) (-2.021,0.657) (-2.075,0.606) (-2.125,0.550) (-2.171,0.491) (-2.212,0.428) (-2.247,0.362) (-2.277,0.293) (-2.301,0.222) (-2.318,0.149) (-2.328,0.075) (-2.331,-0.000) (-2.328,-0.075) (-2.318,-0.149) (-2.300,-0.222) (-2.277,-0.294) (-2.247,-0.362) (-2.212,-0.428) (-2.171,-0.491) (-2.125,-0.551) (-2.075,-0.606) (-2.020,-0.658) (-1.962,-0.705) (-1.901,-0.748) (-1.837,-0.787) (-1.770,-0.822) (-1.702,-0.852) (-1.631,-0.879) (-1.560,-0.901) (-1.487,-0.918)}
\LC{(6.535,-10.000) (6.476,-9.962) (6.418,-9.915) (6.360,-9.867) (6.302,-9.820) (6.244,-9.772) (6.186,-9.724) (6.129,-9.676) (6.071,-9.628) (6.014,-9.579) (5.957,-9.531) (5.900,-9.482) (5.843,-9.433) (5.787,-9.384) (5.730,-9.334) (5.674,-9.285) (5.618,-9.235) (5.562,-9.185) (5.506,-9.135) (5.451,-9.084) (5.395,-9.034) (5.340,-8.983) (5.285,-8.932) (5.230,-8.881) (5.175,-8.830) (5.120,-8.779) (5.066,-8.727) (5.012,-8.675) (4.958,-8.623) (4.904,-8.571) (4.850,-8.519) (4.796,-8.467) (4.743,-8.414) (4.690,-8.361) (4.637,-8.308) (4.584,-8.255) (4.531,-8.201) (4.479,-8.148) (4.426,-8.094) (4.374,-8.040) (4.322,-7.986) (4.271,-7.932) (4.219,-7.877) (4.168,-7.823) (4.117,-7.768) (4.066,-7.713) (4.015,-7.657) (3.964,-7.602) (3.914,-7.547) (3.864,-7.491) (3.814,-7.435) (3.764,-7.379) (3.714,-7.323) (3.665,-7.266) (3.616,-7.210) (3.567,-7.153) (3.518,-7.096) (3.469,-7.039) (3.421,-6.981) (3.373,-6.924) (3.325,-6.866) (3.277,-6.808) (3.229,-6.750) (3.182,-6.692) (3.135,-6.634) (3.088,-6.575) (3.041,-6.517) (2.995,-6.458) (2.949,-6.399) (2.903,-6.339) (2.857,-6.280) (2.811,-6.220) (2.766,-6.161) (2.721,-6.101) (2.676,-6.041) (2.632,-5.980) (2.587,-5.920) (2.543,-5.859)}
\LC{(-2.806,6.393) (-2.851,6.453) (-2.896,6.513) (-2.941,6.573) (-2.986,6.633) (-3.031,6.693) (-3.077,6.752) (-3.123,6.812) (-3.169,6.871) (-3.215,6.930) (-3.261,6.989) (-3.308,7.048) (-3.355,7.106) (-3.402,7.165) (-3.449,7.223) (-3.496,7.281) (-3.544,7.339) (-3.592,7.397) (-3.640,7.454) (-3.688,7.512) (-3.736,7.569) (-3.785,7.626) (-3.834,7.683) (-3.883,7.740) (-3.932,7.797) (-3.981,7.853) (-4.031,7.909) (-4.081,7.966) (-4.131,8.021) (-4.181,8.077) (-4.231,8.133) (-4.282,8.188) (-4.332,8.244) (-4.383,8.299) (-4.434,8.354) (-4.485,8.409) (-4.537,8.463) (-4.588,8.518) (-4.640,8.572) (-4.692,8.626) (-4.744,8.680) (-4.796,8.734) (-4.849,8.788) (-4.901,8.841) (-4.954,8.894) (-5.007,8.947) (-5.060,9.000) (-5.113,9.053) (-5.167,9.106) (-5.220,9.158) (-5.274,9.211) (-5.328,9.263) (-5.382,9.315) (-5.436,9.366) (-5.491,9.418) (-5.545,9.470) (-5.600,9.521) (-5.655,9.572) (-5.710,9.623) (-5.765,9.674) (-5.821,9.724) (-5.876,9.775) (-5.932,9.825) (-5.988,9.875) (-6.044,9.925) (-6.100,9.975) (-6.137,10.000)}
\LC{(-1.501,-10.000) (-1.485,-9.935) (-1.468,-9.862) (-1.451,-9.789) (-1.435,-9.716) (-1.418,-9.643) (-1.402,-9.569) (-1.386,-9.496) (-1.370,-9.423) (-1.353,-9.350) (-1.337,-9.276) (-1.321,-9.203) (-1.305,-9.130) (-1.289,-9.057) (-1.274,-8.983) (-1.258,-8.910) (-1.242,-8.837) (-1.226,-8.763) (-1.211,-8.690) (-1.195,-8.617) (-1.180,-8.543) (-1.164,-8.470) (-1.149,-8.396) (-1.134,-8.323) (-1.119,-8.249) (-1.104,-8.176) (-1.089,-8.103) (-1.074,-8.029) (-1.059,-7.956) (-1.044,-7.882) (-1.029,-7.808) (-1.014,-7.735) (-1.000,-7.661) (-0.985,-7.588) (-0.971,-7.514) (-0.956,-7.441) (-0.942,-7.367) (-0.928,-7.293) (-0.913,-7.220) (-0.899,-7.146) (-0.885,-7.072) (-0.871,-6.999) (-0.857,-6.925) (-0.843,-6.851) (-0.830,-6.778) (-0.816,-6.704) (-0.802,-6.630) (-0.789,-6.556) (-0.775,-6.483) (-0.762,-6.409) (-0.748,-6.335) (-0.735,-6.261) (-0.722,-6.187) (-0.709,-6.113) (-0.696,-6.040) (-0.683,-5.966) (-0.670,-5.892) (-0.657,-5.818) (-0.645,-5.744) (-0.632,-5.670) (-0.620,-5.596) (-0.607,-5.522) (-0.595,-5.448) (-0.582,-5.374) (-0.570,-5.300) (-0.558,-5.226) (-0.546,-5.152) (-0.534,-5.078) (-0.522,-5.004)}
\LC{(-5.630,4.299) (-5.705,4.304) (-5.780,4.308) (-5.855,4.311) (-5.929,4.313) (-6.004,4.314) (-6.079,4.314) (-6.154,4.314) (-6.229,4.312) (-6.304,4.310) (-6.379,4.307) (-6.454,4.303) (-6.529,4.298) (-6.604,4.292) (-6.679,4.285) (-6.753,4.278) (-6.828,4.269) (-6.902,4.260) (-6.976,4.249) (-7.051,4.238) (-7.125,4.226) (-7.198,4.213) (-7.272,4.199) (-7.346,4.184) (-7.419,4.168) (-7.492,4.151) (-7.565,4.134) (-7.638,4.115) (-7.710,4.096) (-7.782,4.075) (-7.854,4.054) (-7.926,4.031) (-7.997,4.008) (-8.068,3.984) (-8.138,3.959) (-8.209,3.932) (-8.279,3.905) (-8.348,3.877) (-8.417,3.848) (-8.486,3.818) (-8.554,3.787) (-8.622,3.755) (-8.690,3.723) (-8.757,3.689) (-8.823,3.654) (-8.889,3.618) (-8.955,3.582) (-9.019,3.544) (-9.084,3.506) (-9.148,3.466) (-9.211,3.426) (-9.273,3.384) (-9.335,3.342) (-9.396,3.299) (-9.457,3.254) (-9.517,3.209) (-9.576,3.163) (-9.634,3.116) (-9.692,3.068) (-9.749,3.019) (-9.805,2.969) (-9.860,2.919) (-9.915,2.867) (-9.968,2.815) (-10.000,2.779)}

    % Dipolo físico - puntas de flecha
% Generado por scripts/generar_datos.py (no editar a mano).
\FL{1.382}{0.003}{-89.7}
\FL{1.658}{2.024}{26.2}
\FL{4.571}{2.223}{-20.4}
\FL{6.266}{-0.023}{-90.6}
\FL{4.528}{-2.238}{-160.4}
\FL{1.617}{-2.004}{153.2}
\FL{-1.658}{2.024}{153.8}
\FL{-4.571}{2.223}{-159.6}
\FL{-6.266}{-0.023}{-89.4}
\FL{-4.528}{-2.238}{-19.6}
\FL{-1.617}{-2.004}{26.8}
\FL{-1.382}{0.003}{-90.3}
\FL{3.873}{-8.712}{126.6}
\FL{2.098}{-5.992}{119.5}
\FL{0.715}{-3.054}{110.5}
\FL{-4.032}{-8.737}{52.1}
\FL{-2.189}{-6.033}{59.5}
\FL{-0.751}{-3.094}{68.8}
\FL{8.276}{-6.267}{176.4}
\FL{4.938}{-5.465}{156.7}
\FL{2.064}{-3.588}{136.7}
\FL{0.802}{2.873}{65.6}
\FL{2.098}{5.181}{56.2}
\FL{3.739}{7.291}{48.3}
\FL{5.620}{9.155}{41.3}
\FL{4.821}{2.884}{-9.5}
\FL{7.202}{1.384}{-58.2}
\FL{7.199}{-1.389}{-121.9}
\FL{4.816}{-2.885}{-170.6}
\FL{-2.682}{3.198}{153.6}
\FL{-5.281}{3.843}{178.8}
\FL{-7.883}{3.242}{-151.6}
\FL{-9.692}{1.308}{-112.7}
\FL{-9.687}{-1.320}{-67.1}
\FL{-7.871}{-3.248}{-28.2}
\FL{-5.268}{-3.843}{1.3}
\FL{-2.670}{-3.192}{26.5}
\FL{-0.035}{-8.072}{89.6}
\FL{-0.013}{-4.247}{89.7}
\FL{-1.003}{3.033}{118.6}
\FL{-2.708}{5.495}{130.4}
\FL{-4.844}{7.595}{140.4}
\FL{-7.300}{9.313}{149.6}
\FL{-3.060}{4.764}{141.0}
\FL{-5.630}{6.344}{155.7}
\FL{-8.515}{7.227}{170.3}
\FL{-8.690}{-7.904}{13.4}
\FL{-6.190}{-7.039}{24.8}
\FL{-3.910}{-5.697}{36.3}
\FL{-1.957}{-3.913}{48.8}
\FL{-7.291}{2.794}{-148.2}
\FL{-8.947}{-0.031}{-89.4}
\FL{-7.260}{-2.813}{-31.3}
\FL{8.390}{-7.354}{168.6}
\FL{5.326}{-6.297}{153.4}
\FL{2.988}{3.361}{23.6}
\FL{6.362}{3.813}{-9.2}
\FL{9.300}{2.161}{-52.1}
\FL{-4.856}{3.051}{-173.2}
\FL{-7.529}{1.485}{-122.6}
\FL{-7.511}{-1.513}{-56.8}
\FL{-4.823}{-3.055}{-6.4}
\FL{3.610}{6.131}{42.8}
\FL{5.966}{7.941}{32.3}
\FL{8.606}{9.304}{22.3}
\FL{3.458}{1.326}{-32.4}
\FL{3.422}{-1.349}{-148.8}
\FL{2.162}{-8.591}{111.4}
\FL{1.164}{-5.763}{107.4}
\FL{-4.971}{6.769}{146.7}
\FL{-8.260}{8.387}{160.9}
\FL{9.417}{0.793}{-75.7}
\FL{8.741}{-2.115}{-129.3}
\FL{6.090}{-3.595}{-170.0}
\FL{3.046}{-3.272}{158.7}
\FL{7.948}{-9.304}{153.7}
\FL{5.021}{-7.488}{142.5}
\FL{2.520}{-5.122}{130.4}
\FL{-8.680}{-5.785}{-3.3}
\FL{-6.026}{-5.537}{13.8}
\FL{-3.553}{-4.540}{30.1}
\FL{8.604}{-5.222}{-171.6}
\FL{5.771}{-5.124}{168.1}
\FL{-0.956}{4.794}{107.2}
\FL{-2.251}{8.312}{113.0}
\FL{-8.277}{-8.959}{22.4}
\FL{-5.083}{-7.207}{35.2}
\FL{-2.318}{-8.626}{67.2}
\FL{-1.257}{-5.820}{71.5}
\FL{-0.004}{4.753}{90.1}
\FL{-0.010}{8.253}{90.1}
\FL{6.061}{4.483}{2.2}
\FL{8.816}{3.979}{-23.8}
\FL{-4.746}{4.560}{164.0}
\FL{-8.332}{4.699}{-167.7}
\FL{4.558}{6.008}{31.8}
\FL{8.109}{7.522}{14.4}
\FL{-2.869}{1.182}{-150.5}
\FL{-2.846}{-1.194}{-28.6}
\FL{-4.231}{1.380}{-140.6}
\FL{-4.219}{-1.390}{-39.0}
\FL{-5.708}{5.696}{161.9}
\FL{-8.556}{6.180}{179.0}
\FL{3.472}{4.284}{28.4}
\FL{6.008}{5.179}{10.4}
\FL{8.695}{5.222}{-8.9}
\FL{2.008}{6.220}{62.7}
\FL{3.508}{8.788}{56.8}
\FL{-2.133}{6.128}{119.3}
\FL{-3.837}{8.775}{126.1}
\FL{-8.700}{-4.059}{-22.1}
\FL{-5.758}{-4.481}{4.9}
\FL{0.826}{5.128}{76.1}
\FL{1.763}{8.396}{72.1}
\FL{8.583}{-8.612}{161.5}
\FL{5.938}{-7.436}{150.5}
\FL{5.658}{5.921}{20.5}
\FL{8.529}{6.559}{4.4}
\FL{0.668}{8.131}{82.9}
\FL{-0.752}{7.848}{98.3}
\FL{2.730}{-0.006}{-90.4}
\FL{0.986}{-7.685}{101.0}
\FL{4.841}{1.247}{-48.7}
\FL{4.815}{-1.276}{-132.4}
\FL{-5.728}{-8.996}{39.7}
\FL{-3.401}{-6.732}{48.8}
\FL{8.823}{-3.571}{-150.9}
\FL{5.940}{-4.272}{179.4}
\FL{-7.426}{-6.771}{13.5}
\FL{-2.331}{-0.025}{-88.2}
\FL{5.432}{-9.068}{137.7}
\FL{3.421}{-6.981}{130.1}
\FL{-4.383}{8.299}{132.8}
\FL{-0.966}{-7.490}{78.9}
\FL{-7.997}{4.008}{-161.5}
}
  \fill[mmarojo] (0,0.5) circle (0.2) (0,-0.5) circle (0.2);
  \node[font=\scriptsize, text=mmarojo, fill=white, fill opacity=0.85, text opacity=1, inner sep=0.6pt, anchor=west] at (0.4,0.62) {$+q$};
  \node[font=\scriptsize, text=mmarojo, fill=white, fill opacity=0.85, text opacity=1, inner sep=0.6pt, anchor=west] at (0.4,-0.62) {$-q$};
  \marcoMMA{10}{5}{$x/\ell$}{$z/\ell$}
\end{tikzpicture}\\[1pt]{\small\scshape (b)\; dipolo físico}
\end{minipage}
\caption{Equipotenciales (gris) y líneas de campo (azul) de (a)~un dipolo puntual $\vp=q\ell\,\uz$ y (b)~un dipolo físico con $+q$ en
$z=\ell/2$ y $-q$ en $z=-\ell/2$, ambos con los mismos niveles de potencial. Fuera del disco central, de radio $1{,}2\,\ell$, las dos
representaciones coinciden a simple vista; las pequeñas diferencias se concentran cerca de las cargas y en los «haces» de equipotenciales que
convergen hacia el plano ecuatorial. Recreación de~\cite{FIGURAS-T2}.}
\label{fig:fisico-puntual}
\end{figure}
